
\documentclass{article} % For LaTeX2e
\usepackage{iclr2026_conference,times}

% Optional math commands from https://github.com/goodfeli/dlbook_notation.
\input{math_commands.tex}

\usepackage{hyperref}
\usepackage{url}
\usepackage{algorithm}
\usepackage{algpseudocode}

%%%%%%%%%%%%%%%SAHANA ADDED THIS %%%%%%%%%%%
\usepackage{amsthm}
\usepackage{cleveref}
\usepackage{graphicx}
\newtheorem{theorem}{Theorem}[section]
\newtheorem{lemma}[theorem]{Lemma}
\newtheorem{claim}[theorem]{Claim}
\newtheorem{proposition}[theorem]{Proposition}
\newtheorem{corollary}[theorem]{Corollary}
\newtheorem{fact}[theorem]{Fact}
\newtheorem{example}[theorem]{Example}
\newtheorem{notation}[theorem]{Notation}
\newtheorem{observation}[theorem]{Observation}
\newtheorem{conjecture}[theorem]{Conjecture}

\newtheorem{innercustomgeneric}{\customgenericname}
\providecommand{\customgenericname}{}
\newcommand{\newcustomtheorem}[2]{%
  \newenvironment{#1}[1]
  {%
   \renewcommand\customgenericname{#2}%
   \renewcommand\theinnercustomgeneric{##1}%
   \innercustomgeneric
  }
  {\endinnercustomgeneric}
}

\newcustomtheorem{customtheorem}{Theorem}
\newcustomtheorem{customlemma}{Lemma}
\newcustomtheorem{customclaim}{Claim}
\newcustomtheorem{customproposition}{Proposition}
\newcustomtheorem{customcorollary}{Corollary}
\newcustomtheorem{customfact}{Fact}
\newcustomtheorem{customexample}{Example}
\newcustomtheorem{customnotation}{Notation}
\newcustomtheorem{customobservation}{Observation}
\newcustomtheorem{customconjecture}{Conjecture}
%%%%%%%%%%%%%%%%%%%%%%%%%%%%%%%%%%%%%%%%


\title{Greedy Hybrid PDE Solver with \\ Neural Operators and Traditional solvers}

% Authors must not appear in the submitted version. They should be hidden
% as long as the \iclrfinalcopy macro remains commented out below.
% Non-anonymous submissions will be rejected without review.

\author{Sahana Rayan \\
Department of Statistics\\
University of Michigan\\
Ann Arbor, MI 48104, USA \\
\texttt{\{srayan\}@umich.edu} \\
\And
Yash Patel \\
Department of Statistics\\
University of Michigan\\
Ann Arbor, MI 48104, USA \\
\texttt{\{yppatel\}@umich.edu} \\
\And
Ambuj Tewari \\
Department of Statistics\\
University of Michigan\\
Ann Arbor, MI 48104, USA \\
\texttt{\{tewaria\}@umich.edu} \\
}
% \author{Antiquus S.~Hippocampus, Natalia Cerebro \& Amelie P. Amygdale \thanks{ Use footnote for providing further information
% about author (webpage, alternative address)---\emph{not} for acknowledging
% funding agencies.  Funding acknowledgements go at the end of the paper.} \\
% Department of Computer Science\\
% Cranberry-Lemon University\\
% Pittsburgh, PA 15213, USA \\
% \texttt{\{hippo,brain,jen\}@cs.cranberry-lemon.edu} \\
% \And
% Ji Q. Ren \& Yevgeny LeNet \\
% Department of Computational Neuroscience \\
% University of the Witwatersrand \\
% Joburg, South Africa \\
% \texttt{\{robot,net\}@wits.ac.za} \\
% \AND
% Coauthor \\
% Affiliation \\
% Address \\
% \texttt{email}
% }

% The \author macro works with any number of authors. There are two commands
% used to separate the names and addresses of multiple authors: \And and \AND.
%
% Using \And between authors leaves it to \LaTeX{} to determine where to break
% the lines. Using \AND forces a linebreak at that point. So, if \LaTeX{}
% puts 3 of 4 authors names on the first line, and the last on the second
% line, try using \AND instead of \And before the third author name.

\newcommand{\fix}{\marginpar{FIX}}
\newcommand{\new}{\marginpar{NEW}}

% \iclrfinalcopy % Uncomment for camera-ready version, but NOT for submission.
\begin{document}


\maketitle

\begin{abstract}

When solving PDEs, classical numerical solvers are often computationally expensive, while machine learning methods can suffer from spectral bias, failing to capture high-frequency components of the solution. Designing an optimal hybrid iterative solver--where, at each iteration, a solver is selected from an ensemble of solvers to leverage their complementary strengths --poses a challenging combinatorial problem. While the greedy selection strategy is desirable for its constant-factor approximation guarantee to the optimal solution, it requires knowledge of the true error at each step, which is generally unavailable in practice. We address this by proposing an approximate greedy router that efficiently mimics a greedy approach to solver selection. Empirical results on the Poisson and Helmholtz equations, demonstrate that our method outperforms existing hybrid solver approaches, such as HINTS, achieving more uniform convergence across a range of error frequencies.

\end{abstract}

\section{Introduction}
% \begin{verbatim}
%     Use command \yash{}, \sahana{}, \ambuj{} for comments
% \end{verbatim}

Natural phenomena and engineered systems are often governed by ordinary and partial differential equations (PDEs). Solving these equations enables a variety of tasks - predicting the evolution of the system through simulation (e.x., forecasting weather using Navier Stokes Equation) \citep{kalnay2003atmospheric,bauer2015quiet,staniforth1991semi}, addressing control and optimization problems (e.x., optimizing heat shield design for spacecrafts using heat transfer equations) \citep{anderson1989hypersonic,troltzsch2010optimal}, and tackling inverse problems based on external measurements (e.x. reconstructing brain activity from EEG data using electrophysiological models) \citep{baillet2001electromagnetic,grech2008review}. 

Traditionally, PDEs are solved using finite difference methods \citep{smith1985numerical,leveque2007finite}, which discretize the spatial and/or temporal domain and approximate the solution by solving a system of equations at the discrete grid points. Similarly, finite element methods \citep{hughes2003finite,bathe2006finite,brenner2008mathematical} construct a system of equations based on fitted curves for each finite element and then rely on numerical linear algebra techniques--such as the Jacobi and Gauss-Seidel method \citep{saad2003iterative}--to compute the solution. Such iterative methods, however, suffer from high computational costs, especially when finely resolved grids are required--regardless of the sparsity of the coefficient matrix. Moreover, these solvers lack generalization across different initial conditions, boundary conditions, or forcing functions, as even minor changes to any of these parameters require re-solving the entire system of equations from scratch. 

These challenges have motivated the development of neural operators \citep{kovachki2023neural} - a class of machine learning models that aim to learn the solution operator directly, enabling fast inference across a range of parameters. Neural operators lift the classical linear layer in neural networks to a linear operator--typically a kernel integral operator--between infinite-dimensional function spaces. While the universal approximation theorem of operators \citep{chen1995universal} suggests that neural operators can approximate solution operators with high accuracy, they are prone to spectral bias \citep{liu2021multiscale,liu2024mitigating,you2024mscalefno,khodakarami2025mitigating,xu2025understanding}--similar to traditional neural networks \citep{rahaman2019spectral,xu2019frequency,luo2019theory}--and favor learning the low-frequency components of the solution operator, often struggling to capture high-frequency features that iterative solvers excel at. 

Recognizing this complementarity, HINTS \citep{zhang2024blending} proposes a hybrid solver that leverages the advantages of traditional iterative methods and neural operators  to learn a wide range of eigenmodes more uniformly and quickly. It achieves this by periodically incorporating predictions from a neural operator into iterative updates of the traditional solver. While HINTS observes significant empirical performance improvements in convergence and accuracy when compared to using the standalone neural operator and Jacobi solver, the algorithm's fixed number of solver iterations between neural operator passes can be viewed as wasteful if the high frequency errors are overly dampened, or insufficient if not dampened enough. 

% Furthermore, when HINTS employs multigrid methods, they are often restricted to cycle patterns, i.e. V cycle or W cycle, of fixed depths making it difficult to adapt cycle patterns in response to the distribution of errors across the learned eigenmodes.  
% We propose a model-based router that
% We instead introduce an alternative approach, where the solver to call is chosen dynamically per step, allowing for variable length interleaving between the traditional solver and neural operator. This framework has the additional benefit of seamlessly expanding to multiple solvers, instead of just being restricted to two solvers as HINTS was.

While dynamic schedules of solvers are desirable, identifying the optimal sequence of solvers to minimize final error presents a challenging combinatorial problem, as the search space grows exponentially with the number of available solvers and iterations. Consequently, a practical solution is to employ a greedy algorithm, selecting at each iteration the solver that minimizes prediction error at the current step. To support this greedy approach, we make the following contributions:
\begin{itemize}
    \item Assuming that the true error is known at every iteration, we show that a greedy approach to selecting a sequence of solvers, aimed at minimizing prediction error, achieves a provable, constant-factor approximation to the optimal strategy.
    \item  Given that true error values are typically unavailable at inference time, a convex surrogate loss is introduced using principles of learning to defer, which, when minimized, enables the learned model to imitate the greedy solver without access to true error information. This alignment is guaranteed through the property of Bayes consistency.
    \item We empirically demonstrate that our approximate greedy solver outperforms HINTS and individual solvers in terms of uniform convergence across frequencies, on benchmark PDEs including the Poisson Equation, Helmholtz Equation, and the diffusion equation.
\end{itemize}

% We present a Greedy Hybrid Solver that, at each iteration, myopically selects a solver from a suite of available methods to reduce the error between the true solution and prediction. Assuming that the true error is known at every iteration, we show that a greedy approach to selecting a sequence of solvers, aimed at minimizing prediction error, achieves at least a $1- e^{-\alpha}$-approximation to the optimal strategy. Given that true error values are typically unavailable at inference time, a convex surrogate loss is introduced using principles of learning to defer,  which, when minimized, enables the learned model to imitate the greedy solver without access to true error information. This alignment is guaranteed through the property of Bayes consistency. We empirically demonstrate that our approximate greedy solver outperforms HINTS and individual solvers in terms of uniform convergence across frequencies, on benchmark PDEs including the Poisson Equation, Helmholtz Equation, and the diffusion equation.


\section{Problem Setting and Background} \label{sec:background}

% Let $D \subseteq \mathbb{R}^n$ denote the spatial domain, and let $\mathcal{A} = \{a: D \to \mathbb{R}^d\}$, $\mathcal{F} = \{f: D \to \mathbb{R}^o\}$, and $\mathcal{U} = \{u: D \to \mathbb{R}^m\}$ be the Banach spaces of coefficient, forcing, and solution functions, respectively.
Consider the following Linear PDE:
\begin{equation} \label{eq:mainpde}
    \begin{split}
    \mathcal{L}_{x}^a\left(u\right) &= f ,x \in D \\
    \mathcal{B}_{x}\left(u\right) &= g ,x \in \partial D
    \end{split}
\end{equation}
where $\mathcal{L}_x^a$ is a differential operator with respect to the spatial variable $x \in D$ parameterized by the coefficient function $a$,  $f$ is an external forcing function,  $u$ is the solution to the PDE, $\mathcal{B}_{x}$ is the boundary operator, and $g$ is the boundary term.  

PDEs of this form can be solved numerically using various discretization strategies, such as finite difference, finite element, and spectral methods \citep{boyd2001chebyshev,canuto2006spectral}. In this paper, however, we restrict ourselves to problems discretized by finite difference methods. Finite difference methods \citep{smith1985numerical,leveque2007finite} approximate solutions to PDEs by replacing partial derivatives with discrete differential quotients. For example, the first derivative can be approximated as $\frac{\partial u(x)}{\partial x} \approx \frac{u(x + h) - u(x)}{h}$. When the entire domain is discretized,  evaluating the solution to the PDE at these grid points reduces to solving a system of algebraic equations.

To make this concrete, let $h$ be the grid size and let $G_h(D)$ be a uniform grid with spacing $h$ over the bounded domain $D \subseteq \mathbb{R}^n$. Given a function $g: D \to \mathbb{R}$, we represent its restriction to $G_h(D)$ by $g_h$, i.e. vector of $g$ evaluated at the points of $G_h(D)$.
If $N = |G_h(D)|$, the discretized version of \Cref{eq:mainpde} is expressed as:
\begin{equation} \label{eq:mainpdediscrete}
    \mathcal{L}_h^au_h = f_h
\end{equation}
where $\mathcal{L}_h^a$ is the discretized operator represented as an $N \times N$ matrix, reducing the PDE reduced to a linear system of equations. Unlike much of the neural operator literature, which retains function space formulations to maintain discretization invariance, we follow the conventions of classical numerical analysis and represent all functions and operators as finite-dimensional vectors and matrices. Discretization invariance is not essential in our setting, since we focus on fixed-grid problems and solver collaborations.

\subsection{Iterative PDE Solvers} \label{sec:finitedifference}

While direct methods like Gauss Elimination and Thomas Algorithm are often employed to approximate the solution to systems of equations like \Cref{eq:mainpdediscrete}, they can be computationally expensive in high-dimensional domains. In contrast, iterative methods like Jacobi and Gauss-Seidel provide computational speedups by iteratively updating the solution, gradually converging to the true solution for the PDE. The general iterative update can be written as:
\begin{equation}\label{eq:iterativeupdate}
    u^{(t + 1)} = u^{(t)} + C\left(b - Au^{(t)}\right)
\end{equation}
where $u^{(t + 1)}$ is the $(t + 1)$ iterate of the solution and $C$ is a preconditioning matrix. For the Jacobi method, $C_\mathrm{Jac} = D^{-1}$, where $D$ is the diagonal of $A$; for Gauss-Seidel, $C_\mathrm{GS} = (D + L)^{-1}$, where $L$ denotes the strict lower triangular part. 

However, convergence of these basic iterative methods is often slower for the low frequency components of the solution. Multigrid solvers \citep{briggs2000multigrid,trottenberg2001multigrid,hackbusch2013multi} are an advanced class of iterative methods that tackle this problem with specialized preconditioning matrices. Specifically, multigrid methods interleave solving problems on coarse and fine grid resolutions to dampen more uniformly across frequencies.
% the solvers allow high frequency error components to be sufficiently dampened on a finer discretization and leaving the smooth parts of the error to a coarser grid where low frequency function appear high frequency. 
We defer the reader to \Cref{sec:multigrid} for more background on Multigrid solvers and the particular form of the preconditioning matrix.

% \sahana{Might move the next two paragraphs to the appendix}



\subsection{Neural Operators} \label{sec:neuraloperators}

Suppose we have some observations $\{(a_i, f_i, u_i)\}_{i = 1}^N$ such that $(a_i, f_i) \sim \mathcal{P}_{\mathcal{A} \times \mathcal{F}}$ are i.i.d. samples drawn from some distribution. Let $u_i$ is generated using a deterministic solution operator called $\mathcal{G}^{*} : \mathcal{A} \times \mathcal{F} \to \mathcal{U}$ such that $u_i = \mathcal{G}^{*}(a_i, f_i)$. Assuming a Lebesgue measure over the domain $D$, for a given pair of $(a_i, u_i, f_i)$, neural operators  $\mathcal{G}_{\theta}$ seek to approximate the true solution operator $\mathcal{G}^*$ by considering the following loss:
\begin{equation} \label{eq:NOloss}
    l_{\text{NO}}(a_i, f_i, u_i, \mathcal{G}_{\theta}) = \left\|u_i - \mathcal{G}_{\theta}(a_i, f_i)\right\|^2_{\mathcal{U}} = \int_{D} \left\|u_i(x) - \mathcal{G}_{\theta}(a_i, f_i)(x)\right\|_2^2dx
\end{equation}

The risk of this loss function is denoted by $\mathcal{R}_{\text{NO}}(\mathcal{G}_{\theta}) = \mathbb{E}_{(a, f) \sim  \mathcal{P}_{\mathcal{A} \times \mathcal{F}}}\left[l_{\text{NO}}(a, f, \mathcal{G}^{*}(a, f), \mathcal{G}_{\theta})\right] = \mathbb{E}_{(a, f) \sim  \mathcal{P}_{\mathcal{A} \times \mathcal{U}}}\left[l_{\text{NO}}(a, f, u, \mathcal{G}_{\theta})\right]$. Ideally, $\mathcal{G}_{\theta}$ would minimize $\mathcal{R}_{\text{NO}}(\mathcal{G}_{\theta})$.

In classical PDE theory, if $\mathcal{L}_{x}^a$ is a linear operator and $m = 1$, the Green's function \citep{stakgold2011greens} $G: D \times D \to \mathbb{R}$ is the solution for $\mathcal{L}_{x}^aG(x, .) = \delta_x$ and the solution function takes the following form: $u(x) =\int_D G(x, y) f(y)dy$. Inspired by this, neural operators \citep{kovachki2023neural} replace the traditional linear layer with a linear operator in the form of a kernel integral transform, $(\mathcal{K}v)(x) = \int_{D}\kappa(x, y)v(y)d\nu(y)$, mimicking the actions of Green's functions in a learnable way.  A single neural operator layer that maps a function $v_t: D \to \mathbb{R}^{d_t}$ to function $v_{t + 1}: D \to \mathbb{R}^{d_{t + 1}}$, where $t$ denotes the layer index. This mapping is defined as: $$v_{t + 1}(x) = \sigma_t\left(W_tv_t(x) + (\mathcal{K}_tv_t)(x) + b_t(x)\right)$$ for every $x \in D$ where $W_t \in \mathbb{R}^{d_{t + 1} \times d_{t}}$ are matrices performing a local linear operation (discretization invariant), $\mathcal{K}_t: \{v_t: D \to \mathbb{R}^{d_{t + 1}}\} \to \{v_{t + 1}: D \to \mathbb{R}^{d_{t + 1}}\}$ are global kernel integral operators, $b_t: D \to \mathbb{R}^{d_{t + 1}}$ are bias functions which are often constant, and $\sigma_t$ is a fixed pointwise activation function mapping from $\mathbb{R}^{d_{t + 1}}$ to $\mathbb{R}^{d_{t + 1}}$.


The architecture of the neural operator is discretization-invariant, a property that is clear in the parameterization of $W_t$ and $b_t$. However, a variety of discretization-invariant kernel parameterizations and kernel integral approximation strategies have given rise to many different variants of neural operators.  
% Different variants of neural operators arise from the choice of kernel parameterization and the approximation strategy used to compute the integral. 
For example, Fourier Neural Operators \citep{li2020fourier} parameterize the kernel in the Fourier domain and use convolutions to characterize the integral. Graph Neural Operators \citep{li2020multipole} compute the integral over a subdomain using the Nystrom's approximation. Prior to the rise of neural operators, DeepONets \citep{lu2021learning} were a popular deep learning approach and can also be viewed a neural operator that uses a pointwise kernel with a discretized integral operator. 

\subsection{Greedy Optimization} \label{sec:bggreedy}

Consider the following optimization problem:
$$\max_{S \subseteq \Omega: |S| \leq T} g(S)$$
where $g$ is a set function defined on subsets of a ground set $\Omega$. While an exhaustive search of all possible subsets of $\Omega$ can yield  optimal solution, this approach quickly becomes computationally infeasible as $|\Omega|$ grows. Instead, an efficient alternative is the greedy algorithm which builds the solution subset iteratively by, at each step, adding the element $\omega^*$ that yields the largest marginal gain, i.e. $\omega^* =  \argmax_{\omega \in \Omega \setminus S} g(S \cup \omega)$, and updating $S \leftarrow S \ \cup \{\omega^*\}$. Under certain conditions—most notably when $g$ is non-negative, monotone (the function value never decreases when more elements are added), and submodular (adding an element to a smaller set provides at least as much gain as adding it to a larger set)—the greedy algorithm is guaranteed to achieve a $(1 - e^{-1})$ approximation to the optimal solution \citep{nemhauser1978analysis}. Formally, a function $g$ is submodular if $g(A \ \cup \{\omega\}) - g(A) \geq g(B \ \cup \{\omega\}) - g(B)$ for all $A \subseteq B \subseteq \Omega$ and $\omega \in \Omega \setminus B$ and supermodular if $g(A \ \cup \{\omega\}) - g(A) \leq g(B \ \cup \{\omega\}) - g(B)$ for all $A \subseteq B \subseteq \Omega$ and $\omega \in \Omega \setminus B$. Supermodularity is an essential property for the greedy algorithm to achieve constant-factor approximation guarantees in set function minimization problems \citep{liberty2017greedy}.
% a greedy approach where at each step $S \ \cup \argmax_{\omega \in \Omega / S} g(S \cup \omega)$ is performed can achieve constant factor approximation if certain conditions are satisfied.
% Supermodularity, also known as the increasing differences property, characterizes the idea that appending an element earlier in a sequence has a larger marginal effect than doing so later. 

The suboptimality of the greedy algorithm has been extensively studied for both minimization \citep{bounia2023approximating}  and maximization \citep{das2018approximate,feige2011maximizing,bian2017guarantees,harshaw2019submodular} of set functions under settings where non-negativity, monotonicity, submodularity, or supermodularity are relaxed or only approximately satisfied. 
% However, these classical results are limited to permutation-invariant objective functions, where the function value doesn't depend on the order of the elements.
In contrast, results on the suboptimality of greedy sequence maximization \citep{streeter2008online,alaei2021maximizing,zhang2015string,bernardini2020through,van2024performance,tschiatschek2017selecting} -- where the ordering of the elements affects the function -- have led to various definitions of submodularity and monotonicity in sequential settings. In this work, we leverage some of those tools to analyze the suboptimality of the greedy solution to minimizing the final iterative error.

% \sahana{We can include a background section on greedy set optimization or consistency instead??}

\section{General framework for Hybrid Solvers}

% \sahana{Feel like I should include this in \Cref{sec:background} instead}

To solve \Cref{eq:mainpdediscrete}, consider the following general hybrid solver: at each time step $t + 1$, the update is given by
\begin{equation}
     u^{\left(t + 1\right)}_h = u^{\left(t\right)}_h + C_{S_t}\left(f_h - \mathcal{L}_h^au^{\left(t \right)}_h\right)
\end{equation}
where $\mathcal{C} =\{C_{j}\}_{j = 1}^K$ denotes the set of $K$ possible preconditioning functions and $S_t \in [K] = \{1, \dots, K\}$ indicates the chosen function at step $t$. Here, we use ``preconditioning function'' broadly to refer to any update rule, encompassing both classical approaches--such as Jacobi, Gauss-Seidel, or multigrid, where $C_j(x) = C_jx$ is a preconditioning matrix--and learned or nonlinear models, such as neural networks or operators. 
% Thus, each $C \in \mathcal{C}$ corresponds to a particular solver or update mechanism.
This update naturally generalizes the form of classical iterative methods, as seen in \Cref{eq:iterativeupdate}, by (1) allowing $C_j$ to be non-linear and (2) 
enabling the preconditioning function to be adaptively chosen at each iteration. Furthermore, HINTS \citep{zhang2024blending} can be viewed as a special case of this hybrid solver where the hybrid solver has access to two preconditioning functions ($K = 2$):  $C_1$, a neural operator, and $C_2$, a preconditioning matrix for a classical solver. The routing decision for HINTS is given by $S_t = \1_{t \bmod \tau > 0} + 1$, which selects $C_1$ every $\tau$ steps and $C_2$ otherwise.

% indexed by $j \in [K] = \{1, \dots, K\}$, and the index $j$ may change with time. 
% In the case of HINTS , t

If the error at time step $t + 1$ is denoted by $e^{(t + 1)}_h = u_h - u^{\left(t + 1\right)}_h$, the update rule can be written in terms of the error from the previous step:
\begin{align*}
    e^{(t + 1)}_h = u_h -  u^{\left(t\right)}_h -C_{S_t}\left(\mathcal{L}_h^au_h - \mathcal{L}_h^au^{\left(t\right)}_h\right) = \left(I- C_{S_t} \circ \mathcal{L}_h^a\right)(e^{(t)}_h)
    % &= e^{(t)} - C_{S_t}\left(\mathcal{L}_h^a\left(u_h-u^{\left(t\right)}_h\right)\right) \\
    %&= e^{(t)} - C_{S_t}\mathcal{L}_h^ae^{(t)}
\end{align*}
where $I : \mathbb{R}^N \to \mathbb{R}^N$ denotes an identity map, i.e., $I(x) = x$, and $I - C_{j} \circ\mathcal{L}_h^a$ is the error propagation function for the $j^{\text{th}}$ solver. Here, we use composition notation for generality: when $C_{j}$ is a matrix, this coincides with matrix multiplication; for nonlinear $C_j$, it simply means function application, i.e., $(C_j \circ \mathcal{L}_h^a)(x) = C_j\left(\mathcal{L}_h^ax\right)$.
% The error after $T$ steps becomes:
% \begin{align*}
%     e^{(T)}_h = \left(I_N - C_{S_T} \mathcal{L}_h^a\right) \cdots \left(I_N - C_{S_1} \mathcal{L}_h^a\right)  e^{(0)}_h 
% \end{align*}
% where $I_N$ is the identity matrix. 
The objective of the hybrid solver is to select a sequence of solvers at each time that minimize the error norm after $T$ steps or $\left\|e^{(T)}_h\right\|_2^2$. Formally, we seek to solve 
\begin{equation} \label{eq:errorobjective}
    \min_{S_t \in [K], |S| \leq T} h(S) := \left\|\left(I - C_{S_{|S|}} \circ \mathcal{L}_h^a\right) \circ \dots \circ \left(I - C_{S_1} \circ \mathcal{L}_h^a\right)  \left(e^{(0)}_h \right)\right\|^2_2
\end{equation}
where $S = (S_1, \dots, S_T)$ is a sequence of solver indices 
% and the objective function $h(S)$ is
% \begin{equation} \label{eq:errorobjective2}
%     h(S) = \left\|\left(I - C_{S_{|S|}} \circ \mathcal{L}_h^a\right) \circ \dots \circ \left(I - C_{S_1} \circ \mathcal{L}_h^a\right)  \left(e^{(0)}_h \right)\right\|^2_2
% \end{equation}
with compositions applied from right to left, so that the $S_1$ update acts on the initial error.
% Here, the product is ordered from right to left (i.e. from $t = 1$ to $t = |S|$), reflecting the sequential application of the preconditioners. 

% If running each solver $C_{j}$ incurs a time cost $c_j$, we can formulate a budget constrained optimization problem
% \begin{align*}
%     \min_{S \in [1, K]^*} & h(S) \\
%     \text{subject to} \quad &\sum_{t = 1}^T c_{i_t} \leq B
% \end{align*}

% where $B$ is the total computational budget.

\section{Greedy Algorithm} \label{sec:greedy}

\Cref{eq:errorobjective} defines a non-convex and combinatorial optimization problem whose search space grows exponentially in the time horizon $T$, rendering exact solutions computationally intractable for large $T$ or $K$. As a result, we begin by proposing and characterizing an ``omniscient'' greedy algorithm, in which the error of the initial iterate $e^{(0)}_h$ is known. Clearly, such a residual would be unknown in any problem of practical interest, else the solution function could be trivially obtained by simply predicting $u^{(0)}_h + e^{(0)}_h$. Yet, we demonstrate that this omniscience assumption can be relaxed to produce a practically usable learning strategy that is Bayes consistent with this omniscient approach in \Cref{sec:approxgreedy}, from which we practically recover the suboptimality guarantees shown below.

% \sahana{this Para below feels like a backgroundy paragraph}



\begin{algorithm} 
\caption{Greedy Algorithm for a Hybrid PDE solver}\label{alg:greedy}
\begin{algorithmic}
\Require $\{C_{j}\}_{j = 1}^{K}, T, \mathcal{L}_h^a, e^{(0)}_h$
\State $S^{(0)} \gets \emptyset$
\For{$t \leq T$}
    \State $S^{(t + 1)} \gets S^{(t)} \ \oplus \argmin_{j \in [K]}\left\|\left( I - C_{j} \circ \mathcal{L}_{h}^a\right)\left(e^{(t)}_h\right)\right\|^2_2$
    \State $e^{(t + 1)}_h \gets \left(I - C_{S_{t+ 1}} \circ \mathcal{L}_{h}^a\right) \left(e^{(t)}_h\right)$
\EndFor
\State \Return $S^{(T)}$
\end{algorithmic}

\end{algorithm}

% \subsection{Suboptimality of Greedy Sequence Minimization}

As discussed in \Cref{sec:bggreedy}, if the objective function satisfies supermodularity and monotonicity, a greedy algorithm -- such as the one described in \Cref{alg:greedy} -- can achieve solutions provably close to the optimal strategy. As mentioned there, classical supermodularity analyses focused on set functions, with recent extensions proposed for sequence analysis. We, therefore, continue along this line of extension and introduce the novel concept of weak-$\alpha$-supermodularity below.

Let $\Omega^*$ represent the space of sequences with each element in $\Omega$. Let $S = (S_1, \dots, S_n) \in \Omega^*$ and $S' = (S_1', \dots, S_m') \in \Omega^*$ be two sequences. We denote their concatenation as $S \oplus S' = (S_1, \dots, S_n, S_1', \dots, S_m')$. If $S' = S \ \oplus L$ where $L \in \Omega^*$, then $S$ is a prefix of $S'$ which is denoted as $S \preceq S'$. A sequence function $g: \Omega^* \to \mathbb{R}$ is considered prefix monotonically non-increasing if $g(S \ \oplus S') \leq g(S)$ for all $S, S' \in \Omega^*$ and postfix monotonically non-increasing if $g(S' \ \oplus S) \leq g(S)$ for all $S, S' \in \Omega^*$. 
% \sahana{I may not need this. will discuss later}
% $g$ is considered weakly-$\mu$-postfix monotonic if for all $S, S' \in \Omega^*$, 
% \begin{equation*}
%     g(S' \ \oplus S) \leq \mu g(S)
% \end{equation*}
% Intuitively, $\mu$ is a ``slack'' for the factor by which $g$ violates the standard postfix monotonicity; hence, $\mu\ge 1$, with $\mu=1$ recovering standard postfix monotonicity. 
In addition, a prefix non-increasing function $g$ is considered sequence supermodular if, for all $S', S \in \Omega^*$ such that $S \preceq S'$, it holds that
\begin{equation}\label{eqn:supermod}
    g(S) - g(S \ \oplus \omega ) \geq g(S') - g(S' \ \oplus \omega )
\end{equation}
% for any $\omega \in \Omega$. However, the objective function $h(S)$ described in \Cref{eq:errorobjective} may not, in general, satisfy this property. Therefore, we introduce weakly-$\alpha$-supermodularity for sequences. A prefix non-increasing function $g$ is weakly-$\alpha$-supermodular for sequence length $T$ if, for any $S, S' \in \Omega^T$ such that $|S'| = T$, 
% \begin{equation}\label{eqn:alpha_supermod}
%      g(S) - g(S\ \oplus S') \leq \alpha T \max_{i \in [|S'|]} g(S) - g(S \ \oplus S'_i)
% \end{equation}
for any $\omega \in \Omega$. However, the objective function $h(S)$ described in \Cref{eq:errorobjective} may not, in general, satisfy this property. Therefore, we introduce weak supermodularity for sequences. A prefix non-increasing function $g$ is weakly supermodular with respect to a sequence $S'$ of length $T$ if, for any $S \in \Omega^T$, there exists $\alpha(S') \geq 1$ such that
\begin{equation}\label{eqn:alpha_supermod}
     g(S) - g(S\ \oplus S') \leq \alpha(S') \sum_{i \in [|S'|]} g(S) - g(S \ \oplus S'_i)
\end{equation}
The parameter $\alpha(S')$, also called supermodularity ratio, measures how much $g$ deviates from the standard notion of sequence supermodularity.
% , similar to $\mu$ with postfix monotonicity. 
To see this, we can expand the left-hand side as a telescoping sum $$g(S) - g(S\ \oplus S')\sum_{i = 1}^{|S'|} g(S \ \oplus \left(S'_1, \dots, S'_{i - 1}\right)) - g(S \ \oplus \left(S'_1, \dots, S'_{i}\right))$$ Each term represents the marginal decreased from appending $S_i'$ after the first $i-1$ elements of $S'$. Weak supermodularity ensures that each of these marginal terms is controlled by the effect of appending $S_i'$ directly to $S$ ($g(S) - g(S \ \oplus S'_i)$). This encodes the idea that postponing the inclusion of $S_i'$ cannot yield a significantly larger benefit compared to adding it earlier. The supermodularity ratio $\alpha(S)$ quantifies the extent to which ``future gains'' from delays may exceed immediate gains. 

% See \Cref{sec:proofalpha=1} for a formal proof of this equivalence between \Cref{eqn:alpha_supermod} and \Cref{eqn:supermod} for the $\alpha=1$ case
% If we ignore the summation, by satisifying weak supermodularity, we can expect $g(S \ \oplus \left(S'_1, \dots, S'_{i - 1}\right)) - g(S \ \oplus \left(S'_1, \dots, S'_{i}\right)) \leq \alpha(S')\left(g(S) - g(S \ \oplus S'_i)\right)$ which basically says that there is no benefit of appending $S_i'$ in the future as opposed to a earlier timestep

% To see this, notice that for $\alpha=1$, \Cref{eqn:alpha_supermod} captures the notion that the total effect of concatenating a $T$-length sequence $S'$ cannot improve upon $T$ repeats of the single best step. In other words, there is no way for consecutive steps to ``complement'' one another to accrue a better improvement than the two steps individually combined. See \Cref{sec:proofalpha=1} for a formal proof of this equivalence between \Cref{eqn:alpha_supermod} and \Cref{eqn:supermod} for the $\alpha=1$ case. Thus, for $\alpha > 1$, such complementarity of steps is permitted but only within a bounded factor. Intuitively, therefore, the higher $\alpha$, the more it is possible there will be some long-range complementarity of the selection of actions that render greedy selection suboptimal.

% where $\alpha \geq 1$. When $\alpha = 1$, the functions is sequence supermodular (See \Cref{sec:proofalpha=1} for proof)

Having introduced these notions, we now characterize the suboptimality of greedy solutions of weakly supermodular and 
% weakly-$\mu$-
postfix monotonic sequence functions in \Cref{th:suboptgreedy} and subsequently discuss the interpretation of such results.

% \Cref{th:suboptgreedy} characterizes the suboptimality of greedy solutions to weakly-$\alpha$-supermodular, prefix monotonic, and weakly-$\mu$-postfix monotonic sequence functions. 


\begin{theorem} \label{th:suboptgreedy}
    Let $g: \Omega^* \to \mathbb{R}$ be a weakly supermodular function with respect to the optimal solution $O = \argmin_{S \in \Omega^T} h(S)$ with a supermodularity ratio of $\alpha(O)$ and 
    % weak-$\mu$-
    postfix monotonicity. Let the greedy solution of length $T$ be $S^T$. Then,
    \begin{equation*}
        g(S^T) \leq  \left(1 - \left(1 - \frac{1}{\alpha(O) T}\right)^T\right)
        g(O) +  \left(1 - \frac{1}{\alpha(O) T}\right)^Tf(\emptyset)
    \end{equation*}
    % Equivalently, %\sahana{stuff below is wrong}
    % \begin{equation*}
    %     \frac{g(\emptyset) - g(S^T)}{g(\emptyset) - \mu g(O)} \geq \left(1 - \left(1 - \frac{1}{\alpha T}\right)^T\right) \geq \left(1 - \frac{1}{e^{\alpha}}\right)
    % \end{equation*}
\end{theorem}

The proof of the above theorem can be found in \Cref{sec:proofsuboptgreedy}. As $T \to \infty$, the factor $\left(1 - \left(1 - \frac{1}{\alpha(O) T}\right)^T\right)$ in \Cref{th:suboptgreedy} decreases to $1 - e^{-\alpha(O)}$, showing that the worst-case performance of the greedy algorithm saturates and does not degrade further with longer sequences. Additionally, as discussed previously, as $\alpha(O) \to\infty$, the bound becomes looser: larger values of $\alpha(O)$ reflect stronger long-range dependencies, where delaying beneficial inclusions is rewarded more.

% the class of weakly-$\alpha$-supermodular functions includes more and more functions where the greedy strategy can be suboptimal due to potential long-range complementarity, hence the loosening of the bound in this setting.

The applicability of \Cref{th:suboptgreedy} to designing a greedy hybrid solver depends on verifying that the sequence function $f$ satisfies both weak supermodularity with respect to the optimal solution and postfix monotonicity-- properties established in Proposition \ref{th:weaklyalphasupermodular}. We defer the full proof of Proposition \ref{th:weaklyalphasupermodular} to \Cref{sec:proofweaklyalphasupermodular}.

\begin{proposition} \label{th:weaklyalphasupermodular}
    Suppose that for all $j \in [K]$, the error propagation function $I - C_j \circ \mathcal{L}_h^a$ is $\rho_j$-Lipschitz continuous with $\rho_j < 1$, and that $(I - C_j \circ \mathcal{L}_h^a) (0_N) = 0_N$. Then, the function $h$ is weakly supermodular with respect to the optimal solution $O$, with $$\alpha(O) = \max\left\{\frac{4}{T - \sum_{i = 1}^{T} \rho_{O_i}^2} , 1\right\}$$Furthermore, if $I - C_j \circ \mathcal{L}_h^a$ is invertible for all $j \in [K]$, $h$ is also postfix monotonically non-increasing.
\end{proposition}
The conditions of Proposition \ref{th:weaklyalphasupermodular} are natural in many applications. For classical iterative solvers, the lipschitz constant is typically $\|I_N - C_j \mathcal{L}^a_h\|$, which is less than $1$ whenever the update dampens errors--a property that holds for well-posed linear elliptic PDEs. Similarly, a neural network can be trained to be Lipschitz continuous in its inputs by applying appropriate regularization techniques to its weights \citep{gouk2021regularisation}. A small Lipschitz constant implies that the preconditioning function closely approximates the inverse of $\mathcal{L}^a_h$, so a sufficiently well-trained model should satisfy the Lipschitz constant inequality. The condition that each preconditioning function maps $0_N$ to $0_N$ is both natural and desirable: it ensures that no unnecessary updates are made when the residual $f_h - \mathcal{L}^a_h u_h^{(t)}$ is $0$. Classical solvers satisfy this by design, and any learned neural network can enforce this condition by excluding bias terms. Invertibility of the error propagation functions with classical solvers since the preconditioning 


The form of $\alpha(O)$ in Proposition \ref{th:weaklyalphasupermodular} highlights that the suboptimality gap described in \Cref{th:suboptgreedy} is governed by the collective contraction factors of the solvers chosen by the optimal solution: as $\sum_{i = 1}^{T} \rho_{O_i}^2$ approaches $T$, $\alpha(O)$ increases. This reflects the fact that when all solvers in the optimal sequence are close to the boundary of stability ($\rho_{O_i} \approx 1$), weak supermodularity degrades and greedy performance suffers.

% in classical solvers, when the residual is $0$, the additive term in \Cref{eq:iterativeupdate} vanishes, ensuring no further correction when convergence to the true solution is achieved. For neural networks, this property holds as long as there are no bias terms included.
% Proposition \ref{th:weaklyalphasupermodular} thus ensures that all the conditions required by \Cref{th:suboptgreedy} are satisfied as long as no solver in $\mathcal{C}$ amplifies errors.
% Notably, any well-trained neural operator is expected to reduce the error norm, and traditional solvers such as Jacobi and Gauss-Seidel are known to exhibit this error-reducing property for linear elliptic PDEs under standard conditions. 
% \subsection{Suboptimality of Greedy Sequence Minimization}

% \sahana{do we need this stuff below?}
\Cref{th:suboptgreedy} thus indicates that the approximation guarantee is the strongest when the sequence exhibits near-supermodular behavior, i.e. $\alpha \approx 1$, achieving the best lower bound. Such favorable conditions often arise, for instance, when the PDE of interest is linear with constant coefficients and periodic boundary conditions. Under these conditions, the error propagation matrices of classical solvers such as Jacobi and Gauss-Seidel exhibit a circulant structure making them both diagonalizable by the discrete Fourier transform matrix. Furthermore, when the learned neural operator is a single-layer Fourier Neural Operator without non-linearities, it too is diagonalizable by the discrete Fourier basis. More generally, when the the error propagation matrices $\left(I - C_j \mathcal{L}_{h}^a\right)$ are simultaneously diagonalizable, i.e each admits a spectral decomposition with a shared eigenbasis, the original problem reduces to an allocation problem. This reduction is critical in analyzing the supermodularity of the objective with simultaneously diagonalizable error propagation matrices.
\begin{proposition} \label{th:simultaneoussupermodular}
    Let $\|I_N - C_j\mathcal{L}_h^a\| \leq 1$ for all $j \in [K]$ and $\left(I - C_j \mathcal{L}_{h}^a\right) = P\Lambda_jP^{-1}$. Then, $h$ is supermodular. 
\end{proposition}
The proof of \Cref{th:simultaneoussupermodular} can be found in \Cref{sec:proofsimultaneoussupermodular}.



% \begin{corollary}[Suboptimality of HINTS]
    
% \end{corollary}


\section{Approximate Greedy Router} \label{sec:approxgreedy}

As discussed, the results from \Cref{sec:greedy} indicate the error reduction achieved with the greedy solution closely matches that of the optimal sequence, but such results are predicated on having an accurate estimate of the initial error, $e^{(0)}_h$. A poor approximation of $e^{(0)}_h$ can result in miscalibrated decisions, with the approximation errors propagating and amplifying over subsequent steps. To remedy this, we wish to design a router that learns to select solvers myopically, as described in \Cref{alg:greedy}, without access to the true initial error $e^{(0)}_h$. 

Let $\mathcal{A}, \mathcal{F}, \mathcal{U} \subseteq \mathbb{R}^N$ denote the spaces of coefficient, forcing, and solution functions, which are restricted to the grid $G_h$. Let $\mathcal{P}_{\mathcal{A} \times \mathcal{F}}$ denote the data-generating distribution over $\mathcal{A} \times \mathcal{F}$. At each time step $t$, the router $r$ selects a solver based on the coefficient function $a \in \mathcal{A}$, forcing function $f \in \mathcal{F}$, and the current iterate of the solution function $u^{(t)} \in \mathcal{U}$. 
% $u^{(t)}$ is a function of the initial guess of the solution function $u^{(0)} \sim \mathcal{P}_{\mathcal{U}}$ and 
When $r(a_h, f_h, u^{(t)}_h) = j$, the solution function is updated by the $j^{\text{th}}$ solver, resulting in an error of $\left\|\left( I - C_{j} \circ  \mathcal{L}_{h}^a\right)\left(e^{(t)}_h\right)\right\|^2_2$ at that step. Learning such a router requires minimizing the following loss objective
\begin{equation} \label{eq:routerloss}
    l_{\text{route}}\left(r, a_h, f_h, u^{(t)}_h, u_h\right) = \sum_{j = 1}^K \left\|\left( I - C_{j} \circ \mathcal{L}_{h}^a\right)\left(u_h - u^{(t)}_h\right)\right\|^2_2\1_{r(a_h, f_h, u^{(t)}_h) = j}
\end{equation}
The corresponding $l_{\text{route}}$-risk is defined as $\mathcal{R}_{\text{route}}\left(r\right) = \mathbb{E}_{a, f \sim \mathcal{P}_{\mathcal{A} \times \mathcal{F}}}\left[l_{\text{route}}\left(r, a_h, f_h, u^{(t)}_h, u_h\right)\right]$.  Although in practice the iterate sequence $\{u^{(t)}_h\}$ depends on $r$, for theoretical analysis, we fix the generation of iterates by assuming a teacher-forced greedy rollout \citep{williams1989learning,lamb2016professor}. Formally, with $u_h^{(0)} = 0_N$, for $t \geq 1$, 
% assume that the input $u^{(t)}_h$ is the true greedy decision at the $t^{th}$ step. Formally,
$$j^*_t = \text{argmin}_{j \in [K]} \left\|\left( I - C_{j} \circ \mathcal{L}_{h}^a\right)\left(e^{(t - 1)}_h\right)\right\|^2_2, \quad u^{(t)}_h = u^{(t - 1)}_h + C_{j^{*}_t}\left(f_h - \mathcal{L}_h^a u_h^{(t - 1)}\right)$$
Hence, the inputs $\{u^{(t)}_h\}$ are a deterministic function of $(a_h, f_h)$. At each $t$, $r(a_h, f_h, u^{(t)}_h)$ is computed on the teacher-forced iterate to evaluate $l_{\text{route}}$, while the next iterate $u^{(t + 1)}_h$ is advanced using the greedy choice $j^*_{t + 1}$. This divorces the iterate generation from the router's previous selections and yields a risk that depends only on $(a_h, f_h)$. However, full teacher forcing induces a distributional mismatch at test time (exposure bias). To mitigate this, our experiments employ a scheduled sampling curriculum \citep{bengio2015scheduled} and the details of this training strategy can be found in \Cref{sec:training details}.

% The router then makes a decision using the greedy decision from the previous step and $a_h, f_h$, the loss is evaluated and the next time step is again teacher forcer
% We emphasize that a distribution over all solution functions is not needed here, since each solution is fully determined by $(a_h, f_h)$. By construction, the entire sequence $\{u^{(t)}_h\}$ is fully determined by $a_h, f_h$, and the initial guess $u^{(0)}_h$. In this work, we assume that $u^{(0)}_h$ is generated by a neural operator, making it a deterministic function of $a_h$ and $f_h$. 

Minimizing $\mathcal{R}_{\text{route}}\left(r\right)$ yields the following Bayes-Optimal Router:
\begin{equation}\label{eq:bayesoptimalrouter}
    r^*(a_h, f_h, u^{(t)}_h) = \text{argmin}_{j \in [K]} \left\|\left( I - C_{j} \circ \mathcal{L}_{h}^a\right)\left(e^{(t)}_h\right)\right\|^2_2
\end{equation} 
% \sahana{is the above equation clear or does this need to be proved somehow}
which aligns with the decision rule in \Cref{alg:greedy}, confirming that $l_{\text{route}}$ is consistent with learning the greedy algorithm.

\subsection{Surrogate Loss}

Since \Cref{eq:routerloss} is discontinuous and non-convex, direct minimization is intractable in practice. To overcome this, we introduce a surrogate loss that satisfies three key properties: (1) convexity, enabling efficient optimization; (2) serving as an upper bound on the original loss; and (3) Bayes consistency, ensuring the Bayes optimal decision of the original loss is preserved upon minimization. Formally, a surrogate $\phi$ is considered to be Bayes consistent with respect to the loss $l$ if 
\begin{equation*}
    \lim_{n \to \infty} \mathcal{R}_{\phi}(f_n) - \mathcal{R}_{\phi}^* \implies \lim_{n \to \infty} \mathcal{R}_{l}(f_n) - \mathcal{R}_{l}^* 
\end{equation*}
where $\mathcal{R}_{l}(f) = \mathbb{E}\left[l(f(X), Y)\right]$ and $\mathcal{R}_{l}^* = \inf_f\mathcal{R}_{l}(f)$. In other words, in the limit of infinite data, if the risk of a sequence of learned hypotheses $\{f_n\}$ converges to the optimal risk under $\phi$, it also converges to the optimal risk with respect to the original loss $l$.

To define a surrogate loss for the routing problem, consider a set of scoring functions $\mathbf{g} =\{g_j\}_{j = 1}^K$ with $g_j: \mathcal{A} \times \mathcal{F} \times \mathcal{U} \to \mathbb{R}$ and we define the router as $r(a, f, u^{(t)}) = \text{argmax}_{j \in [K]} g_j(a, f, u^{(t)})$. Then, minimizing the following surrogate loss yields the same decision as minimizing \Cref{eq:routerloss}:
\begin{equation} \label{eq:routersurrogate}
    % \Psi\left(\mathbf{g}, a_h, f_h, u^{(t)}_h, u_h\right) = -\sum_{j = 1}^K \left(\Tilde{e} -\left\|\left( I - C_{j} \circ \mathcal{L}_{h}^a\right)\left(e^{(t)}_h\right)\right\|^2_2\right) \log\left(\frac{\exp\left(g_j(a, f, u^{(t)})\right)}{\sum_{k = 1}^K\exp\left(g_k(a, f, u^{(t)})\right)}\right)% \psi\left(\mathbf{g}(a, f, u^{(t)}),j\right)
    \Psi\left(\mathbf{g}, a_h, f_h, u^{(t)}_h, u_h\right) = -\sum_{j = 1}^K \sum_{k = 1}^K\left\|\left( I - C_{k} \circ \mathcal{L}_{h}^a\right)\left(u_h - u^{(t)}_h\right)\right\|^2_2\1_{k\neq j}\log\left(\frac{\exp\left(g_j(a, f, u^{(t)})\right)}{\sum_{m = 1}^K\exp\left(g_m(a, f, u^{(t)})\right)}\right)
\end{equation}

% where 
% $\psi\left(\mathbf{g}(a, f, u^{(t)}),j\right)$ is a consistent surrogate for the multiclass $0/1$ loss and 
% $\Tilde{e} = \max_{j \in [1, K]} \left\|\left( I_N - C_{j} \mathcal{L}_{h}\right)e^{(t)}\right\|^2_2$. 
The $\Psi$-risk is denoted by $\mathcal{R}_{\Psi}\left(\mathbf{g}\right) = \mathbb{E}_{a, f \sim \mathcal{P}_{\mathcal{A} \times \mathcal{F}}}\left[\Psi\left(\mathbf{g}, a, f, u^{(t)}, u\right)\right]$. The convexity of $\Psi$ with respect to $\mathbf{g}$ follows from the convexity of log-softmax function in its inputs. Moreover, $\Psi$ upper bounds $l_{\text{route}}$ up to a constant factor, and we refer the reader to \Cref{sec:proofupperbound} for the proof. Finally, \Cref{th:consistency} shows that $\Psi$ achieves Bayes consistency with respect to $l$; the proof can be found in \Cref{sec:proofconsistency}.

\begin{theorem} \label{th:consistency}
    Let $\left\|\left( I - C_{j} \circ \mathcal{L}_{h}^a\right)\left(e^{(t)}_h\right)\right\|^2_2 < \Bar{E} < \infty$  for all $j \in [K]$. If there exists $j\in [K]$ such that $\left\|\left( I - C_{j} \circ \mathcal{L}_{h}^a\right)\left(e^{(t)}_h\right)\right\|^2_2 > E_{min} > 0$, then, for any collection of solvers $\{C_j\}_{j= 1}^K$ and linear discrete operator $\mathcal{L}_h^a$, $\Psi$ is Bayes consistent surrogate for $l_{\text{route}}$.
\end{theorem}
In \Cref{th:consistency}, the assumption of a uniform upper bound on the errors is satisfied as long as all preconditioning functions are guaranteed to dampen errors. It is also realistic to assume that there is at least one solver that cannot drive the error to $0$, as many solver ensembles may include methods that can only provide approximate solutions. Thus, \Cref{th:consistency} ensures that, by minimizing the surrogate loss $\Psi$, one can recover the desired Bayes optimal decision rule described in \Cref{eq:bayesoptimalrouter}, which corresponds to the greedy solution of \Cref{eq:errorobjective}.





% \sahana{maybe a generalization bound here}

\section{Related Works}

\paragraph{Hybrid PDE Solvers:}
As highlighted in \Cref{sec:neuraloperators}, surrogate models are becoming of interest in settings where the solution of PDEs across initial conditions or parametric specifications is of interest. Yet, direct prediction of the solution field by such trained surrogates is often insufficient, due to the lack of guarantees around such predictions. Early works in the space of data-driven solvers, therefore, sought convergence guarantees by adapting the known guarantees of iterative numerical schemes and predicting parameters of such schemes, such as preconditioning matrices $C$ or multi-grid smoothing and coarsening matrices \citep{taghibakhshi2021optimization,caldana2024deep,kopanivcakova2025deeponet,huang2022learning,katrutsa2020black}. 

These initial works, however, failed to exploit the full expressivity of neural surrogates, thus offering only modest speedups over classic solvers. For this reason, deep learning hybrid iterative methods are becoming of greater interest, in which solutions are iteratively refined in the manner of \Cref{eq:iterativeupdate}, with every $M$-th step invoking a trained surrogate $\mathcal{G}\left(b - Au^{(t)}\right)$ to replace $C\left(b - Au^{(t)}\right)$.

% such trained surrogates is often undesirable, due to their lack of convergence guarantees and poor performance on high-frequency modes, unlike classical numerical solvers. For this reason, recent interest has emerged in the coupling of data-driven and numerical solvers to exploit the complementary strengths of these approaches. 

In particular, the HINTS approach of \citep{zhang2024blending} served as the foundational work in this space and similarly focused on the problem formulation of \Cref{eq:mainpde}. In particular, the HINTS method couples a single, fixed numerical method, such as Jacobi or Gauss-Seidel, with a pre-trained DeepONet $\mathcal{G}$ to perform iterative updating on the discretized solution $u^{\left(t\right)}_h$. 
% Iterations using the numerical solver perform the appropriately parameterized update given by \Cref{eq:iterativeupdate}; those using the DeepONet would predict the next iterate using the currently measured residual, formally as $u^{\left(t+1\right)}_h := \mathcal{G}u^{\left(t\right)}_h$. 
Notably, the schedule responsible for interleaving the numerical and neural solvers was fixed for HINTS; for instance, the Jacobi-DeepONet hybrid solver tested in their experiments repeatedly invoked 24 update steps of the Jacobi solver followed by a call to the DeepONet. Even in this preliminary work, the interleaving was demonstrated to significantly outperform numerical solvers in terms of speed of convergence. 

Given the initial promise of such results, many works have emerged to expand upon the HINTS methodology. For instance, \citep{kahana2023geometry} extended HINTS to enable its use on geometries distinct from, but related to, those on which the neural operator was originally trained. In parallel, works have emerged to complement the empirical findings presented in the HINTS paper with stronger theoretical grounding, often considering slight variants of the HINTS methodology for ease of analysis. In particular, both \citep{hu2025hybrid} and \citep{cui2022fourier} characterized the dampening of error modes over iterates, with the former replacing the DeepONet of HINTS with an MIONet \citep{jin2022mionet} and the latter a Fourier Neural Operator \citep{li2020fourier}.

Critically, all hybrid solvers proposed heretofore are limited in their expressivity in two ways: they, firstly, only consider a \textit{single} numerical solver to complement the trained surrogate and secondly interleave the numerical and neural solvers with a schedule fixed a priori as guided by user intuition. These were the primary shortcomings that we addressed in our proposed method, which result in significant empirical improvements, as we demonstrate in \Cref{section:experiments}.

\paragraph{Model Routing:}
With large language models (LLMs) increasingly serving as a unified endpoint for a diverse collection of queries, significant interest has emerged around the problem of ``model routing'' \cite{shnitzer2023large,hu2024routerbench,ding2024hybrid,huang2025routereval}. Formally, the ``routing'' problem seeks to find the optimal selection $f^{*}$ amongst some collection of models $f_{1},...,f_{M} : \mathcal{X}\to\mathcal{Y}$ for an input query $x$, where each $f_i$ represents a fixed model. Notably, the framework of model routing is agnostic to the choice of what $f_i$ represents; that is, while $f_i$ could be a trained LLM, it could also represent dispatching to a human expert.

Such routing decisions often seek to balance a combination of anticipated model correctness, inference cost, and latency requirements. In the simplest cases, such decisions are based on hand-crafted heuristics, such as thresholding uncertainty estimates for the cheaper models to determine whether routing to more expensive models is necessary \cite{chuang2025confident,chuang2024learning}. The rigidity of such rules, however, can fail to capture subtleties in the optimal routing. As a result, routing is increasingly handled by training a separate router network $r$, which classifies a query $x$ as being best served by model $r(x)\in[M]$ amongst the full collection of models, such that the final prediction is $f_{r(x)}(x)$ \cite{mohammadshahi2024routoo,hari2023tryage,vsakota2024fly}.

Abstractly, our proposed method can be seen as learning a router for the collection of numerical and neural solver methods. However, there are several novel aspects of our work that distinguish it from others of this literature. In particular, our work is the first to demonstrate that such a framework is applicable to the problem of learning hybrid PDE solvers, differing from prior hybrid solver works, which simply fixed the ``router'' a priori to a predetermined schedule. This novel application further enabled us to provide the detailed error characterizations presented in \Cref{sec:greedy}, which is uniquely possible in our setting, given the formal characterizations possible for PDE numerical solvers compared to LLM responses. In addition, our proposed surrogate loss of \Cref{eq:routersurrogate} directly optimizes the iterate residuals, whereas alternative works on model routing for LLMs are often trained against a cross-entropy loss on labels identifying the optimal model.

\section{Experiments}\label{section:experiments}

In this section, we empirically demonstrate the fast, uniform convergence of the approximate greedy router on Poisson and Helmholtz equations posed on the unit domain $D = [0, 1]^d$ with $d\in \{1, 2\}$. For Poisson we use the constant-coefficient model
$$- \Delta u(x) = f(x), \quad x\in D$$ and for Helmholtz, we adopt the sign convention 
$$-\Delta u(x) - a^2(x)u(x) = f(x), \quad x\in D$$ where $a^2(x) \geq 0$. Furthermore, we impose periodic boundary conditions for both equations; for periodic Poisson the right-hand side is centered to satisfy the compatibility condition $\int_D f(x)dx = 0$.

The spatial domain $D$ is discretized into uniform grids. In one dimension, we use $N = 65$ points and in two dimensions, we use a $33 \times 33$ grid. The discrete Laplacian $\Delta_h$ are the standard three-point (1D) 
% $\begin{bmatrix}
%     1 & -2 & 1
% \end{bmatrix}$ 
stencil and five-point (2D) 
% $\begin{bmatrix}
%     0 & -1 & 0 \\
%     1 & -4 & 1 \\
%     0 & -1 & 0
% \end{bmatrix}$ 
stencils. Data is sampled from a zero-mean Gaussian Random Field on the periodic domain with covariance operator $(- \Delta + 9I)^{-2}$. We generate samples in Fourier space: for each non-zero mode $k$, we draw an independent complex coefficient from a Gaussian distribution with mean 0 and variance $(4\pi^2 \|k\|^2_2 + 9)^{-2}$, enforce a Hermitian symmetry to obtain a real-valued field, set the DC mode to $0$ to ensure zero mean for Poisson, and apply inverse Discrete Fourier Transform to obtain the field in physical space. For each sample, we compute reference solutions with a least squares solver and treat them as ground truth.  

The set of solvers or preconditioning functions compirses a weighted Jacobi relaxation, a Gauss-Seidel (GS) sweep, one two-grid  multigrid (MG) step, and a learned DeepONet residual correction. As baselines, we evaluate (i) single-solver schedules -- Jacobi only, GS only , MG only-- as well as (ii) HINTS variants instantiated with each classical solver (HINTS-Jacobi, HINTS-GS, HINTS-MG), where the DeepONet correction is interleaved every 24 iterations for Jacobi/GS and every 14 V-cycles for MG. We then train LSTM-based routers using the loss in \Cref{eq:routersurrogate} under four different solver-access sets which include Jacobi+DeepONet (Greedy-Jacobi), GS+DeepONet (Greedy-GS), GS+Jacobi+DeepONet (Greedy-GS,Jacobi), and MG+DeepONet (Greedy-MG). Training Details of the DeepONet and the routers can be found in \Cref{sec:training details}.

A key comparison principle is that methods are only comparable within the same solver-access sets. For example, Greedy-Jacobi and HINTS-GS are not directly comparable because the former has no access to GS. Accordingly, we partition results by solver family (Jacobi-related, GS-related, and MG-related) and compare methods only within each family. Each method is run for 300 iterations for Jacobi and GS, and 100 V cycles for MG. \Cref{fig:errorcomparison} shows the convergence history of a representative solution function, and \Cref{tab:errorcomparison} summarizes such convergence plots across $64$ test samples by reporting the mean final error $\left\|e^{(T)}_h\right\|^2_2$ and the mean area under the curve (AUC), where $AUC_T = \sum_{t = 1}^T\left\|e^{(t)}_h\right\|^2_2$. AUC is meaningful because it measures anytime performance: smaller AUC reflects lower error across the run rather than only at the final step. It also naturally penalizes non-monotone spikes that undo progress. 

Across all solver families, the greedy router achieves the lowest final error and lowest mean AUC: Greedy-Jacobi, Greedy-GS, and Greedy-MG outperform both single-solver baselines and their HINTS counterparts. While HINTS improves over single-solver schedules, it often exhibits sawtooth error traces as seend in \Cref{fig:errorcomparison} as the fixed schedules often invoke the DeepONet when it's suboptimal. By contrast, the greedy routers only routes to solvers that yield an immediate error drop, which directly translates into the reported AUC gains and near-monotone error decay. Additional empirical results are provided in \Cref{sec:moreexperimentalresults} where all solvers are compared under different metrics.


\begin{table}[]
    \centering
    \resizebox{\textwidth}{!}{%
    \begin{tabular}{|c|c|c|c|c|c|c|}
        \hline
         Equation & \multicolumn{2}{|c|}{1D Poisson} & \multicolumn{2}{|c|}{2D Poisson} & \multicolumn{2}{|c|}{2D Helmholtz} \\
         \hline 
         Methods & $\|e^{(T)}_h\|$ & AUC & $\|e^{(T)}_h\|$ & AUC& $\|e^{(T)}_h\|$ & AUC\\
         \hline
         \multicolumn{7}{|c|}{Jacobi-related Solvers} \\
        \hline
        Jacobi Only & 0.19 (0.1) & 128.66 (69.98) & 0.07 (0.03) & 115.5 (42.96) & NaN & NaN \\
        HINTS-Jacobi & 0.01  & 22.71 (11.36) & 0.18 & 65.75 (8.61) & NaN & NaN \\
        Greedy-Jacobi &\textbf{0.001} & \textbf{1.12 (0.44)} & \textbf{0.01}  & \textbf{18.48 (4.06)} & NaN & NaN \\
        % Jacobi Only & 128.662 (69.985) & 115.505 (42.964) & NaN \\
        % HINTS-Jacobi & 22.708 (11.36) & 65.746 (8.607) & NaN \\
        % Greedy-Jacobi & \textbf{1.122 (0.442)} & \textbf{18.482 (4.064)} & NaN \\
        \hline
         \multicolumn{7}{|c|}{GS-related Solvers} \\
         \hline
        % GS Only & 80.554 (43.835) & 61.671 (23.022) & NaN \\
        % HINTS-GS & 17.8 (9.208) & 56.627 (7.75) & NaN \\
        % Greedy-GS & \textbf{0.674 (0.274)} & \textbf{9.898 (2.263)} & NaN \\
        % Greedy-GS,Jacobi & 1.07 (0.431) & 62.575 (23.426) & NaN \\
        GS Only & 0.05 (0.03) & 80.55 (43.83) & 0.005  & 61.67 (23.02) & NaN & NaN \\
    HINTS-GS & 0.003  & 17.8 (9.208) & 0.174 & 56.63 (7.75) & NaN & NaN \\
    Greedy-GS & $\mathbf{<10^{-3}}$ & \textbf{0.674 (0.274)} & \textbf{0.001} & \textbf{9.9 (2.26)} & NaN & NaN \\
    Greedy-GS,Jacobi & $<10^{-3}$ & 1.07 (0.43) & 0.013  & 62.57 (23.43) & NaN & NaN \\
        \hline
         \multicolumn{7}{|c|}{MG-related Solvers} \\
         \hline
         MG Only & 0.05 (0.01) & 20.51 (7.46) & 0.002 & 10.97 (3.03) & NaN & NaN \\
        HINTS-MG & 0.002  & 7.64 (3.663) & 0.079 & 15.02 (2.03) & NaN & NaN \\
        Greedy-MG & $\mathbf{<10^{-3}}$ & \textbf{0.24 (0.1)} & \textbf{0.001} & \textbf{2.58 (0.54)} & NaN & NaN \\
        % MG Only & 20.513 (7.463) & 10.97 (3.035) & NaN \\
        % HINTS-MG & 7.64 (3.663) & 15.022 (2.027) & NaN \\
        % Greedy-MG & \textbf{0.24 (0.102)} & \textbf{2.582 (0.542)} & NaN \\
         \hline
    \end{tabular}
    }
    \caption{Final error and Area under the error curve (AUC) of squared $L^2$ error (lower is better). Values are mean ($\pm$ standard error (s.e.) over 64 test instances; both mean and s.e. are reported in $\times 10^{-3}$. AUC is $\sum_{t=0}^{T-1}|e_h^{(t)}|_2^2$ over a fixed budget (Jacobi/GS: $T{=}300$ iterations; MG: $T{=}100$ V-cycles). If a standard error is not shown, it is $< 10^{-3}$ in the reported units (raw $< 10^{-6}$). Bold indicates the best method within each solver family.}
    \label{tab:errorcomparison}
\end{table}




\begin{figure}
    \centering
    %%% Creator: Matplotlib, PGF backend
%%
%% To include the figure in your LaTeX document, write
%%   \input{<filename>.pgf}
%%
%% Make sure the required packages are loaded in your preamble
%%   \usepackage{pgf}
%%
%% Also ensure that all the required font packages are loaded; for instance,
%% the lmodern package is sometimes necessary when using math font.
%%   \usepackage{lmodern}
%%
%% Figures using additional raster images can only be included by \input if
%% they are in the same directory as the main LaTeX file. For loading figures
%% from other directories you can use the `import` package
%%   \usepackage{import}
%%
%% and then include the figures with
%%   \import{<path to file>}{<filename>.pgf}
%%
%% Matplotlib used the following preamble
%%   \def\mathdefault#1{#1}
%%   \everymath=\expandafter{\the\everymath\displaystyle}
%%   \IfFileExists{scrextend.sty}{
%%     \usepackage[fontsize=10.000000pt]{scrextend}
%%   }{
%%     \renewcommand{\normalsize}{\fontsize{10.000000}{12.000000}\selectfont}
%%     \normalsize
%%   }
%%   
%%   \ifdefined\pdftexversion\else  % non-pdftex case.
%%     \usepackage{fontspec}
%%     \setmainfont{DejaVuSerif.ttf}[Path=\detokenize{/home/srayan/miniconda3/envs/greedy/lib/python3.13/site-packages/matplotlib/mpl-data/fonts/ttf/}]
%%     \setsansfont{DejaVuSans.ttf}[Path=\detokenize{/home/srayan/miniconda3/envs/greedy/lib/python3.13/site-packages/matplotlib/mpl-data/fonts/ttf/}]
%%     \setmonofont{DejaVuSansMono.ttf}[Path=\detokenize{/home/srayan/miniconda3/envs/greedy/lib/python3.13/site-packages/matplotlib/mpl-data/fonts/ttf/}]
%%   \fi
%%   \makeatletter\@ifpackageloaded{underscore}{}{\usepackage[strings]{underscore}}\makeatother
%%
\begingroup%
\makeatletter%
\begin{pgfpicture}%
\pgfpathrectangle{\pgfpointorigin}{\pgfqpoint{15.000000in}{10.000000in}}%
\pgfusepath{use as bounding box, clip}%
\begin{pgfscope}%
\pgfsetbuttcap%
\pgfsetmiterjoin%
\definecolor{currentfill}{rgb}{1.000000,1.000000,1.000000}%
\pgfsetfillcolor{currentfill}%
\pgfsetlinewidth{0.000000pt}%
\definecolor{currentstroke}{rgb}{1.000000,1.000000,1.000000}%
\pgfsetstrokecolor{currentstroke}%
\pgfsetdash{}{0pt}%
\pgfpathmoveto{\pgfqpoint{0.000000in}{0.000000in}}%
\pgfpathlineto{\pgfqpoint{15.000000in}{0.000000in}}%
\pgfpathlineto{\pgfqpoint{15.000000in}{10.000000in}}%
\pgfpathlineto{\pgfqpoint{0.000000in}{10.000000in}}%
\pgfpathlineto{\pgfqpoint{0.000000in}{0.000000in}}%
\pgfpathclose%
\pgfusepath{fill}%
\end{pgfscope}%
\begin{pgfscope}%
\pgfsetbuttcap%
\pgfsetmiterjoin%
\definecolor{currentfill}{rgb}{1.000000,1.000000,1.000000}%
\pgfsetfillcolor{currentfill}%
\pgfsetlinewidth{0.000000pt}%
\definecolor{currentstroke}{rgb}{0.000000,0.000000,0.000000}%
\pgfsetstrokecolor{currentstroke}%
\pgfsetstrokeopacity{0.000000}%
\pgfsetdash{}{0pt}%
\pgfpathmoveto{\pgfqpoint{1.875000in}{5.300000in}}%
\pgfpathlineto{\pgfqpoint{5.294118in}{5.300000in}}%
\pgfpathlineto{\pgfqpoint{5.294118in}{8.800000in}}%
\pgfpathlineto{\pgfqpoint{1.875000in}{8.800000in}}%
\pgfpathlineto{\pgfqpoint{1.875000in}{5.300000in}}%
\pgfpathclose%
\pgfusepath{fill}%
\end{pgfscope}%
\begin{pgfscope}%
\pgfsetbuttcap%
\pgfsetroundjoin%
\definecolor{currentfill}{rgb}{0.000000,0.000000,0.000000}%
\pgfsetfillcolor{currentfill}%
\pgfsetlinewidth{0.803000pt}%
\definecolor{currentstroke}{rgb}{0.000000,0.000000,0.000000}%
\pgfsetstrokecolor{currentstroke}%
\pgfsetdash{}{0pt}%
\pgfsys@defobject{currentmarker}{\pgfqpoint{0.000000in}{-0.048611in}}{\pgfqpoint{0.000000in}{0.000000in}}{%
\pgfpathmoveto{\pgfqpoint{0.000000in}{0.000000in}}%
\pgfpathlineto{\pgfqpoint{0.000000in}{-0.048611in}}%
\pgfusepath{stroke,fill}%
}%
\begin{pgfscope}%
\pgfsys@transformshift{2.030414in}{5.300000in}%
\pgfsys@useobject{currentmarker}{}%
\end{pgfscope}%
\end{pgfscope}%
\begin{pgfscope}%
\definecolor{textcolor}{rgb}{0.000000,0.000000,0.000000}%
\pgfsetstrokecolor{textcolor}%
\pgfsetfillcolor{textcolor}%
\pgftext[x=2.030414in,y=5.202778in,,top]{\color{textcolor}{\sffamily\fontsize{10.000000}{12.000000}\selectfont\catcode`\^=\active\def^{\ifmmode\sp\else\^{}\fi}\catcode`\%=\active\def%{\%}0}}%
\end{pgfscope}%
\begin{pgfscope}%
\pgfsetbuttcap%
\pgfsetroundjoin%
\definecolor{currentfill}{rgb}{0.000000,0.000000,0.000000}%
\pgfsetfillcolor{currentfill}%
\pgfsetlinewidth{0.803000pt}%
\definecolor{currentstroke}{rgb}{0.000000,0.000000,0.000000}%
\pgfsetstrokecolor{currentstroke}%
\pgfsetdash{}{0pt}%
\pgfsys@defobject{currentmarker}{\pgfqpoint{0.000000in}{-0.048611in}}{\pgfqpoint{0.000000in}{0.000000in}}{%
\pgfpathmoveto{\pgfqpoint{0.000000in}{0.000000in}}%
\pgfpathlineto{\pgfqpoint{0.000000in}{-0.048611in}}%
\pgfusepath{stroke,fill}%
}%
\begin{pgfscope}%
\pgfsys@transformshift{2.550195in}{5.300000in}%
\pgfsys@useobject{currentmarker}{}%
\end{pgfscope}%
\end{pgfscope}%
\begin{pgfscope}%
\definecolor{textcolor}{rgb}{0.000000,0.000000,0.000000}%
\pgfsetstrokecolor{textcolor}%
\pgfsetfillcolor{textcolor}%
\pgftext[x=2.550195in,y=5.202778in,,top]{\color{textcolor}{\sffamily\fontsize{10.000000}{12.000000}\selectfont\catcode`\^=\active\def^{\ifmmode\sp\else\^{}\fi}\catcode`\%=\active\def%{\%}50}}%
\end{pgfscope}%
\begin{pgfscope}%
\pgfsetbuttcap%
\pgfsetroundjoin%
\definecolor{currentfill}{rgb}{0.000000,0.000000,0.000000}%
\pgfsetfillcolor{currentfill}%
\pgfsetlinewidth{0.803000pt}%
\definecolor{currentstroke}{rgb}{0.000000,0.000000,0.000000}%
\pgfsetstrokecolor{currentstroke}%
\pgfsetdash{}{0pt}%
\pgfsys@defobject{currentmarker}{\pgfqpoint{0.000000in}{-0.048611in}}{\pgfqpoint{0.000000in}{0.000000in}}{%
\pgfpathmoveto{\pgfqpoint{0.000000in}{0.000000in}}%
\pgfpathlineto{\pgfqpoint{0.000000in}{-0.048611in}}%
\pgfusepath{stroke,fill}%
}%
\begin{pgfscope}%
\pgfsys@transformshift{3.069976in}{5.300000in}%
\pgfsys@useobject{currentmarker}{}%
\end{pgfscope}%
\end{pgfscope}%
\begin{pgfscope}%
\definecolor{textcolor}{rgb}{0.000000,0.000000,0.000000}%
\pgfsetstrokecolor{textcolor}%
\pgfsetfillcolor{textcolor}%
\pgftext[x=3.069976in,y=5.202778in,,top]{\color{textcolor}{\sffamily\fontsize{10.000000}{12.000000}\selectfont\catcode`\^=\active\def^{\ifmmode\sp\else\^{}\fi}\catcode`\%=\active\def%{\%}100}}%
\end{pgfscope}%
\begin{pgfscope}%
\pgfsetbuttcap%
\pgfsetroundjoin%
\definecolor{currentfill}{rgb}{0.000000,0.000000,0.000000}%
\pgfsetfillcolor{currentfill}%
\pgfsetlinewidth{0.803000pt}%
\definecolor{currentstroke}{rgb}{0.000000,0.000000,0.000000}%
\pgfsetstrokecolor{currentstroke}%
\pgfsetdash{}{0pt}%
\pgfsys@defobject{currentmarker}{\pgfqpoint{0.000000in}{-0.048611in}}{\pgfqpoint{0.000000in}{0.000000in}}{%
\pgfpathmoveto{\pgfqpoint{0.000000in}{0.000000in}}%
\pgfpathlineto{\pgfqpoint{0.000000in}{-0.048611in}}%
\pgfusepath{stroke,fill}%
}%
\begin{pgfscope}%
\pgfsys@transformshift{3.589757in}{5.300000in}%
\pgfsys@useobject{currentmarker}{}%
\end{pgfscope}%
\end{pgfscope}%
\begin{pgfscope}%
\definecolor{textcolor}{rgb}{0.000000,0.000000,0.000000}%
\pgfsetstrokecolor{textcolor}%
\pgfsetfillcolor{textcolor}%
\pgftext[x=3.589757in,y=5.202778in,,top]{\color{textcolor}{\sffamily\fontsize{10.000000}{12.000000}\selectfont\catcode`\^=\active\def^{\ifmmode\sp\else\^{}\fi}\catcode`\%=\active\def%{\%}150}}%
\end{pgfscope}%
\begin{pgfscope}%
\pgfsetbuttcap%
\pgfsetroundjoin%
\definecolor{currentfill}{rgb}{0.000000,0.000000,0.000000}%
\pgfsetfillcolor{currentfill}%
\pgfsetlinewidth{0.803000pt}%
\definecolor{currentstroke}{rgb}{0.000000,0.000000,0.000000}%
\pgfsetstrokecolor{currentstroke}%
\pgfsetdash{}{0pt}%
\pgfsys@defobject{currentmarker}{\pgfqpoint{0.000000in}{-0.048611in}}{\pgfqpoint{0.000000in}{0.000000in}}{%
\pgfpathmoveto{\pgfqpoint{0.000000in}{0.000000in}}%
\pgfpathlineto{\pgfqpoint{0.000000in}{-0.048611in}}%
\pgfusepath{stroke,fill}%
}%
\begin{pgfscope}%
\pgfsys@transformshift{4.109537in}{5.300000in}%
\pgfsys@useobject{currentmarker}{}%
\end{pgfscope}%
\end{pgfscope}%
\begin{pgfscope}%
\definecolor{textcolor}{rgb}{0.000000,0.000000,0.000000}%
\pgfsetstrokecolor{textcolor}%
\pgfsetfillcolor{textcolor}%
\pgftext[x=4.109537in,y=5.202778in,,top]{\color{textcolor}{\sffamily\fontsize{10.000000}{12.000000}\selectfont\catcode`\^=\active\def^{\ifmmode\sp\else\^{}\fi}\catcode`\%=\active\def%{\%}200}}%
\end{pgfscope}%
\begin{pgfscope}%
\pgfsetbuttcap%
\pgfsetroundjoin%
\definecolor{currentfill}{rgb}{0.000000,0.000000,0.000000}%
\pgfsetfillcolor{currentfill}%
\pgfsetlinewidth{0.803000pt}%
\definecolor{currentstroke}{rgb}{0.000000,0.000000,0.000000}%
\pgfsetstrokecolor{currentstroke}%
\pgfsetdash{}{0pt}%
\pgfsys@defobject{currentmarker}{\pgfqpoint{0.000000in}{-0.048611in}}{\pgfqpoint{0.000000in}{0.000000in}}{%
\pgfpathmoveto{\pgfqpoint{0.000000in}{0.000000in}}%
\pgfpathlineto{\pgfqpoint{0.000000in}{-0.048611in}}%
\pgfusepath{stroke,fill}%
}%
\begin{pgfscope}%
\pgfsys@transformshift{4.629318in}{5.300000in}%
\pgfsys@useobject{currentmarker}{}%
\end{pgfscope}%
\end{pgfscope}%
\begin{pgfscope}%
\definecolor{textcolor}{rgb}{0.000000,0.000000,0.000000}%
\pgfsetstrokecolor{textcolor}%
\pgfsetfillcolor{textcolor}%
\pgftext[x=4.629318in,y=5.202778in,,top]{\color{textcolor}{\sffamily\fontsize{10.000000}{12.000000}\selectfont\catcode`\^=\active\def^{\ifmmode\sp\else\^{}\fi}\catcode`\%=\active\def%{\%}250}}%
\end{pgfscope}%
\begin{pgfscope}%
\pgfsetbuttcap%
\pgfsetroundjoin%
\definecolor{currentfill}{rgb}{0.000000,0.000000,0.000000}%
\pgfsetfillcolor{currentfill}%
\pgfsetlinewidth{0.803000pt}%
\definecolor{currentstroke}{rgb}{0.000000,0.000000,0.000000}%
\pgfsetstrokecolor{currentstroke}%
\pgfsetdash{}{0pt}%
\pgfsys@defobject{currentmarker}{\pgfqpoint{0.000000in}{-0.048611in}}{\pgfqpoint{0.000000in}{0.000000in}}{%
\pgfpathmoveto{\pgfqpoint{0.000000in}{0.000000in}}%
\pgfpathlineto{\pgfqpoint{0.000000in}{-0.048611in}}%
\pgfusepath{stroke,fill}%
}%
\begin{pgfscope}%
\pgfsys@transformshift{5.149099in}{5.300000in}%
\pgfsys@useobject{currentmarker}{}%
\end{pgfscope}%
\end{pgfscope}%
\begin{pgfscope}%
\definecolor{textcolor}{rgb}{0.000000,0.000000,0.000000}%
\pgfsetstrokecolor{textcolor}%
\pgfsetfillcolor{textcolor}%
\pgftext[x=5.149099in,y=5.202778in,,top]{\color{textcolor}{\sffamily\fontsize{10.000000}{12.000000}\selectfont\catcode`\^=\active\def^{\ifmmode\sp\else\^{}\fi}\catcode`\%=\active\def%{\%}300}}%
\end{pgfscope}%
\begin{pgfscope}%
\definecolor{textcolor}{rgb}{0.000000,0.000000,0.000000}%
\pgfsetstrokecolor{textcolor}%
\pgfsetfillcolor{textcolor}%
\pgftext[x=3.584559in,y=5.012809in,,top]{\color{textcolor}{\sffamily\fontsize{10.000000}{12.000000}\selectfont\catcode`\^=\active\def^{\ifmmode\sp\else\^{}\fi}\catcode`\%=\active\def%{\%}Iterations}}%
\end{pgfscope}%
\begin{pgfscope}%
\pgfsetbuttcap%
\pgfsetroundjoin%
\definecolor{currentfill}{rgb}{0.000000,0.000000,0.000000}%
\pgfsetfillcolor{currentfill}%
\pgfsetlinewidth{0.803000pt}%
\definecolor{currentstroke}{rgb}{0.000000,0.000000,0.000000}%
\pgfsetstrokecolor{currentstroke}%
\pgfsetdash{}{0pt}%
\pgfsys@defobject{currentmarker}{\pgfqpoint{-0.048611in}{0.000000in}}{\pgfqpoint{-0.000000in}{0.000000in}}{%
\pgfpathmoveto{\pgfqpoint{-0.000000in}{0.000000in}}%
\pgfpathlineto{\pgfqpoint{-0.048611in}{0.000000in}}%
\pgfusepath{stroke,fill}%
}%
\begin{pgfscope}%
\pgfsys@transformshift{1.875000in}{5.806344in}%
\pgfsys@useobject{currentmarker}{}%
\end{pgfscope}%
\end{pgfscope}%
\begin{pgfscope}%
\definecolor{textcolor}{rgb}{0.000000,0.000000,0.000000}%
\pgfsetstrokecolor{textcolor}%
\pgfsetfillcolor{textcolor}%
\pgftext[x=1.489775in, y=5.753582in, left, base]{\color{textcolor}{\sffamily\fontsize{10.000000}{12.000000}\selectfont\catcode`\^=\active\def^{\ifmmode\sp\else\^{}\fi}\catcode`\%=\active\def%{\%}$\mathdefault{10^{-6}}$}}%
\end{pgfscope}%
\begin{pgfscope}%
\pgfsetbuttcap%
\pgfsetroundjoin%
\definecolor{currentfill}{rgb}{0.000000,0.000000,0.000000}%
\pgfsetfillcolor{currentfill}%
\pgfsetlinewidth{0.803000pt}%
\definecolor{currentstroke}{rgb}{0.000000,0.000000,0.000000}%
\pgfsetstrokecolor{currentstroke}%
\pgfsetdash{}{0pt}%
\pgfsys@defobject{currentmarker}{\pgfqpoint{-0.048611in}{0.000000in}}{\pgfqpoint{-0.000000in}{0.000000in}}{%
\pgfpathmoveto{\pgfqpoint{-0.000000in}{0.000000in}}%
\pgfpathlineto{\pgfqpoint{-0.048611in}{0.000000in}}%
\pgfusepath{stroke,fill}%
}%
\begin{pgfscope}%
\pgfsys@transformshift{1.875000in}{6.687896in}%
\pgfsys@useobject{currentmarker}{}%
\end{pgfscope}%
\end{pgfscope}%
\begin{pgfscope}%
\definecolor{textcolor}{rgb}{0.000000,0.000000,0.000000}%
\pgfsetstrokecolor{textcolor}%
\pgfsetfillcolor{textcolor}%
\pgftext[x=1.489775in, y=6.635134in, left, base]{\color{textcolor}{\sffamily\fontsize{10.000000}{12.000000}\selectfont\catcode`\^=\active\def^{\ifmmode\sp\else\^{}\fi}\catcode`\%=\active\def%{\%}$\mathdefault{10^{-5}}$}}%
\end{pgfscope}%
\begin{pgfscope}%
\pgfsetbuttcap%
\pgfsetroundjoin%
\definecolor{currentfill}{rgb}{0.000000,0.000000,0.000000}%
\pgfsetfillcolor{currentfill}%
\pgfsetlinewidth{0.803000pt}%
\definecolor{currentstroke}{rgb}{0.000000,0.000000,0.000000}%
\pgfsetstrokecolor{currentstroke}%
\pgfsetdash{}{0pt}%
\pgfsys@defobject{currentmarker}{\pgfqpoint{-0.048611in}{0.000000in}}{\pgfqpoint{-0.000000in}{0.000000in}}{%
\pgfpathmoveto{\pgfqpoint{-0.000000in}{0.000000in}}%
\pgfpathlineto{\pgfqpoint{-0.048611in}{0.000000in}}%
\pgfusepath{stroke,fill}%
}%
\begin{pgfscope}%
\pgfsys@transformshift{1.875000in}{7.569448in}%
\pgfsys@useobject{currentmarker}{}%
\end{pgfscope}%
\end{pgfscope}%
\begin{pgfscope}%
\definecolor{textcolor}{rgb}{0.000000,0.000000,0.000000}%
\pgfsetstrokecolor{textcolor}%
\pgfsetfillcolor{textcolor}%
\pgftext[x=1.489775in, y=7.516686in, left, base]{\color{textcolor}{\sffamily\fontsize{10.000000}{12.000000}\selectfont\catcode`\^=\active\def^{\ifmmode\sp\else\^{}\fi}\catcode`\%=\active\def%{\%}$\mathdefault{10^{-4}}$}}%
\end{pgfscope}%
\begin{pgfscope}%
\pgfsetbuttcap%
\pgfsetroundjoin%
\definecolor{currentfill}{rgb}{0.000000,0.000000,0.000000}%
\pgfsetfillcolor{currentfill}%
\pgfsetlinewidth{0.803000pt}%
\definecolor{currentstroke}{rgb}{0.000000,0.000000,0.000000}%
\pgfsetstrokecolor{currentstroke}%
\pgfsetdash{}{0pt}%
\pgfsys@defobject{currentmarker}{\pgfqpoint{-0.048611in}{0.000000in}}{\pgfqpoint{-0.000000in}{0.000000in}}{%
\pgfpathmoveto{\pgfqpoint{-0.000000in}{0.000000in}}%
\pgfpathlineto{\pgfqpoint{-0.048611in}{0.000000in}}%
\pgfusepath{stroke,fill}%
}%
\begin{pgfscope}%
\pgfsys@transformshift{1.875000in}{8.451000in}%
\pgfsys@useobject{currentmarker}{}%
\end{pgfscope}%
\end{pgfscope}%
\begin{pgfscope}%
\definecolor{textcolor}{rgb}{0.000000,0.000000,0.000000}%
\pgfsetstrokecolor{textcolor}%
\pgfsetfillcolor{textcolor}%
\pgftext[x=1.489775in, y=8.398238in, left, base]{\color{textcolor}{\sffamily\fontsize{10.000000}{12.000000}\selectfont\catcode`\^=\active\def^{\ifmmode\sp\else\^{}\fi}\catcode`\%=\active\def%{\%}$\mathdefault{10^{-3}}$}}%
\end{pgfscope}%
\begin{pgfscope}%
\pgfsetbuttcap%
\pgfsetroundjoin%
\definecolor{currentfill}{rgb}{0.000000,0.000000,0.000000}%
\pgfsetfillcolor{currentfill}%
\pgfsetlinewidth{0.602250pt}%
\definecolor{currentstroke}{rgb}{0.000000,0.000000,0.000000}%
\pgfsetstrokecolor{currentstroke}%
\pgfsetdash{}{0pt}%
\pgfsys@defobject{currentmarker}{\pgfqpoint{-0.027778in}{0.000000in}}{\pgfqpoint{-0.000000in}{0.000000in}}{%
\pgfpathmoveto{\pgfqpoint{-0.000000in}{0.000000in}}%
\pgfpathlineto{\pgfqpoint{-0.027778in}{0.000000in}}%
\pgfusepath{stroke,fill}%
}%
\begin{pgfscope}%
\pgfsys@transformshift{1.875000in}{5.345399in}%
\pgfsys@useobject{currentmarker}{}%
\end{pgfscope}%
\end{pgfscope}%
\begin{pgfscope}%
\pgfsetbuttcap%
\pgfsetroundjoin%
\definecolor{currentfill}{rgb}{0.000000,0.000000,0.000000}%
\pgfsetfillcolor{currentfill}%
\pgfsetlinewidth{0.602250pt}%
\definecolor{currentstroke}{rgb}{0.000000,0.000000,0.000000}%
\pgfsetstrokecolor{currentstroke}%
\pgfsetdash{}{0pt}%
\pgfsys@defobject{currentmarker}{\pgfqpoint{-0.027778in}{0.000000in}}{\pgfqpoint{-0.000000in}{0.000000in}}{%
\pgfpathmoveto{\pgfqpoint{-0.000000in}{0.000000in}}%
\pgfpathlineto{\pgfqpoint{-0.027778in}{0.000000in}}%
\pgfusepath{stroke,fill}%
}%
\begin{pgfscope}%
\pgfsys@transformshift{1.875000in}{5.455539in}%
\pgfsys@useobject{currentmarker}{}%
\end{pgfscope}%
\end{pgfscope}%
\begin{pgfscope}%
\pgfsetbuttcap%
\pgfsetroundjoin%
\definecolor{currentfill}{rgb}{0.000000,0.000000,0.000000}%
\pgfsetfillcolor{currentfill}%
\pgfsetlinewidth{0.602250pt}%
\definecolor{currentstroke}{rgb}{0.000000,0.000000,0.000000}%
\pgfsetstrokecolor{currentstroke}%
\pgfsetdash{}{0pt}%
\pgfsys@defobject{currentmarker}{\pgfqpoint{-0.027778in}{0.000000in}}{\pgfqpoint{-0.000000in}{0.000000in}}{%
\pgfpathmoveto{\pgfqpoint{-0.000000in}{0.000000in}}%
\pgfpathlineto{\pgfqpoint{-0.027778in}{0.000000in}}%
\pgfusepath{stroke,fill}%
}%
\begin{pgfscope}%
\pgfsys@transformshift{1.875000in}{5.540970in}%
\pgfsys@useobject{currentmarker}{}%
\end{pgfscope}%
\end{pgfscope}%
\begin{pgfscope}%
\pgfsetbuttcap%
\pgfsetroundjoin%
\definecolor{currentfill}{rgb}{0.000000,0.000000,0.000000}%
\pgfsetfillcolor{currentfill}%
\pgfsetlinewidth{0.602250pt}%
\definecolor{currentstroke}{rgb}{0.000000,0.000000,0.000000}%
\pgfsetstrokecolor{currentstroke}%
\pgfsetdash{}{0pt}%
\pgfsys@defobject{currentmarker}{\pgfqpoint{-0.027778in}{0.000000in}}{\pgfqpoint{-0.000000in}{0.000000in}}{%
\pgfpathmoveto{\pgfqpoint{-0.000000in}{0.000000in}}%
\pgfpathlineto{\pgfqpoint{-0.027778in}{0.000000in}}%
\pgfusepath{stroke,fill}%
}%
\begin{pgfscope}%
\pgfsys@transformshift{1.875000in}{5.610773in}%
\pgfsys@useobject{currentmarker}{}%
\end{pgfscope}%
\end{pgfscope}%
\begin{pgfscope}%
\pgfsetbuttcap%
\pgfsetroundjoin%
\definecolor{currentfill}{rgb}{0.000000,0.000000,0.000000}%
\pgfsetfillcolor{currentfill}%
\pgfsetlinewidth{0.602250pt}%
\definecolor{currentstroke}{rgb}{0.000000,0.000000,0.000000}%
\pgfsetstrokecolor{currentstroke}%
\pgfsetdash{}{0pt}%
\pgfsys@defobject{currentmarker}{\pgfqpoint{-0.027778in}{0.000000in}}{\pgfqpoint{-0.000000in}{0.000000in}}{%
\pgfpathmoveto{\pgfqpoint{-0.000000in}{0.000000in}}%
\pgfpathlineto{\pgfqpoint{-0.027778in}{0.000000in}}%
\pgfusepath{stroke,fill}%
}%
\begin{pgfscope}%
\pgfsys@transformshift{1.875000in}{5.669790in}%
\pgfsys@useobject{currentmarker}{}%
\end{pgfscope}%
\end{pgfscope}%
\begin{pgfscope}%
\pgfsetbuttcap%
\pgfsetroundjoin%
\definecolor{currentfill}{rgb}{0.000000,0.000000,0.000000}%
\pgfsetfillcolor{currentfill}%
\pgfsetlinewidth{0.602250pt}%
\definecolor{currentstroke}{rgb}{0.000000,0.000000,0.000000}%
\pgfsetstrokecolor{currentstroke}%
\pgfsetdash{}{0pt}%
\pgfsys@defobject{currentmarker}{\pgfqpoint{-0.027778in}{0.000000in}}{\pgfqpoint{-0.000000in}{0.000000in}}{%
\pgfpathmoveto{\pgfqpoint{-0.000000in}{0.000000in}}%
\pgfpathlineto{\pgfqpoint{-0.027778in}{0.000000in}}%
\pgfusepath{stroke,fill}%
}%
\begin{pgfscope}%
\pgfsys@transformshift{1.875000in}{5.720913in}%
\pgfsys@useobject{currentmarker}{}%
\end{pgfscope}%
\end{pgfscope}%
\begin{pgfscope}%
\pgfsetbuttcap%
\pgfsetroundjoin%
\definecolor{currentfill}{rgb}{0.000000,0.000000,0.000000}%
\pgfsetfillcolor{currentfill}%
\pgfsetlinewidth{0.602250pt}%
\definecolor{currentstroke}{rgb}{0.000000,0.000000,0.000000}%
\pgfsetstrokecolor{currentstroke}%
\pgfsetdash{}{0pt}%
\pgfsys@defobject{currentmarker}{\pgfqpoint{-0.027778in}{0.000000in}}{\pgfqpoint{-0.000000in}{0.000000in}}{%
\pgfpathmoveto{\pgfqpoint{-0.000000in}{0.000000in}}%
\pgfpathlineto{\pgfqpoint{-0.027778in}{0.000000in}}%
\pgfusepath{stroke,fill}%
}%
\begin{pgfscope}%
\pgfsys@transformshift{1.875000in}{5.766006in}%
\pgfsys@useobject{currentmarker}{}%
\end{pgfscope}%
\end{pgfscope}%
\begin{pgfscope}%
\pgfsetbuttcap%
\pgfsetroundjoin%
\definecolor{currentfill}{rgb}{0.000000,0.000000,0.000000}%
\pgfsetfillcolor{currentfill}%
\pgfsetlinewidth{0.602250pt}%
\definecolor{currentstroke}{rgb}{0.000000,0.000000,0.000000}%
\pgfsetstrokecolor{currentstroke}%
\pgfsetdash{}{0pt}%
\pgfsys@defobject{currentmarker}{\pgfqpoint{-0.027778in}{0.000000in}}{\pgfqpoint{-0.000000in}{0.000000in}}{%
\pgfpathmoveto{\pgfqpoint{-0.000000in}{0.000000in}}%
\pgfpathlineto{\pgfqpoint{-0.027778in}{0.000000in}}%
\pgfusepath{stroke,fill}%
}%
\begin{pgfscope}%
\pgfsys@transformshift{1.875000in}{6.071717in}%
\pgfsys@useobject{currentmarker}{}%
\end{pgfscope}%
\end{pgfscope}%
\begin{pgfscope}%
\pgfsetbuttcap%
\pgfsetroundjoin%
\definecolor{currentfill}{rgb}{0.000000,0.000000,0.000000}%
\pgfsetfillcolor{currentfill}%
\pgfsetlinewidth{0.602250pt}%
\definecolor{currentstroke}{rgb}{0.000000,0.000000,0.000000}%
\pgfsetstrokecolor{currentstroke}%
\pgfsetdash{}{0pt}%
\pgfsys@defobject{currentmarker}{\pgfqpoint{-0.027778in}{0.000000in}}{\pgfqpoint{-0.000000in}{0.000000in}}{%
\pgfpathmoveto{\pgfqpoint{-0.000000in}{0.000000in}}%
\pgfpathlineto{\pgfqpoint{-0.027778in}{0.000000in}}%
\pgfusepath{stroke,fill}%
}%
\begin{pgfscope}%
\pgfsys@transformshift{1.875000in}{6.226951in}%
\pgfsys@useobject{currentmarker}{}%
\end{pgfscope}%
\end{pgfscope}%
\begin{pgfscope}%
\pgfsetbuttcap%
\pgfsetroundjoin%
\definecolor{currentfill}{rgb}{0.000000,0.000000,0.000000}%
\pgfsetfillcolor{currentfill}%
\pgfsetlinewidth{0.602250pt}%
\definecolor{currentstroke}{rgb}{0.000000,0.000000,0.000000}%
\pgfsetstrokecolor{currentstroke}%
\pgfsetdash{}{0pt}%
\pgfsys@defobject{currentmarker}{\pgfqpoint{-0.027778in}{0.000000in}}{\pgfqpoint{-0.000000in}{0.000000in}}{%
\pgfpathmoveto{\pgfqpoint{-0.000000in}{0.000000in}}%
\pgfpathlineto{\pgfqpoint{-0.027778in}{0.000000in}}%
\pgfusepath{stroke,fill}%
}%
\begin{pgfscope}%
\pgfsys@transformshift{1.875000in}{6.337091in}%
\pgfsys@useobject{currentmarker}{}%
\end{pgfscope}%
\end{pgfscope}%
\begin{pgfscope}%
\pgfsetbuttcap%
\pgfsetroundjoin%
\definecolor{currentfill}{rgb}{0.000000,0.000000,0.000000}%
\pgfsetfillcolor{currentfill}%
\pgfsetlinewidth{0.602250pt}%
\definecolor{currentstroke}{rgb}{0.000000,0.000000,0.000000}%
\pgfsetstrokecolor{currentstroke}%
\pgfsetdash{}{0pt}%
\pgfsys@defobject{currentmarker}{\pgfqpoint{-0.027778in}{0.000000in}}{\pgfqpoint{-0.000000in}{0.000000in}}{%
\pgfpathmoveto{\pgfqpoint{-0.000000in}{0.000000in}}%
\pgfpathlineto{\pgfqpoint{-0.027778in}{0.000000in}}%
\pgfusepath{stroke,fill}%
}%
\begin{pgfscope}%
\pgfsys@transformshift{1.875000in}{6.422522in}%
\pgfsys@useobject{currentmarker}{}%
\end{pgfscope}%
\end{pgfscope}%
\begin{pgfscope}%
\pgfsetbuttcap%
\pgfsetroundjoin%
\definecolor{currentfill}{rgb}{0.000000,0.000000,0.000000}%
\pgfsetfillcolor{currentfill}%
\pgfsetlinewidth{0.602250pt}%
\definecolor{currentstroke}{rgb}{0.000000,0.000000,0.000000}%
\pgfsetstrokecolor{currentstroke}%
\pgfsetdash{}{0pt}%
\pgfsys@defobject{currentmarker}{\pgfqpoint{-0.027778in}{0.000000in}}{\pgfqpoint{-0.000000in}{0.000000in}}{%
\pgfpathmoveto{\pgfqpoint{-0.000000in}{0.000000in}}%
\pgfpathlineto{\pgfqpoint{-0.027778in}{0.000000in}}%
\pgfusepath{stroke,fill}%
}%
\begin{pgfscope}%
\pgfsys@transformshift{1.875000in}{6.492324in}%
\pgfsys@useobject{currentmarker}{}%
\end{pgfscope}%
\end{pgfscope}%
\begin{pgfscope}%
\pgfsetbuttcap%
\pgfsetroundjoin%
\definecolor{currentfill}{rgb}{0.000000,0.000000,0.000000}%
\pgfsetfillcolor{currentfill}%
\pgfsetlinewidth{0.602250pt}%
\definecolor{currentstroke}{rgb}{0.000000,0.000000,0.000000}%
\pgfsetstrokecolor{currentstroke}%
\pgfsetdash{}{0pt}%
\pgfsys@defobject{currentmarker}{\pgfqpoint{-0.027778in}{0.000000in}}{\pgfqpoint{-0.000000in}{0.000000in}}{%
\pgfpathmoveto{\pgfqpoint{-0.000000in}{0.000000in}}%
\pgfpathlineto{\pgfqpoint{-0.027778in}{0.000000in}}%
\pgfusepath{stroke,fill}%
}%
\begin{pgfscope}%
\pgfsys@transformshift{1.875000in}{6.551342in}%
\pgfsys@useobject{currentmarker}{}%
\end{pgfscope}%
\end{pgfscope}%
\begin{pgfscope}%
\pgfsetbuttcap%
\pgfsetroundjoin%
\definecolor{currentfill}{rgb}{0.000000,0.000000,0.000000}%
\pgfsetfillcolor{currentfill}%
\pgfsetlinewidth{0.602250pt}%
\definecolor{currentstroke}{rgb}{0.000000,0.000000,0.000000}%
\pgfsetstrokecolor{currentstroke}%
\pgfsetdash{}{0pt}%
\pgfsys@defobject{currentmarker}{\pgfqpoint{-0.027778in}{0.000000in}}{\pgfqpoint{-0.000000in}{0.000000in}}{%
\pgfpathmoveto{\pgfqpoint{-0.000000in}{0.000000in}}%
\pgfpathlineto{\pgfqpoint{-0.027778in}{0.000000in}}%
\pgfusepath{stroke,fill}%
}%
\begin{pgfscope}%
\pgfsys@transformshift{1.875000in}{6.602464in}%
\pgfsys@useobject{currentmarker}{}%
\end{pgfscope}%
\end{pgfscope}%
\begin{pgfscope}%
\pgfsetbuttcap%
\pgfsetroundjoin%
\definecolor{currentfill}{rgb}{0.000000,0.000000,0.000000}%
\pgfsetfillcolor{currentfill}%
\pgfsetlinewidth{0.602250pt}%
\definecolor{currentstroke}{rgb}{0.000000,0.000000,0.000000}%
\pgfsetstrokecolor{currentstroke}%
\pgfsetdash{}{0pt}%
\pgfsys@defobject{currentmarker}{\pgfqpoint{-0.027778in}{0.000000in}}{\pgfqpoint{-0.000000in}{0.000000in}}{%
\pgfpathmoveto{\pgfqpoint{-0.000000in}{0.000000in}}%
\pgfpathlineto{\pgfqpoint{-0.027778in}{0.000000in}}%
\pgfusepath{stroke,fill}%
}%
\begin{pgfscope}%
\pgfsys@transformshift{1.875000in}{6.647558in}%
\pgfsys@useobject{currentmarker}{}%
\end{pgfscope}%
\end{pgfscope}%
\begin{pgfscope}%
\pgfsetbuttcap%
\pgfsetroundjoin%
\definecolor{currentfill}{rgb}{0.000000,0.000000,0.000000}%
\pgfsetfillcolor{currentfill}%
\pgfsetlinewidth{0.602250pt}%
\definecolor{currentstroke}{rgb}{0.000000,0.000000,0.000000}%
\pgfsetstrokecolor{currentstroke}%
\pgfsetdash{}{0pt}%
\pgfsys@defobject{currentmarker}{\pgfqpoint{-0.027778in}{0.000000in}}{\pgfqpoint{-0.000000in}{0.000000in}}{%
\pgfpathmoveto{\pgfqpoint{-0.000000in}{0.000000in}}%
\pgfpathlineto{\pgfqpoint{-0.027778in}{0.000000in}}%
\pgfusepath{stroke,fill}%
}%
\begin{pgfscope}%
\pgfsys@transformshift{1.875000in}{6.953269in}%
\pgfsys@useobject{currentmarker}{}%
\end{pgfscope}%
\end{pgfscope}%
\begin{pgfscope}%
\pgfsetbuttcap%
\pgfsetroundjoin%
\definecolor{currentfill}{rgb}{0.000000,0.000000,0.000000}%
\pgfsetfillcolor{currentfill}%
\pgfsetlinewidth{0.602250pt}%
\definecolor{currentstroke}{rgb}{0.000000,0.000000,0.000000}%
\pgfsetstrokecolor{currentstroke}%
\pgfsetdash{}{0pt}%
\pgfsys@defobject{currentmarker}{\pgfqpoint{-0.027778in}{0.000000in}}{\pgfqpoint{-0.000000in}{0.000000in}}{%
\pgfpathmoveto{\pgfqpoint{-0.000000in}{0.000000in}}%
\pgfpathlineto{\pgfqpoint{-0.027778in}{0.000000in}}%
\pgfusepath{stroke,fill}%
}%
\begin{pgfscope}%
\pgfsys@transformshift{1.875000in}{7.108503in}%
\pgfsys@useobject{currentmarker}{}%
\end{pgfscope}%
\end{pgfscope}%
\begin{pgfscope}%
\pgfsetbuttcap%
\pgfsetroundjoin%
\definecolor{currentfill}{rgb}{0.000000,0.000000,0.000000}%
\pgfsetfillcolor{currentfill}%
\pgfsetlinewidth{0.602250pt}%
\definecolor{currentstroke}{rgb}{0.000000,0.000000,0.000000}%
\pgfsetstrokecolor{currentstroke}%
\pgfsetdash{}{0pt}%
\pgfsys@defobject{currentmarker}{\pgfqpoint{-0.027778in}{0.000000in}}{\pgfqpoint{-0.000000in}{0.000000in}}{%
\pgfpathmoveto{\pgfqpoint{-0.000000in}{0.000000in}}%
\pgfpathlineto{\pgfqpoint{-0.027778in}{0.000000in}}%
\pgfusepath{stroke,fill}%
}%
\begin{pgfscope}%
\pgfsys@transformshift{1.875000in}{7.218643in}%
\pgfsys@useobject{currentmarker}{}%
\end{pgfscope}%
\end{pgfscope}%
\begin{pgfscope}%
\pgfsetbuttcap%
\pgfsetroundjoin%
\definecolor{currentfill}{rgb}{0.000000,0.000000,0.000000}%
\pgfsetfillcolor{currentfill}%
\pgfsetlinewidth{0.602250pt}%
\definecolor{currentstroke}{rgb}{0.000000,0.000000,0.000000}%
\pgfsetstrokecolor{currentstroke}%
\pgfsetdash{}{0pt}%
\pgfsys@defobject{currentmarker}{\pgfqpoint{-0.027778in}{0.000000in}}{\pgfqpoint{-0.000000in}{0.000000in}}{%
\pgfpathmoveto{\pgfqpoint{-0.000000in}{0.000000in}}%
\pgfpathlineto{\pgfqpoint{-0.027778in}{0.000000in}}%
\pgfusepath{stroke,fill}%
}%
\begin{pgfscope}%
\pgfsys@transformshift{1.875000in}{7.304074in}%
\pgfsys@useobject{currentmarker}{}%
\end{pgfscope}%
\end{pgfscope}%
\begin{pgfscope}%
\pgfsetbuttcap%
\pgfsetroundjoin%
\definecolor{currentfill}{rgb}{0.000000,0.000000,0.000000}%
\pgfsetfillcolor{currentfill}%
\pgfsetlinewidth{0.602250pt}%
\definecolor{currentstroke}{rgb}{0.000000,0.000000,0.000000}%
\pgfsetstrokecolor{currentstroke}%
\pgfsetdash{}{0pt}%
\pgfsys@defobject{currentmarker}{\pgfqpoint{-0.027778in}{0.000000in}}{\pgfqpoint{-0.000000in}{0.000000in}}{%
\pgfpathmoveto{\pgfqpoint{-0.000000in}{0.000000in}}%
\pgfpathlineto{\pgfqpoint{-0.027778in}{0.000000in}}%
\pgfusepath{stroke,fill}%
}%
\begin{pgfscope}%
\pgfsys@transformshift{1.875000in}{7.373876in}%
\pgfsys@useobject{currentmarker}{}%
\end{pgfscope}%
\end{pgfscope}%
\begin{pgfscope}%
\pgfsetbuttcap%
\pgfsetroundjoin%
\definecolor{currentfill}{rgb}{0.000000,0.000000,0.000000}%
\pgfsetfillcolor{currentfill}%
\pgfsetlinewidth{0.602250pt}%
\definecolor{currentstroke}{rgb}{0.000000,0.000000,0.000000}%
\pgfsetstrokecolor{currentstroke}%
\pgfsetdash{}{0pt}%
\pgfsys@defobject{currentmarker}{\pgfqpoint{-0.027778in}{0.000000in}}{\pgfqpoint{-0.000000in}{0.000000in}}{%
\pgfpathmoveto{\pgfqpoint{-0.000000in}{0.000000in}}%
\pgfpathlineto{\pgfqpoint{-0.027778in}{0.000000in}}%
\pgfusepath{stroke,fill}%
}%
\begin{pgfscope}%
\pgfsys@transformshift{1.875000in}{7.432894in}%
\pgfsys@useobject{currentmarker}{}%
\end{pgfscope}%
\end{pgfscope}%
\begin{pgfscope}%
\pgfsetbuttcap%
\pgfsetroundjoin%
\definecolor{currentfill}{rgb}{0.000000,0.000000,0.000000}%
\pgfsetfillcolor{currentfill}%
\pgfsetlinewidth{0.602250pt}%
\definecolor{currentstroke}{rgb}{0.000000,0.000000,0.000000}%
\pgfsetstrokecolor{currentstroke}%
\pgfsetdash{}{0pt}%
\pgfsys@defobject{currentmarker}{\pgfqpoint{-0.027778in}{0.000000in}}{\pgfqpoint{-0.000000in}{0.000000in}}{%
\pgfpathmoveto{\pgfqpoint{-0.000000in}{0.000000in}}%
\pgfpathlineto{\pgfqpoint{-0.027778in}{0.000000in}}%
\pgfusepath{stroke,fill}%
}%
\begin{pgfscope}%
\pgfsys@transformshift{1.875000in}{7.484016in}%
\pgfsys@useobject{currentmarker}{}%
\end{pgfscope}%
\end{pgfscope}%
\begin{pgfscope}%
\pgfsetbuttcap%
\pgfsetroundjoin%
\definecolor{currentfill}{rgb}{0.000000,0.000000,0.000000}%
\pgfsetfillcolor{currentfill}%
\pgfsetlinewidth{0.602250pt}%
\definecolor{currentstroke}{rgb}{0.000000,0.000000,0.000000}%
\pgfsetstrokecolor{currentstroke}%
\pgfsetdash{}{0pt}%
\pgfsys@defobject{currentmarker}{\pgfqpoint{-0.027778in}{0.000000in}}{\pgfqpoint{-0.000000in}{0.000000in}}{%
\pgfpathmoveto{\pgfqpoint{-0.000000in}{0.000000in}}%
\pgfpathlineto{\pgfqpoint{-0.027778in}{0.000000in}}%
\pgfusepath{stroke,fill}%
}%
\begin{pgfscope}%
\pgfsys@transformshift{1.875000in}{7.529110in}%
\pgfsys@useobject{currentmarker}{}%
\end{pgfscope}%
\end{pgfscope}%
\begin{pgfscope}%
\pgfsetbuttcap%
\pgfsetroundjoin%
\definecolor{currentfill}{rgb}{0.000000,0.000000,0.000000}%
\pgfsetfillcolor{currentfill}%
\pgfsetlinewidth{0.602250pt}%
\definecolor{currentstroke}{rgb}{0.000000,0.000000,0.000000}%
\pgfsetstrokecolor{currentstroke}%
\pgfsetdash{}{0pt}%
\pgfsys@defobject{currentmarker}{\pgfqpoint{-0.027778in}{0.000000in}}{\pgfqpoint{-0.000000in}{0.000000in}}{%
\pgfpathmoveto{\pgfqpoint{-0.000000in}{0.000000in}}%
\pgfpathlineto{\pgfqpoint{-0.027778in}{0.000000in}}%
\pgfusepath{stroke,fill}%
}%
\begin{pgfscope}%
\pgfsys@transformshift{1.875000in}{7.834821in}%
\pgfsys@useobject{currentmarker}{}%
\end{pgfscope}%
\end{pgfscope}%
\begin{pgfscope}%
\pgfsetbuttcap%
\pgfsetroundjoin%
\definecolor{currentfill}{rgb}{0.000000,0.000000,0.000000}%
\pgfsetfillcolor{currentfill}%
\pgfsetlinewidth{0.602250pt}%
\definecolor{currentstroke}{rgb}{0.000000,0.000000,0.000000}%
\pgfsetstrokecolor{currentstroke}%
\pgfsetdash{}{0pt}%
\pgfsys@defobject{currentmarker}{\pgfqpoint{-0.027778in}{0.000000in}}{\pgfqpoint{-0.000000in}{0.000000in}}{%
\pgfpathmoveto{\pgfqpoint{-0.000000in}{0.000000in}}%
\pgfpathlineto{\pgfqpoint{-0.027778in}{0.000000in}}%
\pgfusepath{stroke,fill}%
}%
\begin{pgfscope}%
\pgfsys@transformshift{1.875000in}{7.990055in}%
\pgfsys@useobject{currentmarker}{}%
\end{pgfscope}%
\end{pgfscope}%
\begin{pgfscope}%
\pgfsetbuttcap%
\pgfsetroundjoin%
\definecolor{currentfill}{rgb}{0.000000,0.000000,0.000000}%
\pgfsetfillcolor{currentfill}%
\pgfsetlinewidth{0.602250pt}%
\definecolor{currentstroke}{rgb}{0.000000,0.000000,0.000000}%
\pgfsetstrokecolor{currentstroke}%
\pgfsetdash{}{0pt}%
\pgfsys@defobject{currentmarker}{\pgfqpoint{-0.027778in}{0.000000in}}{\pgfqpoint{-0.000000in}{0.000000in}}{%
\pgfpathmoveto{\pgfqpoint{-0.000000in}{0.000000in}}%
\pgfpathlineto{\pgfqpoint{-0.027778in}{0.000000in}}%
\pgfusepath{stroke,fill}%
}%
\begin{pgfscope}%
\pgfsys@transformshift{1.875000in}{8.100195in}%
\pgfsys@useobject{currentmarker}{}%
\end{pgfscope}%
\end{pgfscope}%
\begin{pgfscope}%
\pgfsetbuttcap%
\pgfsetroundjoin%
\definecolor{currentfill}{rgb}{0.000000,0.000000,0.000000}%
\pgfsetfillcolor{currentfill}%
\pgfsetlinewidth{0.602250pt}%
\definecolor{currentstroke}{rgb}{0.000000,0.000000,0.000000}%
\pgfsetstrokecolor{currentstroke}%
\pgfsetdash{}{0pt}%
\pgfsys@defobject{currentmarker}{\pgfqpoint{-0.027778in}{0.000000in}}{\pgfqpoint{-0.000000in}{0.000000in}}{%
\pgfpathmoveto{\pgfqpoint{-0.000000in}{0.000000in}}%
\pgfpathlineto{\pgfqpoint{-0.027778in}{0.000000in}}%
\pgfusepath{stroke,fill}%
}%
\begin{pgfscope}%
\pgfsys@transformshift{1.875000in}{8.185626in}%
\pgfsys@useobject{currentmarker}{}%
\end{pgfscope}%
\end{pgfscope}%
\begin{pgfscope}%
\pgfsetbuttcap%
\pgfsetroundjoin%
\definecolor{currentfill}{rgb}{0.000000,0.000000,0.000000}%
\pgfsetfillcolor{currentfill}%
\pgfsetlinewidth{0.602250pt}%
\definecolor{currentstroke}{rgb}{0.000000,0.000000,0.000000}%
\pgfsetstrokecolor{currentstroke}%
\pgfsetdash{}{0pt}%
\pgfsys@defobject{currentmarker}{\pgfqpoint{-0.027778in}{0.000000in}}{\pgfqpoint{-0.000000in}{0.000000in}}{%
\pgfpathmoveto{\pgfqpoint{-0.000000in}{0.000000in}}%
\pgfpathlineto{\pgfqpoint{-0.027778in}{0.000000in}}%
\pgfusepath{stroke,fill}%
}%
\begin{pgfscope}%
\pgfsys@transformshift{1.875000in}{8.255428in}%
\pgfsys@useobject{currentmarker}{}%
\end{pgfscope}%
\end{pgfscope}%
\begin{pgfscope}%
\pgfsetbuttcap%
\pgfsetroundjoin%
\definecolor{currentfill}{rgb}{0.000000,0.000000,0.000000}%
\pgfsetfillcolor{currentfill}%
\pgfsetlinewidth{0.602250pt}%
\definecolor{currentstroke}{rgb}{0.000000,0.000000,0.000000}%
\pgfsetstrokecolor{currentstroke}%
\pgfsetdash{}{0pt}%
\pgfsys@defobject{currentmarker}{\pgfqpoint{-0.027778in}{0.000000in}}{\pgfqpoint{-0.000000in}{0.000000in}}{%
\pgfpathmoveto{\pgfqpoint{-0.000000in}{0.000000in}}%
\pgfpathlineto{\pgfqpoint{-0.027778in}{0.000000in}}%
\pgfusepath{stroke,fill}%
}%
\begin{pgfscope}%
\pgfsys@transformshift{1.875000in}{8.314446in}%
\pgfsys@useobject{currentmarker}{}%
\end{pgfscope}%
\end{pgfscope}%
\begin{pgfscope}%
\pgfsetbuttcap%
\pgfsetroundjoin%
\definecolor{currentfill}{rgb}{0.000000,0.000000,0.000000}%
\pgfsetfillcolor{currentfill}%
\pgfsetlinewidth{0.602250pt}%
\definecolor{currentstroke}{rgb}{0.000000,0.000000,0.000000}%
\pgfsetstrokecolor{currentstroke}%
\pgfsetdash{}{0pt}%
\pgfsys@defobject{currentmarker}{\pgfqpoint{-0.027778in}{0.000000in}}{\pgfqpoint{-0.000000in}{0.000000in}}{%
\pgfpathmoveto{\pgfqpoint{-0.000000in}{0.000000in}}%
\pgfpathlineto{\pgfqpoint{-0.027778in}{0.000000in}}%
\pgfusepath{stroke,fill}%
}%
\begin{pgfscope}%
\pgfsys@transformshift{1.875000in}{8.365568in}%
\pgfsys@useobject{currentmarker}{}%
\end{pgfscope}%
\end{pgfscope}%
\begin{pgfscope}%
\pgfsetbuttcap%
\pgfsetroundjoin%
\definecolor{currentfill}{rgb}{0.000000,0.000000,0.000000}%
\pgfsetfillcolor{currentfill}%
\pgfsetlinewidth{0.602250pt}%
\definecolor{currentstroke}{rgb}{0.000000,0.000000,0.000000}%
\pgfsetstrokecolor{currentstroke}%
\pgfsetdash{}{0pt}%
\pgfsys@defobject{currentmarker}{\pgfqpoint{-0.027778in}{0.000000in}}{\pgfqpoint{-0.000000in}{0.000000in}}{%
\pgfpathmoveto{\pgfqpoint{-0.000000in}{0.000000in}}%
\pgfpathlineto{\pgfqpoint{-0.027778in}{0.000000in}}%
\pgfusepath{stroke,fill}%
}%
\begin{pgfscope}%
\pgfsys@transformshift{1.875000in}{8.410662in}%
\pgfsys@useobject{currentmarker}{}%
\end{pgfscope}%
\end{pgfscope}%
\begin{pgfscope}%
\pgfsetbuttcap%
\pgfsetroundjoin%
\definecolor{currentfill}{rgb}{0.000000,0.000000,0.000000}%
\pgfsetfillcolor{currentfill}%
\pgfsetlinewidth{0.602250pt}%
\definecolor{currentstroke}{rgb}{0.000000,0.000000,0.000000}%
\pgfsetstrokecolor{currentstroke}%
\pgfsetdash{}{0pt}%
\pgfsys@defobject{currentmarker}{\pgfqpoint{-0.027778in}{0.000000in}}{\pgfqpoint{-0.000000in}{0.000000in}}{%
\pgfpathmoveto{\pgfqpoint{-0.000000in}{0.000000in}}%
\pgfpathlineto{\pgfqpoint{-0.027778in}{0.000000in}}%
\pgfusepath{stroke,fill}%
}%
\begin{pgfscope}%
\pgfsys@transformshift{1.875000in}{8.716373in}%
\pgfsys@useobject{currentmarker}{}%
\end{pgfscope}%
\end{pgfscope}%
\begin{pgfscope}%
\definecolor{textcolor}{rgb}{0.000000,0.000000,0.000000}%
\pgfsetstrokecolor{textcolor}%
\pgfsetfillcolor{textcolor}%
\pgftext[x=1.434220in,y=7.050000in,,bottom,rotate=90.000000]{\color{textcolor}{\sffamily\fontsize{10.000000}{12.000000}\selectfont\catcode`\^=\active\def^{\ifmmode\sp\else\^{}\fi}\catcode`\%=\active\def%{\%}Error}}%
\end{pgfscope}%
\begin{pgfscope}%
\pgfpathrectangle{\pgfqpoint{1.875000in}{5.300000in}}{\pgfqpoint{3.419118in}{3.500000in}}%
\pgfusepath{clip}%
\pgfsetrectcap%
\pgfsetroundjoin%
\pgfsetlinewidth{1.505625pt}%
\definecolor{currentstroke}{rgb}{0.121569,0.466667,0.705882}%
\pgfsetstrokecolor{currentstroke}%
\pgfsetdash{}{0pt}%
\pgfpathmoveto{\pgfqpoint{2.030414in}{8.127663in}}%
\pgfpathlineto{\pgfqpoint{2.207140in}{8.092774in}}%
\pgfpathlineto{\pgfqpoint{2.435843in}{8.050040in}}%
\pgfpathlineto{\pgfqpoint{2.758107in}{7.992243in}}%
\pgfpathlineto{\pgfqpoint{3.298679in}{7.897836in}}%
\pgfpathlineto{\pgfqpoint{4.785252in}{7.641292in}}%
\pgfpathlineto{\pgfqpoint{5.138703in}{7.580379in}}%
\pgfpathlineto{\pgfqpoint{5.138703in}{7.580379in}}%
\pgfusepath{stroke}%
\end{pgfscope}%
\begin{pgfscope}%
\pgfpathrectangle{\pgfqpoint{1.875000in}{5.300000in}}{\pgfqpoint{3.419118in}{3.500000in}}%
\pgfusepath{clip}%
\pgfsetrectcap%
\pgfsetroundjoin%
\pgfsetlinewidth{1.505625pt}%
\definecolor{currentstroke}{rgb}{1.000000,0.498039,0.054902}%
\pgfsetstrokecolor{currentstroke}%
\pgfsetdash{}{0pt}%
\pgfpathmoveto{\pgfqpoint{2.030414in}{8.307443in}}%
\pgfpathlineto{\pgfqpoint{2.259118in}{8.264864in}}%
\pgfpathlineto{\pgfqpoint{2.269514in}{8.262972in}}%
\pgfpathlineto{\pgfqpoint{2.279909in}{6.947992in}}%
\pgfpathlineto{\pgfqpoint{2.290305in}{6.913400in}}%
\pgfpathlineto{\pgfqpoint{2.300700in}{6.889766in}}%
\pgfpathlineto{\pgfqpoint{2.321492in}{6.851668in}}%
\pgfpathlineto{\pgfqpoint{2.342283in}{6.819558in}}%
\pgfpathlineto{\pgfqpoint{2.363074in}{6.791070in}}%
\pgfpathlineto{\pgfqpoint{2.394261in}{6.753030in}}%
\pgfpathlineto{\pgfqpoint{2.425448in}{6.719113in}}%
\pgfpathlineto{\pgfqpoint{2.456635in}{6.688290in}}%
\pgfpathlineto{\pgfqpoint{2.498217in}{6.650825in}}%
\pgfpathlineto{\pgfqpoint{2.529404in}{6.624810in}}%
\pgfpathlineto{\pgfqpoint{2.539800in}{6.778814in}}%
\pgfpathlineto{\pgfqpoint{2.550195in}{6.771203in}}%
\pgfpathlineto{\pgfqpoint{2.570986in}{6.761873in}}%
\pgfpathlineto{\pgfqpoint{2.602173in}{6.750908in}}%
\pgfpathlineto{\pgfqpoint{2.654151in}{6.735283in}}%
\pgfpathlineto{\pgfqpoint{2.737316in}{6.713072in}}%
\pgfpathlineto{\pgfqpoint{2.789294in}{6.700183in}}%
\pgfpathlineto{\pgfqpoint{2.799690in}{6.666587in}}%
\pgfpathlineto{\pgfqpoint{2.841272in}{6.656128in}}%
\pgfpathlineto{\pgfqpoint{2.924437in}{6.637970in}}%
\pgfpathlineto{\pgfqpoint{3.049185in}{6.613210in}}%
\pgfpathlineto{\pgfqpoint{3.059580in}{6.657502in}}%
\pgfpathlineto{\pgfqpoint{3.101163in}{6.647573in}}%
\pgfpathlineto{\pgfqpoint{3.194723in}{6.627946in}}%
\pgfpathlineto{\pgfqpoint{3.309075in}{6.605697in}}%
\pgfpathlineto{\pgfqpoint{3.319471in}{6.686693in}}%
\pgfpathlineto{\pgfqpoint{3.371449in}{6.674780in}}%
\pgfpathlineto{\pgfqpoint{3.475405in}{6.653533in}}%
\pgfpathlineto{\pgfqpoint{3.568965in}{6.635559in}}%
\pgfpathlineto{\pgfqpoint{3.579361in}{6.690981in}}%
\pgfpathlineto{\pgfqpoint{3.631339in}{6.678935in}}%
\pgfpathlineto{\pgfqpoint{3.735295in}{6.657480in}}%
\pgfpathlineto{\pgfqpoint{3.828856in}{6.639381in}}%
\pgfpathlineto{\pgfqpoint{3.839251in}{6.688004in}}%
\pgfpathlineto{\pgfqpoint{3.891229in}{6.675749in}}%
\pgfpathlineto{\pgfqpoint{3.995186in}{6.654007in}}%
\pgfpathlineto{\pgfqpoint{4.088746in}{6.635739in}}%
\pgfpathlineto{\pgfqpoint{4.099142in}{6.679027in}}%
\pgfpathlineto{\pgfqpoint{4.151120in}{6.666495in}}%
\pgfpathlineto{\pgfqpoint{4.244680in}{6.646514in}}%
\pgfpathlineto{\pgfqpoint{4.348636in}{6.625918in}}%
\pgfpathlineto{\pgfqpoint{4.359032in}{6.666464in}}%
\pgfpathlineto{\pgfqpoint{4.400615in}{6.656033in}}%
\pgfpathlineto{\pgfqpoint{4.483779in}{6.637619in}}%
\pgfpathlineto{\pgfqpoint{4.608527in}{6.612404in}}%
\pgfpathlineto{\pgfqpoint{4.618922in}{6.652081in}}%
\pgfpathlineto{\pgfqpoint{4.660505in}{6.641364in}}%
\pgfpathlineto{\pgfqpoint{4.743670in}{6.622573in}}%
\pgfpathlineto{\pgfqpoint{4.868417in}{6.597007in}}%
\pgfpathlineto{\pgfqpoint{4.878813in}{6.637160in}}%
\pgfpathlineto{\pgfqpoint{4.920395in}{6.626156in}}%
\pgfpathlineto{\pgfqpoint{5.003560in}{6.606990in}}%
\pgfpathlineto{\pgfqpoint{5.117912in}{6.583144in}}%
\pgfpathlineto{\pgfqpoint{5.128308in}{6.581066in}}%
\pgfpathlineto{\pgfqpoint{5.138703in}{6.622715in}}%
\pgfpathlineto{\pgfqpoint{5.138703in}{6.622715in}}%
\pgfusepath{stroke}%
\end{pgfscope}%
\begin{pgfscope}%
\pgfpathrectangle{\pgfqpoint{1.875000in}{5.300000in}}{\pgfqpoint{3.419118in}{3.500000in}}%
\pgfusepath{clip}%
\pgfsetrectcap%
\pgfsetroundjoin%
\pgfsetlinewidth{1.505625pt}%
\definecolor{currentstroke}{rgb}{0.172549,0.627451,0.172549}%
\pgfsetstrokecolor{currentstroke}%
\pgfsetdash{}{0pt}%
\pgfpathmoveto{\pgfqpoint{2.030414in}{6.886577in}}%
\pgfpathlineto{\pgfqpoint{2.040810in}{6.865139in}}%
\pgfpathlineto{\pgfqpoint{2.061601in}{6.830008in}}%
\pgfpathlineto{\pgfqpoint{2.092788in}{6.782231in}}%
\pgfpathlineto{\pgfqpoint{2.123975in}{6.737941in}}%
\pgfpathlineto{\pgfqpoint{2.165557in}{6.683295in}}%
\pgfpathlineto{\pgfqpoint{2.207140in}{6.633333in}}%
\pgfpathlineto{\pgfqpoint{2.238327in}{6.598877in}}%
\pgfpathlineto{\pgfqpoint{2.269514in}{6.566991in}}%
\pgfpathlineto{\pgfqpoint{2.300700in}{6.537646in}}%
\pgfpathlineto{\pgfqpoint{2.331887in}{6.510784in}}%
\pgfpathlineto{\pgfqpoint{2.363074in}{6.486313in}}%
\pgfpathlineto{\pgfqpoint{2.394261in}{6.464108in}}%
\pgfpathlineto{\pgfqpoint{2.425448in}{6.444013in}}%
\pgfpathlineto{\pgfqpoint{2.467030in}{6.420194in}}%
\pgfpathlineto{\pgfqpoint{2.508613in}{6.399364in}}%
\pgfpathlineto{\pgfqpoint{2.550195in}{6.381071in}}%
\pgfpathlineto{\pgfqpoint{2.602173in}{6.361129in}}%
\pgfpathlineto{\pgfqpoint{2.654151in}{6.343756in}}%
\pgfpathlineto{\pgfqpoint{2.716525in}{6.325455in}}%
\pgfpathlineto{\pgfqpoint{2.789294in}{6.306619in}}%
\pgfpathlineto{\pgfqpoint{2.882855in}{6.285035in}}%
\pgfpathlineto{\pgfqpoint{2.997207in}{6.261102in}}%
\pgfpathlineto{\pgfqpoint{3.153141in}{6.230847in}}%
\pgfpathlineto{\pgfqpoint{3.381844in}{6.188892in}}%
\pgfpathlineto{\pgfqpoint{3.745691in}{6.124637in}}%
\pgfpathlineto{\pgfqpoint{4.400615in}{6.011565in}}%
\pgfpathlineto{\pgfqpoint{5.138703in}{5.885447in}}%
\pgfpathlineto{\pgfqpoint{5.138703in}{5.885447in}}%
\pgfusepath{stroke}%
\end{pgfscope}%
\begin{pgfscope}%
\pgfsetrectcap%
\pgfsetmiterjoin%
\pgfsetlinewidth{0.803000pt}%
\definecolor{currentstroke}{rgb}{0.000000,0.000000,0.000000}%
\pgfsetstrokecolor{currentstroke}%
\pgfsetdash{}{0pt}%
\pgfpathmoveto{\pgfqpoint{1.875000in}{5.300000in}}%
\pgfpathlineto{\pgfqpoint{1.875000in}{8.800000in}}%
\pgfusepath{stroke}%
\end{pgfscope}%
\begin{pgfscope}%
\pgfsetrectcap%
\pgfsetmiterjoin%
\pgfsetlinewidth{0.803000pt}%
\definecolor{currentstroke}{rgb}{0.000000,0.000000,0.000000}%
\pgfsetstrokecolor{currentstroke}%
\pgfsetdash{}{0pt}%
\pgfpathmoveto{\pgfqpoint{5.294118in}{5.300000in}}%
\pgfpathlineto{\pgfqpoint{5.294118in}{8.800000in}}%
\pgfusepath{stroke}%
\end{pgfscope}%
\begin{pgfscope}%
\pgfsetrectcap%
\pgfsetmiterjoin%
\pgfsetlinewidth{0.803000pt}%
\definecolor{currentstroke}{rgb}{0.000000,0.000000,0.000000}%
\pgfsetstrokecolor{currentstroke}%
\pgfsetdash{}{0pt}%
\pgfpathmoveto{\pgfqpoint{1.875000in}{5.300000in}}%
\pgfpathlineto{\pgfqpoint{5.294118in}{5.300000in}}%
\pgfusepath{stroke}%
\end{pgfscope}%
\begin{pgfscope}%
\pgfsetrectcap%
\pgfsetmiterjoin%
\pgfsetlinewidth{0.803000pt}%
\definecolor{currentstroke}{rgb}{0.000000,0.000000,0.000000}%
\pgfsetstrokecolor{currentstroke}%
\pgfsetdash{}{0pt}%
\pgfpathmoveto{\pgfqpoint{1.875000in}{8.800000in}}%
\pgfpathlineto{\pgfqpoint{5.294118in}{8.800000in}}%
\pgfusepath{stroke}%
\end{pgfscope}%
\begin{pgfscope}%
\definecolor{textcolor}{rgb}{0.000000,0.000000,0.000000}%
\pgfsetstrokecolor{textcolor}%
\pgfsetfillcolor{textcolor}%
\pgftext[x=3.584559in,y=8.883333in,,base]{\color{textcolor}{\sffamily\fontsize{12.000000}{14.400000}\selectfont\catcode`\^=\active\def^{\ifmmode\sp\else\^{}\fi}\catcode`\%=\active\def%{\%}Jacobi}}%
\end{pgfscope}%
\begin{pgfscope}%
\pgfsetbuttcap%
\pgfsetmiterjoin%
\definecolor{currentfill}{rgb}{1.000000,1.000000,1.000000}%
\pgfsetfillcolor{currentfill}%
\pgfsetlinewidth{0.000000pt}%
\definecolor{currentstroke}{rgb}{0.000000,0.000000,0.000000}%
\pgfsetstrokecolor{currentstroke}%
\pgfsetstrokeopacity{0.000000}%
\pgfsetdash{}{0pt}%
\pgfpathmoveto{\pgfqpoint{5.977941in}{5.300000in}}%
\pgfpathlineto{\pgfqpoint{9.397059in}{5.300000in}}%
\pgfpathlineto{\pgfqpoint{9.397059in}{8.800000in}}%
\pgfpathlineto{\pgfqpoint{5.977941in}{8.800000in}}%
\pgfpathlineto{\pgfqpoint{5.977941in}{5.300000in}}%
\pgfpathclose%
\pgfusepath{fill}%
\end{pgfscope}%
\begin{pgfscope}%
\pgfsetbuttcap%
\pgfsetroundjoin%
\definecolor{currentfill}{rgb}{0.000000,0.000000,0.000000}%
\pgfsetfillcolor{currentfill}%
\pgfsetlinewidth{0.803000pt}%
\definecolor{currentstroke}{rgb}{0.000000,0.000000,0.000000}%
\pgfsetstrokecolor{currentstroke}%
\pgfsetdash{}{0pt}%
\pgfsys@defobject{currentmarker}{\pgfqpoint{0.000000in}{-0.048611in}}{\pgfqpoint{0.000000in}{0.000000in}}{%
\pgfpathmoveto{\pgfqpoint{0.000000in}{0.000000in}}%
\pgfpathlineto{\pgfqpoint{0.000000in}{-0.048611in}}%
\pgfusepath{stroke,fill}%
}%
\begin{pgfscope}%
\pgfsys@transformshift{6.133356in}{5.300000in}%
\pgfsys@useobject{currentmarker}{}%
\end{pgfscope}%
\end{pgfscope}%
\begin{pgfscope}%
\definecolor{textcolor}{rgb}{0.000000,0.000000,0.000000}%
\pgfsetstrokecolor{textcolor}%
\pgfsetfillcolor{textcolor}%
\pgftext[x=6.133356in,y=5.202778in,,top]{\color{textcolor}{\sffamily\fontsize{10.000000}{12.000000}\selectfont\catcode`\^=\active\def^{\ifmmode\sp\else\^{}\fi}\catcode`\%=\active\def%{\%}0}}%
\end{pgfscope}%
\begin{pgfscope}%
\pgfsetbuttcap%
\pgfsetroundjoin%
\definecolor{currentfill}{rgb}{0.000000,0.000000,0.000000}%
\pgfsetfillcolor{currentfill}%
\pgfsetlinewidth{0.803000pt}%
\definecolor{currentstroke}{rgb}{0.000000,0.000000,0.000000}%
\pgfsetstrokecolor{currentstroke}%
\pgfsetdash{}{0pt}%
\pgfsys@defobject{currentmarker}{\pgfqpoint{0.000000in}{-0.048611in}}{\pgfqpoint{0.000000in}{0.000000in}}{%
\pgfpathmoveto{\pgfqpoint{0.000000in}{0.000000in}}%
\pgfpathlineto{\pgfqpoint{0.000000in}{-0.048611in}}%
\pgfusepath{stroke,fill}%
}%
\begin{pgfscope}%
\pgfsys@transformshift{6.653136in}{5.300000in}%
\pgfsys@useobject{currentmarker}{}%
\end{pgfscope}%
\end{pgfscope}%
\begin{pgfscope}%
\definecolor{textcolor}{rgb}{0.000000,0.000000,0.000000}%
\pgfsetstrokecolor{textcolor}%
\pgfsetfillcolor{textcolor}%
\pgftext[x=6.653136in,y=5.202778in,,top]{\color{textcolor}{\sffamily\fontsize{10.000000}{12.000000}\selectfont\catcode`\^=\active\def^{\ifmmode\sp\else\^{}\fi}\catcode`\%=\active\def%{\%}50}}%
\end{pgfscope}%
\begin{pgfscope}%
\pgfsetbuttcap%
\pgfsetroundjoin%
\definecolor{currentfill}{rgb}{0.000000,0.000000,0.000000}%
\pgfsetfillcolor{currentfill}%
\pgfsetlinewidth{0.803000pt}%
\definecolor{currentstroke}{rgb}{0.000000,0.000000,0.000000}%
\pgfsetstrokecolor{currentstroke}%
\pgfsetdash{}{0pt}%
\pgfsys@defobject{currentmarker}{\pgfqpoint{0.000000in}{-0.048611in}}{\pgfqpoint{0.000000in}{0.000000in}}{%
\pgfpathmoveto{\pgfqpoint{0.000000in}{0.000000in}}%
\pgfpathlineto{\pgfqpoint{0.000000in}{-0.048611in}}%
\pgfusepath{stroke,fill}%
}%
\begin{pgfscope}%
\pgfsys@transformshift{7.172917in}{5.300000in}%
\pgfsys@useobject{currentmarker}{}%
\end{pgfscope}%
\end{pgfscope}%
\begin{pgfscope}%
\definecolor{textcolor}{rgb}{0.000000,0.000000,0.000000}%
\pgfsetstrokecolor{textcolor}%
\pgfsetfillcolor{textcolor}%
\pgftext[x=7.172917in,y=5.202778in,,top]{\color{textcolor}{\sffamily\fontsize{10.000000}{12.000000}\selectfont\catcode`\^=\active\def^{\ifmmode\sp\else\^{}\fi}\catcode`\%=\active\def%{\%}100}}%
\end{pgfscope}%
\begin{pgfscope}%
\pgfsetbuttcap%
\pgfsetroundjoin%
\definecolor{currentfill}{rgb}{0.000000,0.000000,0.000000}%
\pgfsetfillcolor{currentfill}%
\pgfsetlinewidth{0.803000pt}%
\definecolor{currentstroke}{rgb}{0.000000,0.000000,0.000000}%
\pgfsetstrokecolor{currentstroke}%
\pgfsetdash{}{0pt}%
\pgfsys@defobject{currentmarker}{\pgfqpoint{0.000000in}{-0.048611in}}{\pgfqpoint{0.000000in}{0.000000in}}{%
\pgfpathmoveto{\pgfqpoint{0.000000in}{0.000000in}}%
\pgfpathlineto{\pgfqpoint{0.000000in}{-0.048611in}}%
\pgfusepath{stroke,fill}%
}%
\begin{pgfscope}%
\pgfsys@transformshift{7.692698in}{5.300000in}%
\pgfsys@useobject{currentmarker}{}%
\end{pgfscope}%
\end{pgfscope}%
\begin{pgfscope}%
\definecolor{textcolor}{rgb}{0.000000,0.000000,0.000000}%
\pgfsetstrokecolor{textcolor}%
\pgfsetfillcolor{textcolor}%
\pgftext[x=7.692698in,y=5.202778in,,top]{\color{textcolor}{\sffamily\fontsize{10.000000}{12.000000}\selectfont\catcode`\^=\active\def^{\ifmmode\sp\else\^{}\fi}\catcode`\%=\active\def%{\%}150}}%
\end{pgfscope}%
\begin{pgfscope}%
\pgfsetbuttcap%
\pgfsetroundjoin%
\definecolor{currentfill}{rgb}{0.000000,0.000000,0.000000}%
\pgfsetfillcolor{currentfill}%
\pgfsetlinewidth{0.803000pt}%
\definecolor{currentstroke}{rgb}{0.000000,0.000000,0.000000}%
\pgfsetstrokecolor{currentstroke}%
\pgfsetdash{}{0pt}%
\pgfsys@defobject{currentmarker}{\pgfqpoint{0.000000in}{-0.048611in}}{\pgfqpoint{0.000000in}{0.000000in}}{%
\pgfpathmoveto{\pgfqpoint{0.000000in}{0.000000in}}%
\pgfpathlineto{\pgfqpoint{0.000000in}{-0.048611in}}%
\pgfusepath{stroke,fill}%
}%
\begin{pgfscope}%
\pgfsys@transformshift{8.212479in}{5.300000in}%
\pgfsys@useobject{currentmarker}{}%
\end{pgfscope}%
\end{pgfscope}%
\begin{pgfscope}%
\definecolor{textcolor}{rgb}{0.000000,0.000000,0.000000}%
\pgfsetstrokecolor{textcolor}%
\pgfsetfillcolor{textcolor}%
\pgftext[x=8.212479in,y=5.202778in,,top]{\color{textcolor}{\sffamily\fontsize{10.000000}{12.000000}\selectfont\catcode`\^=\active\def^{\ifmmode\sp\else\^{}\fi}\catcode`\%=\active\def%{\%}200}}%
\end{pgfscope}%
\begin{pgfscope}%
\pgfsetbuttcap%
\pgfsetroundjoin%
\definecolor{currentfill}{rgb}{0.000000,0.000000,0.000000}%
\pgfsetfillcolor{currentfill}%
\pgfsetlinewidth{0.803000pt}%
\definecolor{currentstroke}{rgb}{0.000000,0.000000,0.000000}%
\pgfsetstrokecolor{currentstroke}%
\pgfsetdash{}{0pt}%
\pgfsys@defobject{currentmarker}{\pgfqpoint{0.000000in}{-0.048611in}}{\pgfqpoint{0.000000in}{0.000000in}}{%
\pgfpathmoveto{\pgfqpoint{0.000000in}{0.000000in}}%
\pgfpathlineto{\pgfqpoint{0.000000in}{-0.048611in}}%
\pgfusepath{stroke,fill}%
}%
\begin{pgfscope}%
\pgfsys@transformshift{8.732259in}{5.300000in}%
\pgfsys@useobject{currentmarker}{}%
\end{pgfscope}%
\end{pgfscope}%
\begin{pgfscope}%
\definecolor{textcolor}{rgb}{0.000000,0.000000,0.000000}%
\pgfsetstrokecolor{textcolor}%
\pgfsetfillcolor{textcolor}%
\pgftext[x=8.732259in,y=5.202778in,,top]{\color{textcolor}{\sffamily\fontsize{10.000000}{12.000000}\selectfont\catcode`\^=\active\def^{\ifmmode\sp\else\^{}\fi}\catcode`\%=\active\def%{\%}250}}%
\end{pgfscope}%
\begin{pgfscope}%
\pgfsetbuttcap%
\pgfsetroundjoin%
\definecolor{currentfill}{rgb}{0.000000,0.000000,0.000000}%
\pgfsetfillcolor{currentfill}%
\pgfsetlinewidth{0.803000pt}%
\definecolor{currentstroke}{rgb}{0.000000,0.000000,0.000000}%
\pgfsetstrokecolor{currentstroke}%
\pgfsetdash{}{0pt}%
\pgfsys@defobject{currentmarker}{\pgfqpoint{0.000000in}{-0.048611in}}{\pgfqpoint{0.000000in}{0.000000in}}{%
\pgfpathmoveto{\pgfqpoint{0.000000in}{0.000000in}}%
\pgfpathlineto{\pgfqpoint{0.000000in}{-0.048611in}}%
\pgfusepath{stroke,fill}%
}%
\begin{pgfscope}%
\pgfsys@transformshift{9.252040in}{5.300000in}%
\pgfsys@useobject{currentmarker}{}%
\end{pgfscope}%
\end{pgfscope}%
\begin{pgfscope}%
\definecolor{textcolor}{rgb}{0.000000,0.000000,0.000000}%
\pgfsetstrokecolor{textcolor}%
\pgfsetfillcolor{textcolor}%
\pgftext[x=9.252040in,y=5.202778in,,top]{\color{textcolor}{\sffamily\fontsize{10.000000}{12.000000}\selectfont\catcode`\^=\active\def^{\ifmmode\sp\else\^{}\fi}\catcode`\%=\active\def%{\%}300}}%
\end{pgfscope}%
\begin{pgfscope}%
\definecolor{textcolor}{rgb}{0.000000,0.000000,0.000000}%
\pgfsetstrokecolor{textcolor}%
\pgfsetfillcolor{textcolor}%
\pgftext[x=7.687500in,y=5.012809in,,top]{\color{textcolor}{\sffamily\fontsize{10.000000}{12.000000}\selectfont\catcode`\^=\active\def^{\ifmmode\sp\else\^{}\fi}\catcode`\%=\active\def%{\%}Iterations}}%
\end{pgfscope}%
\begin{pgfscope}%
\pgfsetbuttcap%
\pgfsetroundjoin%
\definecolor{currentfill}{rgb}{0.000000,0.000000,0.000000}%
\pgfsetfillcolor{currentfill}%
\pgfsetlinewidth{0.803000pt}%
\definecolor{currentstroke}{rgb}{0.000000,0.000000,0.000000}%
\pgfsetstrokecolor{currentstroke}%
\pgfsetdash{}{0pt}%
\pgfsys@defobject{currentmarker}{\pgfqpoint{-0.048611in}{0.000000in}}{\pgfqpoint{-0.000000in}{0.000000in}}{%
\pgfpathmoveto{\pgfqpoint{-0.000000in}{0.000000in}}%
\pgfpathlineto{\pgfqpoint{-0.048611in}{0.000000in}}%
\pgfusepath{stroke,fill}%
}%
\begin{pgfscope}%
\pgfsys@transformshift{5.977941in}{5.806344in}%
\pgfsys@useobject{currentmarker}{}%
\end{pgfscope}%
\end{pgfscope}%
\begin{pgfscope}%
\pgfsetbuttcap%
\pgfsetroundjoin%
\definecolor{currentfill}{rgb}{0.000000,0.000000,0.000000}%
\pgfsetfillcolor{currentfill}%
\pgfsetlinewidth{0.803000pt}%
\definecolor{currentstroke}{rgb}{0.000000,0.000000,0.000000}%
\pgfsetstrokecolor{currentstroke}%
\pgfsetdash{}{0pt}%
\pgfsys@defobject{currentmarker}{\pgfqpoint{-0.048611in}{0.000000in}}{\pgfqpoint{-0.000000in}{0.000000in}}{%
\pgfpathmoveto{\pgfqpoint{-0.000000in}{0.000000in}}%
\pgfpathlineto{\pgfqpoint{-0.048611in}{0.000000in}}%
\pgfusepath{stroke,fill}%
}%
\begin{pgfscope}%
\pgfsys@transformshift{5.977941in}{6.687896in}%
\pgfsys@useobject{currentmarker}{}%
\end{pgfscope}%
\end{pgfscope}%
\begin{pgfscope}%
\pgfsetbuttcap%
\pgfsetroundjoin%
\definecolor{currentfill}{rgb}{0.000000,0.000000,0.000000}%
\pgfsetfillcolor{currentfill}%
\pgfsetlinewidth{0.803000pt}%
\definecolor{currentstroke}{rgb}{0.000000,0.000000,0.000000}%
\pgfsetstrokecolor{currentstroke}%
\pgfsetdash{}{0pt}%
\pgfsys@defobject{currentmarker}{\pgfqpoint{-0.048611in}{0.000000in}}{\pgfqpoint{-0.000000in}{0.000000in}}{%
\pgfpathmoveto{\pgfqpoint{-0.000000in}{0.000000in}}%
\pgfpathlineto{\pgfqpoint{-0.048611in}{0.000000in}}%
\pgfusepath{stroke,fill}%
}%
\begin{pgfscope}%
\pgfsys@transformshift{5.977941in}{7.569448in}%
\pgfsys@useobject{currentmarker}{}%
\end{pgfscope}%
\end{pgfscope}%
\begin{pgfscope}%
\pgfsetbuttcap%
\pgfsetroundjoin%
\definecolor{currentfill}{rgb}{0.000000,0.000000,0.000000}%
\pgfsetfillcolor{currentfill}%
\pgfsetlinewidth{0.803000pt}%
\definecolor{currentstroke}{rgb}{0.000000,0.000000,0.000000}%
\pgfsetstrokecolor{currentstroke}%
\pgfsetdash{}{0pt}%
\pgfsys@defobject{currentmarker}{\pgfqpoint{-0.048611in}{0.000000in}}{\pgfqpoint{-0.000000in}{0.000000in}}{%
\pgfpathmoveto{\pgfqpoint{-0.000000in}{0.000000in}}%
\pgfpathlineto{\pgfqpoint{-0.048611in}{0.000000in}}%
\pgfusepath{stroke,fill}%
}%
\begin{pgfscope}%
\pgfsys@transformshift{5.977941in}{8.451000in}%
\pgfsys@useobject{currentmarker}{}%
\end{pgfscope}%
\end{pgfscope}%
\begin{pgfscope}%
\pgfsetbuttcap%
\pgfsetroundjoin%
\definecolor{currentfill}{rgb}{0.000000,0.000000,0.000000}%
\pgfsetfillcolor{currentfill}%
\pgfsetlinewidth{0.602250pt}%
\definecolor{currentstroke}{rgb}{0.000000,0.000000,0.000000}%
\pgfsetstrokecolor{currentstroke}%
\pgfsetdash{}{0pt}%
\pgfsys@defobject{currentmarker}{\pgfqpoint{-0.027778in}{0.000000in}}{\pgfqpoint{-0.000000in}{0.000000in}}{%
\pgfpathmoveto{\pgfqpoint{-0.000000in}{0.000000in}}%
\pgfpathlineto{\pgfqpoint{-0.027778in}{0.000000in}}%
\pgfusepath{stroke,fill}%
}%
\begin{pgfscope}%
\pgfsys@transformshift{5.977941in}{5.345399in}%
\pgfsys@useobject{currentmarker}{}%
\end{pgfscope}%
\end{pgfscope}%
\begin{pgfscope}%
\pgfsetbuttcap%
\pgfsetroundjoin%
\definecolor{currentfill}{rgb}{0.000000,0.000000,0.000000}%
\pgfsetfillcolor{currentfill}%
\pgfsetlinewidth{0.602250pt}%
\definecolor{currentstroke}{rgb}{0.000000,0.000000,0.000000}%
\pgfsetstrokecolor{currentstroke}%
\pgfsetdash{}{0pt}%
\pgfsys@defobject{currentmarker}{\pgfqpoint{-0.027778in}{0.000000in}}{\pgfqpoint{-0.000000in}{0.000000in}}{%
\pgfpathmoveto{\pgfqpoint{-0.000000in}{0.000000in}}%
\pgfpathlineto{\pgfqpoint{-0.027778in}{0.000000in}}%
\pgfusepath{stroke,fill}%
}%
\begin{pgfscope}%
\pgfsys@transformshift{5.977941in}{5.455539in}%
\pgfsys@useobject{currentmarker}{}%
\end{pgfscope}%
\end{pgfscope}%
\begin{pgfscope}%
\pgfsetbuttcap%
\pgfsetroundjoin%
\definecolor{currentfill}{rgb}{0.000000,0.000000,0.000000}%
\pgfsetfillcolor{currentfill}%
\pgfsetlinewidth{0.602250pt}%
\definecolor{currentstroke}{rgb}{0.000000,0.000000,0.000000}%
\pgfsetstrokecolor{currentstroke}%
\pgfsetdash{}{0pt}%
\pgfsys@defobject{currentmarker}{\pgfqpoint{-0.027778in}{0.000000in}}{\pgfqpoint{-0.000000in}{0.000000in}}{%
\pgfpathmoveto{\pgfqpoint{-0.000000in}{0.000000in}}%
\pgfpathlineto{\pgfqpoint{-0.027778in}{0.000000in}}%
\pgfusepath{stroke,fill}%
}%
\begin{pgfscope}%
\pgfsys@transformshift{5.977941in}{5.540970in}%
\pgfsys@useobject{currentmarker}{}%
\end{pgfscope}%
\end{pgfscope}%
\begin{pgfscope}%
\pgfsetbuttcap%
\pgfsetroundjoin%
\definecolor{currentfill}{rgb}{0.000000,0.000000,0.000000}%
\pgfsetfillcolor{currentfill}%
\pgfsetlinewidth{0.602250pt}%
\definecolor{currentstroke}{rgb}{0.000000,0.000000,0.000000}%
\pgfsetstrokecolor{currentstroke}%
\pgfsetdash{}{0pt}%
\pgfsys@defobject{currentmarker}{\pgfqpoint{-0.027778in}{0.000000in}}{\pgfqpoint{-0.000000in}{0.000000in}}{%
\pgfpathmoveto{\pgfqpoint{-0.000000in}{0.000000in}}%
\pgfpathlineto{\pgfqpoint{-0.027778in}{0.000000in}}%
\pgfusepath{stroke,fill}%
}%
\begin{pgfscope}%
\pgfsys@transformshift{5.977941in}{5.610773in}%
\pgfsys@useobject{currentmarker}{}%
\end{pgfscope}%
\end{pgfscope}%
\begin{pgfscope}%
\pgfsetbuttcap%
\pgfsetroundjoin%
\definecolor{currentfill}{rgb}{0.000000,0.000000,0.000000}%
\pgfsetfillcolor{currentfill}%
\pgfsetlinewidth{0.602250pt}%
\definecolor{currentstroke}{rgb}{0.000000,0.000000,0.000000}%
\pgfsetstrokecolor{currentstroke}%
\pgfsetdash{}{0pt}%
\pgfsys@defobject{currentmarker}{\pgfqpoint{-0.027778in}{0.000000in}}{\pgfqpoint{-0.000000in}{0.000000in}}{%
\pgfpathmoveto{\pgfqpoint{-0.000000in}{0.000000in}}%
\pgfpathlineto{\pgfqpoint{-0.027778in}{0.000000in}}%
\pgfusepath{stroke,fill}%
}%
\begin{pgfscope}%
\pgfsys@transformshift{5.977941in}{5.669790in}%
\pgfsys@useobject{currentmarker}{}%
\end{pgfscope}%
\end{pgfscope}%
\begin{pgfscope}%
\pgfsetbuttcap%
\pgfsetroundjoin%
\definecolor{currentfill}{rgb}{0.000000,0.000000,0.000000}%
\pgfsetfillcolor{currentfill}%
\pgfsetlinewidth{0.602250pt}%
\definecolor{currentstroke}{rgb}{0.000000,0.000000,0.000000}%
\pgfsetstrokecolor{currentstroke}%
\pgfsetdash{}{0pt}%
\pgfsys@defobject{currentmarker}{\pgfqpoint{-0.027778in}{0.000000in}}{\pgfqpoint{-0.000000in}{0.000000in}}{%
\pgfpathmoveto{\pgfqpoint{-0.000000in}{0.000000in}}%
\pgfpathlineto{\pgfqpoint{-0.027778in}{0.000000in}}%
\pgfusepath{stroke,fill}%
}%
\begin{pgfscope}%
\pgfsys@transformshift{5.977941in}{5.720913in}%
\pgfsys@useobject{currentmarker}{}%
\end{pgfscope}%
\end{pgfscope}%
\begin{pgfscope}%
\pgfsetbuttcap%
\pgfsetroundjoin%
\definecolor{currentfill}{rgb}{0.000000,0.000000,0.000000}%
\pgfsetfillcolor{currentfill}%
\pgfsetlinewidth{0.602250pt}%
\definecolor{currentstroke}{rgb}{0.000000,0.000000,0.000000}%
\pgfsetstrokecolor{currentstroke}%
\pgfsetdash{}{0pt}%
\pgfsys@defobject{currentmarker}{\pgfqpoint{-0.027778in}{0.000000in}}{\pgfqpoint{-0.000000in}{0.000000in}}{%
\pgfpathmoveto{\pgfqpoint{-0.000000in}{0.000000in}}%
\pgfpathlineto{\pgfqpoint{-0.027778in}{0.000000in}}%
\pgfusepath{stroke,fill}%
}%
\begin{pgfscope}%
\pgfsys@transformshift{5.977941in}{5.766006in}%
\pgfsys@useobject{currentmarker}{}%
\end{pgfscope}%
\end{pgfscope}%
\begin{pgfscope}%
\pgfsetbuttcap%
\pgfsetroundjoin%
\definecolor{currentfill}{rgb}{0.000000,0.000000,0.000000}%
\pgfsetfillcolor{currentfill}%
\pgfsetlinewidth{0.602250pt}%
\definecolor{currentstroke}{rgb}{0.000000,0.000000,0.000000}%
\pgfsetstrokecolor{currentstroke}%
\pgfsetdash{}{0pt}%
\pgfsys@defobject{currentmarker}{\pgfqpoint{-0.027778in}{0.000000in}}{\pgfqpoint{-0.000000in}{0.000000in}}{%
\pgfpathmoveto{\pgfqpoint{-0.000000in}{0.000000in}}%
\pgfpathlineto{\pgfqpoint{-0.027778in}{0.000000in}}%
\pgfusepath{stroke,fill}%
}%
\begin{pgfscope}%
\pgfsys@transformshift{5.977941in}{6.071717in}%
\pgfsys@useobject{currentmarker}{}%
\end{pgfscope}%
\end{pgfscope}%
\begin{pgfscope}%
\pgfsetbuttcap%
\pgfsetroundjoin%
\definecolor{currentfill}{rgb}{0.000000,0.000000,0.000000}%
\pgfsetfillcolor{currentfill}%
\pgfsetlinewidth{0.602250pt}%
\definecolor{currentstroke}{rgb}{0.000000,0.000000,0.000000}%
\pgfsetstrokecolor{currentstroke}%
\pgfsetdash{}{0pt}%
\pgfsys@defobject{currentmarker}{\pgfqpoint{-0.027778in}{0.000000in}}{\pgfqpoint{-0.000000in}{0.000000in}}{%
\pgfpathmoveto{\pgfqpoint{-0.000000in}{0.000000in}}%
\pgfpathlineto{\pgfqpoint{-0.027778in}{0.000000in}}%
\pgfusepath{stroke,fill}%
}%
\begin{pgfscope}%
\pgfsys@transformshift{5.977941in}{6.226951in}%
\pgfsys@useobject{currentmarker}{}%
\end{pgfscope}%
\end{pgfscope}%
\begin{pgfscope}%
\pgfsetbuttcap%
\pgfsetroundjoin%
\definecolor{currentfill}{rgb}{0.000000,0.000000,0.000000}%
\pgfsetfillcolor{currentfill}%
\pgfsetlinewidth{0.602250pt}%
\definecolor{currentstroke}{rgb}{0.000000,0.000000,0.000000}%
\pgfsetstrokecolor{currentstroke}%
\pgfsetdash{}{0pt}%
\pgfsys@defobject{currentmarker}{\pgfqpoint{-0.027778in}{0.000000in}}{\pgfqpoint{-0.000000in}{0.000000in}}{%
\pgfpathmoveto{\pgfqpoint{-0.000000in}{0.000000in}}%
\pgfpathlineto{\pgfqpoint{-0.027778in}{0.000000in}}%
\pgfusepath{stroke,fill}%
}%
\begin{pgfscope}%
\pgfsys@transformshift{5.977941in}{6.337091in}%
\pgfsys@useobject{currentmarker}{}%
\end{pgfscope}%
\end{pgfscope}%
\begin{pgfscope}%
\pgfsetbuttcap%
\pgfsetroundjoin%
\definecolor{currentfill}{rgb}{0.000000,0.000000,0.000000}%
\pgfsetfillcolor{currentfill}%
\pgfsetlinewidth{0.602250pt}%
\definecolor{currentstroke}{rgb}{0.000000,0.000000,0.000000}%
\pgfsetstrokecolor{currentstroke}%
\pgfsetdash{}{0pt}%
\pgfsys@defobject{currentmarker}{\pgfqpoint{-0.027778in}{0.000000in}}{\pgfqpoint{-0.000000in}{0.000000in}}{%
\pgfpathmoveto{\pgfqpoint{-0.000000in}{0.000000in}}%
\pgfpathlineto{\pgfqpoint{-0.027778in}{0.000000in}}%
\pgfusepath{stroke,fill}%
}%
\begin{pgfscope}%
\pgfsys@transformshift{5.977941in}{6.422522in}%
\pgfsys@useobject{currentmarker}{}%
\end{pgfscope}%
\end{pgfscope}%
\begin{pgfscope}%
\pgfsetbuttcap%
\pgfsetroundjoin%
\definecolor{currentfill}{rgb}{0.000000,0.000000,0.000000}%
\pgfsetfillcolor{currentfill}%
\pgfsetlinewidth{0.602250pt}%
\definecolor{currentstroke}{rgb}{0.000000,0.000000,0.000000}%
\pgfsetstrokecolor{currentstroke}%
\pgfsetdash{}{0pt}%
\pgfsys@defobject{currentmarker}{\pgfqpoint{-0.027778in}{0.000000in}}{\pgfqpoint{-0.000000in}{0.000000in}}{%
\pgfpathmoveto{\pgfqpoint{-0.000000in}{0.000000in}}%
\pgfpathlineto{\pgfqpoint{-0.027778in}{0.000000in}}%
\pgfusepath{stroke,fill}%
}%
\begin{pgfscope}%
\pgfsys@transformshift{5.977941in}{6.492324in}%
\pgfsys@useobject{currentmarker}{}%
\end{pgfscope}%
\end{pgfscope}%
\begin{pgfscope}%
\pgfsetbuttcap%
\pgfsetroundjoin%
\definecolor{currentfill}{rgb}{0.000000,0.000000,0.000000}%
\pgfsetfillcolor{currentfill}%
\pgfsetlinewidth{0.602250pt}%
\definecolor{currentstroke}{rgb}{0.000000,0.000000,0.000000}%
\pgfsetstrokecolor{currentstroke}%
\pgfsetdash{}{0pt}%
\pgfsys@defobject{currentmarker}{\pgfqpoint{-0.027778in}{0.000000in}}{\pgfqpoint{-0.000000in}{0.000000in}}{%
\pgfpathmoveto{\pgfqpoint{-0.000000in}{0.000000in}}%
\pgfpathlineto{\pgfqpoint{-0.027778in}{0.000000in}}%
\pgfusepath{stroke,fill}%
}%
\begin{pgfscope}%
\pgfsys@transformshift{5.977941in}{6.551342in}%
\pgfsys@useobject{currentmarker}{}%
\end{pgfscope}%
\end{pgfscope}%
\begin{pgfscope}%
\pgfsetbuttcap%
\pgfsetroundjoin%
\definecolor{currentfill}{rgb}{0.000000,0.000000,0.000000}%
\pgfsetfillcolor{currentfill}%
\pgfsetlinewidth{0.602250pt}%
\definecolor{currentstroke}{rgb}{0.000000,0.000000,0.000000}%
\pgfsetstrokecolor{currentstroke}%
\pgfsetdash{}{0pt}%
\pgfsys@defobject{currentmarker}{\pgfqpoint{-0.027778in}{0.000000in}}{\pgfqpoint{-0.000000in}{0.000000in}}{%
\pgfpathmoveto{\pgfqpoint{-0.000000in}{0.000000in}}%
\pgfpathlineto{\pgfqpoint{-0.027778in}{0.000000in}}%
\pgfusepath{stroke,fill}%
}%
\begin{pgfscope}%
\pgfsys@transformshift{5.977941in}{6.602464in}%
\pgfsys@useobject{currentmarker}{}%
\end{pgfscope}%
\end{pgfscope}%
\begin{pgfscope}%
\pgfsetbuttcap%
\pgfsetroundjoin%
\definecolor{currentfill}{rgb}{0.000000,0.000000,0.000000}%
\pgfsetfillcolor{currentfill}%
\pgfsetlinewidth{0.602250pt}%
\definecolor{currentstroke}{rgb}{0.000000,0.000000,0.000000}%
\pgfsetstrokecolor{currentstroke}%
\pgfsetdash{}{0pt}%
\pgfsys@defobject{currentmarker}{\pgfqpoint{-0.027778in}{0.000000in}}{\pgfqpoint{-0.000000in}{0.000000in}}{%
\pgfpathmoveto{\pgfqpoint{-0.000000in}{0.000000in}}%
\pgfpathlineto{\pgfqpoint{-0.027778in}{0.000000in}}%
\pgfusepath{stroke,fill}%
}%
\begin{pgfscope}%
\pgfsys@transformshift{5.977941in}{6.647558in}%
\pgfsys@useobject{currentmarker}{}%
\end{pgfscope}%
\end{pgfscope}%
\begin{pgfscope}%
\pgfsetbuttcap%
\pgfsetroundjoin%
\definecolor{currentfill}{rgb}{0.000000,0.000000,0.000000}%
\pgfsetfillcolor{currentfill}%
\pgfsetlinewidth{0.602250pt}%
\definecolor{currentstroke}{rgb}{0.000000,0.000000,0.000000}%
\pgfsetstrokecolor{currentstroke}%
\pgfsetdash{}{0pt}%
\pgfsys@defobject{currentmarker}{\pgfqpoint{-0.027778in}{0.000000in}}{\pgfqpoint{-0.000000in}{0.000000in}}{%
\pgfpathmoveto{\pgfqpoint{-0.000000in}{0.000000in}}%
\pgfpathlineto{\pgfqpoint{-0.027778in}{0.000000in}}%
\pgfusepath{stroke,fill}%
}%
\begin{pgfscope}%
\pgfsys@transformshift{5.977941in}{6.953269in}%
\pgfsys@useobject{currentmarker}{}%
\end{pgfscope}%
\end{pgfscope}%
\begin{pgfscope}%
\pgfsetbuttcap%
\pgfsetroundjoin%
\definecolor{currentfill}{rgb}{0.000000,0.000000,0.000000}%
\pgfsetfillcolor{currentfill}%
\pgfsetlinewidth{0.602250pt}%
\definecolor{currentstroke}{rgb}{0.000000,0.000000,0.000000}%
\pgfsetstrokecolor{currentstroke}%
\pgfsetdash{}{0pt}%
\pgfsys@defobject{currentmarker}{\pgfqpoint{-0.027778in}{0.000000in}}{\pgfqpoint{-0.000000in}{0.000000in}}{%
\pgfpathmoveto{\pgfqpoint{-0.000000in}{0.000000in}}%
\pgfpathlineto{\pgfqpoint{-0.027778in}{0.000000in}}%
\pgfusepath{stroke,fill}%
}%
\begin{pgfscope}%
\pgfsys@transformshift{5.977941in}{7.108503in}%
\pgfsys@useobject{currentmarker}{}%
\end{pgfscope}%
\end{pgfscope}%
\begin{pgfscope}%
\pgfsetbuttcap%
\pgfsetroundjoin%
\definecolor{currentfill}{rgb}{0.000000,0.000000,0.000000}%
\pgfsetfillcolor{currentfill}%
\pgfsetlinewidth{0.602250pt}%
\definecolor{currentstroke}{rgb}{0.000000,0.000000,0.000000}%
\pgfsetstrokecolor{currentstroke}%
\pgfsetdash{}{0pt}%
\pgfsys@defobject{currentmarker}{\pgfqpoint{-0.027778in}{0.000000in}}{\pgfqpoint{-0.000000in}{0.000000in}}{%
\pgfpathmoveto{\pgfqpoint{-0.000000in}{0.000000in}}%
\pgfpathlineto{\pgfqpoint{-0.027778in}{0.000000in}}%
\pgfusepath{stroke,fill}%
}%
\begin{pgfscope}%
\pgfsys@transformshift{5.977941in}{7.218643in}%
\pgfsys@useobject{currentmarker}{}%
\end{pgfscope}%
\end{pgfscope}%
\begin{pgfscope}%
\pgfsetbuttcap%
\pgfsetroundjoin%
\definecolor{currentfill}{rgb}{0.000000,0.000000,0.000000}%
\pgfsetfillcolor{currentfill}%
\pgfsetlinewidth{0.602250pt}%
\definecolor{currentstroke}{rgb}{0.000000,0.000000,0.000000}%
\pgfsetstrokecolor{currentstroke}%
\pgfsetdash{}{0pt}%
\pgfsys@defobject{currentmarker}{\pgfqpoint{-0.027778in}{0.000000in}}{\pgfqpoint{-0.000000in}{0.000000in}}{%
\pgfpathmoveto{\pgfqpoint{-0.000000in}{0.000000in}}%
\pgfpathlineto{\pgfqpoint{-0.027778in}{0.000000in}}%
\pgfusepath{stroke,fill}%
}%
\begin{pgfscope}%
\pgfsys@transformshift{5.977941in}{7.304074in}%
\pgfsys@useobject{currentmarker}{}%
\end{pgfscope}%
\end{pgfscope}%
\begin{pgfscope}%
\pgfsetbuttcap%
\pgfsetroundjoin%
\definecolor{currentfill}{rgb}{0.000000,0.000000,0.000000}%
\pgfsetfillcolor{currentfill}%
\pgfsetlinewidth{0.602250pt}%
\definecolor{currentstroke}{rgb}{0.000000,0.000000,0.000000}%
\pgfsetstrokecolor{currentstroke}%
\pgfsetdash{}{0pt}%
\pgfsys@defobject{currentmarker}{\pgfqpoint{-0.027778in}{0.000000in}}{\pgfqpoint{-0.000000in}{0.000000in}}{%
\pgfpathmoveto{\pgfqpoint{-0.000000in}{0.000000in}}%
\pgfpathlineto{\pgfqpoint{-0.027778in}{0.000000in}}%
\pgfusepath{stroke,fill}%
}%
\begin{pgfscope}%
\pgfsys@transformshift{5.977941in}{7.373876in}%
\pgfsys@useobject{currentmarker}{}%
\end{pgfscope}%
\end{pgfscope}%
\begin{pgfscope}%
\pgfsetbuttcap%
\pgfsetroundjoin%
\definecolor{currentfill}{rgb}{0.000000,0.000000,0.000000}%
\pgfsetfillcolor{currentfill}%
\pgfsetlinewidth{0.602250pt}%
\definecolor{currentstroke}{rgb}{0.000000,0.000000,0.000000}%
\pgfsetstrokecolor{currentstroke}%
\pgfsetdash{}{0pt}%
\pgfsys@defobject{currentmarker}{\pgfqpoint{-0.027778in}{0.000000in}}{\pgfqpoint{-0.000000in}{0.000000in}}{%
\pgfpathmoveto{\pgfqpoint{-0.000000in}{0.000000in}}%
\pgfpathlineto{\pgfqpoint{-0.027778in}{0.000000in}}%
\pgfusepath{stroke,fill}%
}%
\begin{pgfscope}%
\pgfsys@transformshift{5.977941in}{7.432894in}%
\pgfsys@useobject{currentmarker}{}%
\end{pgfscope}%
\end{pgfscope}%
\begin{pgfscope}%
\pgfsetbuttcap%
\pgfsetroundjoin%
\definecolor{currentfill}{rgb}{0.000000,0.000000,0.000000}%
\pgfsetfillcolor{currentfill}%
\pgfsetlinewidth{0.602250pt}%
\definecolor{currentstroke}{rgb}{0.000000,0.000000,0.000000}%
\pgfsetstrokecolor{currentstroke}%
\pgfsetdash{}{0pt}%
\pgfsys@defobject{currentmarker}{\pgfqpoint{-0.027778in}{0.000000in}}{\pgfqpoint{-0.000000in}{0.000000in}}{%
\pgfpathmoveto{\pgfqpoint{-0.000000in}{0.000000in}}%
\pgfpathlineto{\pgfqpoint{-0.027778in}{0.000000in}}%
\pgfusepath{stroke,fill}%
}%
\begin{pgfscope}%
\pgfsys@transformshift{5.977941in}{7.484016in}%
\pgfsys@useobject{currentmarker}{}%
\end{pgfscope}%
\end{pgfscope}%
\begin{pgfscope}%
\pgfsetbuttcap%
\pgfsetroundjoin%
\definecolor{currentfill}{rgb}{0.000000,0.000000,0.000000}%
\pgfsetfillcolor{currentfill}%
\pgfsetlinewidth{0.602250pt}%
\definecolor{currentstroke}{rgb}{0.000000,0.000000,0.000000}%
\pgfsetstrokecolor{currentstroke}%
\pgfsetdash{}{0pt}%
\pgfsys@defobject{currentmarker}{\pgfqpoint{-0.027778in}{0.000000in}}{\pgfqpoint{-0.000000in}{0.000000in}}{%
\pgfpathmoveto{\pgfqpoint{-0.000000in}{0.000000in}}%
\pgfpathlineto{\pgfqpoint{-0.027778in}{0.000000in}}%
\pgfusepath{stroke,fill}%
}%
\begin{pgfscope}%
\pgfsys@transformshift{5.977941in}{7.529110in}%
\pgfsys@useobject{currentmarker}{}%
\end{pgfscope}%
\end{pgfscope}%
\begin{pgfscope}%
\pgfsetbuttcap%
\pgfsetroundjoin%
\definecolor{currentfill}{rgb}{0.000000,0.000000,0.000000}%
\pgfsetfillcolor{currentfill}%
\pgfsetlinewidth{0.602250pt}%
\definecolor{currentstroke}{rgb}{0.000000,0.000000,0.000000}%
\pgfsetstrokecolor{currentstroke}%
\pgfsetdash{}{0pt}%
\pgfsys@defobject{currentmarker}{\pgfqpoint{-0.027778in}{0.000000in}}{\pgfqpoint{-0.000000in}{0.000000in}}{%
\pgfpathmoveto{\pgfqpoint{-0.000000in}{0.000000in}}%
\pgfpathlineto{\pgfqpoint{-0.027778in}{0.000000in}}%
\pgfusepath{stroke,fill}%
}%
\begin{pgfscope}%
\pgfsys@transformshift{5.977941in}{7.834821in}%
\pgfsys@useobject{currentmarker}{}%
\end{pgfscope}%
\end{pgfscope}%
\begin{pgfscope}%
\pgfsetbuttcap%
\pgfsetroundjoin%
\definecolor{currentfill}{rgb}{0.000000,0.000000,0.000000}%
\pgfsetfillcolor{currentfill}%
\pgfsetlinewidth{0.602250pt}%
\definecolor{currentstroke}{rgb}{0.000000,0.000000,0.000000}%
\pgfsetstrokecolor{currentstroke}%
\pgfsetdash{}{0pt}%
\pgfsys@defobject{currentmarker}{\pgfqpoint{-0.027778in}{0.000000in}}{\pgfqpoint{-0.000000in}{0.000000in}}{%
\pgfpathmoveto{\pgfqpoint{-0.000000in}{0.000000in}}%
\pgfpathlineto{\pgfqpoint{-0.027778in}{0.000000in}}%
\pgfusepath{stroke,fill}%
}%
\begin{pgfscope}%
\pgfsys@transformshift{5.977941in}{7.990055in}%
\pgfsys@useobject{currentmarker}{}%
\end{pgfscope}%
\end{pgfscope}%
\begin{pgfscope}%
\pgfsetbuttcap%
\pgfsetroundjoin%
\definecolor{currentfill}{rgb}{0.000000,0.000000,0.000000}%
\pgfsetfillcolor{currentfill}%
\pgfsetlinewidth{0.602250pt}%
\definecolor{currentstroke}{rgb}{0.000000,0.000000,0.000000}%
\pgfsetstrokecolor{currentstroke}%
\pgfsetdash{}{0pt}%
\pgfsys@defobject{currentmarker}{\pgfqpoint{-0.027778in}{0.000000in}}{\pgfqpoint{-0.000000in}{0.000000in}}{%
\pgfpathmoveto{\pgfqpoint{-0.000000in}{0.000000in}}%
\pgfpathlineto{\pgfqpoint{-0.027778in}{0.000000in}}%
\pgfusepath{stroke,fill}%
}%
\begin{pgfscope}%
\pgfsys@transformshift{5.977941in}{8.100195in}%
\pgfsys@useobject{currentmarker}{}%
\end{pgfscope}%
\end{pgfscope}%
\begin{pgfscope}%
\pgfsetbuttcap%
\pgfsetroundjoin%
\definecolor{currentfill}{rgb}{0.000000,0.000000,0.000000}%
\pgfsetfillcolor{currentfill}%
\pgfsetlinewidth{0.602250pt}%
\definecolor{currentstroke}{rgb}{0.000000,0.000000,0.000000}%
\pgfsetstrokecolor{currentstroke}%
\pgfsetdash{}{0pt}%
\pgfsys@defobject{currentmarker}{\pgfqpoint{-0.027778in}{0.000000in}}{\pgfqpoint{-0.000000in}{0.000000in}}{%
\pgfpathmoveto{\pgfqpoint{-0.000000in}{0.000000in}}%
\pgfpathlineto{\pgfqpoint{-0.027778in}{0.000000in}}%
\pgfusepath{stroke,fill}%
}%
\begin{pgfscope}%
\pgfsys@transformshift{5.977941in}{8.185626in}%
\pgfsys@useobject{currentmarker}{}%
\end{pgfscope}%
\end{pgfscope}%
\begin{pgfscope}%
\pgfsetbuttcap%
\pgfsetroundjoin%
\definecolor{currentfill}{rgb}{0.000000,0.000000,0.000000}%
\pgfsetfillcolor{currentfill}%
\pgfsetlinewidth{0.602250pt}%
\definecolor{currentstroke}{rgb}{0.000000,0.000000,0.000000}%
\pgfsetstrokecolor{currentstroke}%
\pgfsetdash{}{0pt}%
\pgfsys@defobject{currentmarker}{\pgfqpoint{-0.027778in}{0.000000in}}{\pgfqpoint{-0.000000in}{0.000000in}}{%
\pgfpathmoveto{\pgfqpoint{-0.000000in}{0.000000in}}%
\pgfpathlineto{\pgfqpoint{-0.027778in}{0.000000in}}%
\pgfusepath{stroke,fill}%
}%
\begin{pgfscope}%
\pgfsys@transformshift{5.977941in}{8.255428in}%
\pgfsys@useobject{currentmarker}{}%
\end{pgfscope}%
\end{pgfscope}%
\begin{pgfscope}%
\pgfsetbuttcap%
\pgfsetroundjoin%
\definecolor{currentfill}{rgb}{0.000000,0.000000,0.000000}%
\pgfsetfillcolor{currentfill}%
\pgfsetlinewidth{0.602250pt}%
\definecolor{currentstroke}{rgb}{0.000000,0.000000,0.000000}%
\pgfsetstrokecolor{currentstroke}%
\pgfsetdash{}{0pt}%
\pgfsys@defobject{currentmarker}{\pgfqpoint{-0.027778in}{0.000000in}}{\pgfqpoint{-0.000000in}{0.000000in}}{%
\pgfpathmoveto{\pgfqpoint{-0.000000in}{0.000000in}}%
\pgfpathlineto{\pgfqpoint{-0.027778in}{0.000000in}}%
\pgfusepath{stroke,fill}%
}%
\begin{pgfscope}%
\pgfsys@transformshift{5.977941in}{8.314446in}%
\pgfsys@useobject{currentmarker}{}%
\end{pgfscope}%
\end{pgfscope}%
\begin{pgfscope}%
\pgfsetbuttcap%
\pgfsetroundjoin%
\definecolor{currentfill}{rgb}{0.000000,0.000000,0.000000}%
\pgfsetfillcolor{currentfill}%
\pgfsetlinewidth{0.602250pt}%
\definecolor{currentstroke}{rgb}{0.000000,0.000000,0.000000}%
\pgfsetstrokecolor{currentstroke}%
\pgfsetdash{}{0pt}%
\pgfsys@defobject{currentmarker}{\pgfqpoint{-0.027778in}{0.000000in}}{\pgfqpoint{-0.000000in}{0.000000in}}{%
\pgfpathmoveto{\pgfqpoint{-0.000000in}{0.000000in}}%
\pgfpathlineto{\pgfqpoint{-0.027778in}{0.000000in}}%
\pgfusepath{stroke,fill}%
}%
\begin{pgfscope}%
\pgfsys@transformshift{5.977941in}{8.365568in}%
\pgfsys@useobject{currentmarker}{}%
\end{pgfscope}%
\end{pgfscope}%
\begin{pgfscope}%
\pgfsetbuttcap%
\pgfsetroundjoin%
\definecolor{currentfill}{rgb}{0.000000,0.000000,0.000000}%
\pgfsetfillcolor{currentfill}%
\pgfsetlinewidth{0.602250pt}%
\definecolor{currentstroke}{rgb}{0.000000,0.000000,0.000000}%
\pgfsetstrokecolor{currentstroke}%
\pgfsetdash{}{0pt}%
\pgfsys@defobject{currentmarker}{\pgfqpoint{-0.027778in}{0.000000in}}{\pgfqpoint{-0.000000in}{0.000000in}}{%
\pgfpathmoveto{\pgfqpoint{-0.000000in}{0.000000in}}%
\pgfpathlineto{\pgfqpoint{-0.027778in}{0.000000in}}%
\pgfusepath{stroke,fill}%
}%
\begin{pgfscope}%
\pgfsys@transformshift{5.977941in}{8.410662in}%
\pgfsys@useobject{currentmarker}{}%
\end{pgfscope}%
\end{pgfscope}%
\begin{pgfscope}%
\pgfsetbuttcap%
\pgfsetroundjoin%
\definecolor{currentfill}{rgb}{0.000000,0.000000,0.000000}%
\pgfsetfillcolor{currentfill}%
\pgfsetlinewidth{0.602250pt}%
\definecolor{currentstroke}{rgb}{0.000000,0.000000,0.000000}%
\pgfsetstrokecolor{currentstroke}%
\pgfsetdash{}{0pt}%
\pgfsys@defobject{currentmarker}{\pgfqpoint{-0.027778in}{0.000000in}}{\pgfqpoint{-0.000000in}{0.000000in}}{%
\pgfpathmoveto{\pgfqpoint{-0.000000in}{0.000000in}}%
\pgfpathlineto{\pgfqpoint{-0.027778in}{0.000000in}}%
\pgfusepath{stroke,fill}%
}%
\begin{pgfscope}%
\pgfsys@transformshift{5.977941in}{8.716373in}%
\pgfsys@useobject{currentmarker}{}%
\end{pgfscope}%
\end{pgfscope}%
\begin{pgfscope}%
\pgfpathrectangle{\pgfqpoint{5.977941in}{5.300000in}}{\pgfqpoint{3.419118in}{3.500000in}}%
\pgfusepath{clip}%
\pgfsetrectcap%
\pgfsetroundjoin%
\pgfsetlinewidth{1.505625pt}%
\definecolor{currentstroke}{rgb}{0.121569,0.466667,0.705882}%
\pgfsetstrokecolor{currentstroke}%
\pgfsetdash{}{0pt}%
\pgfpathmoveto{\pgfqpoint{6.133356in}{8.640909in}}%
\pgfpathlineto{\pgfqpoint{7.495181in}{8.176507in}}%
\pgfpathlineto{\pgfqpoint{9.241644in}{7.581242in}}%
\pgfpathlineto{\pgfqpoint{9.241644in}{7.581242in}}%
\pgfusepath{stroke}%
\end{pgfscope}%
\begin{pgfscope}%
\pgfpathrectangle{\pgfqpoint{5.977941in}{5.300000in}}{\pgfqpoint{3.419118in}{3.500000in}}%
\pgfusepath{clip}%
\pgfsetrectcap%
\pgfsetroundjoin%
\pgfsetlinewidth{1.505625pt}%
\definecolor{currentstroke}{rgb}{1.000000,0.498039,0.054902}%
\pgfsetstrokecolor{currentstroke}%
\pgfsetdash{}{0pt}%
\pgfpathmoveto{\pgfqpoint{6.133356in}{8.640909in}}%
\pgfpathlineto{\pgfqpoint{6.372455in}{8.559274in}}%
\pgfpathlineto{\pgfqpoint{6.382850in}{6.796347in}}%
\pgfpathlineto{\pgfqpoint{6.393246in}{6.682512in}}%
\pgfpathlineto{\pgfqpoint{6.403642in}{6.599572in}}%
\pgfpathlineto{\pgfqpoint{6.414037in}{6.535025in}}%
\pgfpathlineto{\pgfqpoint{6.424433in}{6.482181in}}%
\pgfpathlineto{\pgfqpoint{6.434828in}{6.437556in}}%
\pgfpathlineto{\pgfqpoint{6.445224in}{6.398951in}}%
\pgfpathlineto{\pgfqpoint{6.466015in}{6.333918in}}%
\pgfpathlineto{\pgfqpoint{6.486807in}{6.279555in}}%
\pgfpathlineto{\pgfqpoint{6.507598in}{6.232127in}}%
\pgfpathlineto{\pgfqpoint{6.528389in}{6.189586in}}%
\pgfpathlineto{\pgfqpoint{6.559576in}{6.132456in}}%
\pgfpathlineto{\pgfqpoint{6.590763in}{6.081403in}}%
\pgfpathlineto{\pgfqpoint{6.621950in}{6.035136in}}%
\pgfpathlineto{\pgfqpoint{6.632345in}{6.020642in}}%
\pgfpathlineto{\pgfqpoint{6.642741in}{6.308847in}}%
\pgfpathlineto{\pgfqpoint{6.673928in}{6.282861in}}%
\pgfpathlineto{\pgfqpoint{6.705114in}{6.261346in}}%
\pgfpathlineto{\pgfqpoint{6.736301in}{6.242396in}}%
\pgfpathlineto{\pgfqpoint{6.777884in}{6.219749in}}%
\pgfpathlineto{\pgfqpoint{6.829862in}{6.194373in}}%
\pgfpathlineto{\pgfqpoint{6.892235in}{6.166781in}}%
\pgfpathlineto{\pgfqpoint{6.902631in}{6.286731in}}%
\pgfpathlineto{\pgfqpoint{6.965005in}{6.245785in}}%
\pgfpathlineto{\pgfqpoint{7.006587in}{6.223381in}}%
\pgfpathlineto{\pgfqpoint{7.058565in}{6.198166in}}%
\pgfpathlineto{\pgfqpoint{7.120939in}{6.170632in}}%
\pgfpathlineto{\pgfqpoint{7.152126in}{6.157648in}}%
\pgfpathlineto{\pgfqpoint{7.162521in}{6.283495in}}%
\pgfpathlineto{\pgfqpoint{7.235291in}{6.237601in}}%
\pgfpathlineto{\pgfqpoint{7.276873in}{6.216111in}}%
\pgfpathlineto{\pgfqpoint{7.328851in}{6.191665in}}%
\pgfpathlineto{\pgfqpoint{7.391225in}{6.164727in}}%
\pgfpathlineto{\pgfqpoint{7.412016in}{6.156173in}}%
\pgfpathlineto{\pgfqpoint{7.422412in}{6.282909in}}%
\pgfpathlineto{\pgfqpoint{7.495181in}{6.237182in}}%
\pgfpathlineto{\pgfqpoint{7.536764in}{6.215751in}}%
\pgfpathlineto{\pgfqpoint{7.588742in}{6.191364in}}%
\pgfpathlineto{\pgfqpoint{7.651115in}{6.164479in}}%
\pgfpathlineto{\pgfqpoint{7.671907in}{6.155933in}}%
\pgfpathlineto{\pgfqpoint{7.682302in}{6.282826in}}%
\pgfpathlineto{\pgfqpoint{7.755071in}{6.237138in}}%
\pgfpathlineto{\pgfqpoint{7.796654in}{6.215726in}}%
\pgfpathlineto{\pgfqpoint{7.848632in}{6.191343in}}%
\pgfpathlineto{\pgfqpoint{7.911006in}{6.164459in}}%
\pgfpathlineto{\pgfqpoint{7.931797in}{6.155917in}}%
\pgfpathlineto{\pgfqpoint{7.942193in}{6.282825in}}%
\pgfpathlineto{\pgfqpoint{8.014962in}{6.237128in}}%
\pgfpathlineto{\pgfqpoint{8.056544in}{6.215709in}}%
\pgfpathlineto{\pgfqpoint{8.108522in}{6.191324in}}%
\pgfpathlineto{\pgfqpoint{8.170896in}{6.164441in}}%
\pgfpathlineto{\pgfqpoint{8.191687in}{6.155894in}}%
\pgfpathlineto{\pgfqpoint{8.202083in}{6.282978in}}%
\pgfpathlineto{\pgfqpoint{8.264457in}{6.243022in}}%
\pgfpathlineto{\pgfqpoint{8.306039in}{6.221048in}}%
\pgfpathlineto{\pgfqpoint{8.358017in}{6.196213in}}%
\pgfpathlineto{\pgfqpoint{8.420391in}{6.168968in}}%
\pgfpathlineto{\pgfqpoint{8.451578in}{6.156082in}}%
\pgfpathlineto{\pgfqpoint{8.461973in}{6.282853in}}%
\pgfpathlineto{\pgfqpoint{8.534743in}{6.237156in}}%
\pgfpathlineto{\pgfqpoint{8.576325in}{6.215732in}}%
\pgfpathlineto{\pgfqpoint{8.628303in}{6.191349in}}%
\pgfpathlineto{\pgfqpoint{8.690677in}{6.164466in}}%
\pgfpathlineto{\pgfqpoint{8.711468in}{6.155916in}}%
\pgfpathlineto{\pgfqpoint{8.721864in}{6.282739in}}%
\pgfpathlineto{\pgfqpoint{8.794633in}{6.237027in}}%
\pgfpathlineto{\pgfqpoint{8.836215in}{6.215603in}}%
\pgfpathlineto{\pgfqpoint{8.888193in}{6.191216in}}%
\pgfpathlineto{\pgfqpoint{8.950567in}{6.164329in}}%
\pgfpathlineto{\pgfqpoint{8.971358in}{6.155786in}}%
\pgfpathlineto{\pgfqpoint{8.981754in}{6.282917in}}%
\pgfpathlineto{\pgfqpoint{9.054523in}{6.237236in}}%
\pgfpathlineto{\pgfqpoint{9.096106in}{6.215821in}}%
\pgfpathlineto{\pgfqpoint{9.148084in}{6.191446in}}%
\pgfpathlineto{\pgfqpoint{9.210458in}{6.164567in}}%
\pgfpathlineto{\pgfqpoint{9.231249in}{6.156024in}}%
\pgfpathlineto{\pgfqpoint{9.241644in}{6.283029in}}%
\pgfpathlineto{\pgfqpoint{9.241644in}{6.283029in}}%
\pgfusepath{stroke}%
\end{pgfscope}%
\begin{pgfscope}%
\pgfpathrectangle{\pgfqpoint{5.977941in}{5.300000in}}{\pgfqpoint{3.419118in}{3.500000in}}%
\pgfusepath{clip}%
\pgfsetrectcap%
\pgfsetroundjoin%
\pgfsetlinewidth{1.505625pt}%
\definecolor{currentstroke}{rgb}{0.172549,0.627451,0.172549}%
\pgfsetstrokecolor{currentstroke}%
\pgfsetdash{}{0pt}%
\pgfpathmoveto{\pgfqpoint{6.133356in}{7.073895in}}%
\pgfpathlineto{\pgfqpoint{6.143751in}{7.030419in}}%
\pgfpathlineto{\pgfqpoint{6.154147in}{6.992677in}}%
\pgfpathlineto{\pgfqpoint{6.174938in}{6.932733in}}%
\pgfpathlineto{\pgfqpoint{6.206125in}{6.852582in}}%
\pgfpathlineto{\pgfqpoint{6.237312in}{6.778194in}}%
\pgfpathlineto{\pgfqpoint{6.268499in}{6.710360in}}%
\pgfpathlineto{\pgfqpoint{6.289290in}{6.669148in}}%
\pgfpathlineto{\pgfqpoint{6.310081in}{6.631331in}}%
\pgfpathlineto{\pgfqpoint{6.330872in}{6.596992in}}%
\pgfpathlineto{\pgfqpoint{6.351664in}{6.566117in}}%
\pgfpathlineto{\pgfqpoint{6.372455in}{6.538588in}}%
\pgfpathlineto{\pgfqpoint{6.393246in}{6.514189in}}%
\pgfpathlineto{\pgfqpoint{6.414037in}{6.492635in}}%
\pgfpathlineto{\pgfqpoint{6.434828in}{6.473599in}}%
\pgfpathlineto{\pgfqpoint{6.455620in}{6.456733in}}%
\pgfpathlineto{\pgfqpoint{6.486807in}{6.434781in}}%
\pgfpathlineto{\pgfqpoint{6.517993in}{6.415953in}}%
\pgfpathlineto{\pgfqpoint{6.559576in}{6.394300in}}%
\pgfpathlineto{\pgfqpoint{6.601158in}{6.375285in}}%
\pgfpathlineto{\pgfqpoint{6.663532in}{6.349698in}}%
\pgfpathlineto{\pgfqpoint{6.746697in}{6.318410in}}%
\pgfpathlineto{\pgfqpoint{6.881840in}{6.270291in}}%
\pgfpathlineto{\pgfqpoint{7.162521in}{6.173241in}}%
\pgfpathlineto{\pgfqpoint{9.241644in}{5.459091in}}%
\pgfpathlineto{\pgfqpoint{9.241644in}{5.459091in}}%
\pgfusepath{stroke}%
\end{pgfscope}%
\begin{pgfscope}%
\pgfpathrectangle{\pgfqpoint{5.977941in}{5.300000in}}{\pgfqpoint{3.419118in}{3.500000in}}%
\pgfusepath{clip}%
\pgfsetrectcap%
\pgfsetroundjoin%
\pgfsetlinewidth{1.505625pt}%
\definecolor{currentstroke}{rgb}{0.839216,0.152941,0.156863}%
\pgfsetstrokecolor{currentstroke}%
\pgfsetdash{}{0pt}%
\pgfpathmoveto{\pgfqpoint{6.133356in}{6.956891in}}%
\pgfpathlineto{\pgfqpoint{6.143751in}{6.901936in}}%
\pgfpathlineto{\pgfqpoint{6.154147in}{6.862045in}}%
\pgfpathlineto{\pgfqpoint{6.164542in}{6.827727in}}%
\pgfpathlineto{\pgfqpoint{6.185334in}{6.791769in}}%
\pgfpathlineto{\pgfqpoint{6.206125in}{6.760090in}}%
\pgfpathlineto{\pgfqpoint{6.237312in}{6.717580in}}%
\pgfpathlineto{\pgfqpoint{6.268499in}{6.679647in}}%
\pgfpathlineto{\pgfqpoint{6.299685in}{6.645454in}}%
\pgfpathlineto{\pgfqpoint{6.330872in}{6.614473in}}%
\pgfpathlineto{\pgfqpoint{6.362059in}{6.586311in}}%
\pgfpathlineto{\pgfqpoint{6.393246in}{6.560640in}}%
\pgfpathlineto{\pgfqpoint{6.434828in}{6.529801in}}%
\pgfpathlineto{\pgfqpoint{6.476411in}{6.502304in}}%
\pgfpathlineto{\pgfqpoint{6.517993in}{6.477645in}}%
\pgfpathlineto{\pgfqpoint{6.559576in}{6.455396in}}%
\pgfpathlineto{\pgfqpoint{6.611554in}{6.430418in}}%
\pgfpathlineto{\pgfqpoint{6.663532in}{6.408042in}}%
\pgfpathlineto{\pgfqpoint{6.725906in}{6.383952in}}%
\pgfpathlineto{\pgfqpoint{6.788279in}{6.362270in}}%
\pgfpathlineto{\pgfqpoint{6.861049in}{6.339359in}}%
\pgfpathlineto{\pgfqpoint{6.944214in}{6.315609in}}%
\pgfpathlineto{\pgfqpoint{7.037774in}{6.291250in}}%
\pgfpathlineto{\pgfqpoint{7.152126in}{6.263993in}}%
\pgfpathlineto{\pgfqpoint{7.287269in}{6.234347in}}%
\pgfpathlineto{\pgfqpoint{7.443203in}{6.202532in}}%
\pgfpathlineto{\pgfqpoint{7.640720in}{6.164631in}}%
\pgfpathlineto{\pgfqpoint{7.900610in}{6.117186in}}%
\pgfpathlineto{\pgfqpoint{8.274852in}{6.051278in}}%
\pgfpathlineto{\pgfqpoint{8.950567in}{5.934895in}}%
\pgfpathlineto{\pgfqpoint{9.241644in}{5.885143in}}%
\pgfpathlineto{\pgfqpoint{9.241644in}{5.885143in}}%
\pgfusepath{stroke}%
\end{pgfscope}%
\begin{pgfscope}%
\pgfsetrectcap%
\pgfsetmiterjoin%
\pgfsetlinewidth{0.803000pt}%
\definecolor{currentstroke}{rgb}{0.000000,0.000000,0.000000}%
\pgfsetstrokecolor{currentstroke}%
\pgfsetdash{}{0pt}%
\pgfpathmoveto{\pgfqpoint{5.977941in}{5.300000in}}%
\pgfpathlineto{\pgfqpoint{5.977941in}{8.800000in}}%
\pgfusepath{stroke}%
\end{pgfscope}%
\begin{pgfscope}%
\pgfsetrectcap%
\pgfsetmiterjoin%
\pgfsetlinewidth{0.803000pt}%
\definecolor{currentstroke}{rgb}{0.000000,0.000000,0.000000}%
\pgfsetstrokecolor{currentstroke}%
\pgfsetdash{}{0pt}%
\pgfpathmoveto{\pgfqpoint{9.397059in}{5.300000in}}%
\pgfpathlineto{\pgfqpoint{9.397059in}{8.800000in}}%
\pgfusepath{stroke}%
\end{pgfscope}%
\begin{pgfscope}%
\pgfsetrectcap%
\pgfsetmiterjoin%
\pgfsetlinewidth{0.803000pt}%
\definecolor{currentstroke}{rgb}{0.000000,0.000000,0.000000}%
\pgfsetstrokecolor{currentstroke}%
\pgfsetdash{}{0pt}%
\pgfpathmoveto{\pgfqpoint{5.977941in}{5.300000in}}%
\pgfpathlineto{\pgfqpoint{9.397059in}{5.300000in}}%
\pgfusepath{stroke}%
\end{pgfscope}%
\begin{pgfscope}%
\pgfsetrectcap%
\pgfsetmiterjoin%
\pgfsetlinewidth{0.803000pt}%
\definecolor{currentstroke}{rgb}{0.000000,0.000000,0.000000}%
\pgfsetstrokecolor{currentstroke}%
\pgfsetdash{}{0pt}%
\pgfpathmoveto{\pgfqpoint{5.977941in}{8.800000in}}%
\pgfpathlineto{\pgfqpoint{9.397059in}{8.800000in}}%
\pgfusepath{stroke}%
\end{pgfscope}%
\begin{pgfscope}%
\definecolor{textcolor}{rgb}{0.000000,0.000000,0.000000}%
\pgfsetstrokecolor{textcolor}%
\pgfsetfillcolor{textcolor}%
\pgftext[x=7.687500in,y=8.883333in,,base]{\color{textcolor}{\sffamily\fontsize{12.000000}{14.400000}\selectfont\catcode`\^=\active\def^{\ifmmode\sp\else\^{}\fi}\catcode`\%=\active\def%{\%}GS}}%
\end{pgfscope}%
\begin{pgfscope}%
\pgfsetbuttcap%
\pgfsetmiterjoin%
\definecolor{currentfill}{rgb}{1.000000,1.000000,1.000000}%
\pgfsetfillcolor{currentfill}%
\pgfsetlinewidth{0.000000pt}%
\definecolor{currentstroke}{rgb}{0.000000,0.000000,0.000000}%
\pgfsetstrokecolor{currentstroke}%
\pgfsetstrokeopacity{0.000000}%
\pgfsetdash{}{0pt}%
\pgfpathmoveto{\pgfqpoint{10.080882in}{5.300000in}}%
\pgfpathlineto{\pgfqpoint{13.500000in}{5.300000in}}%
\pgfpathlineto{\pgfqpoint{13.500000in}{8.800000in}}%
\pgfpathlineto{\pgfqpoint{10.080882in}{8.800000in}}%
\pgfpathlineto{\pgfqpoint{10.080882in}{5.300000in}}%
\pgfpathclose%
\pgfusepath{fill}%
\end{pgfscope}%
\begin{pgfscope}%
\pgfsetbuttcap%
\pgfsetroundjoin%
\definecolor{currentfill}{rgb}{0.000000,0.000000,0.000000}%
\pgfsetfillcolor{currentfill}%
\pgfsetlinewidth{0.803000pt}%
\definecolor{currentstroke}{rgb}{0.000000,0.000000,0.000000}%
\pgfsetstrokecolor{currentstroke}%
\pgfsetdash{}{0pt}%
\pgfsys@defobject{currentmarker}{\pgfqpoint{0.000000in}{-0.048611in}}{\pgfqpoint{0.000000in}{0.000000in}}{%
\pgfpathmoveto{\pgfqpoint{0.000000in}{0.000000in}}%
\pgfpathlineto{\pgfqpoint{0.000000in}{-0.048611in}}%
\pgfusepath{stroke,fill}%
}%
\begin{pgfscope}%
\pgfsys@transformshift{10.236297in}{5.300000in}%
\pgfsys@useobject{currentmarker}{}%
\end{pgfscope}%
\end{pgfscope}%
\begin{pgfscope}%
\definecolor{textcolor}{rgb}{0.000000,0.000000,0.000000}%
\pgfsetstrokecolor{textcolor}%
\pgfsetfillcolor{textcolor}%
\pgftext[x=10.236297in,y=5.202778in,,top]{\color{textcolor}{\sffamily\fontsize{10.000000}{12.000000}\selectfont\catcode`\^=\active\def^{\ifmmode\sp\else\^{}\fi}\catcode`\%=\active\def%{\%}0}}%
\end{pgfscope}%
\begin{pgfscope}%
\pgfsetbuttcap%
\pgfsetroundjoin%
\definecolor{currentfill}{rgb}{0.000000,0.000000,0.000000}%
\pgfsetfillcolor{currentfill}%
\pgfsetlinewidth{0.803000pt}%
\definecolor{currentstroke}{rgb}{0.000000,0.000000,0.000000}%
\pgfsetstrokecolor{currentstroke}%
\pgfsetdash{}{0pt}%
\pgfsys@defobject{currentmarker}{\pgfqpoint{0.000000in}{-0.048611in}}{\pgfqpoint{0.000000in}{0.000000in}}{%
\pgfpathmoveto{\pgfqpoint{0.000000in}{0.000000in}}%
\pgfpathlineto{\pgfqpoint{0.000000in}{-0.048611in}}%
\pgfusepath{stroke,fill}%
}%
\begin{pgfscope}%
\pgfsys@transformshift{10.864234in}{5.300000in}%
\pgfsys@useobject{currentmarker}{}%
\end{pgfscope}%
\end{pgfscope}%
\begin{pgfscope}%
\definecolor{textcolor}{rgb}{0.000000,0.000000,0.000000}%
\pgfsetstrokecolor{textcolor}%
\pgfsetfillcolor{textcolor}%
\pgftext[x=10.864234in,y=5.202778in,,top]{\color{textcolor}{\sffamily\fontsize{10.000000}{12.000000}\selectfont\catcode`\^=\active\def^{\ifmmode\sp\else\^{}\fi}\catcode`\%=\active\def%{\%}20}}%
\end{pgfscope}%
\begin{pgfscope}%
\pgfsetbuttcap%
\pgfsetroundjoin%
\definecolor{currentfill}{rgb}{0.000000,0.000000,0.000000}%
\pgfsetfillcolor{currentfill}%
\pgfsetlinewidth{0.803000pt}%
\definecolor{currentstroke}{rgb}{0.000000,0.000000,0.000000}%
\pgfsetstrokecolor{currentstroke}%
\pgfsetdash{}{0pt}%
\pgfsys@defobject{currentmarker}{\pgfqpoint{0.000000in}{-0.048611in}}{\pgfqpoint{0.000000in}{0.000000in}}{%
\pgfpathmoveto{\pgfqpoint{0.000000in}{0.000000in}}%
\pgfpathlineto{\pgfqpoint{0.000000in}{-0.048611in}}%
\pgfusepath{stroke,fill}%
}%
\begin{pgfscope}%
\pgfsys@transformshift{11.492171in}{5.300000in}%
\pgfsys@useobject{currentmarker}{}%
\end{pgfscope}%
\end{pgfscope}%
\begin{pgfscope}%
\definecolor{textcolor}{rgb}{0.000000,0.000000,0.000000}%
\pgfsetstrokecolor{textcolor}%
\pgfsetfillcolor{textcolor}%
\pgftext[x=11.492171in,y=5.202778in,,top]{\color{textcolor}{\sffamily\fontsize{10.000000}{12.000000}\selectfont\catcode`\^=\active\def^{\ifmmode\sp\else\^{}\fi}\catcode`\%=\active\def%{\%}40}}%
\end{pgfscope}%
\begin{pgfscope}%
\pgfsetbuttcap%
\pgfsetroundjoin%
\definecolor{currentfill}{rgb}{0.000000,0.000000,0.000000}%
\pgfsetfillcolor{currentfill}%
\pgfsetlinewidth{0.803000pt}%
\definecolor{currentstroke}{rgb}{0.000000,0.000000,0.000000}%
\pgfsetstrokecolor{currentstroke}%
\pgfsetdash{}{0pt}%
\pgfsys@defobject{currentmarker}{\pgfqpoint{0.000000in}{-0.048611in}}{\pgfqpoint{0.000000in}{0.000000in}}{%
\pgfpathmoveto{\pgfqpoint{0.000000in}{0.000000in}}%
\pgfpathlineto{\pgfqpoint{0.000000in}{-0.048611in}}%
\pgfusepath{stroke,fill}%
}%
\begin{pgfscope}%
\pgfsys@transformshift{12.120108in}{5.300000in}%
\pgfsys@useobject{currentmarker}{}%
\end{pgfscope}%
\end{pgfscope}%
\begin{pgfscope}%
\definecolor{textcolor}{rgb}{0.000000,0.000000,0.000000}%
\pgfsetstrokecolor{textcolor}%
\pgfsetfillcolor{textcolor}%
\pgftext[x=12.120108in,y=5.202778in,,top]{\color{textcolor}{\sffamily\fontsize{10.000000}{12.000000}\selectfont\catcode`\^=\active\def^{\ifmmode\sp\else\^{}\fi}\catcode`\%=\active\def%{\%}60}}%
\end{pgfscope}%
\begin{pgfscope}%
\pgfsetbuttcap%
\pgfsetroundjoin%
\definecolor{currentfill}{rgb}{0.000000,0.000000,0.000000}%
\pgfsetfillcolor{currentfill}%
\pgfsetlinewidth{0.803000pt}%
\definecolor{currentstroke}{rgb}{0.000000,0.000000,0.000000}%
\pgfsetstrokecolor{currentstroke}%
\pgfsetdash{}{0pt}%
\pgfsys@defobject{currentmarker}{\pgfqpoint{0.000000in}{-0.048611in}}{\pgfqpoint{0.000000in}{0.000000in}}{%
\pgfpathmoveto{\pgfqpoint{0.000000in}{0.000000in}}%
\pgfpathlineto{\pgfqpoint{0.000000in}{-0.048611in}}%
\pgfusepath{stroke,fill}%
}%
\begin{pgfscope}%
\pgfsys@transformshift{12.748045in}{5.300000in}%
\pgfsys@useobject{currentmarker}{}%
\end{pgfscope}%
\end{pgfscope}%
\begin{pgfscope}%
\definecolor{textcolor}{rgb}{0.000000,0.000000,0.000000}%
\pgfsetstrokecolor{textcolor}%
\pgfsetfillcolor{textcolor}%
\pgftext[x=12.748045in,y=5.202778in,,top]{\color{textcolor}{\sffamily\fontsize{10.000000}{12.000000}\selectfont\catcode`\^=\active\def^{\ifmmode\sp\else\^{}\fi}\catcode`\%=\active\def%{\%}80}}%
\end{pgfscope}%
\begin{pgfscope}%
\pgfsetbuttcap%
\pgfsetroundjoin%
\definecolor{currentfill}{rgb}{0.000000,0.000000,0.000000}%
\pgfsetfillcolor{currentfill}%
\pgfsetlinewidth{0.803000pt}%
\definecolor{currentstroke}{rgb}{0.000000,0.000000,0.000000}%
\pgfsetstrokecolor{currentstroke}%
\pgfsetdash{}{0pt}%
\pgfsys@defobject{currentmarker}{\pgfqpoint{0.000000in}{-0.048611in}}{\pgfqpoint{0.000000in}{0.000000in}}{%
\pgfpathmoveto{\pgfqpoint{0.000000in}{0.000000in}}%
\pgfpathlineto{\pgfqpoint{0.000000in}{-0.048611in}}%
\pgfusepath{stroke,fill}%
}%
\begin{pgfscope}%
\pgfsys@transformshift{13.375982in}{5.300000in}%
\pgfsys@useobject{currentmarker}{}%
\end{pgfscope}%
\end{pgfscope}%
\begin{pgfscope}%
\definecolor{textcolor}{rgb}{0.000000,0.000000,0.000000}%
\pgfsetstrokecolor{textcolor}%
\pgfsetfillcolor{textcolor}%
\pgftext[x=13.375982in,y=5.202778in,,top]{\color{textcolor}{\sffamily\fontsize{10.000000}{12.000000}\selectfont\catcode`\^=\active\def^{\ifmmode\sp\else\^{}\fi}\catcode`\%=\active\def%{\%}100}}%
\end{pgfscope}%
\begin{pgfscope}%
\definecolor{textcolor}{rgb}{0.000000,0.000000,0.000000}%
\pgfsetstrokecolor{textcolor}%
\pgfsetfillcolor{textcolor}%
\pgftext[x=11.790441in,y=5.012809in,,top]{\color{textcolor}{\sffamily\fontsize{10.000000}{12.000000}\selectfont\catcode`\^=\active\def^{\ifmmode\sp\else\^{}\fi}\catcode`\%=\active\def%{\%}V cycles}}%
\end{pgfscope}%
\begin{pgfscope}%
\pgfsetbuttcap%
\pgfsetroundjoin%
\definecolor{currentfill}{rgb}{0.000000,0.000000,0.000000}%
\pgfsetfillcolor{currentfill}%
\pgfsetlinewidth{0.803000pt}%
\definecolor{currentstroke}{rgb}{0.000000,0.000000,0.000000}%
\pgfsetstrokecolor{currentstroke}%
\pgfsetdash{}{0pt}%
\pgfsys@defobject{currentmarker}{\pgfqpoint{-0.048611in}{0.000000in}}{\pgfqpoint{-0.000000in}{0.000000in}}{%
\pgfpathmoveto{\pgfqpoint{-0.000000in}{0.000000in}}%
\pgfpathlineto{\pgfqpoint{-0.048611in}{0.000000in}}%
\pgfusepath{stroke,fill}%
}%
\begin{pgfscope}%
\pgfsys@transformshift{10.080882in}{5.806344in}%
\pgfsys@useobject{currentmarker}{}%
\end{pgfscope}%
\end{pgfscope}%
\begin{pgfscope}%
\pgfsetbuttcap%
\pgfsetroundjoin%
\definecolor{currentfill}{rgb}{0.000000,0.000000,0.000000}%
\pgfsetfillcolor{currentfill}%
\pgfsetlinewidth{0.803000pt}%
\definecolor{currentstroke}{rgb}{0.000000,0.000000,0.000000}%
\pgfsetstrokecolor{currentstroke}%
\pgfsetdash{}{0pt}%
\pgfsys@defobject{currentmarker}{\pgfqpoint{-0.048611in}{0.000000in}}{\pgfqpoint{-0.000000in}{0.000000in}}{%
\pgfpathmoveto{\pgfqpoint{-0.000000in}{0.000000in}}%
\pgfpathlineto{\pgfqpoint{-0.048611in}{0.000000in}}%
\pgfusepath{stroke,fill}%
}%
\begin{pgfscope}%
\pgfsys@transformshift{10.080882in}{6.687896in}%
\pgfsys@useobject{currentmarker}{}%
\end{pgfscope}%
\end{pgfscope}%
\begin{pgfscope}%
\pgfsetbuttcap%
\pgfsetroundjoin%
\definecolor{currentfill}{rgb}{0.000000,0.000000,0.000000}%
\pgfsetfillcolor{currentfill}%
\pgfsetlinewidth{0.803000pt}%
\definecolor{currentstroke}{rgb}{0.000000,0.000000,0.000000}%
\pgfsetstrokecolor{currentstroke}%
\pgfsetdash{}{0pt}%
\pgfsys@defobject{currentmarker}{\pgfqpoint{-0.048611in}{0.000000in}}{\pgfqpoint{-0.000000in}{0.000000in}}{%
\pgfpathmoveto{\pgfqpoint{-0.000000in}{0.000000in}}%
\pgfpathlineto{\pgfqpoint{-0.048611in}{0.000000in}}%
\pgfusepath{stroke,fill}%
}%
\begin{pgfscope}%
\pgfsys@transformshift{10.080882in}{7.569448in}%
\pgfsys@useobject{currentmarker}{}%
\end{pgfscope}%
\end{pgfscope}%
\begin{pgfscope}%
\pgfsetbuttcap%
\pgfsetroundjoin%
\definecolor{currentfill}{rgb}{0.000000,0.000000,0.000000}%
\pgfsetfillcolor{currentfill}%
\pgfsetlinewidth{0.803000pt}%
\definecolor{currentstroke}{rgb}{0.000000,0.000000,0.000000}%
\pgfsetstrokecolor{currentstroke}%
\pgfsetdash{}{0pt}%
\pgfsys@defobject{currentmarker}{\pgfqpoint{-0.048611in}{0.000000in}}{\pgfqpoint{-0.000000in}{0.000000in}}{%
\pgfpathmoveto{\pgfqpoint{-0.000000in}{0.000000in}}%
\pgfpathlineto{\pgfqpoint{-0.048611in}{0.000000in}}%
\pgfusepath{stroke,fill}%
}%
\begin{pgfscope}%
\pgfsys@transformshift{10.080882in}{8.451000in}%
\pgfsys@useobject{currentmarker}{}%
\end{pgfscope}%
\end{pgfscope}%
\begin{pgfscope}%
\pgfsetbuttcap%
\pgfsetroundjoin%
\definecolor{currentfill}{rgb}{0.000000,0.000000,0.000000}%
\pgfsetfillcolor{currentfill}%
\pgfsetlinewidth{0.602250pt}%
\definecolor{currentstroke}{rgb}{0.000000,0.000000,0.000000}%
\pgfsetstrokecolor{currentstroke}%
\pgfsetdash{}{0pt}%
\pgfsys@defobject{currentmarker}{\pgfqpoint{-0.027778in}{0.000000in}}{\pgfqpoint{-0.000000in}{0.000000in}}{%
\pgfpathmoveto{\pgfqpoint{-0.000000in}{0.000000in}}%
\pgfpathlineto{\pgfqpoint{-0.027778in}{0.000000in}}%
\pgfusepath{stroke,fill}%
}%
\begin{pgfscope}%
\pgfsys@transformshift{10.080882in}{5.345399in}%
\pgfsys@useobject{currentmarker}{}%
\end{pgfscope}%
\end{pgfscope}%
\begin{pgfscope}%
\pgfsetbuttcap%
\pgfsetroundjoin%
\definecolor{currentfill}{rgb}{0.000000,0.000000,0.000000}%
\pgfsetfillcolor{currentfill}%
\pgfsetlinewidth{0.602250pt}%
\definecolor{currentstroke}{rgb}{0.000000,0.000000,0.000000}%
\pgfsetstrokecolor{currentstroke}%
\pgfsetdash{}{0pt}%
\pgfsys@defobject{currentmarker}{\pgfqpoint{-0.027778in}{0.000000in}}{\pgfqpoint{-0.000000in}{0.000000in}}{%
\pgfpathmoveto{\pgfqpoint{-0.000000in}{0.000000in}}%
\pgfpathlineto{\pgfqpoint{-0.027778in}{0.000000in}}%
\pgfusepath{stroke,fill}%
}%
\begin{pgfscope}%
\pgfsys@transformshift{10.080882in}{5.455539in}%
\pgfsys@useobject{currentmarker}{}%
\end{pgfscope}%
\end{pgfscope}%
\begin{pgfscope}%
\pgfsetbuttcap%
\pgfsetroundjoin%
\definecolor{currentfill}{rgb}{0.000000,0.000000,0.000000}%
\pgfsetfillcolor{currentfill}%
\pgfsetlinewidth{0.602250pt}%
\definecolor{currentstroke}{rgb}{0.000000,0.000000,0.000000}%
\pgfsetstrokecolor{currentstroke}%
\pgfsetdash{}{0pt}%
\pgfsys@defobject{currentmarker}{\pgfqpoint{-0.027778in}{0.000000in}}{\pgfqpoint{-0.000000in}{0.000000in}}{%
\pgfpathmoveto{\pgfqpoint{-0.000000in}{0.000000in}}%
\pgfpathlineto{\pgfqpoint{-0.027778in}{0.000000in}}%
\pgfusepath{stroke,fill}%
}%
\begin{pgfscope}%
\pgfsys@transformshift{10.080882in}{5.540970in}%
\pgfsys@useobject{currentmarker}{}%
\end{pgfscope}%
\end{pgfscope}%
\begin{pgfscope}%
\pgfsetbuttcap%
\pgfsetroundjoin%
\definecolor{currentfill}{rgb}{0.000000,0.000000,0.000000}%
\pgfsetfillcolor{currentfill}%
\pgfsetlinewidth{0.602250pt}%
\definecolor{currentstroke}{rgb}{0.000000,0.000000,0.000000}%
\pgfsetstrokecolor{currentstroke}%
\pgfsetdash{}{0pt}%
\pgfsys@defobject{currentmarker}{\pgfqpoint{-0.027778in}{0.000000in}}{\pgfqpoint{-0.000000in}{0.000000in}}{%
\pgfpathmoveto{\pgfqpoint{-0.000000in}{0.000000in}}%
\pgfpathlineto{\pgfqpoint{-0.027778in}{0.000000in}}%
\pgfusepath{stroke,fill}%
}%
\begin{pgfscope}%
\pgfsys@transformshift{10.080882in}{5.610773in}%
\pgfsys@useobject{currentmarker}{}%
\end{pgfscope}%
\end{pgfscope}%
\begin{pgfscope}%
\pgfsetbuttcap%
\pgfsetroundjoin%
\definecolor{currentfill}{rgb}{0.000000,0.000000,0.000000}%
\pgfsetfillcolor{currentfill}%
\pgfsetlinewidth{0.602250pt}%
\definecolor{currentstroke}{rgb}{0.000000,0.000000,0.000000}%
\pgfsetstrokecolor{currentstroke}%
\pgfsetdash{}{0pt}%
\pgfsys@defobject{currentmarker}{\pgfqpoint{-0.027778in}{0.000000in}}{\pgfqpoint{-0.000000in}{0.000000in}}{%
\pgfpathmoveto{\pgfqpoint{-0.000000in}{0.000000in}}%
\pgfpathlineto{\pgfqpoint{-0.027778in}{0.000000in}}%
\pgfusepath{stroke,fill}%
}%
\begin{pgfscope}%
\pgfsys@transformshift{10.080882in}{5.669790in}%
\pgfsys@useobject{currentmarker}{}%
\end{pgfscope}%
\end{pgfscope}%
\begin{pgfscope}%
\pgfsetbuttcap%
\pgfsetroundjoin%
\definecolor{currentfill}{rgb}{0.000000,0.000000,0.000000}%
\pgfsetfillcolor{currentfill}%
\pgfsetlinewidth{0.602250pt}%
\definecolor{currentstroke}{rgb}{0.000000,0.000000,0.000000}%
\pgfsetstrokecolor{currentstroke}%
\pgfsetdash{}{0pt}%
\pgfsys@defobject{currentmarker}{\pgfqpoint{-0.027778in}{0.000000in}}{\pgfqpoint{-0.000000in}{0.000000in}}{%
\pgfpathmoveto{\pgfqpoint{-0.000000in}{0.000000in}}%
\pgfpathlineto{\pgfqpoint{-0.027778in}{0.000000in}}%
\pgfusepath{stroke,fill}%
}%
\begin{pgfscope}%
\pgfsys@transformshift{10.080882in}{5.720913in}%
\pgfsys@useobject{currentmarker}{}%
\end{pgfscope}%
\end{pgfscope}%
\begin{pgfscope}%
\pgfsetbuttcap%
\pgfsetroundjoin%
\definecolor{currentfill}{rgb}{0.000000,0.000000,0.000000}%
\pgfsetfillcolor{currentfill}%
\pgfsetlinewidth{0.602250pt}%
\definecolor{currentstroke}{rgb}{0.000000,0.000000,0.000000}%
\pgfsetstrokecolor{currentstroke}%
\pgfsetdash{}{0pt}%
\pgfsys@defobject{currentmarker}{\pgfqpoint{-0.027778in}{0.000000in}}{\pgfqpoint{-0.000000in}{0.000000in}}{%
\pgfpathmoveto{\pgfqpoint{-0.000000in}{0.000000in}}%
\pgfpathlineto{\pgfqpoint{-0.027778in}{0.000000in}}%
\pgfusepath{stroke,fill}%
}%
\begin{pgfscope}%
\pgfsys@transformshift{10.080882in}{5.766006in}%
\pgfsys@useobject{currentmarker}{}%
\end{pgfscope}%
\end{pgfscope}%
\begin{pgfscope}%
\pgfsetbuttcap%
\pgfsetroundjoin%
\definecolor{currentfill}{rgb}{0.000000,0.000000,0.000000}%
\pgfsetfillcolor{currentfill}%
\pgfsetlinewidth{0.602250pt}%
\definecolor{currentstroke}{rgb}{0.000000,0.000000,0.000000}%
\pgfsetstrokecolor{currentstroke}%
\pgfsetdash{}{0pt}%
\pgfsys@defobject{currentmarker}{\pgfqpoint{-0.027778in}{0.000000in}}{\pgfqpoint{-0.000000in}{0.000000in}}{%
\pgfpathmoveto{\pgfqpoint{-0.000000in}{0.000000in}}%
\pgfpathlineto{\pgfqpoint{-0.027778in}{0.000000in}}%
\pgfusepath{stroke,fill}%
}%
\begin{pgfscope}%
\pgfsys@transformshift{10.080882in}{6.071717in}%
\pgfsys@useobject{currentmarker}{}%
\end{pgfscope}%
\end{pgfscope}%
\begin{pgfscope}%
\pgfsetbuttcap%
\pgfsetroundjoin%
\definecolor{currentfill}{rgb}{0.000000,0.000000,0.000000}%
\pgfsetfillcolor{currentfill}%
\pgfsetlinewidth{0.602250pt}%
\definecolor{currentstroke}{rgb}{0.000000,0.000000,0.000000}%
\pgfsetstrokecolor{currentstroke}%
\pgfsetdash{}{0pt}%
\pgfsys@defobject{currentmarker}{\pgfqpoint{-0.027778in}{0.000000in}}{\pgfqpoint{-0.000000in}{0.000000in}}{%
\pgfpathmoveto{\pgfqpoint{-0.000000in}{0.000000in}}%
\pgfpathlineto{\pgfqpoint{-0.027778in}{0.000000in}}%
\pgfusepath{stroke,fill}%
}%
\begin{pgfscope}%
\pgfsys@transformshift{10.080882in}{6.226951in}%
\pgfsys@useobject{currentmarker}{}%
\end{pgfscope}%
\end{pgfscope}%
\begin{pgfscope}%
\pgfsetbuttcap%
\pgfsetroundjoin%
\definecolor{currentfill}{rgb}{0.000000,0.000000,0.000000}%
\pgfsetfillcolor{currentfill}%
\pgfsetlinewidth{0.602250pt}%
\definecolor{currentstroke}{rgb}{0.000000,0.000000,0.000000}%
\pgfsetstrokecolor{currentstroke}%
\pgfsetdash{}{0pt}%
\pgfsys@defobject{currentmarker}{\pgfqpoint{-0.027778in}{0.000000in}}{\pgfqpoint{-0.000000in}{0.000000in}}{%
\pgfpathmoveto{\pgfqpoint{-0.000000in}{0.000000in}}%
\pgfpathlineto{\pgfqpoint{-0.027778in}{0.000000in}}%
\pgfusepath{stroke,fill}%
}%
\begin{pgfscope}%
\pgfsys@transformshift{10.080882in}{6.337091in}%
\pgfsys@useobject{currentmarker}{}%
\end{pgfscope}%
\end{pgfscope}%
\begin{pgfscope}%
\pgfsetbuttcap%
\pgfsetroundjoin%
\definecolor{currentfill}{rgb}{0.000000,0.000000,0.000000}%
\pgfsetfillcolor{currentfill}%
\pgfsetlinewidth{0.602250pt}%
\definecolor{currentstroke}{rgb}{0.000000,0.000000,0.000000}%
\pgfsetstrokecolor{currentstroke}%
\pgfsetdash{}{0pt}%
\pgfsys@defobject{currentmarker}{\pgfqpoint{-0.027778in}{0.000000in}}{\pgfqpoint{-0.000000in}{0.000000in}}{%
\pgfpathmoveto{\pgfqpoint{-0.000000in}{0.000000in}}%
\pgfpathlineto{\pgfqpoint{-0.027778in}{0.000000in}}%
\pgfusepath{stroke,fill}%
}%
\begin{pgfscope}%
\pgfsys@transformshift{10.080882in}{6.422522in}%
\pgfsys@useobject{currentmarker}{}%
\end{pgfscope}%
\end{pgfscope}%
\begin{pgfscope}%
\pgfsetbuttcap%
\pgfsetroundjoin%
\definecolor{currentfill}{rgb}{0.000000,0.000000,0.000000}%
\pgfsetfillcolor{currentfill}%
\pgfsetlinewidth{0.602250pt}%
\definecolor{currentstroke}{rgb}{0.000000,0.000000,0.000000}%
\pgfsetstrokecolor{currentstroke}%
\pgfsetdash{}{0pt}%
\pgfsys@defobject{currentmarker}{\pgfqpoint{-0.027778in}{0.000000in}}{\pgfqpoint{-0.000000in}{0.000000in}}{%
\pgfpathmoveto{\pgfqpoint{-0.000000in}{0.000000in}}%
\pgfpathlineto{\pgfqpoint{-0.027778in}{0.000000in}}%
\pgfusepath{stroke,fill}%
}%
\begin{pgfscope}%
\pgfsys@transformshift{10.080882in}{6.492324in}%
\pgfsys@useobject{currentmarker}{}%
\end{pgfscope}%
\end{pgfscope}%
\begin{pgfscope}%
\pgfsetbuttcap%
\pgfsetroundjoin%
\definecolor{currentfill}{rgb}{0.000000,0.000000,0.000000}%
\pgfsetfillcolor{currentfill}%
\pgfsetlinewidth{0.602250pt}%
\definecolor{currentstroke}{rgb}{0.000000,0.000000,0.000000}%
\pgfsetstrokecolor{currentstroke}%
\pgfsetdash{}{0pt}%
\pgfsys@defobject{currentmarker}{\pgfqpoint{-0.027778in}{0.000000in}}{\pgfqpoint{-0.000000in}{0.000000in}}{%
\pgfpathmoveto{\pgfqpoint{-0.000000in}{0.000000in}}%
\pgfpathlineto{\pgfqpoint{-0.027778in}{0.000000in}}%
\pgfusepath{stroke,fill}%
}%
\begin{pgfscope}%
\pgfsys@transformshift{10.080882in}{6.551342in}%
\pgfsys@useobject{currentmarker}{}%
\end{pgfscope}%
\end{pgfscope}%
\begin{pgfscope}%
\pgfsetbuttcap%
\pgfsetroundjoin%
\definecolor{currentfill}{rgb}{0.000000,0.000000,0.000000}%
\pgfsetfillcolor{currentfill}%
\pgfsetlinewidth{0.602250pt}%
\definecolor{currentstroke}{rgb}{0.000000,0.000000,0.000000}%
\pgfsetstrokecolor{currentstroke}%
\pgfsetdash{}{0pt}%
\pgfsys@defobject{currentmarker}{\pgfqpoint{-0.027778in}{0.000000in}}{\pgfqpoint{-0.000000in}{0.000000in}}{%
\pgfpathmoveto{\pgfqpoint{-0.000000in}{0.000000in}}%
\pgfpathlineto{\pgfqpoint{-0.027778in}{0.000000in}}%
\pgfusepath{stroke,fill}%
}%
\begin{pgfscope}%
\pgfsys@transformshift{10.080882in}{6.602464in}%
\pgfsys@useobject{currentmarker}{}%
\end{pgfscope}%
\end{pgfscope}%
\begin{pgfscope}%
\pgfsetbuttcap%
\pgfsetroundjoin%
\definecolor{currentfill}{rgb}{0.000000,0.000000,0.000000}%
\pgfsetfillcolor{currentfill}%
\pgfsetlinewidth{0.602250pt}%
\definecolor{currentstroke}{rgb}{0.000000,0.000000,0.000000}%
\pgfsetstrokecolor{currentstroke}%
\pgfsetdash{}{0pt}%
\pgfsys@defobject{currentmarker}{\pgfqpoint{-0.027778in}{0.000000in}}{\pgfqpoint{-0.000000in}{0.000000in}}{%
\pgfpathmoveto{\pgfqpoint{-0.000000in}{0.000000in}}%
\pgfpathlineto{\pgfqpoint{-0.027778in}{0.000000in}}%
\pgfusepath{stroke,fill}%
}%
\begin{pgfscope}%
\pgfsys@transformshift{10.080882in}{6.647558in}%
\pgfsys@useobject{currentmarker}{}%
\end{pgfscope}%
\end{pgfscope}%
\begin{pgfscope}%
\pgfsetbuttcap%
\pgfsetroundjoin%
\definecolor{currentfill}{rgb}{0.000000,0.000000,0.000000}%
\pgfsetfillcolor{currentfill}%
\pgfsetlinewidth{0.602250pt}%
\definecolor{currentstroke}{rgb}{0.000000,0.000000,0.000000}%
\pgfsetstrokecolor{currentstroke}%
\pgfsetdash{}{0pt}%
\pgfsys@defobject{currentmarker}{\pgfqpoint{-0.027778in}{0.000000in}}{\pgfqpoint{-0.000000in}{0.000000in}}{%
\pgfpathmoveto{\pgfqpoint{-0.000000in}{0.000000in}}%
\pgfpathlineto{\pgfqpoint{-0.027778in}{0.000000in}}%
\pgfusepath{stroke,fill}%
}%
\begin{pgfscope}%
\pgfsys@transformshift{10.080882in}{6.953269in}%
\pgfsys@useobject{currentmarker}{}%
\end{pgfscope}%
\end{pgfscope}%
\begin{pgfscope}%
\pgfsetbuttcap%
\pgfsetroundjoin%
\definecolor{currentfill}{rgb}{0.000000,0.000000,0.000000}%
\pgfsetfillcolor{currentfill}%
\pgfsetlinewidth{0.602250pt}%
\definecolor{currentstroke}{rgb}{0.000000,0.000000,0.000000}%
\pgfsetstrokecolor{currentstroke}%
\pgfsetdash{}{0pt}%
\pgfsys@defobject{currentmarker}{\pgfqpoint{-0.027778in}{0.000000in}}{\pgfqpoint{-0.000000in}{0.000000in}}{%
\pgfpathmoveto{\pgfqpoint{-0.000000in}{0.000000in}}%
\pgfpathlineto{\pgfqpoint{-0.027778in}{0.000000in}}%
\pgfusepath{stroke,fill}%
}%
\begin{pgfscope}%
\pgfsys@transformshift{10.080882in}{7.108503in}%
\pgfsys@useobject{currentmarker}{}%
\end{pgfscope}%
\end{pgfscope}%
\begin{pgfscope}%
\pgfsetbuttcap%
\pgfsetroundjoin%
\definecolor{currentfill}{rgb}{0.000000,0.000000,0.000000}%
\pgfsetfillcolor{currentfill}%
\pgfsetlinewidth{0.602250pt}%
\definecolor{currentstroke}{rgb}{0.000000,0.000000,0.000000}%
\pgfsetstrokecolor{currentstroke}%
\pgfsetdash{}{0pt}%
\pgfsys@defobject{currentmarker}{\pgfqpoint{-0.027778in}{0.000000in}}{\pgfqpoint{-0.000000in}{0.000000in}}{%
\pgfpathmoveto{\pgfqpoint{-0.000000in}{0.000000in}}%
\pgfpathlineto{\pgfqpoint{-0.027778in}{0.000000in}}%
\pgfusepath{stroke,fill}%
}%
\begin{pgfscope}%
\pgfsys@transformshift{10.080882in}{7.218643in}%
\pgfsys@useobject{currentmarker}{}%
\end{pgfscope}%
\end{pgfscope}%
\begin{pgfscope}%
\pgfsetbuttcap%
\pgfsetroundjoin%
\definecolor{currentfill}{rgb}{0.000000,0.000000,0.000000}%
\pgfsetfillcolor{currentfill}%
\pgfsetlinewidth{0.602250pt}%
\definecolor{currentstroke}{rgb}{0.000000,0.000000,0.000000}%
\pgfsetstrokecolor{currentstroke}%
\pgfsetdash{}{0pt}%
\pgfsys@defobject{currentmarker}{\pgfqpoint{-0.027778in}{0.000000in}}{\pgfqpoint{-0.000000in}{0.000000in}}{%
\pgfpathmoveto{\pgfqpoint{-0.000000in}{0.000000in}}%
\pgfpathlineto{\pgfqpoint{-0.027778in}{0.000000in}}%
\pgfusepath{stroke,fill}%
}%
\begin{pgfscope}%
\pgfsys@transformshift{10.080882in}{7.304074in}%
\pgfsys@useobject{currentmarker}{}%
\end{pgfscope}%
\end{pgfscope}%
\begin{pgfscope}%
\pgfsetbuttcap%
\pgfsetroundjoin%
\definecolor{currentfill}{rgb}{0.000000,0.000000,0.000000}%
\pgfsetfillcolor{currentfill}%
\pgfsetlinewidth{0.602250pt}%
\definecolor{currentstroke}{rgb}{0.000000,0.000000,0.000000}%
\pgfsetstrokecolor{currentstroke}%
\pgfsetdash{}{0pt}%
\pgfsys@defobject{currentmarker}{\pgfqpoint{-0.027778in}{0.000000in}}{\pgfqpoint{-0.000000in}{0.000000in}}{%
\pgfpathmoveto{\pgfqpoint{-0.000000in}{0.000000in}}%
\pgfpathlineto{\pgfqpoint{-0.027778in}{0.000000in}}%
\pgfusepath{stroke,fill}%
}%
\begin{pgfscope}%
\pgfsys@transformshift{10.080882in}{7.373876in}%
\pgfsys@useobject{currentmarker}{}%
\end{pgfscope}%
\end{pgfscope}%
\begin{pgfscope}%
\pgfsetbuttcap%
\pgfsetroundjoin%
\definecolor{currentfill}{rgb}{0.000000,0.000000,0.000000}%
\pgfsetfillcolor{currentfill}%
\pgfsetlinewidth{0.602250pt}%
\definecolor{currentstroke}{rgb}{0.000000,0.000000,0.000000}%
\pgfsetstrokecolor{currentstroke}%
\pgfsetdash{}{0pt}%
\pgfsys@defobject{currentmarker}{\pgfqpoint{-0.027778in}{0.000000in}}{\pgfqpoint{-0.000000in}{0.000000in}}{%
\pgfpathmoveto{\pgfqpoint{-0.000000in}{0.000000in}}%
\pgfpathlineto{\pgfqpoint{-0.027778in}{0.000000in}}%
\pgfusepath{stroke,fill}%
}%
\begin{pgfscope}%
\pgfsys@transformshift{10.080882in}{7.432894in}%
\pgfsys@useobject{currentmarker}{}%
\end{pgfscope}%
\end{pgfscope}%
\begin{pgfscope}%
\pgfsetbuttcap%
\pgfsetroundjoin%
\definecolor{currentfill}{rgb}{0.000000,0.000000,0.000000}%
\pgfsetfillcolor{currentfill}%
\pgfsetlinewidth{0.602250pt}%
\definecolor{currentstroke}{rgb}{0.000000,0.000000,0.000000}%
\pgfsetstrokecolor{currentstroke}%
\pgfsetdash{}{0pt}%
\pgfsys@defobject{currentmarker}{\pgfqpoint{-0.027778in}{0.000000in}}{\pgfqpoint{-0.000000in}{0.000000in}}{%
\pgfpathmoveto{\pgfqpoint{-0.000000in}{0.000000in}}%
\pgfpathlineto{\pgfqpoint{-0.027778in}{0.000000in}}%
\pgfusepath{stroke,fill}%
}%
\begin{pgfscope}%
\pgfsys@transformshift{10.080882in}{7.484016in}%
\pgfsys@useobject{currentmarker}{}%
\end{pgfscope}%
\end{pgfscope}%
\begin{pgfscope}%
\pgfsetbuttcap%
\pgfsetroundjoin%
\definecolor{currentfill}{rgb}{0.000000,0.000000,0.000000}%
\pgfsetfillcolor{currentfill}%
\pgfsetlinewidth{0.602250pt}%
\definecolor{currentstroke}{rgb}{0.000000,0.000000,0.000000}%
\pgfsetstrokecolor{currentstroke}%
\pgfsetdash{}{0pt}%
\pgfsys@defobject{currentmarker}{\pgfqpoint{-0.027778in}{0.000000in}}{\pgfqpoint{-0.000000in}{0.000000in}}{%
\pgfpathmoveto{\pgfqpoint{-0.000000in}{0.000000in}}%
\pgfpathlineto{\pgfqpoint{-0.027778in}{0.000000in}}%
\pgfusepath{stroke,fill}%
}%
\begin{pgfscope}%
\pgfsys@transformshift{10.080882in}{7.529110in}%
\pgfsys@useobject{currentmarker}{}%
\end{pgfscope}%
\end{pgfscope}%
\begin{pgfscope}%
\pgfsetbuttcap%
\pgfsetroundjoin%
\definecolor{currentfill}{rgb}{0.000000,0.000000,0.000000}%
\pgfsetfillcolor{currentfill}%
\pgfsetlinewidth{0.602250pt}%
\definecolor{currentstroke}{rgb}{0.000000,0.000000,0.000000}%
\pgfsetstrokecolor{currentstroke}%
\pgfsetdash{}{0pt}%
\pgfsys@defobject{currentmarker}{\pgfqpoint{-0.027778in}{0.000000in}}{\pgfqpoint{-0.000000in}{0.000000in}}{%
\pgfpathmoveto{\pgfqpoint{-0.000000in}{0.000000in}}%
\pgfpathlineto{\pgfqpoint{-0.027778in}{0.000000in}}%
\pgfusepath{stroke,fill}%
}%
\begin{pgfscope}%
\pgfsys@transformshift{10.080882in}{7.834821in}%
\pgfsys@useobject{currentmarker}{}%
\end{pgfscope}%
\end{pgfscope}%
\begin{pgfscope}%
\pgfsetbuttcap%
\pgfsetroundjoin%
\definecolor{currentfill}{rgb}{0.000000,0.000000,0.000000}%
\pgfsetfillcolor{currentfill}%
\pgfsetlinewidth{0.602250pt}%
\definecolor{currentstroke}{rgb}{0.000000,0.000000,0.000000}%
\pgfsetstrokecolor{currentstroke}%
\pgfsetdash{}{0pt}%
\pgfsys@defobject{currentmarker}{\pgfqpoint{-0.027778in}{0.000000in}}{\pgfqpoint{-0.000000in}{0.000000in}}{%
\pgfpathmoveto{\pgfqpoint{-0.000000in}{0.000000in}}%
\pgfpathlineto{\pgfqpoint{-0.027778in}{0.000000in}}%
\pgfusepath{stroke,fill}%
}%
\begin{pgfscope}%
\pgfsys@transformshift{10.080882in}{7.990055in}%
\pgfsys@useobject{currentmarker}{}%
\end{pgfscope}%
\end{pgfscope}%
\begin{pgfscope}%
\pgfsetbuttcap%
\pgfsetroundjoin%
\definecolor{currentfill}{rgb}{0.000000,0.000000,0.000000}%
\pgfsetfillcolor{currentfill}%
\pgfsetlinewidth{0.602250pt}%
\definecolor{currentstroke}{rgb}{0.000000,0.000000,0.000000}%
\pgfsetstrokecolor{currentstroke}%
\pgfsetdash{}{0pt}%
\pgfsys@defobject{currentmarker}{\pgfqpoint{-0.027778in}{0.000000in}}{\pgfqpoint{-0.000000in}{0.000000in}}{%
\pgfpathmoveto{\pgfqpoint{-0.000000in}{0.000000in}}%
\pgfpathlineto{\pgfqpoint{-0.027778in}{0.000000in}}%
\pgfusepath{stroke,fill}%
}%
\begin{pgfscope}%
\pgfsys@transformshift{10.080882in}{8.100195in}%
\pgfsys@useobject{currentmarker}{}%
\end{pgfscope}%
\end{pgfscope}%
\begin{pgfscope}%
\pgfsetbuttcap%
\pgfsetroundjoin%
\definecolor{currentfill}{rgb}{0.000000,0.000000,0.000000}%
\pgfsetfillcolor{currentfill}%
\pgfsetlinewidth{0.602250pt}%
\definecolor{currentstroke}{rgb}{0.000000,0.000000,0.000000}%
\pgfsetstrokecolor{currentstroke}%
\pgfsetdash{}{0pt}%
\pgfsys@defobject{currentmarker}{\pgfqpoint{-0.027778in}{0.000000in}}{\pgfqpoint{-0.000000in}{0.000000in}}{%
\pgfpathmoveto{\pgfqpoint{-0.000000in}{0.000000in}}%
\pgfpathlineto{\pgfqpoint{-0.027778in}{0.000000in}}%
\pgfusepath{stroke,fill}%
}%
\begin{pgfscope}%
\pgfsys@transformshift{10.080882in}{8.185626in}%
\pgfsys@useobject{currentmarker}{}%
\end{pgfscope}%
\end{pgfscope}%
\begin{pgfscope}%
\pgfsetbuttcap%
\pgfsetroundjoin%
\definecolor{currentfill}{rgb}{0.000000,0.000000,0.000000}%
\pgfsetfillcolor{currentfill}%
\pgfsetlinewidth{0.602250pt}%
\definecolor{currentstroke}{rgb}{0.000000,0.000000,0.000000}%
\pgfsetstrokecolor{currentstroke}%
\pgfsetdash{}{0pt}%
\pgfsys@defobject{currentmarker}{\pgfqpoint{-0.027778in}{0.000000in}}{\pgfqpoint{-0.000000in}{0.000000in}}{%
\pgfpathmoveto{\pgfqpoint{-0.000000in}{0.000000in}}%
\pgfpathlineto{\pgfqpoint{-0.027778in}{0.000000in}}%
\pgfusepath{stroke,fill}%
}%
\begin{pgfscope}%
\pgfsys@transformshift{10.080882in}{8.255428in}%
\pgfsys@useobject{currentmarker}{}%
\end{pgfscope}%
\end{pgfscope}%
\begin{pgfscope}%
\pgfsetbuttcap%
\pgfsetroundjoin%
\definecolor{currentfill}{rgb}{0.000000,0.000000,0.000000}%
\pgfsetfillcolor{currentfill}%
\pgfsetlinewidth{0.602250pt}%
\definecolor{currentstroke}{rgb}{0.000000,0.000000,0.000000}%
\pgfsetstrokecolor{currentstroke}%
\pgfsetdash{}{0pt}%
\pgfsys@defobject{currentmarker}{\pgfqpoint{-0.027778in}{0.000000in}}{\pgfqpoint{-0.000000in}{0.000000in}}{%
\pgfpathmoveto{\pgfqpoint{-0.000000in}{0.000000in}}%
\pgfpathlineto{\pgfqpoint{-0.027778in}{0.000000in}}%
\pgfusepath{stroke,fill}%
}%
\begin{pgfscope}%
\pgfsys@transformshift{10.080882in}{8.314446in}%
\pgfsys@useobject{currentmarker}{}%
\end{pgfscope}%
\end{pgfscope}%
\begin{pgfscope}%
\pgfsetbuttcap%
\pgfsetroundjoin%
\definecolor{currentfill}{rgb}{0.000000,0.000000,0.000000}%
\pgfsetfillcolor{currentfill}%
\pgfsetlinewidth{0.602250pt}%
\definecolor{currentstroke}{rgb}{0.000000,0.000000,0.000000}%
\pgfsetstrokecolor{currentstroke}%
\pgfsetdash{}{0pt}%
\pgfsys@defobject{currentmarker}{\pgfqpoint{-0.027778in}{0.000000in}}{\pgfqpoint{-0.000000in}{0.000000in}}{%
\pgfpathmoveto{\pgfqpoint{-0.000000in}{0.000000in}}%
\pgfpathlineto{\pgfqpoint{-0.027778in}{0.000000in}}%
\pgfusepath{stroke,fill}%
}%
\begin{pgfscope}%
\pgfsys@transformshift{10.080882in}{8.365568in}%
\pgfsys@useobject{currentmarker}{}%
\end{pgfscope}%
\end{pgfscope}%
\begin{pgfscope}%
\pgfsetbuttcap%
\pgfsetroundjoin%
\definecolor{currentfill}{rgb}{0.000000,0.000000,0.000000}%
\pgfsetfillcolor{currentfill}%
\pgfsetlinewidth{0.602250pt}%
\definecolor{currentstroke}{rgb}{0.000000,0.000000,0.000000}%
\pgfsetstrokecolor{currentstroke}%
\pgfsetdash{}{0pt}%
\pgfsys@defobject{currentmarker}{\pgfqpoint{-0.027778in}{0.000000in}}{\pgfqpoint{-0.000000in}{0.000000in}}{%
\pgfpathmoveto{\pgfqpoint{-0.000000in}{0.000000in}}%
\pgfpathlineto{\pgfqpoint{-0.027778in}{0.000000in}}%
\pgfusepath{stroke,fill}%
}%
\begin{pgfscope}%
\pgfsys@transformshift{10.080882in}{8.410662in}%
\pgfsys@useobject{currentmarker}{}%
\end{pgfscope}%
\end{pgfscope}%
\begin{pgfscope}%
\pgfsetbuttcap%
\pgfsetroundjoin%
\definecolor{currentfill}{rgb}{0.000000,0.000000,0.000000}%
\pgfsetfillcolor{currentfill}%
\pgfsetlinewidth{0.602250pt}%
\definecolor{currentstroke}{rgb}{0.000000,0.000000,0.000000}%
\pgfsetstrokecolor{currentstroke}%
\pgfsetdash{}{0pt}%
\pgfsys@defobject{currentmarker}{\pgfqpoint{-0.027778in}{0.000000in}}{\pgfqpoint{-0.000000in}{0.000000in}}{%
\pgfpathmoveto{\pgfqpoint{-0.000000in}{0.000000in}}%
\pgfpathlineto{\pgfqpoint{-0.027778in}{0.000000in}}%
\pgfusepath{stroke,fill}%
}%
\begin{pgfscope}%
\pgfsys@transformshift{10.080882in}{8.716373in}%
\pgfsys@useobject{currentmarker}{}%
\end{pgfscope}%
\end{pgfscope}%
\begin{pgfscope}%
\pgfpathrectangle{\pgfqpoint{10.080882in}{5.300000in}}{\pgfqpoint{3.419118in}{3.500000in}}%
\pgfusepath{clip}%
\pgfsetrectcap%
\pgfsetroundjoin%
\pgfsetlinewidth{1.505625pt}%
\definecolor{currentstroke}{rgb}{0.121569,0.466667,0.705882}%
\pgfsetstrokecolor{currentstroke}%
\pgfsetdash{}{0pt}%
\pgfpathmoveto{\pgfqpoint{10.236297in}{8.244732in}}%
\pgfpathlineto{\pgfqpoint{10.267694in}{8.224241in}}%
\pgfpathlineto{\pgfqpoint{10.299091in}{8.203841in}}%
\pgfpathlineto{\pgfqpoint{10.330487in}{8.183496in}}%
\pgfpathlineto{\pgfqpoint{10.361884in}{8.163187in}}%
\pgfpathlineto{\pgfqpoint{10.393281in}{8.142903in}}%
\pgfpathlineto{\pgfqpoint{10.424678in}{8.122635in}}%
\pgfpathlineto{\pgfqpoint{10.456075in}{8.102379in}}%
\pgfpathlineto{\pgfqpoint{10.487472in}{8.082132in}}%
\pgfpathlineto{\pgfqpoint{10.518868in}{8.061890in}}%
\pgfpathlineto{\pgfqpoint{10.550265in}{8.047403in}}%
\pgfpathlineto{\pgfqpoint{10.581662in}{8.038670in}}%
\pgfpathlineto{\pgfqpoint{10.613059in}{8.029944in}}%
\pgfpathlineto{\pgfqpoint{10.644456in}{8.021218in}}%
\pgfpathlineto{\pgfqpoint{10.675853in}{8.012493in}}%
\pgfpathlineto{\pgfqpoint{10.707250in}{8.003767in}}%
\pgfpathlineto{\pgfqpoint{10.738646in}{7.995042in}}%
\pgfpathlineto{\pgfqpoint{10.770043in}{7.986317in}}%
\pgfpathlineto{\pgfqpoint{10.801440in}{7.977593in}}%
\pgfpathlineto{\pgfqpoint{10.832837in}{7.968868in}}%
\pgfpathlineto{\pgfqpoint{10.864234in}{7.960143in}}%
\pgfpathlineto{\pgfqpoint{10.895631in}{7.951419in}}%
\pgfpathlineto{\pgfqpoint{10.927028in}{7.942694in}}%
\pgfpathlineto{\pgfqpoint{10.958424in}{7.933970in}}%
\pgfpathlineto{\pgfqpoint{10.989821in}{7.925245in}}%
\pgfpathlineto{\pgfqpoint{11.021218in}{7.916520in}}%
\pgfpathlineto{\pgfqpoint{11.052615in}{7.907796in}}%
\pgfpathlineto{\pgfqpoint{11.084012in}{7.899071in}}%
\pgfpathlineto{\pgfqpoint{11.115409in}{7.890346in}}%
\pgfpathlineto{\pgfqpoint{11.146806in}{7.881622in}}%
\pgfpathlineto{\pgfqpoint{11.178202in}{7.872897in}}%
\pgfpathlineto{\pgfqpoint{11.209599in}{7.864172in}}%
\pgfpathlineto{\pgfqpoint{11.240996in}{7.855448in}}%
\pgfpathlineto{\pgfqpoint{11.272393in}{7.846723in}}%
\pgfpathlineto{\pgfqpoint{11.303790in}{7.837998in}}%
\pgfpathlineto{\pgfqpoint{11.335187in}{7.829273in}}%
\pgfpathlineto{\pgfqpoint{11.366584in}{7.820548in}}%
\pgfpathlineto{\pgfqpoint{11.397980in}{7.811823in}}%
\pgfpathlineto{\pgfqpoint{11.429377in}{7.803098in}}%
\pgfpathlineto{\pgfqpoint{11.460774in}{7.794373in}}%
\pgfpathlineto{\pgfqpoint{11.492171in}{7.785648in}}%
\pgfpathlineto{\pgfqpoint{11.523568in}{7.776924in}}%
\pgfpathlineto{\pgfqpoint{11.554965in}{7.768198in}}%
\pgfpathlineto{\pgfqpoint{11.586362in}{7.759473in}}%
\pgfpathlineto{\pgfqpoint{11.617758in}{7.750748in}}%
\pgfpathlineto{\pgfqpoint{11.649155in}{7.742023in}}%
\pgfpathlineto{\pgfqpoint{11.680552in}{7.733298in}}%
\pgfpathlineto{\pgfqpoint{11.711949in}{7.724573in}}%
\pgfpathlineto{\pgfqpoint{11.743346in}{7.715848in}}%
\pgfpathlineto{\pgfqpoint{11.774743in}{7.707123in}}%
\pgfpathlineto{\pgfqpoint{11.806140in}{7.698398in}}%
\pgfpathlineto{\pgfqpoint{11.837536in}{7.689673in}}%
\pgfpathlineto{\pgfqpoint{11.868933in}{7.680948in}}%
\pgfpathlineto{\pgfqpoint{11.900330in}{7.672223in}}%
\pgfpathlineto{\pgfqpoint{11.931727in}{7.663498in}}%
\pgfpathlineto{\pgfqpoint{11.963124in}{7.654773in}}%
\pgfpathlineto{\pgfqpoint{11.994521in}{7.646047in}}%
\pgfpathlineto{\pgfqpoint{12.025918in}{7.637322in}}%
\pgfpathlineto{\pgfqpoint{12.057314in}{7.628597in}}%
\pgfpathlineto{\pgfqpoint{12.088711in}{7.619872in}}%
\pgfpathlineto{\pgfqpoint{12.120108in}{7.611147in}}%
\pgfpathlineto{\pgfqpoint{12.151505in}{7.602422in}}%
\pgfpathlineto{\pgfqpoint{12.182902in}{7.593697in}}%
\pgfpathlineto{\pgfqpoint{12.214299in}{7.584972in}}%
\pgfpathlineto{\pgfqpoint{12.245696in}{7.576247in}}%
\pgfpathlineto{\pgfqpoint{12.277092in}{7.567522in}}%
\pgfpathlineto{\pgfqpoint{12.308489in}{7.558796in}}%
\pgfpathlineto{\pgfqpoint{12.339886in}{7.550071in}}%
\pgfpathlineto{\pgfqpoint{12.371283in}{7.541346in}}%
\pgfpathlineto{\pgfqpoint{12.402680in}{7.532621in}}%
\pgfpathlineto{\pgfqpoint{12.434077in}{7.523896in}}%
\pgfpathlineto{\pgfqpoint{12.465474in}{7.515171in}}%
\pgfpathlineto{\pgfqpoint{12.496870in}{7.506446in}}%
\pgfpathlineto{\pgfqpoint{12.528267in}{7.497721in}}%
\pgfpathlineto{\pgfqpoint{12.559664in}{7.488996in}}%
\pgfpathlineto{\pgfqpoint{12.591061in}{7.480271in}}%
\pgfpathlineto{\pgfqpoint{12.622458in}{7.471546in}}%
\pgfpathlineto{\pgfqpoint{12.653855in}{7.462821in}}%
\pgfpathlineto{\pgfqpoint{12.685252in}{7.454096in}}%
\pgfpathlineto{\pgfqpoint{12.716648in}{7.445371in}}%
\pgfpathlineto{\pgfqpoint{12.748045in}{7.436646in}}%
\pgfpathlineto{\pgfqpoint{12.779442in}{7.427921in}}%
\pgfpathlineto{\pgfqpoint{12.810839in}{7.419196in}}%
\pgfpathlineto{\pgfqpoint{12.842236in}{7.410471in}}%
\pgfpathlineto{\pgfqpoint{12.873633in}{7.401746in}}%
\pgfpathlineto{\pgfqpoint{12.905030in}{7.393022in}}%
\pgfpathlineto{\pgfqpoint{12.936426in}{7.384297in}}%
\pgfpathlineto{\pgfqpoint{12.967823in}{7.375572in}}%
\pgfpathlineto{\pgfqpoint{12.999220in}{7.366847in}}%
\pgfpathlineto{\pgfqpoint{13.030617in}{7.358122in}}%
\pgfpathlineto{\pgfqpoint{13.062014in}{7.349397in}}%
\pgfpathlineto{\pgfqpoint{13.093411in}{7.340672in}}%
\pgfpathlineto{\pgfqpoint{13.124808in}{7.331947in}}%
\pgfpathlineto{\pgfqpoint{13.156204in}{7.323223in}}%
\pgfpathlineto{\pgfqpoint{13.187601in}{7.314498in}}%
\pgfpathlineto{\pgfqpoint{13.218998in}{7.305773in}}%
\pgfpathlineto{\pgfqpoint{13.250395in}{7.297048in}}%
\pgfpathlineto{\pgfqpoint{13.281792in}{7.288324in}}%
\pgfpathlineto{\pgfqpoint{13.313189in}{7.279599in}}%
\pgfpathlineto{\pgfqpoint{13.344586in}{7.270874in}}%
\pgfusepath{stroke}%
\end{pgfscope}%
\begin{pgfscope}%
\pgfpathrectangle{\pgfqpoint{10.080882in}{5.300000in}}{\pgfqpoint{3.419118in}{3.500000in}}%
\pgfusepath{clip}%
\pgfsetrectcap%
\pgfsetroundjoin%
\pgfsetlinewidth{1.505625pt}%
\definecolor{currentstroke}{rgb}{1.000000,0.498039,0.054902}%
\pgfsetstrokecolor{currentstroke}%
\pgfsetdash{}{0pt}%
\pgfpathmoveto{\pgfqpoint{10.236297in}{7.965941in}}%
\pgfpathlineto{\pgfqpoint{10.267694in}{7.943024in}}%
\pgfpathlineto{\pgfqpoint{10.299091in}{7.920912in}}%
\pgfpathlineto{\pgfqpoint{10.330487in}{7.905425in}}%
\pgfpathlineto{\pgfqpoint{10.361884in}{7.896196in}}%
\pgfpathlineto{\pgfqpoint{10.393281in}{7.887045in}}%
\pgfpathlineto{\pgfqpoint{10.424678in}{7.877956in}}%
\pgfpathlineto{\pgfqpoint{10.456075in}{7.868921in}}%
\pgfpathlineto{\pgfqpoint{10.487472in}{7.859934in}}%
\pgfpathlineto{\pgfqpoint{10.518868in}{7.850988in}}%
\pgfpathlineto{\pgfqpoint{10.550265in}{7.842078in}}%
\pgfpathlineto{\pgfqpoint{10.581662in}{7.833199in}}%
\pgfpathlineto{\pgfqpoint{10.613059in}{7.824348in}}%
\pgfpathlineto{\pgfqpoint{10.644456in}{7.815521in}}%
\pgfpathlineto{\pgfqpoint{10.675853in}{6.500799in}}%
\pgfpathlineto{\pgfqpoint{10.707250in}{6.454800in}}%
\pgfpathlineto{\pgfqpoint{10.738646in}{6.435830in}}%
\pgfpathlineto{\pgfqpoint{10.770043in}{6.420268in}}%
\pgfpathlineto{\pgfqpoint{10.801440in}{6.406511in}}%
\pgfpathlineto{\pgfqpoint{10.832837in}{6.393914in}}%
\pgfpathlineto{\pgfqpoint{10.864234in}{6.382135in}}%
\pgfpathlineto{\pgfqpoint{10.895631in}{6.370959in}}%
\pgfpathlineto{\pgfqpoint{10.927028in}{6.360240in}}%
\pgfpathlineto{\pgfqpoint{10.958424in}{6.349876in}}%
\pgfpathlineto{\pgfqpoint{10.989821in}{6.339793in}}%
\pgfpathlineto{\pgfqpoint{11.021218in}{6.329938in}}%
\pgfpathlineto{\pgfqpoint{11.052615in}{6.320267in}}%
\pgfpathlineto{\pgfqpoint{11.084012in}{6.310750in}}%
\pgfpathlineto{\pgfqpoint{11.115409in}{6.301361in}}%
\pgfpathlineto{\pgfqpoint{11.146806in}{6.284620in}}%
\pgfpathlineto{\pgfqpoint{11.178202in}{6.263939in}}%
\pgfpathlineto{\pgfqpoint{11.209599in}{6.248669in}}%
\pgfpathlineto{\pgfqpoint{11.240996in}{6.234765in}}%
\pgfpathlineto{\pgfqpoint{11.272393in}{6.221733in}}%
\pgfpathlineto{\pgfqpoint{11.303790in}{6.209375in}}%
\pgfpathlineto{\pgfqpoint{11.335187in}{6.197562in}}%
\pgfpathlineto{\pgfqpoint{11.366584in}{6.186200in}}%
\pgfpathlineto{\pgfqpoint{11.397980in}{6.175216in}}%
\pgfpathlineto{\pgfqpoint{11.429377in}{6.164549in}}%
\pgfpathlineto{\pgfqpoint{11.460774in}{6.154152in}}%
\pgfpathlineto{\pgfqpoint{11.492171in}{6.143985in}}%
\pgfpathlineto{\pgfqpoint{11.523568in}{6.134016in}}%
\pgfpathlineto{\pgfqpoint{11.554965in}{6.124219in}}%
\pgfpathlineto{\pgfqpoint{11.586362in}{6.114569in}}%
\pgfpathlineto{\pgfqpoint{11.617758in}{6.285094in}}%
\pgfpathlineto{\pgfqpoint{11.649155in}{6.264039in}}%
\pgfpathlineto{\pgfqpoint{11.680552in}{6.248745in}}%
\pgfpathlineto{\pgfqpoint{11.711949in}{6.234860in}}%
\pgfpathlineto{\pgfqpoint{11.743346in}{6.221856in}}%
\pgfpathlineto{\pgfqpoint{11.774743in}{6.209529in}}%
\pgfpathlineto{\pgfqpoint{11.806140in}{6.197749in}}%
\pgfpathlineto{\pgfqpoint{11.837536in}{6.186419in}}%
\pgfpathlineto{\pgfqpoint{11.868933in}{6.175465in}}%
\pgfpathlineto{\pgfqpoint{11.900330in}{6.164827in}}%
\pgfpathlineto{\pgfqpoint{11.931727in}{6.154457in}}%
\pgfpathlineto{\pgfqpoint{11.963124in}{6.144315in}}%
\pgfpathlineto{\pgfqpoint{11.994521in}{6.134369in}}%
\pgfpathlineto{\pgfqpoint{12.025918in}{6.124592in}}%
\pgfpathlineto{\pgfqpoint{12.057314in}{6.114961in}}%
\pgfpathlineto{\pgfqpoint{12.088711in}{6.285009in}}%
\pgfpathlineto{\pgfqpoint{12.120108in}{6.263930in}}%
\pgfpathlineto{\pgfqpoint{12.151505in}{6.248631in}}%
\pgfpathlineto{\pgfqpoint{12.182902in}{6.234744in}}%
\pgfpathlineto{\pgfqpoint{12.214299in}{6.221738in}}%
\pgfpathlineto{\pgfqpoint{12.245696in}{6.209410in}}%
\pgfpathlineto{\pgfqpoint{12.277092in}{6.197630in}}%
\pgfpathlineto{\pgfqpoint{12.308489in}{6.186301in}}%
\pgfpathlineto{\pgfqpoint{12.339886in}{6.175348in}}%
\pgfpathlineto{\pgfqpoint{12.371283in}{6.164709in}}%
\pgfpathlineto{\pgfqpoint{12.402680in}{6.154340in}}%
\pgfpathlineto{\pgfqpoint{12.434077in}{6.144198in}}%
\pgfpathlineto{\pgfqpoint{12.465474in}{6.134253in}}%
\pgfpathlineto{\pgfqpoint{12.496870in}{6.124476in}}%
\pgfpathlineto{\pgfqpoint{12.528267in}{6.114845in}}%
\pgfpathlineto{\pgfqpoint{12.559664in}{6.285018in}}%
\pgfpathlineto{\pgfqpoint{12.591061in}{6.263939in}}%
\pgfpathlineto{\pgfqpoint{12.622458in}{6.248640in}}%
\pgfpathlineto{\pgfqpoint{12.653855in}{6.234751in}}%
\pgfpathlineto{\pgfqpoint{12.685252in}{6.221745in}}%
\pgfpathlineto{\pgfqpoint{12.716648in}{6.209417in}}%
\pgfpathlineto{\pgfqpoint{12.748045in}{6.197636in}}%
\pgfpathlineto{\pgfqpoint{12.779442in}{6.186307in}}%
\pgfpathlineto{\pgfqpoint{12.810839in}{6.175353in}}%
\pgfpathlineto{\pgfqpoint{12.842236in}{6.164715in}}%
\pgfpathlineto{\pgfqpoint{12.873633in}{6.154346in}}%
\pgfpathlineto{\pgfqpoint{12.905030in}{6.144205in}}%
\pgfpathlineto{\pgfqpoint{12.936426in}{6.134260in}}%
\pgfpathlineto{\pgfqpoint{12.967823in}{6.124483in}}%
\pgfpathlineto{\pgfqpoint{12.999220in}{6.114852in}}%
\pgfpathlineto{\pgfqpoint{13.030617in}{6.285014in}}%
\pgfpathlineto{\pgfqpoint{13.062014in}{6.263935in}}%
\pgfpathlineto{\pgfqpoint{13.093411in}{6.248636in}}%
\pgfpathlineto{\pgfqpoint{13.124808in}{6.234749in}}%
\pgfpathlineto{\pgfqpoint{13.156204in}{6.221744in}}%
\pgfpathlineto{\pgfqpoint{13.187601in}{6.209417in}}%
\pgfpathlineto{\pgfqpoint{13.218998in}{6.197637in}}%
\pgfpathlineto{\pgfqpoint{13.250395in}{6.186308in}}%
\pgfpathlineto{\pgfqpoint{13.281792in}{6.175355in}}%
\pgfpathlineto{\pgfqpoint{13.313189in}{6.164717in}}%
\pgfpathlineto{\pgfqpoint{13.344586in}{6.154347in}}%
\pgfusepath{stroke}%
\end{pgfscope}%
\begin{pgfscope}%
\pgfpathrectangle{\pgfqpoint{10.080882in}{5.300000in}}{\pgfqpoint{3.419118in}{3.500000in}}%
\pgfusepath{clip}%
\pgfsetrectcap%
\pgfsetroundjoin%
\pgfsetlinewidth{1.505625pt}%
\definecolor{currentstroke}{rgb}{0.172549,0.627451,0.172549}%
\pgfsetstrokecolor{currentstroke}%
\pgfsetdash{}{0pt}%
\pgfpathmoveto{\pgfqpoint{10.236297in}{6.929213in}}%
\pgfpathlineto{\pgfqpoint{10.267694in}{6.620083in}}%
\pgfpathlineto{\pgfqpoint{10.299091in}{6.555783in}}%
\pgfpathlineto{\pgfqpoint{10.330487in}{6.503792in}}%
\pgfpathlineto{\pgfqpoint{10.361884in}{6.461068in}}%
\pgfpathlineto{\pgfqpoint{10.393281in}{6.425953in}}%
\pgfpathlineto{\pgfqpoint{10.424678in}{6.396993in}}%
\pgfpathlineto{\pgfqpoint{10.456075in}{6.372881in}}%
\pgfpathlineto{\pgfqpoint{10.487472in}{6.352507in}}%
\pgfpathlineto{\pgfqpoint{10.518868in}{6.334971in}}%
\pgfpathlineto{\pgfqpoint{10.550265in}{6.319568in}}%
\pgfpathlineto{\pgfqpoint{10.581662in}{6.305755in}}%
\pgfpathlineto{\pgfqpoint{10.613059in}{6.293127in}}%
\pgfpathlineto{\pgfqpoint{10.644456in}{6.281380in}}%
\pgfpathlineto{\pgfqpoint{10.675853in}{6.270294in}}%
\pgfpathlineto{\pgfqpoint{10.707250in}{6.259707in}}%
\pgfpathlineto{\pgfqpoint{10.738646in}{6.249491in}}%
\pgfpathlineto{\pgfqpoint{10.770043in}{6.239566in}}%
\pgfpathlineto{\pgfqpoint{10.801440in}{6.229865in}}%
\pgfpathlineto{\pgfqpoint{10.832837in}{6.220336in}}%
\pgfpathlineto{\pgfqpoint{10.864234in}{6.210948in}}%
\pgfpathlineto{\pgfqpoint{10.895631in}{6.201667in}}%
\pgfpathlineto{\pgfqpoint{10.927028in}{6.192477in}}%
\pgfpathlineto{\pgfqpoint{10.958424in}{6.183360in}}%
\pgfpathlineto{\pgfqpoint{10.989821in}{6.174301in}}%
\pgfpathlineto{\pgfqpoint{11.021218in}{6.165292in}}%
\pgfpathlineto{\pgfqpoint{11.052615in}{6.156329in}}%
\pgfpathlineto{\pgfqpoint{11.084012in}{6.147399in}}%
\pgfpathlineto{\pgfqpoint{11.115409in}{6.138503in}}%
\pgfpathlineto{\pgfqpoint{11.146806in}{6.129627in}}%
\pgfpathlineto{\pgfqpoint{11.178202in}{6.120774in}}%
\pgfpathlineto{\pgfqpoint{11.209599in}{6.111945in}}%
\pgfpathlineto{\pgfqpoint{11.240996in}{6.103129in}}%
\pgfpathlineto{\pgfqpoint{11.272393in}{6.094330in}}%
\pgfpathlineto{\pgfqpoint{11.303790in}{6.085541in}}%
\pgfpathlineto{\pgfqpoint{11.335187in}{6.076764in}}%
\pgfpathlineto{\pgfqpoint{11.366584in}{6.067995in}}%
\pgfpathlineto{\pgfqpoint{11.397980in}{6.059243in}}%
\pgfpathlineto{\pgfqpoint{11.429377in}{6.050490in}}%
\pgfpathlineto{\pgfqpoint{11.460774in}{6.041747in}}%
\pgfpathlineto{\pgfqpoint{11.492171in}{6.033008in}}%
\pgfpathlineto{\pgfqpoint{11.523568in}{6.024273in}}%
\pgfpathlineto{\pgfqpoint{11.554965in}{6.015537in}}%
\pgfpathlineto{\pgfqpoint{11.586362in}{6.006810in}}%
\pgfpathlineto{\pgfqpoint{11.617758in}{5.998089in}}%
\pgfpathlineto{\pgfqpoint{11.649155in}{5.989370in}}%
\pgfpathlineto{\pgfqpoint{11.680552in}{5.980654in}}%
\pgfpathlineto{\pgfqpoint{11.711949in}{5.971937in}}%
\pgfpathlineto{\pgfqpoint{11.743346in}{5.963225in}}%
\pgfpathlineto{\pgfqpoint{11.774743in}{5.954517in}}%
\pgfpathlineto{\pgfqpoint{11.806140in}{5.945812in}}%
\pgfpathlineto{\pgfqpoint{11.837536in}{5.937107in}}%
\pgfpathlineto{\pgfqpoint{11.868933in}{5.928401in}}%
\pgfpathlineto{\pgfqpoint{11.900330in}{5.919699in}}%
\pgfpathlineto{\pgfqpoint{11.931727in}{5.910993in}}%
\pgfpathlineto{\pgfqpoint{11.963124in}{5.902300in}}%
\pgfpathlineto{\pgfqpoint{11.994521in}{5.893601in}}%
\pgfpathlineto{\pgfqpoint{12.025918in}{5.884903in}}%
\pgfpathlineto{\pgfqpoint{12.057314in}{5.876202in}}%
\pgfpathlineto{\pgfqpoint{12.088711in}{5.867508in}}%
\pgfpathlineto{\pgfqpoint{12.120108in}{5.858805in}}%
\pgfpathlineto{\pgfqpoint{12.151505in}{5.850117in}}%
\pgfpathlineto{\pgfqpoint{12.182902in}{5.841429in}}%
\pgfpathlineto{\pgfqpoint{12.214299in}{5.832747in}}%
\pgfpathlineto{\pgfqpoint{12.245696in}{5.824052in}}%
\pgfpathlineto{\pgfqpoint{12.277092in}{5.815368in}}%
\pgfpathlineto{\pgfqpoint{12.308489in}{5.806681in}}%
\pgfpathlineto{\pgfqpoint{12.339886in}{5.797998in}}%
\pgfpathlineto{\pgfqpoint{12.371283in}{5.789311in}}%
\pgfpathlineto{\pgfqpoint{12.402680in}{5.780625in}}%
\pgfpathlineto{\pgfqpoint{12.434077in}{5.771946in}}%
\pgfpathlineto{\pgfqpoint{12.465474in}{5.763261in}}%
\pgfpathlineto{\pgfqpoint{12.496870in}{5.754578in}}%
\pgfpathlineto{\pgfqpoint{12.528267in}{5.745893in}}%
\pgfpathlineto{\pgfqpoint{12.559664in}{5.737206in}}%
\pgfpathlineto{\pgfqpoint{12.591061in}{5.728526in}}%
\pgfpathlineto{\pgfqpoint{12.622458in}{5.719848in}}%
\pgfpathlineto{\pgfqpoint{12.653855in}{5.711175in}}%
\pgfpathlineto{\pgfqpoint{12.685252in}{5.702508in}}%
\pgfpathlineto{\pgfqpoint{12.716648in}{5.693841in}}%
\pgfpathlineto{\pgfqpoint{12.748045in}{5.685175in}}%
\pgfpathlineto{\pgfqpoint{12.779442in}{5.676509in}}%
\pgfpathlineto{\pgfqpoint{12.810839in}{5.667859in}}%
\pgfpathlineto{\pgfqpoint{12.842236in}{5.659191in}}%
\pgfpathlineto{\pgfqpoint{12.873633in}{5.650532in}}%
\pgfpathlineto{\pgfqpoint{12.905030in}{5.641876in}}%
\pgfpathlineto{\pgfqpoint{12.936426in}{5.633228in}}%
\pgfpathlineto{\pgfqpoint{12.967823in}{5.624567in}}%
\pgfpathlineto{\pgfqpoint{12.999220in}{5.615927in}}%
\pgfpathlineto{\pgfqpoint{13.030617in}{5.607279in}}%
\pgfpathlineto{\pgfqpoint{13.062014in}{5.598624in}}%
\pgfpathlineto{\pgfqpoint{13.093411in}{5.589974in}}%
\pgfpathlineto{\pgfqpoint{13.124808in}{5.581332in}}%
\pgfpathlineto{\pgfqpoint{13.156204in}{5.572678in}}%
\pgfpathlineto{\pgfqpoint{13.187601in}{5.564043in}}%
\pgfpathlineto{\pgfqpoint{13.218998in}{5.555403in}}%
\pgfpathlineto{\pgfqpoint{13.250395in}{5.546786in}}%
\pgfpathlineto{\pgfqpoint{13.281792in}{5.538144in}}%
\pgfpathlineto{\pgfqpoint{13.313189in}{5.529507in}}%
\pgfpathlineto{\pgfqpoint{13.344586in}{5.520866in}}%
\pgfusepath{stroke}%
\end{pgfscope}%
\begin{pgfscope}%
\pgfsetrectcap%
\pgfsetmiterjoin%
\pgfsetlinewidth{0.803000pt}%
\definecolor{currentstroke}{rgb}{0.000000,0.000000,0.000000}%
\pgfsetstrokecolor{currentstroke}%
\pgfsetdash{}{0pt}%
\pgfpathmoveto{\pgfqpoint{10.080882in}{5.300000in}}%
\pgfpathlineto{\pgfqpoint{10.080882in}{8.800000in}}%
\pgfusepath{stroke}%
\end{pgfscope}%
\begin{pgfscope}%
\pgfsetrectcap%
\pgfsetmiterjoin%
\pgfsetlinewidth{0.803000pt}%
\definecolor{currentstroke}{rgb}{0.000000,0.000000,0.000000}%
\pgfsetstrokecolor{currentstroke}%
\pgfsetdash{}{0pt}%
\pgfpathmoveto{\pgfqpoint{13.500000in}{5.300000in}}%
\pgfpathlineto{\pgfqpoint{13.500000in}{8.800000in}}%
\pgfusepath{stroke}%
\end{pgfscope}%
\begin{pgfscope}%
\pgfsetrectcap%
\pgfsetmiterjoin%
\pgfsetlinewidth{0.803000pt}%
\definecolor{currentstroke}{rgb}{0.000000,0.000000,0.000000}%
\pgfsetstrokecolor{currentstroke}%
\pgfsetdash{}{0pt}%
\pgfpathmoveto{\pgfqpoint{10.080882in}{5.300000in}}%
\pgfpathlineto{\pgfqpoint{13.500000in}{5.300000in}}%
\pgfusepath{stroke}%
\end{pgfscope}%
\begin{pgfscope}%
\pgfsetrectcap%
\pgfsetmiterjoin%
\pgfsetlinewidth{0.803000pt}%
\definecolor{currentstroke}{rgb}{0.000000,0.000000,0.000000}%
\pgfsetstrokecolor{currentstroke}%
\pgfsetdash{}{0pt}%
\pgfpathmoveto{\pgfqpoint{10.080882in}{8.800000in}}%
\pgfpathlineto{\pgfqpoint{13.500000in}{8.800000in}}%
\pgfusepath{stroke}%
\end{pgfscope}%
\begin{pgfscope}%
\definecolor{textcolor}{rgb}{0.000000,0.000000,0.000000}%
\pgfsetstrokecolor{textcolor}%
\pgfsetfillcolor{textcolor}%
\pgftext[x=11.790441in,y=8.883333in,,base]{\color{textcolor}{\sffamily\fontsize{12.000000}{14.400000}\selectfont\catcode`\^=\active\def^{\ifmmode\sp\else\^{}\fi}\catcode`\%=\active\def%{\%}MG}}%
\end{pgfscope}%
\begin{pgfscope}%
\pgfsetbuttcap%
\pgfsetmiterjoin%
\definecolor{currentfill}{rgb}{1.000000,1.000000,1.000000}%
\pgfsetfillcolor{currentfill}%
\pgfsetlinewidth{0.000000pt}%
\definecolor{currentstroke}{rgb}{0.000000,0.000000,0.000000}%
\pgfsetstrokecolor{currentstroke}%
\pgfsetstrokeopacity{0.000000}%
\pgfsetdash{}{0pt}%
\pgfpathmoveto{\pgfqpoint{1.875000in}{1.100000in}}%
\pgfpathlineto{\pgfqpoint{5.294118in}{1.100000in}}%
\pgfpathlineto{\pgfqpoint{5.294118in}{4.600000in}}%
\pgfpathlineto{\pgfqpoint{1.875000in}{4.600000in}}%
\pgfpathlineto{\pgfqpoint{1.875000in}{1.100000in}}%
\pgfpathclose%
\pgfusepath{fill}%
\end{pgfscope}%
\begin{pgfscope}%
\pgfsetbuttcap%
\pgfsetroundjoin%
\definecolor{currentfill}{rgb}{0.000000,0.000000,0.000000}%
\pgfsetfillcolor{currentfill}%
\pgfsetlinewidth{0.803000pt}%
\definecolor{currentstroke}{rgb}{0.000000,0.000000,0.000000}%
\pgfsetstrokecolor{currentstroke}%
\pgfsetdash{}{0pt}%
\pgfsys@defobject{currentmarker}{\pgfqpoint{0.000000in}{-0.048611in}}{\pgfqpoint{0.000000in}{0.000000in}}{%
\pgfpathmoveto{\pgfqpoint{0.000000in}{0.000000in}}%
\pgfpathlineto{\pgfqpoint{0.000000in}{-0.048611in}}%
\pgfusepath{stroke,fill}%
}%
\begin{pgfscope}%
\pgfsys@transformshift{2.030414in}{1.100000in}%
\pgfsys@useobject{currentmarker}{}%
\end{pgfscope}%
\end{pgfscope}%
\begin{pgfscope}%
\definecolor{textcolor}{rgb}{0.000000,0.000000,0.000000}%
\pgfsetstrokecolor{textcolor}%
\pgfsetfillcolor{textcolor}%
\pgftext[x=2.030414in,y=1.002778in,,top]{\color{textcolor}{\sffamily\fontsize{10.000000}{12.000000}\selectfont\catcode`\^=\active\def^{\ifmmode\sp\else\^{}\fi}\catcode`\%=\active\def%{\%}0}}%
\end{pgfscope}%
\begin{pgfscope}%
\pgfsetbuttcap%
\pgfsetroundjoin%
\definecolor{currentfill}{rgb}{0.000000,0.000000,0.000000}%
\pgfsetfillcolor{currentfill}%
\pgfsetlinewidth{0.803000pt}%
\definecolor{currentstroke}{rgb}{0.000000,0.000000,0.000000}%
\pgfsetstrokecolor{currentstroke}%
\pgfsetdash{}{0pt}%
\pgfsys@defobject{currentmarker}{\pgfqpoint{0.000000in}{-0.048611in}}{\pgfqpoint{0.000000in}{0.000000in}}{%
\pgfpathmoveto{\pgfqpoint{0.000000in}{0.000000in}}%
\pgfpathlineto{\pgfqpoint{0.000000in}{-0.048611in}}%
\pgfusepath{stroke,fill}%
}%
\begin{pgfscope}%
\pgfsys@transformshift{2.550195in}{1.100000in}%
\pgfsys@useobject{currentmarker}{}%
\end{pgfscope}%
\end{pgfscope}%
\begin{pgfscope}%
\definecolor{textcolor}{rgb}{0.000000,0.000000,0.000000}%
\pgfsetstrokecolor{textcolor}%
\pgfsetfillcolor{textcolor}%
\pgftext[x=2.550195in,y=1.002778in,,top]{\color{textcolor}{\sffamily\fontsize{10.000000}{12.000000}\selectfont\catcode`\^=\active\def^{\ifmmode\sp\else\^{}\fi}\catcode`\%=\active\def%{\%}50}}%
\end{pgfscope}%
\begin{pgfscope}%
\pgfsetbuttcap%
\pgfsetroundjoin%
\definecolor{currentfill}{rgb}{0.000000,0.000000,0.000000}%
\pgfsetfillcolor{currentfill}%
\pgfsetlinewidth{0.803000pt}%
\definecolor{currentstroke}{rgb}{0.000000,0.000000,0.000000}%
\pgfsetstrokecolor{currentstroke}%
\pgfsetdash{}{0pt}%
\pgfsys@defobject{currentmarker}{\pgfqpoint{0.000000in}{-0.048611in}}{\pgfqpoint{0.000000in}{0.000000in}}{%
\pgfpathmoveto{\pgfqpoint{0.000000in}{0.000000in}}%
\pgfpathlineto{\pgfqpoint{0.000000in}{-0.048611in}}%
\pgfusepath{stroke,fill}%
}%
\begin{pgfscope}%
\pgfsys@transformshift{3.069976in}{1.100000in}%
\pgfsys@useobject{currentmarker}{}%
\end{pgfscope}%
\end{pgfscope}%
\begin{pgfscope}%
\definecolor{textcolor}{rgb}{0.000000,0.000000,0.000000}%
\pgfsetstrokecolor{textcolor}%
\pgfsetfillcolor{textcolor}%
\pgftext[x=3.069976in,y=1.002778in,,top]{\color{textcolor}{\sffamily\fontsize{10.000000}{12.000000}\selectfont\catcode`\^=\active\def^{\ifmmode\sp\else\^{}\fi}\catcode`\%=\active\def%{\%}100}}%
\end{pgfscope}%
\begin{pgfscope}%
\pgfsetbuttcap%
\pgfsetroundjoin%
\definecolor{currentfill}{rgb}{0.000000,0.000000,0.000000}%
\pgfsetfillcolor{currentfill}%
\pgfsetlinewidth{0.803000pt}%
\definecolor{currentstroke}{rgb}{0.000000,0.000000,0.000000}%
\pgfsetstrokecolor{currentstroke}%
\pgfsetdash{}{0pt}%
\pgfsys@defobject{currentmarker}{\pgfqpoint{0.000000in}{-0.048611in}}{\pgfqpoint{0.000000in}{0.000000in}}{%
\pgfpathmoveto{\pgfqpoint{0.000000in}{0.000000in}}%
\pgfpathlineto{\pgfqpoint{0.000000in}{-0.048611in}}%
\pgfusepath{stroke,fill}%
}%
\begin{pgfscope}%
\pgfsys@transformshift{3.589757in}{1.100000in}%
\pgfsys@useobject{currentmarker}{}%
\end{pgfscope}%
\end{pgfscope}%
\begin{pgfscope}%
\definecolor{textcolor}{rgb}{0.000000,0.000000,0.000000}%
\pgfsetstrokecolor{textcolor}%
\pgfsetfillcolor{textcolor}%
\pgftext[x=3.589757in,y=1.002778in,,top]{\color{textcolor}{\sffamily\fontsize{10.000000}{12.000000}\selectfont\catcode`\^=\active\def^{\ifmmode\sp\else\^{}\fi}\catcode`\%=\active\def%{\%}150}}%
\end{pgfscope}%
\begin{pgfscope}%
\pgfsetbuttcap%
\pgfsetroundjoin%
\definecolor{currentfill}{rgb}{0.000000,0.000000,0.000000}%
\pgfsetfillcolor{currentfill}%
\pgfsetlinewidth{0.803000pt}%
\definecolor{currentstroke}{rgb}{0.000000,0.000000,0.000000}%
\pgfsetstrokecolor{currentstroke}%
\pgfsetdash{}{0pt}%
\pgfsys@defobject{currentmarker}{\pgfqpoint{0.000000in}{-0.048611in}}{\pgfqpoint{0.000000in}{0.000000in}}{%
\pgfpathmoveto{\pgfqpoint{0.000000in}{0.000000in}}%
\pgfpathlineto{\pgfqpoint{0.000000in}{-0.048611in}}%
\pgfusepath{stroke,fill}%
}%
\begin{pgfscope}%
\pgfsys@transformshift{4.109537in}{1.100000in}%
\pgfsys@useobject{currentmarker}{}%
\end{pgfscope}%
\end{pgfscope}%
\begin{pgfscope}%
\definecolor{textcolor}{rgb}{0.000000,0.000000,0.000000}%
\pgfsetstrokecolor{textcolor}%
\pgfsetfillcolor{textcolor}%
\pgftext[x=4.109537in,y=1.002778in,,top]{\color{textcolor}{\sffamily\fontsize{10.000000}{12.000000}\selectfont\catcode`\^=\active\def^{\ifmmode\sp\else\^{}\fi}\catcode`\%=\active\def%{\%}200}}%
\end{pgfscope}%
\begin{pgfscope}%
\pgfsetbuttcap%
\pgfsetroundjoin%
\definecolor{currentfill}{rgb}{0.000000,0.000000,0.000000}%
\pgfsetfillcolor{currentfill}%
\pgfsetlinewidth{0.803000pt}%
\definecolor{currentstroke}{rgb}{0.000000,0.000000,0.000000}%
\pgfsetstrokecolor{currentstroke}%
\pgfsetdash{}{0pt}%
\pgfsys@defobject{currentmarker}{\pgfqpoint{0.000000in}{-0.048611in}}{\pgfqpoint{0.000000in}{0.000000in}}{%
\pgfpathmoveto{\pgfqpoint{0.000000in}{0.000000in}}%
\pgfpathlineto{\pgfqpoint{0.000000in}{-0.048611in}}%
\pgfusepath{stroke,fill}%
}%
\begin{pgfscope}%
\pgfsys@transformshift{4.629318in}{1.100000in}%
\pgfsys@useobject{currentmarker}{}%
\end{pgfscope}%
\end{pgfscope}%
\begin{pgfscope}%
\definecolor{textcolor}{rgb}{0.000000,0.000000,0.000000}%
\pgfsetstrokecolor{textcolor}%
\pgfsetfillcolor{textcolor}%
\pgftext[x=4.629318in,y=1.002778in,,top]{\color{textcolor}{\sffamily\fontsize{10.000000}{12.000000}\selectfont\catcode`\^=\active\def^{\ifmmode\sp\else\^{}\fi}\catcode`\%=\active\def%{\%}250}}%
\end{pgfscope}%
\begin{pgfscope}%
\pgfsetbuttcap%
\pgfsetroundjoin%
\definecolor{currentfill}{rgb}{0.000000,0.000000,0.000000}%
\pgfsetfillcolor{currentfill}%
\pgfsetlinewidth{0.803000pt}%
\definecolor{currentstroke}{rgb}{0.000000,0.000000,0.000000}%
\pgfsetstrokecolor{currentstroke}%
\pgfsetdash{}{0pt}%
\pgfsys@defobject{currentmarker}{\pgfqpoint{0.000000in}{-0.048611in}}{\pgfqpoint{0.000000in}{0.000000in}}{%
\pgfpathmoveto{\pgfqpoint{0.000000in}{0.000000in}}%
\pgfpathlineto{\pgfqpoint{0.000000in}{-0.048611in}}%
\pgfusepath{stroke,fill}%
}%
\begin{pgfscope}%
\pgfsys@transformshift{5.149099in}{1.100000in}%
\pgfsys@useobject{currentmarker}{}%
\end{pgfscope}%
\end{pgfscope}%
\begin{pgfscope}%
\definecolor{textcolor}{rgb}{0.000000,0.000000,0.000000}%
\pgfsetstrokecolor{textcolor}%
\pgfsetfillcolor{textcolor}%
\pgftext[x=5.149099in,y=1.002778in,,top]{\color{textcolor}{\sffamily\fontsize{10.000000}{12.000000}\selectfont\catcode`\^=\active\def^{\ifmmode\sp\else\^{}\fi}\catcode`\%=\active\def%{\%}300}}%
\end{pgfscope}%
\begin{pgfscope}%
\definecolor{textcolor}{rgb}{0.000000,0.000000,0.000000}%
\pgfsetstrokecolor{textcolor}%
\pgfsetfillcolor{textcolor}%
\pgftext[x=3.584559in,y=0.812809in,,top]{\color{textcolor}{\sffamily\fontsize{10.000000}{12.000000}\selectfont\catcode`\^=\active\def^{\ifmmode\sp\else\^{}\fi}\catcode`\%=\active\def%{\%}Iterations}}%
\end{pgfscope}%
\begin{pgfscope}%
\pgfsetbuttcap%
\pgfsetroundjoin%
\definecolor{currentfill}{rgb}{0.000000,0.000000,0.000000}%
\pgfsetfillcolor{currentfill}%
\pgfsetlinewidth{0.803000pt}%
\definecolor{currentstroke}{rgb}{0.000000,0.000000,0.000000}%
\pgfsetstrokecolor{currentstroke}%
\pgfsetdash{}{0pt}%
\pgfsys@defobject{currentmarker}{\pgfqpoint{-0.048611in}{0.000000in}}{\pgfqpoint{-0.000000in}{0.000000in}}{%
\pgfpathmoveto{\pgfqpoint{-0.000000in}{0.000000in}}%
\pgfpathlineto{\pgfqpoint{-0.048611in}{0.000000in}}%
\pgfusepath{stroke,fill}%
}%
\begin{pgfscope}%
\pgfsys@transformshift{1.875000in}{1.660284in}%
\pgfsys@useobject{currentmarker}{}%
\end{pgfscope}%
\end{pgfscope}%
\begin{pgfscope}%
\definecolor{textcolor}{rgb}{0.000000,0.000000,0.000000}%
\pgfsetstrokecolor{textcolor}%
\pgfsetfillcolor{textcolor}%
\pgftext[x=1.489775in, y=1.607522in, left, base]{\color{textcolor}{\sffamily\fontsize{10.000000}{12.000000}\selectfont\catcode`\^=\active\def^{\ifmmode\sp\else\^{}\fi}\catcode`\%=\active\def%{\%}$\mathdefault{10^{-6}}$}}%
\end{pgfscope}%
\begin{pgfscope}%
\pgfsetbuttcap%
\pgfsetroundjoin%
\definecolor{currentfill}{rgb}{0.000000,0.000000,0.000000}%
\pgfsetfillcolor{currentfill}%
\pgfsetlinewidth{0.803000pt}%
\definecolor{currentstroke}{rgb}{0.000000,0.000000,0.000000}%
\pgfsetstrokecolor{currentstroke}%
\pgfsetdash{}{0pt}%
\pgfsys@defobject{currentmarker}{\pgfqpoint{-0.048611in}{0.000000in}}{\pgfqpoint{-0.000000in}{0.000000in}}{%
\pgfpathmoveto{\pgfqpoint{-0.000000in}{0.000000in}}%
\pgfpathlineto{\pgfqpoint{-0.048611in}{0.000000in}}%
\pgfusepath{stroke,fill}%
}%
\begin{pgfscope}%
\pgfsys@transformshift{1.875000in}{2.534463in}%
\pgfsys@useobject{currentmarker}{}%
\end{pgfscope}%
\end{pgfscope}%
\begin{pgfscope}%
\definecolor{textcolor}{rgb}{0.000000,0.000000,0.000000}%
\pgfsetstrokecolor{textcolor}%
\pgfsetfillcolor{textcolor}%
\pgftext[x=1.489775in, y=2.481702in, left, base]{\color{textcolor}{\sffamily\fontsize{10.000000}{12.000000}\selectfont\catcode`\^=\active\def^{\ifmmode\sp\else\^{}\fi}\catcode`\%=\active\def%{\%}$\mathdefault{10^{-5}}$}}%
\end{pgfscope}%
\begin{pgfscope}%
\pgfsetbuttcap%
\pgfsetroundjoin%
\definecolor{currentfill}{rgb}{0.000000,0.000000,0.000000}%
\pgfsetfillcolor{currentfill}%
\pgfsetlinewidth{0.803000pt}%
\definecolor{currentstroke}{rgb}{0.000000,0.000000,0.000000}%
\pgfsetstrokecolor{currentstroke}%
\pgfsetdash{}{0pt}%
\pgfsys@defobject{currentmarker}{\pgfqpoint{-0.048611in}{0.000000in}}{\pgfqpoint{-0.000000in}{0.000000in}}{%
\pgfpathmoveto{\pgfqpoint{-0.000000in}{0.000000in}}%
\pgfpathlineto{\pgfqpoint{-0.048611in}{0.000000in}}%
\pgfusepath{stroke,fill}%
}%
\begin{pgfscope}%
\pgfsys@transformshift{1.875000in}{3.408642in}%
\pgfsys@useobject{currentmarker}{}%
\end{pgfscope}%
\end{pgfscope}%
\begin{pgfscope}%
\definecolor{textcolor}{rgb}{0.000000,0.000000,0.000000}%
\pgfsetstrokecolor{textcolor}%
\pgfsetfillcolor{textcolor}%
\pgftext[x=1.489775in, y=3.355881in, left, base]{\color{textcolor}{\sffamily\fontsize{10.000000}{12.000000}\selectfont\catcode`\^=\active\def^{\ifmmode\sp\else\^{}\fi}\catcode`\%=\active\def%{\%}$\mathdefault{10^{-4}}$}}%
\end{pgfscope}%
\begin{pgfscope}%
\pgfsetbuttcap%
\pgfsetroundjoin%
\definecolor{currentfill}{rgb}{0.000000,0.000000,0.000000}%
\pgfsetfillcolor{currentfill}%
\pgfsetlinewidth{0.803000pt}%
\definecolor{currentstroke}{rgb}{0.000000,0.000000,0.000000}%
\pgfsetstrokecolor{currentstroke}%
\pgfsetdash{}{0pt}%
\pgfsys@defobject{currentmarker}{\pgfqpoint{-0.048611in}{0.000000in}}{\pgfqpoint{-0.000000in}{0.000000in}}{%
\pgfpathmoveto{\pgfqpoint{-0.000000in}{0.000000in}}%
\pgfpathlineto{\pgfqpoint{-0.048611in}{0.000000in}}%
\pgfusepath{stroke,fill}%
}%
\begin{pgfscope}%
\pgfsys@transformshift{1.875000in}{4.282822in}%
\pgfsys@useobject{currentmarker}{}%
\end{pgfscope}%
\end{pgfscope}%
\begin{pgfscope}%
\definecolor{textcolor}{rgb}{0.000000,0.000000,0.000000}%
\pgfsetstrokecolor{textcolor}%
\pgfsetfillcolor{textcolor}%
\pgftext[x=1.489775in, y=4.230060in, left, base]{\color{textcolor}{\sffamily\fontsize{10.000000}{12.000000}\selectfont\catcode`\^=\active\def^{\ifmmode\sp\else\^{}\fi}\catcode`\%=\active\def%{\%}$\mathdefault{10^{-3}}$}}%
\end{pgfscope}%
\begin{pgfscope}%
\pgfsetbuttcap%
\pgfsetroundjoin%
\definecolor{currentfill}{rgb}{0.000000,0.000000,0.000000}%
\pgfsetfillcolor{currentfill}%
\pgfsetlinewidth{0.602250pt}%
\definecolor{currentstroke}{rgb}{0.000000,0.000000,0.000000}%
\pgfsetstrokecolor{currentstroke}%
\pgfsetdash{}{0pt}%
\pgfsys@defobject{currentmarker}{\pgfqpoint{-0.027778in}{0.000000in}}{\pgfqpoint{-0.000000in}{0.000000in}}{%
\pgfpathmoveto{\pgfqpoint{-0.000000in}{0.000000in}}%
\pgfpathlineto{\pgfqpoint{-0.027778in}{0.000000in}}%
\pgfusepath{stroke,fill}%
}%
\begin{pgfscope}%
\pgfsys@transformshift{1.875000in}{1.203194in}%
\pgfsys@useobject{currentmarker}{}%
\end{pgfscope}%
\end{pgfscope}%
\begin{pgfscope}%
\pgfsetbuttcap%
\pgfsetroundjoin%
\definecolor{currentfill}{rgb}{0.000000,0.000000,0.000000}%
\pgfsetfillcolor{currentfill}%
\pgfsetlinewidth{0.602250pt}%
\definecolor{currentstroke}{rgb}{0.000000,0.000000,0.000000}%
\pgfsetstrokecolor{currentstroke}%
\pgfsetdash{}{0pt}%
\pgfsys@defobject{currentmarker}{\pgfqpoint{-0.027778in}{0.000000in}}{\pgfqpoint{-0.000000in}{0.000000in}}{%
\pgfpathmoveto{\pgfqpoint{-0.000000in}{0.000000in}}%
\pgfpathlineto{\pgfqpoint{-0.027778in}{0.000000in}}%
\pgfusepath{stroke,fill}%
}%
\begin{pgfscope}%
\pgfsys@transformshift{1.875000in}{1.312413in}%
\pgfsys@useobject{currentmarker}{}%
\end{pgfscope}%
\end{pgfscope}%
\begin{pgfscope}%
\pgfsetbuttcap%
\pgfsetroundjoin%
\definecolor{currentfill}{rgb}{0.000000,0.000000,0.000000}%
\pgfsetfillcolor{currentfill}%
\pgfsetlinewidth{0.602250pt}%
\definecolor{currentstroke}{rgb}{0.000000,0.000000,0.000000}%
\pgfsetstrokecolor{currentstroke}%
\pgfsetdash{}{0pt}%
\pgfsys@defobject{currentmarker}{\pgfqpoint{-0.027778in}{0.000000in}}{\pgfqpoint{-0.000000in}{0.000000in}}{%
\pgfpathmoveto{\pgfqpoint{-0.000000in}{0.000000in}}%
\pgfpathlineto{\pgfqpoint{-0.027778in}{0.000000in}}%
\pgfusepath{stroke,fill}%
}%
\begin{pgfscope}%
\pgfsys@transformshift{1.875000in}{1.397130in}%
\pgfsys@useobject{currentmarker}{}%
\end{pgfscope}%
\end{pgfscope}%
\begin{pgfscope}%
\pgfsetbuttcap%
\pgfsetroundjoin%
\definecolor{currentfill}{rgb}{0.000000,0.000000,0.000000}%
\pgfsetfillcolor{currentfill}%
\pgfsetlinewidth{0.602250pt}%
\definecolor{currentstroke}{rgb}{0.000000,0.000000,0.000000}%
\pgfsetstrokecolor{currentstroke}%
\pgfsetdash{}{0pt}%
\pgfsys@defobject{currentmarker}{\pgfqpoint{-0.027778in}{0.000000in}}{\pgfqpoint{-0.000000in}{0.000000in}}{%
\pgfpathmoveto{\pgfqpoint{-0.000000in}{0.000000in}}%
\pgfpathlineto{\pgfqpoint{-0.027778in}{0.000000in}}%
\pgfusepath{stroke,fill}%
}%
\begin{pgfscope}%
\pgfsys@transformshift{1.875000in}{1.466348in}%
\pgfsys@useobject{currentmarker}{}%
\end{pgfscope}%
\end{pgfscope}%
\begin{pgfscope}%
\pgfsetbuttcap%
\pgfsetroundjoin%
\definecolor{currentfill}{rgb}{0.000000,0.000000,0.000000}%
\pgfsetfillcolor{currentfill}%
\pgfsetlinewidth{0.602250pt}%
\definecolor{currentstroke}{rgb}{0.000000,0.000000,0.000000}%
\pgfsetstrokecolor{currentstroke}%
\pgfsetdash{}{0pt}%
\pgfsys@defobject{currentmarker}{\pgfqpoint{-0.027778in}{0.000000in}}{\pgfqpoint{-0.000000in}{0.000000in}}{%
\pgfpathmoveto{\pgfqpoint{-0.000000in}{0.000000in}}%
\pgfpathlineto{\pgfqpoint{-0.027778in}{0.000000in}}%
\pgfusepath{stroke,fill}%
}%
\begin{pgfscope}%
\pgfsys@transformshift{1.875000in}{1.524872in}%
\pgfsys@useobject{currentmarker}{}%
\end{pgfscope}%
\end{pgfscope}%
\begin{pgfscope}%
\pgfsetbuttcap%
\pgfsetroundjoin%
\definecolor{currentfill}{rgb}{0.000000,0.000000,0.000000}%
\pgfsetfillcolor{currentfill}%
\pgfsetlinewidth{0.602250pt}%
\definecolor{currentstroke}{rgb}{0.000000,0.000000,0.000000}%
\pgfsetstrokecolor{currentstroke}%
\pgfsetdash{}{0pt}%
\pgfsys@defobject{currentmarker}{\pgfqpoint{-0.027778in}{0.000000in}}{\pgfqpoint{-0.000000in}{0.000000in}}{%
\pgfpathmoveto{\pgfqpoint{-0.000000in}{0.000000in}}%
\pgfpathlineto{\pgfqpoint{-0.027778in}{0.000000in}}%
\pgfusepath{stroke,fill}%
}%
\begin{pgfscope}%
\pgfsys@transformshift{1.875000in}{1.575567in}%
\pgfsys@useobject{currentmarker}{}%
\end{pgfscope}%
\end{pgfscope}%
\begin{pgfscope}%
\pgfsetbuttcap%
\pgfsetroundjoin%
\definecolor{currentfill}{rgb}{0.000000,0.000000,0.000000}%
\pgfsetfillcolor{currentfill}%
\pgfsetlinewidth{0.602250pt}%
\definecolor{currentstroke}{rgb}{0.000000,0.000000,0.000000}%
\pgfsetstrokecolor{currentstroke}%
\pgfsetdash{}{0pt}%
\pgfsys@defobject{currentmarker}{\pgfqpoint{-0.027778in}{0.000000in}}{\pgfqpoint{-0.000000in}{0.000000in}}{%
\pgfpathmoveto{\pgfqpoint{-0.000000in}{0.000000in}}%
\pgfpathlineto{\pgfqpoint{-0.027778in}{0.000000in}}%
\pgfusepath{stroke,fill}%
}%
\begin{pgfscope}%
\pgfsys@transformshift{1.875000in}{1.620284in}%
\pgfsys@useobject{currentmarker}{}%
\end{pgfscope}%
\end{pgfscope}%
\begin{pgfscope}%
\pgfsetbuttcap%
\pgfsetroundjoin%
\definecolor{currentfill}{rgb}{0.000000,0.000000,0.000000}%
\pgfsetfillcolor{currentfill}%
\pgfsetlinewidth{0.602250pt}%
\definecolor{currentstroke}{rgb}{0.000000,0.000000,0.000000}%
\pgfsetstrokecolor{currentstroke}%
\pgfsetdash{}{0pt}%
\pgfsys@defobject{currentmarker}{\pgfqpoint{-0.027778in}{0.000000in}}{\pgfqpoint{-0.000000in}{0.000000in}}{%
\pgfpathmoveto{\pgfqpoint{-0.000000in}{0.000000in}}%
\pgfpathlineto{\pgfqpoint{-0.027778in}{0.000000in}}%
\pgfusepath{stroke,fill}%
}%
\begin{pgfscope}%
\pgfsys@transformshift{1.875000in}{1.923438in}%
\pgfsys@useobject{currentmarker}{}%
\end{pgfscope}%
\end{pgfscope}%
\begin{pgfscope}%
\pgfsetbuttcap%
\pgfsetroundjoin%
\definecolor{currentfill}{rgb}{0.000000,0.000000,0.000000}%
\pgfsetfillcolor{currentfill}%
\pgfsetlinewidth{0.602250pt}%
\definecolor{currentstroke}{rgb}{0.000000,0.000000,0.000000}%
\pgfsetstrokecolor{currentstroke}%
\pgfsetdash{}{0pt}%
\pgfsys@defobject{currentmarker}{\pgfqpoint{-0.027778in}{0.000000in}}{\pgfqpoint{-0.000000in}{0.000000in}}{%
\pgfpathmoveto{\pgfqpoint{-0.000000in}{0.000000in}}%
\pgfpathlineto{\pgfqpoint{-0.027778in}{0.000000in}}%
\pgfusepath{stroke,fill}%
}%
\begin{pgfscope}%
\pgfsys@transformshift{1.875000in}{2.077373in}%
\pgfsys@useobject{currentmarker}{}%
\end{pgfscope}%
\end{pgfscope}%
\begin{pgfscope}%
\pgfsetbuttcap%
\pgfsetroundjoin%
\definecolor{currentfill}{rgb}{0.000000,0.000000,0.000000}%
\pgfsetfillcolor{currentfill}%
\pgfsetlinewidth{0.602250pt}%
\definecolor{currentstroke}{rgb}{0.000000,0.000000,0.000000}%
\pgfsetstrokecolor{currentstroke}%
\pgfsetdash{}{0pt}%
\pgfsys@defobject{currentmarker}{\pgfqpoint{-0.027778in}{0.000000in}}{\pgfqpoint{-0.000000in}{0.000000in}}{%
\pgfpathmoveto{\pgfqpoint{-0.000000in}{0.000000in}}%
\pgfpathlineto{\pgfqpoint{-0.027778in}{0.000000in}}%
\pgfusepath{stroke,fill}%
}%
\begin{pgfscope}%
\pgfsys@transformshift{1.875000in}{2.186592in}%
\pgfsys@useobject{currentmarker}{}%
\end{pgfscope}%
\end{pgfscope}%
\begin{pgfscope}%
\pgfsetbuttcap%
\pgfsetroundjoin%
\definecolor{currentfill}{rgb}{0.000000,0.000000,0.000000}%
\pgfsetfillcolor{currentfill}%
\pgfsetlinewidth{0.602250pt}%
\definecolor{currentstroke}{rgb}{0.000000,0.000000,0.000000}%
\pgfsetstrokecolor{currentstroke}%
\pgfsetdash{}{0pt}%
\pgfsys@defobject{currentmarker}{\pgfqpoint{-0.027778in}{0.000000in}}{\pgfqpoint{-0.000000in}{0.000000in}}{%
\pgfpathmoveto{\pgfqpoint{-0.000000in}{0.000000in}}%
\pgfpathlineto{\pgfqpoint{-0.027778in}{0.000000in}}%
\pgfusepath{stroke,fill}%
}%
\begin{pgfscope}%
\pgfsys@transformshift{1.875000in}{2.271309in}%
\pgfsys@useobject{currentmarker}{}%
\end{pgfscope}%
\end{pgfscope}%
\begin{pgfscope}%
\pgfsetbuttcap%
\pgfsetroundjoin%
\definecolor{currentfill}{rgb}{0.000000,0.000000,0.000000}%
\pgfsetfillcolor{currentfill}%
\pgfsetlinewidth{0.602250pt}%
\definecolor{currentstroke}{rgb}{0.000000,0.000000,0.000000}%
\pgfsetstrokecolor{currentstroke}%
\pgfsetdash{}{0pt}%
\pgfsys@defobject{currentmarker}{\pgfqpoint{-0.027778in}{0.000000in}}{\pgfqpoint{-0.000000in}{0.000000in}}{%
\pgfpathmoveto{\pgfqpoint{-0.000000in}{0.000000in}}%
\pgfpathlineto{\pgfqpoint{-0.027778in}{0.000000in}}%
\pgfusepath{stroke,fill}%
}%
\begin{pgfscope}%
\pgfsys@transformshift{1.875000in}{2.340528in}%
\pgfsys@useobject{currentmarker}{}%
\end{pgfscope}%
\end{pgfscope}%
\begin{pgfscope}%
\pgfsetbuttcap%
\pgfsetroundjoin%
\definecolor{currentfill}{rgb}{0.000000,0.000000,0.000000}%
\pgfsetfillcolor{currentfill}%
\pgfsetlinewidth{0.602250pt}%
\definecolor{currentstroke}{rgb}{0.000000,0.000000,0.000000}%
\pgfsetstrokecolor{currentstroke}%
\pgfsetdash{}{0pt}%
\pgfsys@defobject{currentmarker}{\pgfqpoint{-0.027778in}{0.000000in}}{\pgfqpoint{-0.000000in}{0.000000in}}{%
\pgfpathmoveto{\pgfqpoint{-0.000000in}{0.000000in}}%
\pgfpathlineto{\pgfqpoint{-0.027778in}{0.000000in}}%
\pgfusepath{stroke,fill}%
}%
\begin{pgfscope}%
\pgfsys@transformshift{1.875000in}{2.399051in}%
\pgfsys@useobject{currentmarker}{}%
\end{pgfscope}%
\end{pgfscope}%
\begin{pgfscope}%
\pgfsetbuttcap%
\pgfsetroundjoin%
\definecolor{currentfill}{rgb}{0.000000,0.000000,0.000000}%
\pgfsetfillcolor{currentfill}%
\pgfsetlinewidth{0.602250pt}%
\definecolor{currentstroke}{rgb}{0.000000,0.000000,0.000000}%
\pgfsetstrokecolor{currentstroke}%
\pgfsetdash{}{0pt}%
\pgfsys@defobject{currentmarker}{\pgfqpoint{-0.027778in}{0.000000in}}{\pgfqpoint{-0.000000in}{0.000000in}}{%
\pgfpathmoveto{\pgfqpoint{-0.000000in}{0.000000in}}%
\pgfpathlineto{\pgfqpoint{-0.027778in}{0.000000in}}%
\pgfusepath{stroke,fill}%
}%
\begin{pgfscope}%
\pgfsys@transformshift{1.875000in}{2.449746in}%
\pgfsys@useobject{currentmarker}{}%
\end{pgfscope}%
\end{pgfscope}%
\begin{pgfscope}%
\pgfsetbuttcap%
\pgfsetroundjoin%
\definecolor{currentfill}{rgb}{0.000000,0.000000,0.000000}%
\pgfsetfillcolor{currentfill}%
\pgfsetlinewidth{0.602250pt}%
\definecolor{currentstroke}{rgb}{0.000000,0.000000,0.000000}%
\pgfsetstrokecolor{currentstroke}%
\pgfsetdash{}{0pt}%
\pgfsys@defobject{currentmarker}{\pgfqpoint{-0.027778in}{0.000000in}}{\pgfqpoint{-0.000000in}{0.000000in}}{%
\pgfpathmoveto{\pgfqpoint{-0.000000in}{0.000000in}}%
\pgfpathlineto{\pgfqpoint{-0.027778in}{0.000000in}}%
\pgfusepath{stroke,fill}%
}%
\begin{pgfscope}%
\pgfsys@transformshift{1.875000in}{2.494463in}%
\pgfsys@useobject{currentmarker}{}%
\end{pgfscope}%
\end{pgfscope}%
\begin{pgfscope}%
\pgfsetbuttcap%
\pgfsetroundjoin%
\definecolor{currentfill}{rgb}{0.000000,0.000000,0.000000}%
\pgfsetfillcolor{currentfill}%
\pgfsetlinewidth{0.602250pt}%
\definecolor{currentstroke}{rgb}{0.000000,0.000000,0.000000}%
\pgfsetstrokecolor{currentstroke}%
\pgfsetdash{}{0pt}%
\pgfsys@defobject{currentmarker}{\pgfqpoint{-0.027778in}{0.000000in}}{\pgfqpoint{-0.000000in}{0.000000in}}{%
\pgfpathmoveto{\pgfqpoint{-0.000000in}{0.000000in}}%
\pgfpathlineto{\pgfqpoint{-0.027778in}{0.000000in}}%
\pgfusepath{stroke,fill}%
}%
\begin{pgfscope}%
\pgfsys@transformshift{1.875000in}{2.797617in}%
\pgfsys@useobject{currentmarker}{}%
\end{pgfscope}%
\end{pgfscope}%
\begin{pgfscope}%
\pgfsetbuttcap%
\pgfsetroundjoin%
\definecolor{currentfill}{rgb}{0.000000,0.000000,0.000000}%
\pgfsetfillcolor{currentfill}%
\pgfsetlinewidth{0.602250pt}%
\definecolor{currentstroke}{rgb}{0.000000,0.000000,0.000000}%
\pgfsetstrokecolor{currentstroke}%
\pgfsetdash{}{0pt}%
\pgfsys@defobject{currentmarker}{\pgfqpoint{-0.027778in}{0.000000in}}{\pgfqpoint{-0.000000in}{0.000000in}}{%
\pgfpathmoveto{\pgfqpoint{-0.000000in}{0.000000in}}%
\pgfpathlineto{\pgfqpoint{-0.027778in}{0.000000in}}%
\pgfusepath{stroke,fill}%
}%
\begin{pgfscope}%
\pgfsys@transformshift{1.875000in}{2.951553in}%
\pgfsys@useobject{currentmarker}{}%
\end{pgfscope}%
\end{pgfscope}%
\begin{pgfscope}%
\pgfsetbuttcap%
\pgfsetroundjoin%
\definecolor{currentfill}{rgb}{0.000000,0.000000,0.000000}%
\pgfsetfillcolor{currentfill}%
\pgfsetlinewidth{0.602250pt}%
\definecolor{currentstroke}{rgb}{0.000000,0.000000,0.000000}%
\pgfsetstrokecolor{currentstroke}%
\pgfsetdash{}{0pt}%
\pgfsys@defobject{currentmarker}{\pgfqpoint{-0.027778in}{0.000000in}}{\pgfqpoint{-0.000000in}{0.000000in}}{%
\pgfpathmoveto{\pgfqpoint{-0.000000in}{0.000000in}}%
\pgfpathlineto{\pgfqpoint{-0.027778in}{0.000000in}}%
\pgfusepath{stroke,fill}%
}%
\begin{pgfscope}%
\pgfsys@transformshift{1.875000in}{3.060771in}%
\pgfsys@useobject{currentmarker}{}%
\end{pgfscope}%
\end{pgfscope}%
\begin{pgfscope}%
\pgfsetbuttcap%
\pgfsetroundjoin%
\definecolor{currentfill}{rgb}{0.000000,0.000000,0.000000}%
\pgfsetfillcolor{currentfill}%
\pgfsetlinewidth{0.602250pt}%
\definecolor{currentstroke}{rgb}{0.000000,0.000000,0.000000}%
\pgfsetstrokecolor{currentstroke}%
\pgfsetdash{}{0pt}%
\pgfsys@defobject{currentmarker}{\pgfqpoint{-0.027778in}{0.000000in}}{\pgfqpoint{-0.000000in}{0.000000in}}{%
\pgfpathmoveto{\pgfqpoint{-0.000000in}{0.000000in}}%
\pgfpathlineto{\pgfqpoint{-0.027778in}{0.000000in}}%
\pgfusepath{stroke,fill}%
}%
\begin{pgfscope}%
\pgfsys@transformshift{1.875000in}{3.145488in}%
\pgfsys@useobject{currentmarker}{}%
\end{pgfscope}%
\end{pgfscope}%
\begin{pgfscope}%
\pgfsetbuttcap%
\pgfsetroundjoin%
\definecolor{currentfill}{rgb}{0.000000,0.000000,0.000000}%
\pgfsetfillcolor{currentfill}%
\pgfsetlinewidth{0.602250pt}%
\definecolor{currentstroke}{rgb}{0.000000,0.000000,0.000000}%
\pgfsetstrokecolor{currentstroke}%
\pgfsetdash{}{0pt}%
\pgfsys@defobject{currentmarker}{\pgfqpoint{-0.027778in}{0.000000in}}{\pgfqpoint{-0.000000in}{0.000000in}}{%
\pgfpathmoveto{\pgfqpoint{-0.000000in}{0.000000in}}%
\pgfpathlineto{\pgfqpoint{-0.027778in}{0.000000in}}%
\pgfusepath{stroke,fill}%
}%
\begin{pgfscope}%
\pgfsys@transformshift{1.875000in}{3.214707in}%
\pgfsys@useobject{currentmarker}{}%
\end{pgfscope}%
\end{pgfscope}%
\begin{pgfscope}%
\pgfsetbuttcap%
\pgfsetroundjoin%
\definecolor{currentfill}{rgb}{0.000000,0.000000,0.000000}%
\pgfsetfillcolor{currentfill}%
\pgfsetlinewidth{0.602250pt}%
\definecolor{currentstroke}{rgb}{0.000000,0.000000,0.000000}%
\pgfsetstrokecolor{currentstroke}%
\pgfsetdash{}{0pt}%
\pgfsys@defobject{currentmarker}{\pgfqpoint{-0.027778in}{0.000000in}}{\pgfqpoint{-0.000000in}{0.000000in}}{%
\pgfpathmoveto{\pgfqpoint{-0.000000in}{0.000000in}}%
\pgfpathlineto{\pgfqpoint{-0.027778in}{0.000000in}}%
\pgfusepath{stroke,fill}%
}%
\begin{pgfscope}%
\pgfsys@transformshift{1.875000in}{3.273230in}%
\pgfsys@useobject{currentmarker}{}%
\end{pgfscope}%
\end{pgfscope}%
\begin{pgfscope}%
\pgfsetbuttcap%
\pgfsetroundjoin%
\definecolor{currentfill}{rgb}{0.000000,0.000000,0.000000}%
\pgfsetfillcolor{currentfill}%
\pgfsetlinewidth{0.602250pt}%
\definecolor{currentstroke}{rgb}{0.000000,0.000000,0.000000}%
\pgfsetstrokecolor{currentstroke}%
\pgfsetdash{}{0pt}%
\pgfsys@defobject{currentmarker}{\pgfqpoint{-0.027778in}{0.000000in}}{\pgfqpoint{-0.000000in}{0.000000in}}{%
\pgfpathmoveto{\pgfqpoint{-0.000000in}{0.000000in}}%
\pgfpathlineto{\pgfqpoint{-0.027778in}{0.000000in}}%
\pgfusepath{stroke,fill}%
}%
\begin{pgfscope}%
\pgfsys@transformshift{1.875000in}{3.323926in}%
\pgfsys@useobject{currentmarker}{}%
\end{pgfscope}%
\end{pgfscope}%
\begin{pgfscope}%
\pgfsetbuttcap%
\pgfsetroundjoin%
\definecolor{currentfill}{rgb}{0.000000,0.000000,0.000000}%
\pgfsetfillcolor{currentfill}%
\pgfsetlinewidth{0.602250pt}%
\definecolor{currentstroke}{rgb}{0.000000,0.000000,0.000000}%
\pgfsetstrokecolor{currentstroke}%
\pgfsetdash{}{0pt}%
\pgfsys@defobject{currentmarker}{\pgfqpoint{-0.027778in}{0.000000in}}{\pgfqpoint{-0.000000in}{0.000000in}}{%
\pgfpathmoveto{\pgfqpoint{-0.000000in}{0.000000in}}%
\pgfpathlineto{\pgfqpoint{-0.027778in}{0.000000in}}%
\pgfusepath{stroke,fill}%
}%
\begin{pgfscope}%
\pgfsys@transformshift{1.875000in}{3.368642in}%
\pgfsys@useobject{currentmarker}{}%
\end{pgfscope}%
\end{pgfscope}%
\begin{pgfscope}%
\pgfsetbuttcap%
\pgfsetroundjoin%
\definecolor{currentfill}{rgb}{0.000000,0.000000,0.000000}%
\pgfsetfillcolor{currentfill}%
\pgfsetlinewidth{0.602250pt}%
\definecolor{currentstroke}{rgb}{0.000000,0.000000,0.000000}%
\pgfsetstrokecolor{currentstroke}%
\pgfsetdash{}{0pt}%
\pgfsys@defobject{currentmarker}{\pgfqpoint{-0.027778in}{0.000000in}}{\pgfqpoint{-0.000000in}{0.000000in}}{%
\pgfpathmoveto{\pgfqpoint{-0.000000in}{0.000000in}}%
\pgfpathlineto{\pgfqpoint{-0.027778in}{0.000000in}}%
\pgfusepath{stroke,fill}%
}%
\begin{pgfscope}%
\pgfsys@transformshift{1.875000in}{3.671797in}%
\pgfsys@useobject{currentmarker}{}%
\end{pgfscope}%
\end{pgfscope}%
\begin{pgfscope}%
\pgfsetbuttcap%
\pgfsetroundjoin%
\definecolor{currentfill}{rgb}{0.000000,0.000000,0.000000}%
\pgfsetfillcolor{currentfill}%
\pgfsetlinewidth{0.602250pt}%
\definecolor{currentstroke}{rgb}{0.000000,0.000000,0.000000}%
\pgfsetstrokecolor{currentstroke}%
\pgfsetdash{}{0pt}%
\pgfsys@defobject{currentmarker}{\pgfqpoint{-0.027778in}{0.000000in}}{\pgfqpoint{-0.000000in}{0.000000in}}{%
\pgfpathmoveto{\pgfqpoint{-0.000000in}{0.000000in}}%
\pgfpathlineto{\pgfqpoint{-0.027778in}{0.000000in}}%
\pgfusepath{stroke,fill}%
}%
\begin{pgfscope}%
\pgfsys@transformshift{1.875000in}{3.825732in}%
\pgfsys@useobject{currentmarker}{}%
\end{pgfscope}%
\end{pgfscope}%
\begin{pgfscope}%
\pgfsetbuttcap%
\pgfsetroundjoin%
\definecolor{currentfill}{rgb}{0.000000,0.000000,0.000000}%
\pgfsetfillcolor{currentfill}%
\pgfsetlinewidth{0.602250pt}%
\definecolor{currentstroke}{rgb}{0.000000,0.000000,0.000000}%
\pgfsetstrokecolor{currentstroke}%
\pgfsetdash{}{0pt}%
\pgfsys@defobject{currentmarker}{\pgfqpoint{-0.027778in}{0.000000in}}{\pgfqpoint{-0.000000in}{0.000000in}}{%
\pgfpathmoveto{\pgfqpoint{-0.000000in}{0.000000in}}%
\pgfpathlineto{\pgfqpoint{-0.027778in}{0.000000in}}%
\pgfusepath{stroke,fill}%
}%
\begin{pgfscope}%
\pgfsys@transformshift{1.875000in}{3.934951in}%
\pgfsys@useobject{currentmarker}{}%
\end{pgfscope}%
\end{pgfscope}%
\begin{pgfscope}%
\pgfsetbuttcap%
\pgfsetroundjoin%
\definecolor{currentfill}{rgb}{0.000000,0.000000,0.000000}%
\pgfsetfillcolor{currentfill}%
\pgfsetlinewidth{0.602250pt}%
\definecolor{currentstroke}{rgb}{0.000000,0.000000,0.000000}%
\pgfsetstrokecolor{currentstroke}%
\pgfsetdash{}{0pt}%
\pgfsys@defobject{currentmarker}{\pgfqpoint{-0.027778in}{0.000000in}}{\pgfqpoint{-0.000000in}{0.000000in}}{%
\pgfpathmoveto{\pgfqpoint{-0.000000in}{0.000000in}}%
\pgfpathlineto{\pgfqpoint{-0.027778in}{0.000000in}}%
\pgfusepath{stroke,fill}%
}%
\begin{pgfscope}%
\pgfsys@transformshift{1.875000in}{4.019667in}%
\pgfsys@useobject{currentmarker}{}%
\end{pgfscope}%
\end{pgfscope}%
\begin{pgfscope}%
\pgfsetbuttcap%
\pgfsetroundjoin%
\definecolor{currentfill}{rgb}{0.000000,0.000000,0.000000}%
\pgfsetfillcolor{currentfill}%
\pgfsetlinewidth{0.602250pt}%
\definecolor{currentstroke}{rgb}{0.000000,0.000000,0.000000}%
\pgfsetstrokecolor{currentstroke}%
\pgfsetdash{}{0pt}%
\pgfsys@defobject{currentmarker}{\pgfqpoint{-0.027778in}{0.000000in}}{\pgfqpoint{-0.000000in}{0.000000in}}{%
\pgfpathmoveto{\pgfqpoint{-0.000000in}{0.000000in}}%
\pgfpathlineto{\pgfqpoint{-0.027778in}{0.000000in}}%
\pgfusepath{stroke,fill}%
}%
\begin{pgfscope}%
\pgfsys@transformshift{1.875000in}{4.088886in}%
\pgfsys@useobject{currentmarker}{}%
\end{pgfscope}%
\end{pgfscope}%
\begin{pgfscope}%
\pgfsetbuttcap%
\pgfsetroundjoin%
\definecolor{currentfill}{rgb}{0.000000,0.000000,0.000000}%
\pgfsetfillcolor{currentfill}%
\pgfsetlinewidth{0.602250pt}%
\definecolor{currentstroke}{rgb}{0.000000,0.000000,0.000000}%
\pgfsetstrokecolor{currentstroke}%
\pgfsetdash{}{0pt}%
\pgfsys@defobject{currentmarker}{\pgfqpoint{-0.027778in}{0.000000in}}{\pgfqpoint{-0.000000in}{0.000000in}}{%
\pgfpathmoveto{\pgfqpoint{-0.000000in}{0.000000in}}%
\pgfpathlineto{\pgfqpoint{-0.027778in}{0.000000in}}%
\pgfusepath{stroke,fill}%
}%
\begin{pgfscope}%
\pgfsys@transformshift{1.875000in}{4.147409in}%
\pgfsys@useobject{currentmarker}{}%
\end{pgfscope}%
\end{pgfscope}%
\begin{pgfscope}%
\pgfsetbuttcap%
\pgfsetroundjoin%
\definecolor{currentfill}{rgb}{0.000000,0.000000,0.000000}%
\pgfsetfillcolor{currentfill}%
\pgfsetlinewidth{0.602250pt}%
\definecolor{currentstroke}{rgb}{0.000000,0.000000,0.000000}%
\pgfsetstrokecolor{currentstroke}%
\pgfsetdash{}{0pt}%
\pgfsys@defobject{currentmarker}{\pgfqpoint{-0.027778in}{0.000000in}}{\pgfqpoint{-0.000000in}{0.000000in}}{%
\pgfpathmoveto{\pgfqpoint{-0.000000in}{0.000000in}}%
\pgfpathlineto{\pgfqpoint{-0.027778in}{0.000000in}}%
\pgfusepath{stroke,fill}%
}%
\begin{pgfscope}%
\pgfsys@transformshift{1.875000in}{4.198105in}%
\pgfsys@useobject{currentmarker}{}%
\end{pgfscope}%
\end{pgfscope}%
\begin{pgfscope}%
\pgfsetbuttcap%
\pgfsetroundjoin%
\definecolor{currentfill}{rgb}{0.000000,0.000000,0.000000}%
\pgfsetfillcolor{currentfill}%
\pgfsetlinewidth{0.602250pt}%
\definecolor{currentstroke}{rgb}{0.000000,0.000000,0.000000}%
\pgfsetstrokecolor{currentstroke}%
\pgfsetdash{}{0pt}%
\pgfsys@defobject{currentmarker}{\pgfqpoint{-0.027778in}{0.000000in}}{\pgfqpoint{-0.000000in}{0.000000in}}{%
\pgfpathmoveto{\pgfqpoint{-0.000000in}{0.000000in}}%
\pgfpathlineto{\pgfqpoint{-0.027778in}{0.000000in}}%
\pgfusepath{stroke,fill}%
}%
\begin{pgfscope}%
\pgfsys@transformshift{1.875000in}{4.242821in}%
\pgfsys@useobject{currentmarker}{}%
\end{pgfscope}%
\end{pgfscope}%
\begin{pgfscope}%
\pgfsetbuttcap%
\pgfsetroundjoin%
\definecolor{currentfill}{rgb}{0.000000,0.000000,0.000000}%
\pgfsetfillcolor{currentfill}%
\pgfsetlinewidth{0.602250pt}%
\definecolor{currentstroke}{rgb}{0.000000,0.000000,0.000000}%
\pgfsetstrokecolor{currentstroke}%
\pgfsetdash{}{0pt}%
\pgfsys@defobject{currentmarker}{\pgfqpoint{-0.027778in}{0.000000in}}{\pgfqpoint{-0.000000in}{0.000000in}}{%
\pgfpathmoveto{\pgfqpoint{-0.000000in}{0.000000in}}%
\pgfpathlineto{\pgfqpoint{-0.027778in}{0.000000in}}%
\pgfusepath{stroke,fill}%
}%
\begin{pgfscope}%
\pgfsys@transformshift{1.875000in}{4.545976in}%
\pgfsys@useobject{currentmarker}{}%
\end{pgfscope}%
\end{pgfscope}%
\begin{pgfscope}%
\definecolor{textcolor}{rgb}{0.000000,0.000000,0.000000}%
\pgfsetstrokecolor{textcolor}%
\pgfsetfillcolor{textcolor}%
\pgftext[x=1.434220in,y=2.850000in,,bottom,rotate=90.000000]{\color{textcolor}{\sffamily\fontsize{10.000000}{12.000000}\selectfont\catcode`\^=\active\def^{\ifmmode\sp\else\^{}\fi}\catcode`\%=\active\def%{\%}Error}}%
\end{pgfscope}%
\begin{pgfscope}%
\pgfpathrectangle{\pgfqpoint{1.875000in}{1.100000in}}{\pgfqpoint{3.419118in}{3.500000in}}%
\pgfusepath{clip}%
\pgfsetrectcap%
\pgfsetroundjoin%
\pgfsetlinewidth{1.505625pt}%
\definecolor{currentstroke}{rgb}{0.121569,0.466667,0.705882}%
\pgfsetstrokecolor{currentstroke}%
\pgfsetdash{}{0pt}%
\pgfpathmoveto{\pgfqpoint{2.030414in}{4.312459in}}%
\pgfpathlineto{\pgfqpoint{2.155162in}{4.268436in}}%
\pgfpathlineto{\pgfqpoint{2.394261in}{4.186737in}}%
\pgfpathlineto{\pgfqpoint{2.851668in}{4.033153in}}%
\pgfpathlineto{\pgfqpoint{3.756086in}{3.732208in}}%
\pgfpathlineto{\pgfqpoint{5.138703in}{3.273605in}}%
\pgfpathlineto{\pgfqpoint{5.138703in}{3.273605in}}%
\pgfusepath{stroke}%
\end{pgfscope}%
\begin{pgfscope}%
\pgfpathrectangle{\pgfqpoint{1.875000in}{1.100000in}}{\pgfqpoint{3.419118in}{3.500000in}}%
\pgfusepath{clip}%
\pgfsetrectcap%
\pgfsetroundjoin%
\pgfsetlinewidth{1.505625pt}%
\definecolor{currentstroke}{rgb}{1.000000,0.498039,0.054902}%
\pgfsetstrokecolor{currentstroke}%
\pgfsetdash{}{0pt}%
\pgfpathmoveto{\pgfqpoint{2.030414in}{4.312459in}}%
\pgfpathlineto{\pgfqpoint{2.155162in}{4.268436in}}%
\pgfpathlineto{\pgfqpoint{2.269514in}{4.229126in}}%
\pgfpathlineto{\pgfqpoint{2.279909in}{3.646209in}}%
\pgfpathlineto{\pgfqpoint{2.290305in}{3.617094in}}%
\pgfpathlineto{\pgfqpoint{2.300700in}{3.600715in}}%
\pgfpathlineto{\pgfqpoint{2.311096in}{3.587727in}}%
\pgfpathlineto{\pgfqpoint{2.331887in}{3.566806in}}%
\pgfpathlineto{\pgfqpoint{2.352678in}{3.549691in}}%
\pgfpathlineto{\pgfqpoint{2.383865in}{3.528033in}}%
\pgfpathlineto{\pgfqpoint{2.415052in}{3.509314in}}%
\pgfpathlineto{\pgfqpoint{2.456635in}{3.487170in}}%
\pgfpathlineto{\pgfqpoint{2.508613in}{3.462341in}}%
\pgfpathlineto{\pgfqpoint{2.529404in}{3.453024in}}%
\pgfpathlineto{\pgfqpoint{2.539800in}{3.629222in}}%
\pgfpathlineto{\pgfqpoint{2.550195in}{3.614584in}}%
\pgfpathlineto{\pgfqpoint{2.570986in}{3.598602in}}%
\pgfpathlineto{\pgfqpoint{2.602173in}{3.579675in}}%
\pgfpathlineto{\pgfqpoint{2.643756in}{3.557899in}}%
\pgfpathlineto{\pgfqpoint{2.695734in}{3.533515in}}%
\pgfpathlineto{\pgfqpoint{2.758107in}{3.506714in}}%
\pgfpathlineto{\pgfqpoint{2.789294in}{3.494004in}}%
\pgfpathlineto{\pgfqpoint{2.799690in}{3.625267in}}%
\pgfpathlineto{\pgfqpoint{2.810086in}{3.612216in}}%
\pgfpathlineto{\pgfqpoint{2.830877in}{3.596997in}}%
\pgfpathlineto{\pgfqpoint{2.862064in}{3.578441in}}%
\pgfpathlineto{\pgfqpoint{2.903646in}{3.556789in}}%
\pgfpathlineto{\pgfqpoint{2.955624in}{3.532390in}}%
\pgfpathlineto{\pgfqpoint{3.017998in}{3.505514in}}%
\pgfpathlineto{\pgfqpoint{3.049185in}{3.492765in}}%
\pgfpathlineto{\pgfqpoint{3.059580in}{3.624086in}}%
\pgfpathlineto{\pgfqpoint{3.069976in}{3.611401in}}%
\pgfpathlineto{\pgfqpoint{3.090767in}{3.596386in}}%
\pgfpathlineto{\pgfqpoint{3.121954in}{3.577947in}}%
\pgfpathlineto{\pgfqpoint{3.163536in}{3.556349in}}%
\pgfpathlineto{\pgfqpoint{3.215515in}{3.531960in}}%
\pgfpathlineto{\pgfqpoint{3.277888in}{3.505072in}}%
\pgfpathlineto{\pgfqpoint{3.309075in}{3.492316in}}%
\pgfpathlineto{\pgfqpoint{3.319471in}{3.623733in}}%
\pgfpathlineto{\pgfqpoint{3.329866in}{3.611164in}}%
\pgfpathlineto{\pgfqpoint{3.350657in}{3.596230in}}%
\pgfpathlineto{\pgfqpoint{3.381844in}{3.577856in}}%
\pgfpathlineto{\pgfqpoint{3.423427in}{3.556308in}}%
\pgfpathlineto{\pgfqpoint{3.475405in}{3.531957in}}%
\pgfpathlineto{\pgfqpoint{3.537779in}{3.505099in}}%
\pgfpathlineto{\pgfqpoint{3.568965in}{3.492355in}}%
\pgfpathlineto{\pgfqpoint{3.579361in}{3.623676in}}%
\pgfpathlineto{\pgfqpoint{3.589757in}{3.611165in}}%
\pgfpathlineto{\pgfqpoint{3.610548in}{3.596280in}}%
\pgfpathlineto{\pgfqpoint{3.641735in}{3.577954in}}%
\pgfpathlineto{\pgfqpoint{3.683317in}{3.556452in}}%
\pgfpathlineto{\pgfqpoint{3.735295in}{3.532144in}}%
\pgfpathlineto{\pgfqpoint{3.797669in}{3.505324in}}%
\pgfpathlineto{\pgfqpoint{3.828856in}{3.492596in}}%
\pgfpathlineto{\pgfqpoint{3.839251in}{3.623769in}}%
\pgfpathlineto{\pgfqpoint{3.849647in}{3.611298in}}%
\pgfpathlineto{\pgfqpoint{3.870438in}{3.596450in}}%
\pgfpathlineto{\pgfqpoint{3.901625in}{3.578164in}}%
\pgfpathlineto{\pgfqpoint{3.943208in}{3.556702in}}%
\pgfpathlineto{\pgfqpoint{3.995186in}{3.532432in}}%
\pgfpathlineto{\pgfqpoint{4.057559in}{3.505648in}}%
\pgfpathlineto{\pgfqpoint{4.088746in}{3.492934in}}%
\pgfpathlineto{\pgfqpoint{4.099142in}{3.623913in}}%
\pgfpathlineto{\pgfqpoint{4.109537in}{3.611474in}}%
\pgfpathlineto{\pgfqpoint{4.130329in}{3.596658in}}%
\pgfpathlineto{\pgfqpoint{4.161515in}{3.578406in}}%
\pgfpathlineto{\pgfqpoint{4.203098in}{3.556979in}}%
\pgfpathlineto{\pgfqpoint{4.255076in}{3.532743in}}%
\pgfpathlineto{\pgfqpoint{4.317450in}{3.505989in}}%
\pgfpathlineto{\pgfqpoint{4.348636in}{3.493287in}}%
\pgfpathlineto{\pgfqpoint{4.359032in}{3.624083in}}%
\pgfpathlineto{\pgfqpoint{4.369428in}{3.611672in}}%
\pgfpathlineto{\pgfqpoint{4.390219in}{3.596884in}}%
\pgfpathlineto{\pgfqpoint{4.421406in}{3.578661in}}%
\pgfpathlineto{\pgfqpoint{4.462988in}{3.557264in}}%
\pgfpathlineto{\pgfqpoint{4.514966in}{3.533057in}}%
\pgfpathlineto{\pgfqpoint{4.577340in}{3.506330in}}%
\pgfpathlineto{\pgfqpoint{4.608527in}{3.493638in}}%
\pgfpathlineto{\pgfqpoint{4.618922in}{3.624248in}}%
\pgfpathlineto{\pgfqpoint{4.629318in}{3.611863in}}%
\pgfpathlineto{\pgfqpoint{4.650109in}{3.597098in}}%
\pgfpathlineto{\pgfqpoint{4.681296in}{3.578901in}}%
\pgfpathlineto{\pgfqpoint{4.722879in}{3.557530in}}%
\pgfpathlineto{\pgfqpoint{4.774857in}{3.533348in}}%
\pgfpathlineto{\pgfqpoint{4.837230in}{3.506644in}}%
\pgfpathlineto{\pgfqpoint{4.868417in}{3.493962in}}%
\pgfpathlineto{\pgfqpoint{4.878813in}{3.624409in}}%
\pgfpathlineto{\pgfqpoint{4.889208in}{3.612046in}}%
\pgfpathlineto{\pgfqpoint{4.910000in}{3.597303in}}%
\pgfpathlineto{\pgfqpoint{4.941187in}{3.579128in}}%
\pgfpathlineto{\pgfqpoint{4.982769in}{3.557780in}}%
\pgfpathlineto{\pgfqpoint{5.034747in}{3.533621in}}%
\pgfpathlineto{\pgfqpoint{5.107516in}{3.502670in}}%
\pgfpathlineto{\pgfqpoint{5.128308in}{3.494263in}}%
\pgfpathlineto{\pgfqpoint{5.138703in}{3.624563in}}%
\pgfpathlineto{\pgfqpoint{5.138703in}{3.624563in}}%
\pgfusepath{stroke}%
\end{pgfscope}%
\begin{pgfscope}%
\pgfpathrectangle{\pgfqpoint{1.875000in}{1.100000in}}{\pgfqpoint{3.419118in}{3.500000in}}%
\pgfusepath{clip}%
\pgfsetrectcap%
\pgfsetroundjoin%
\pgfsetlinewidth{1.505625pt}%
\definecolor{currentstroke}{rgb}{0.172549,0.627451,0.172549}%
\pgfsetstrokecolor{currentstroke}%
\pgfsetdash{}{0pt}%
\pgfpathmoveto{\pgfqpoint{2.030414in}{3.835976in}}%
\pgfpathlineto{\pgfqpoint{2.040810in}{3.791428in}}%
\pgfpathlineto{\pgfqpoint{2.051206in}{3.765058in}}%
\pgfpathlineto{\pgfqpoint{2.061601in}{3.743744in}}%
\pgfpathlineto{\pgfqpoint{2.082393in}{3.709043in}}%
\pgfpathlineto{\pgfqpoint{2.103184in}{3.680667in}}%
\pgfpathlineto{\pgfqpoint{2.123975in}{3.656356in}}%
\pgfpathlineto{\pgfqpoint{2.144766in}{3.634975in}}%
\pgfpathlineto{\pgfqpoint{2.175953in}{3.606961in}}%
\pgfpathlineto{\pgfqpoint{2.207140in}{3.582579in}}%
\pgfpathlineto{\pgfqpoint{2.238327in}{3.560895in}}%
\pgfpathlineto{\pgfqpoint{2.279909in}{3.535097in}}%
\pgfpathlineto{\pgfqpoint{2.321492in}{3.511915in}}%
\pgfpathlineto{\pgfqpoint{2.373470in}{3.485555in}}%
\pgfpathlineto{\pgfqpoint{2.435843in}{3.456600in}}%
\pgfpathlineto{\pgfqpoint{2.508613in}{3.425294in}}%
\pgfpathlineto{\pgfqpoint{2.602173in}{3.387569in}}%
\pgfpathlineto{\pgfqpoint{2.726921in}{3.339965in}}%
\pgfpathlineto{\pgfqpoint{2.882855in}{3.282962in}}%
\pgfpathlineto{\pgfqpoint{3.101163in}{3.205727in}}%
\pgfpathlineto{\pgfqpoint{3.392240in}{3.105285in}}%
\pgfpathlineto{\pgfqpoint{3.797669in}{2.967931in}}%
\pgfpathlineto{\pgfqpoint{4.400615in}{2.766234in}}%
\pgfpathlineto{\pgfqpoint{5.138703in}{2.520940in}}%
\pgfpathlineto{\pgfqpoint{5.138703in}{2.520940in}}%
\pgfusepath{stroke}%
\end{pgfscope}%
\begin{pgfscope}%
\pgfsetrectcap%
\pgfsetmiterjoin%
\pgfsetlinewidth{0.803000pt}%
\definecolor{currentstroke}{rgb}{0.000000,0.000000,0.000000}%
\pgfsetstrokecolor{currentstroke}%
\pgfsetdash{}{0pt}%
\pgfpathmoveto{\pgfqpoint{1.875000in}{1.100000in}}%
\pgfpathlineto{\pgfqpoint{1.875000in}{4.600000in}}%
\pgfusepath{stroke}%
\end{pgfscope}%
\begin{pgfscope}%
\pgfsetrectcap%
\pgfsetmiterjoin%
\pgfsetlinewidth{0.803000pt}%
\definecolor{currentstroke}{rgb}{0.000000,0.000000,0.000000}%
\pgfsetstrokecolor{currentstroke}%
\pgfsetdash{}{0pt}%
\pgfpathmoveto{\pgfqpoint{5.294118in}{1.100000in}}%
\pgfpathlineto{\pgfqpoint{5.294118in}{4.600000in}}%
\pgfusepath{stroke}%
\end{pgfscope}%
\begin{pgfscope}%
\pgfsetrectcap%
\pgfsetmiterjoin%
\pgfsetlinewidth{0.803000pt}%
\definecolor{currentstroke}{rgb}{0.000000,0.000000,0.000000}%
\pgfsetstrokecolor{currentstroke}%
\pgfsetdash{}{0pt}%
\pgfpathmoveto{\pgfqpoint{1.875000in}{1.100000in}}%
\pgfpathlineto{\pgfqpoint{5.294118in}{1.100000in}}%
\pgfusepath{stroke}%
\end{pgfscope}%
\begin{pgfscope}%
\pgfsetrectcap%
\pgfsetmiterjoin%
\pgfsetlinewidth{0.803000pt}%
\definecolor{currentstroke}{rgb}{0.000000,0.000000,0.000000}%
\pgfsetstrokecolor{currentstroke}%
\pgfsetdash{}{0pt}%
\pgfpathmoveto{\pgfqpoint{1.875000in}{4.600000in}}%
\pgfpathlineto{\pgfqpoint{5.294118in}{4.600000in}}%
\pgfusepath{stroke}%
\end{pgfscope}%
\begin{pgfscope}%
\definecolor{textcolor}{rgb}{0.000000,0.000000,0.000000}%
\pgfsetstrokecolor{textcolor}%
\pgfsetfillcolor{textcolor}%
\pgftext[x=3.584559in,y=4.683333in,,base]{\color{textcolor}{\sffamily\fontsize{12.000000}{14.400000}\selectfont\catcode`\^=\active\def^{\ifmmode\sp\else\^{}\fi}\catcode`\%=\active\def%{\%}Jacobi}}%
\end{pgfscope}%
\begin{pgfscope}%
\pgfsetbuttcap%
\pgfsetmiterjoin%
\definecolor{currentfill}{rgb}{1.000000,1.000000,1.000000}%
\pgfsetfillcolor{currentfill}%
\pgfsetlinewidth{0.000000pt}%
\definecolor{currentstroke}{rgb}{0.000000,0.000000,0.000000}%
\pgfsetstrokecolor{currentstroke}%
\pgfsetstrokeopacity{0.000000}%
\pgfsetdash{}{0pt}%
\pgfpathmoveto{\pgfqpoint{5.977941in}{1.100000in}}%
\pgfpathlineto{\pgfqpoint{9.397059in}{1.100000in}}%
\pgfpathlineto{\pgfqpoint{9.397059in}{4.600000in}}%
\pgfpathlineto{\pgfqpoint{5.977941in}{4.600000in}}%
\pgfpathlineto{\pgfqpoint{5.977941in}{1.100000in}}%
\pgfpathclose%
\pgfusepath{fill}%
\end{pgfscope}%
\begin{pgfscope}%
\pgfsetbuttcap%
\pgfsetroundjoin%
\definecolor{currentfill}{rgb}{0.000000,0.000000,0.000000}%
\pgfsetfillcolor{currentfill}%
\pgfsetlinewidth{0.803000pt}%
\definecolor{currentstroke}{rgb}{0.000000,0.000000,0.000000}%
\pgfsetstrokecolor{currentstroke}%
\pgfsetdash{}{0pt}%
\pgfsys@defobject{currentmarker}{\pgfqpoint{0.000000in}{-0.048611in}}{\pgfqpoint{0.000000in}{0.000000in}}{%
\pgfpathmoveto{\pgfqpoint{0.000000in}{0.000000in}}%
\pgfpathlineto{\pgfqpoint{0.000000in}{-0.048611in}}%
\pgfusepath{stroke,fill}%
}%
\begin{pgfscope}%
\pgfsys@transformshift{6.133356in}{1.100000in}%
\pgfsys@useobject{currentmarker}{}%
\end{pgfscope}%
\end{pgfscope}%
\begin{pgfscope}%
\definecolor{textcolor}{rgb}{0.000000,0.000000,0.000000}%
\pgfsetstrokecolor{textcolor}%
\pgfsetfillcolor{textcolor}%
\pgftext[x=6.133356in,y=1.002778in,,top]{\color{textcolor}{\sffamily\fontsize{10.000000}{12.000000}\selectfont\catcode`\^=\active\def^{\ifmmode\sp\else\^{}\fi}\catcode`\%=\active\def%{\%}0}}%
\end{pgfscope}%
\begin{pgfscope}%
\pgfsetbuttcap%
\pgfsetroundjoin%
\definecolor{currentfill}{rgb}{0.000000,0.000000,0.000000}%
\pgfsetfillcolor{currentfill}%
\pgfsetlinewidth{0.803000pt}%
\definecolor{currentstroke}{rgb}{0.000000,0.000000,0.000000}%
\pgfsetstrokecolor{currentstroke}%
\pgfsetdash{}{0pt}%
\pgfsys@defobject{currentmarker}{\pgfqpoint{0.000000in}{-0.048611in}}{\pgfqpoint{0.000000in}{0.000000in}}{%
\pgfpathmoveto{\pgfqpoint{0.000000in}{0.000000in}}%
\pgfpathlineto{\pgfqpoint{0.000000in}{-0.048611in}}%
\pgfusepath{stroke,fill}%
}%
\begin{pgfscope}%
\pgfsys@transformshift{6.653136in}{1.100000in}%
\pgfsys@useobject{currentmarker}{}%
\end{pgfscope}%
\end{pgfscope}%
\begin{pgfscope}%
\definecolor{textcolor}{rgb}{0.000000,0.000000,0.000000}%
\pgfsetstrokecolor{textcolor}%
\pgfsetfillcolor{textcolor}%
\pgftext[x=6.653136in,y=1.002778in,,top]{\color{textcolor}{\sffamily\fontsize{10.000000}{12.000000}\selectfont\catcode`\^=\active\def^{\ifmmode\sp\else\^{}\fi}\catcode`\%=\active\def%{\%}50}}%
\end{pgfscope}%
\begin{pgfscope}%
\pgfsetbuttcap%
\pgfsetroundjoin%
\definecolor{currentfill}{rgb}{0.000000,0.000000,0.000000}%
\pgfsetfillcolor{currentfill}%
\pgfsetlinewidth{0.803000pt}%
\definecolor{currentstroke}{rgb}{0.000000,0.000000,0.000000}%
\pgfsetstrokecolor{currentstroke}%
\pgfsetdash{}{0pt}%
\pgfsys@defobject{currentmarker}{\pgfqpoint{0.000000in}{-0.048611in}}{\pgfqpoint{0.000000in}{0.000000in}}{%
\pgfpathmoveto{\pgfqpoint{0.000000in}{0.000000in}}%
\pgfpathlineto{\pgfqpoint{0.000000in}{-0.048611in}}%
\pgfusepath{stroke,fill}%
}%
\begin{pgfscope}%
\pgfsys@transformshift{7.172917in}{1.100000in}%
\pgfsys@useobject{currentmarker}{}%
\end{pgfscope}%
\end{pgfscope}%
\begin{pgfscope}%
\definecolor{textcolor}{rgb}{0.000000,0.000000,0.000000}%
\pgfsetstrokecolor{textcolor}%
\pgfsetfillcolor{textcolor}%
\pgftext[x=7.172917in,y=1.002778in,,top]{\color{textcolor}{\sffamily\fontsize{10.000000}{12.000000}\selectfont\catcode`\^=\active\def^{\ifmmode\sp\else\^{}\fi}\catcode`\%=\active\def%{\%}100}}%
\end{pgfscope}%
\begin{pgfscope}%
\pgfsetbuttcap%
\pgfsetroundjoin%
\definecolor{currentfill}{rgb}{0.000000,0.000000,0.000000}%
\pgfsetfillcolor{currentfill}%
\pgfsetlinewidth{0.803000pt}%
\definecolor{currentstroke}{rgb}{0.000000,0.000000,0.000000}%
\pgfsetstrokecolor{currentstroke}%
\pgfsetdash{}{0pt}%
\pgfsys@defobject{currentmarker}{\pgfqpoint{0.000000in}{-0.048611in}}{\pgfqpoint{0.000000in}{0.000000in}}{%
\pgfpathmoveto{\pgfqpoint{0.000000in}{0.000000in}}%
\pgfpathlineto{\pgfqpoint{0.000000in}{-0.048611in}}%
\pgfusepath{stroke,fill}%
}%
\begin{pgfscope}%
\pgfsys@transformshift{7.692698in}{1.100000in}%
\pgfsys@useobject{currentmarker}{}%
\end{pgfscope}%
\end{pgfscope}%
\begin{pgfscope}%
\definecolor{textcolor}{rgb}{0.000000,0.000000,0.000000}%
\pgfsetstrokecolor{textcolor}%
\pgfsetfillcolor{textcolor}%
\pgftext[x=7.692698in,y=1.002778in,,top]{\color{textcolor}{\sffamily\fontsize{10.000000}{12.000000}\selectfont\catcode`\^=\active\def^{\ifmmode\sp\else\^{}\fi}\catcode`\%=\active\def%{\%}150}}%
\end{pgfscope}%
\begin{pgfscope}%
\pgfsetbuttcap%
\pgfsetroundjoin%
\definecolor{currentfill}{rgb}{0.000000,0.000000,0.000000}%
\pgfsetfillcolor{currentfill}%
\pgfsetlinewidth{0.803000pt}%
\definecolor{currentstroke}{rgb}{0.000000,0.000000,0.000000}%
\pgfsetstrokecolor{currentstroke}%
\pgfsetdash{}{0pt}%
\pgfsys@defobject{currentmarker}{\pgfqpoint{0.000000in}{-0.048611in}}{\pgfqpoint{0.000000in}{0.000000in}}{%
\pgfpathmoveto{\pgfqpoint{0.000000in}{0.000000in}}%
\pgfpathlineto{\pgfqpoint{0.000000in}{-0.048611in}}%
\pgfusepath{stroke,fill}%
}%
\begin{pgfscope}%
\pgfsys@transformshift{8.212479in}{1.100000in}%
\pgfsys@useobject{currentmarker}{}%
\end{pgfscope}%
\end{pgfscope}%
\begin{pgfscope}%
\definecolor{textcolor}{rgb}{0.000000,0.000000,0.000000}%
\pgfsetstrokecolor{textcolor}%
\pgfsetfillcolor{textcolor}%
\pgftext[x=8.212479in,y=1.002778in,,top]{\color{textcolor}{\sffamily\fontsize{10.000000}{12.000000}\selectfont\catcode`\^=\active\def^{\ifmmode\sp\else\^{}\fi}\catcode`\%=\active\def%{\%}200}}%
\end{pgfscope}%
\begin{pgfscope}%
\pgfsetbuttcap%
\pgfsetroundjoin%
\definecolor{currentfill}{rgb}{0.000000,0.000000,0.000000}%
\pgfsetfillcolor{currentfill}%
\pgfsetlinewidth{0.803000pt}%
\definecolor{currentstroke}{rgb}{0.000000,0.000000,0.000000}%
\pgfsetstrokecolor{currentstroke}%
\pgfsetdash{}{0pt}%
\pgfsys@defobject{currentmarker}{\pgfqpoint{0.000000in}{-0.048611in}}{\pgfqpoint{0.000000in}{0.000000in}}{%
\pgfpathmoveto{\pgfqpoint{0.000000in}{0.000000in}}%
\pgfpathlineto{\pgfqpoint{0.000000in}{-0.048611in}}%
\pgfusepath{stroke,fill}%
}%
\begin{pgfscope}%
\pgfsys@transformshift{8.732259in}{1.100000in}%
\pgfsys@useobject{currentmarker}{}%
\end{pgfscope}%
\end{pgfscope}%
\begin{pgfscope}%
\definecolor{textcolor}{rgb}{0.000000,0.000000,0.000000}%
\pgfsetstrokecolor{textcolor}%
\pgfsetfillcolor{textcolor}%
\pgftext[x=8.732259in,y=1.002778in,,top]{\color{textcolor}{\sffamily\fontsize{10.000000}{12.000000}\selectfont\catcode`\^=\active\def^{\ifmmode\sp\else\^{}\fi}\catcode`\%=\active\def%{\%}250}}%
\end{pgfscope}%
\begin{pgfscope}%
\pgfsetbuttcap%
\pgfsetroundjoin%
\definecolor{currentfill}{rgb}{0.000000,0.000000,0.000000}%
\pgfsetfillcolor{currentfill}%
\pgfsetlinewidth{0.803000pt}%
\definecolor{currentstroke}{rgb}{0.000000,0.000000,0.000000}%
\pgfsetstrokecolor{currentstroke}%
\pgfsetdash{}{0pt}%
\pgfsys@defobject{currentmarker}{\pgfqpoint{0.000000in}{-0.048611in}}{\pgfqpoint{0.000000in}{0.000000in}}{%
\pgfpathmoveto{\pgfqpoint{0.000000in}{0.000000in}}%
\pgfpathlineto{\pgfqpoint{0.000000in}{-0.048611in}}%
\pgfusepath{stroke,fill}%
}%
\begin{pgfscope}%
\pgfsys@transformshift{9.252040in}{1.100000in}%
\pgfsys@useobject{currentmarker}{}%
\end{pgfscope}%
\end{pgfscope}%
\begin{pgfscope}%
\definecolor{textcolor}{rgb}{0.000000,0.000000,0.000000}%
\pgfsetstrokecolor{textcolor}%
\pgfsetfillcolor{textcolor}%
\pgftext[x=9.252040in,y=1.002778in,,top]{\color{textcolor}{\sffamily\fontsize{10.000000}{12.000000}\selectfont\catcode`\^=\active\def^{\ifmmode\sp\else\^{}\fi}\catcode`\%=\active\def%{\%}300}}%
\end{pgfscope}%
\begin{pgfscope}%
\definecolor{textcolor}{rgb}{0.000000,0.000000,0.000000}%
\pgfsetstrokecolor{textcolor}%
\pgfsetfillcolor{textcolor}%
\pgftext[x=7.687500in,y=0.812809in,,top]{\color{textcolor}{\sffamily\fontsize{10.000000}{12.000000}\selectfont\catcode`\^=\active\def^{\ifmmode\sp\else\^{}\fi}\catcode`\%=\active\def%{\%}Iterations}}%
\end{pgfscope}%
\begin{pgfscope}%
\pgfsetbuttcap%
\pgfsetroundjoin%
\definecolor{currentfill}{rgb}{0.000000,0.000000,0.000000}%
\pgfsetfillcolor{currentfill}%
\pgfsetlinewidth{0.803000pt}%
\definecolor{currentstroke}{rgb}{0.000000,0.000000,0.000000}%
\pgfsetstrokecolor{currentstroke}%
\pgfsetdash{}{0pt}%
\pgfsys@defobject{currentmarker}{\pgfqpoint{-0.048611in}{0.000000in}}{\pgfqpoint{-0.000000in}{0.000000in}}{%
\pgfpathmoveto{\pgfqpoint{-0.000000in}{0.000000in}}%
\pgfpathlineto{\pgfqpoint{-0.048611in}{0.000000in}}%
\pgfusepath{stroke,fill}%
}%
\begin{pgfscope}%
\pgfsys@transformshift{5.977941in}{1.660284in}%
\pgfsys@useobject{currentmarker}{}%
\end{pgfscope}%
\end{pgfscope}%
\begin{pgfscope}%
\pgfsetbuttcap%
\pgfsetroundjoin%
\definecolor{currentfill}{rgb}{0.000000,0.000000,0.000000}%
\pgfsetfillcolor{currentfill}%
\pgfsetlinewidth{0.803000pt}%
\definecolor{currentstroke}{rgb}{0.000000,0.000000,0.000000}%
\pgfsetstrokecolor{currentstroke}%
\pgfsetdash{}{0pt}%
\pgfsys@defobject{currentmarker}{\pgfqpoint{-0.048611in}{0.000000in}}{\pgfqpoint{-0.000000in}{0.000000in}}{%
\pgfpathmoveto{\pgfqpoint{-0.000000in}{0.000000in}}%
\pgfpathlineto{\pgfqpoint{-0.048611in}{0.000000in}}%
\pgfusepath{stroke,fill}%
}%
\begin{pgfscope}%
\pgfsys@transformshift{5.977941in}{2.534463in}%
\pgfsys@useobject{currentmarker}{}%
\end{pgfscope}%
\end{pgfscope}%
\begin{pgfscope}%
\pgfsetbuttcap%
\pgfsetroundjoin%
\definecolor{currentfill}{rgb}{0.000000,0.000000,0.000000}%
\pgfsetfillcolor{currentfill}%
\pgfsetlinewidth{0.803000pt}%
\definecolor{currentstroke}{rgb}{0.000000,0.000000,0.000000}%
\pgfsetstrokecolor{currentstroke}%
\pgfsetdash{}{0pt}%
\pgfsys@defobject{currentmarker}{\pgfqpoint{-0.048611in}{0.000000in}}{\pgfqpoint{-0.000000in}{0.000000in}}{%
\pgfpathmoveto{\pgfqpoint{-0.000000in}{0.000000in}}%
\pgfpathlineto{\pgfqpoint{-0.048611in}{0.000000in}}%
\pgfusepath{stroke,fill}%
}%
\begin{pgfscope}%
\pgfsys@transformshift{5.977941in}{3.408642in}%
\pgfsys@useobject{currentmarker}{}%
\end{pgfscope}%
\end{pgfscope}%
\begin{pgfscope}%
\pgfsetbuttcap%
\pgfsetroundjoin%
\definecolor{currentfill}{rgb}{0.000000,0.000000,0.000000}%
\pgfsetfillcolor{currentfill}%
\pgfsetlinewidth{0.803000pt}%
\definecolor{currentstroke}{rgb}{0.000000,0.000000,0.000000}%
\pgfsetstrokecolor{currentstroke}%
\pgfsetdash{}{0pt}%
\pgfsys@defobject{currentmarker}{\pgfqpoint{-0.048611in}{0.000000in}}{\pgfqpoint{-0.000000in}{0.000000in}}{%
\pgfpathmoveto{\pgfqpoint{-0.000000in}{0.000000in}}%
\pgfpathlineto{\pgfqpoint{-0.048611in}{0.000000in}}%
\pgfusepath{stroke,fill}%
}%
\begin{pgfscope}%
\pgfsys@transformshift{5.977941in}{4.282822in}%
\pgfsys@useobject{currentmarker}{}%
\end{pgfscope}%
\end{pgfscope}%
\begin{pgfscope}%
\pgfsetbuttcap%
\pgfsetroundjoin%
\definecolor{currentfill}{rgb}{0.000000,0.000000,0.000000}%
\pgfsetfillcolor{currentfill}%
\pgfsetlinewidth{0.602250pt}%
\definecolor{currentstroke}{rgb}{0.000000,0.000000,0.000000}%
\pgfsetstrokecolor{currentstroke}%
\pgfsetdash{}{0pt}%
\pgfsys@defobject{currentmarker}{\pgfqpoint{-0.027778in}{0.000000in}}{\pgfqpoint{-0.000000in}{0.000000in}}{%
\pgfpathmoveto{\pgfqpoint{-0.000000in}{0.000000in}}%
\pgfpathlineto{\pgfqpoint{-0.027778in}{0.000000in}}%
\pgfusepath{stroke,fill}%
}%
\begin{pgfscope}%
\pgfsys@transformshift{5.977941in}{1.203194in}%
\pgfsys@useobject{currentmarker}{}%
\end{pgfscope}%
\end{pgfscope}%
\begin{pgfscope}%
\pgfsetbuttcap%
\pgfsetroundjoin%
\definecolor{currentfill}{rgb}{0.000000,0.000000,0.000000}%
\pgfsetfillcolor{currentfill}%
\pgfsetlinewidth{0.602250pt}%
\definecolor{currentstroke}{rgb}{0.000000,0.000000,0.000000}%
\pgfsetstrokecolor{currentstroke}%
\pgfsetdash{}{0pt}%
\pgfsys@defobject{currentmarker}{\pgfqpoint{-0.027778in}{0.000000in}}{\pgfqpoint{-0.000000in}{0.000000in}}{%
\pgfpathmoveto{\pgfqpoint{-0.000000in}{0.000000in}}%
\pgfpathlineto{\pgfqpoint{-0.027778in}{0.000000in}}%
\pgfusepath{stroke,fill}%
}%
\begin{pgfscope}%
\pgfsys@transformshift{5.977941in}{1.312413in}%
\pgfsys@useobject{currentmarker}{}%
\end{pgfscope}%
\end{pgfscope}%
\begin{pgfscope}%
\pgfsetbuttcap%
\pgfsetroundjoin%
\definecolor{currentfill}{rgb}{0.000000,0.000000,0.000000}%
\pgfsetfillcolor{currentfill}%
\pgfsetlinewidth{0.602250pt}%
\definecolor{currentstroke}{rgb}{0.000000,0.000000,0.000000}%
\pgfsetstrokecolor{currentstroke}%
\pgfsetdash{}{0pt}%
\pgfsys@defobject{currentmarker}{\pgfqpoint{-0.027778in}{0.000000in}}{\pgfqpoint{-0.000000in}{0.000000in}}{%
\pgfpathmoveto{\pgfqpoint{-0.000000in}{0.000000in}}%
\pgfpathlineto{\pgfqpoint{-0.027778in}{0.000000in}}%
\pgfusepath{stroke,fill}%
}%
\begin{pgfscope}%
\pgfsys@transformshift{5.977941in}{1.397130in}%
\pgfsys@useobject{currentmarker}{}%
\end{pgfscope}%
\end{pgfscope}%
\begin{pgfscope}%
\pgfsetbuttcap%
\pgfsetroundjoin%
\definecolor{currentfill}{rgb}{0.000000,0.000000,0.000000}%
\pgfsetfillcolor{currentfill}%
\pgfsetlinewidth{0.602250pt}%
\definecolor{currentstroke}{rgb}{0.000000,0.000000,0.000000}%
\pgfsetstrokecolor{currentstroke}%
\pgfsetdash{}{0pt}%
\pgfsys@defobject{currentmarker}{\pgfqpoint{-0.027778in}{0.000000in}}{\pgfqpoint{-0.000000in}{0.000000in}}{%
\pgfpathmoveto{\pgfqpoint{-0.000000in}{0.000000in}}%
\pgfpathlineto{\pgfqpoint{-0.027778in}{0.000000in}}%
\pgfusepath{stroke,fill}%
}%
\begin{pgfscope}%
\pgfsys@transformshift{5.977941in}{1.466348in}%
\pgfsys@useobject{currentmarker}{}%
\end{pgfscope}%
\end{pgfscope}%
\begin{pgfscope}%
\pgfsetbuttcap%
\pgfsetroundjoin%
\definecolor{currentfill}{rgb}{0.000000,0.000000,0.000000}%
\pgfsetfillcolor{currentfill}%
\pgfsetlinewidth{0.602250pt}%
\definecolor{currentstroke}{rgb}{0.000000,0.000000,0.000000}%
\pgfsetstrokecolor{currentstroke}%
\pgfsetdash{}{0pt}%
\pgfsys@defobject{currentmarker}{\pgfqpoint{-0.027778in}{0.000000in}}{\pgfqpoint{-0.000000in}{0.000000in}}{%
\pgfpathmoveto{\pgfqpoint{-0.000000in}{0.000000in}}%
\pgfpathlineto{\pgfqpoint{-0.027778in}{0.000000in}}%
\pgfusepath{stroke,fill}%
}%
\begin{pgfscope}%
\pgfsys@transformshift{5.977941in}{1.524872in}%
\pgfsys@useobject{currentmarker}{}%
\end{pgfscope}%
\end{pgfscope}%
\begin{pgfscope}%
\pgfsetbuttcap%
\pgfsetroundjoin%
\definecolor{currentfill}{rgb}{0.000000,0.000000,0.000000}%
\pgfsetfillcolor{currentfill}%
\pgfsetlinewidth{0.602250pt}%
\definecolor{currentstroke}{rgb}{0.000000,0.000000,0.000000}%
\pgfsetstrokecolor{currentstroke}%
\pgfsetdash{}{0pt}%
\pgfsys@defobject{currentmarker}{\pgfqpoint{-0.027778in}{0.000000in}}{\pgfqpoint{-0.000000in}{0.000000in}}{%
\pgfpathmoveto{\pgfqpoint{-0.000000in}{0.000000in}}%
\pgfpathlineto{\pgfqpoint{-0.027778in}{0.000000in}}%
\pgfusepath{stroke,fill}%
}%
\begin{pgfscope}%
\pgfsys@transformshift{5.977941in}{1.575567in}%
\pgfsys@useobject{currentmarker}{}%
\end{pgfscope}%
\end{pgfscope}%
\begin{pgfscope}%
\pgfsetbuttcap%
\pgfsetroundjoin%
\definecolor{currentfill}{rgb}{0.000000,0.000000,0.000000}%
\pgfsetfillcolor{currentfill}%
\pgfsetlinewidth{0.602250pt}%
\definecolor{currentstroke}{rgb}{0.000000,0.000000,0.000000}%
\pgfsetstrokecolor{currentstroke}%
\pgfsetdash{}{0pt}%
\pgfsys@defobject{currentmarker}{\pgfqpoint{-0.027778in}{0.000000in}}{\pgfqpoint{-0.000000in}{0.000000in}}{%
\pgfpathmoveto{\pgfqpoint{-0.000000in}{0.000000in}}%
\pgfpathlineto{\pgfqpoint{-0.027778in}{0.000000in}}%
\pgfusepath{stroke,fill}%
}%
\begin{pgfscope}%
\pgfsys@transformshift{5.977941in}{1.620284in}%
\pgfsys@useobject{currentmarker}{}%
\end{pgfscope}%
\end{pgfscope}%
\begin{pgfscope}%
\pgfsetbuttcap%
\pgfsetroundjoin%
\definecolor{currentfill}{rgb}{0.000000,0.000000,0.000000}%
\pgfsetfillcolor{currentfill}%
\pgfsetlinewidth{0.602250pt}%
\definecolor{currentstroke}{rgb}{0.000000,0.000000,0.000000}%
\pgfsetstrokecolor{currentstroke}%
\pgfsetdash{}{0pt}%
\pgfsys@defobject{currentmarker}{\pgfqpoint{-0.027778in}{0.000000in}}{\pgfqpoint{-0.000000in}{0.000000in}}{%
\pgfpathmoveto{\pgfqpoint{-0.000000in}{0.000000in}}%
\pgfpathlineto{\pgfqpoint{-0.027778in}{0.000000in}}%
\pgfusepath{stroke,fill}%
}%
\begin{pgfscope}%
\pgfsys@transformshift{5.977941in}{1.923438in}%
\pgfsys@useobject{currentmarker}{}%
\end{pgfscope}%
\end{pgfscope}%
\begin{pgfscope}%
\pgfsetbuttcap%
\pgfsetroundjoin%
\definecolor{currentfill}{rgb}{0.000000,0.000000,0.000000}%
\pgfsetfillcolor{currentfill}%
\pgfsetlinewidth{0.602250pt}%
\definecolor{currentstroke}{rgb}{0.000000,0.000000,0.000000}%
\pgfsetstrokecolor{currentstroke}%
\pgfsetdash{}{0pt}%
\pgfsys@defobject{currentmarker}{\pgfqpoint{-0.027778in}{0.000000in}}{\pgfqpoint{-0.000000in}{0.000000in}}{%
\pgfpathmoveto{\pgfqpoint{-0.000000in}{0.000000in}}%
\pgfpathlineto{\pgfqpoint{-0.027778in}{0.000000in}}%
\pgfusepath{stroke,fill}%
}%
\begin{pgfscope}%
\pgfsys@transformshift{5.977941in}{2.077373in}%
\pgfsys@useobject{currentmarker}{}%
\end{pgfscope}%
\end{pgfscope}%
\begin{pgfscope}%
\pgfsetbuttcap%
\pgfsetroundjoin%
\definecolor{currentfill}{rgb}{0.000000,0.000000,0.000000}%
\pgfsetfillcolor{currentfill}%
\pgfsetlinewidth{0.602250pt}%
\definecolor{currentstroke}{rgb}{0.000000,0.000000,0.000000}%
\pgfsetstrokecolor{currentstroke}%
\pgfsetdash{}{0pt}%
\pgfsys@defobject{currentmarker}{\pgfqpoint{-0.027778in}{0.000000in}}{\pgfqpoint{-0.000000in}{0.000000in}}{%
\pgfpathmoveto{\pgfqpoint{-0.000000in}{0.000000in}}%
\pgfpathlineto{\pgfqpoint{-0.027778in}{0.000000in}}%
\pgfusepath{stroke,fill}%
}%
\begin{pgfscope}%
\pgfsys@transformshift{5.977941in}{2.186592in}%
\pgfsys@useobject{currentmarker}{}%
\end{pgfscope}%
\end{pgfscope}%
\begin{pgfscope}%
\pgfsetbuttcap%
\pgfsetroundjoin%
\definecolor{currentfill}{rgb}{0.000000,0.000000,0.000000}%
\pgfsetfillcolor{currentfill}%
\pgfsetlinewidth{0.602250pt}%
\definecolor{currentstroke}{rgb}{0.000000,0.000000,0.000000}%
\pgfsetstrokecolor{currentstroke}%
\pgfsetdash{}{0pt}%
\pgfsys@defobject{currentmarker}{\pgfqpoint{-0.027778in}{0.000000in}}{\pgfqpoint{-0.000000in}{0.000000in}}{%
\pgfpathmoveto{\pgfqpoint{-0.000000in}{0.000000in}}%
\pgfpathlineto{\pgfqpoint{-0.027778in}{0.000000in}}%
\pgfusepath{stroke,fill}%
}%
\begin{pgfscope}%
\pgfsys@transformshift{5.977941in}{2.271309in}%
\pgfsys@useobject{currentmarker}{}%
\end{pgfscope}%
\end{pgfscope}%
\begin{pgfscope}%
\pgfsetbuttcap%
\pgfsetroundjoin%
\definecolor{currentfill}{rgb}{0.000000,0.000000,0.000000}%
\pgfsetfillcolor{currentfill}%
\pgfsetlinewidth{0.602250pt}%
\definecolor{currentstroke}{rgb}{0.000000,0.000000,0.000000}%
\pgfsetstrokecolor{currentstroke}%
\pgfsetdash{}{0pt}%
\pgfsys@defobject{currentmarker}{\pgfqpoint{-0.027778in}{0.000000in}}{\pgfqpoint{-0.000000in}{0.000000in}}{%
\pgfpathmoveto{\pgfqpoint{-0.000000in}{0.000000in}}%
\pgfpathlineto{\pgfqpoint{-0.027778in}{0.000000in}}%
\pgfusepath{stroke,fill}%
}%
\begin{pgfscope}%
\pgfsys@transformshift{5.977941in}{2.340528in}%
\pgfsys@useobject{currentmarker}{}%
\end{pgfscope}%
\end{pgfscope}%
\begin{pgfscope}%
\pgfsetbuttcap%
\pgfsetroundjoin%
\definecolor{currentfill}{rgb}{0.000000,0.000000,0.000000}%
\pgfsetfillcolor{currentfill}%
\pgfsetlinewidth{0.602250pt}%
\definecolor{currentstroke}{rgb}{0.000000,0.000000,0.000000}%
\pgfsetstrokecolor{currentstroke}%
\pgfsetdash{}{0pt}%
\pgfsys@defobject{currentmarker}{\pgfqpoint{-0.027778in}{0.000000in}}{\pgfqpoint{-0.000000in}{0.000000in}}{%
\pgfpathmoveto{\pgfqpoint{-0.000000in}{0.000000in}}%
\pgfpathlineto{\pgfqpoint{-0.027778in}{0.000000in}}%
\pgfusepath{stroke,fill}%
}%
\begin{pgfscope}%
\pgfsys@transformshift{5.977941in}{2.399051in}%
\pgfsys@useobject{currentmarker}{}%
\end{pgfscope}%
\end{pgfscope}%
\begin{pgfscope}%
\pgfsetbuttcap%
\pgfsetroundjoin%
\definecolor{currentfill}{rgb}{0.000000,0.000000,0.000000}%
\pgfsetfillcolor{currentfill}%
\pgfsetlinewidth{0.602250pt}%
\definecolor{currentstroke}{rgb}{0.000000,0.000000,0.000000}%
\pgfsetstrokecolor{currentstroke}%
\pgfsetdash{}{0pt}%
\pgfsys@defobject{currentmarker}{\pgfqpoint{-0.027778in}{0.000000in}}{\pgfqpoint{-0.000000in}{0.000000in}}{%
\pgfpathmoveto{\pgfqpoint{-0.000000in}{0.000000in}}%
\pgfpathlineto{\pgfqpoint{-0.027778in}{0.000000in}}%
\pgfusepath{stroke,fill}%
}%
\begin{pgfscope}%
\pgfsys@transformshift{5.977941in}{2.449746in}%
\pgfsys@useobject{currentmarker}{}%
\end{pgfscope}%
\end{pgfscope}%
\begin{pgfscope}%
\pgfsetbuttcap%
\pgfsetroundjoin%
\definecolor{currentfill}{rgb}{0.000000,0.000000,0.000000}%
\pgfsetfillcolor{currentfill}%
\pgfsetlinewidth{0.602250pt}%
\definecolor{currentstroke}{rgb}{0.000000,0.000000,0.000000}%
\pgfsetstrokecolor{currentstroke}%
\pgfsetdash{}{0pt}%
\pgfsys@defobject{currentmarker}{\pgfqpoint{-0.027778in}{0.000000in}}{\pgfqpoint{-0.000000in}{0.000000in}}{%
\pgfpathmoveto{\pgfqpoint{-0.000000in}{0.000000in}}%
\pgfpathlineto{\pgfqpoint{-0.027778in}{0.000000in}}%
\pgfusepath{stroke,fill}%
}%
\begin{pgfscope}%
\pgfsys@transformshift{5.977941in}{2.494463in}%
\pgfsys@useobject{currentmarker}{}%
\end{pgfscope}%
\end{pgfscope}%
\begin{pgfscope}%
\pgfsetbuttcap%
\pgfsetroundjoin%
\definecolor{currentfill}{rgb}{0.000000,0.000000,0.000000}%
\pgfsetfillcolor{currentfill}%
\pgfsetlinewidth{0.602250pt}%
\definecolor{currentstroke}{rgb}{0.000000,0.000000,0.000000}%
\pgfsetstrokecolor{currentstroke}%
\pgfsetdash{}{0pt}%
\pgfsys@defobject{currentmarker}{\pgfqpoint{-0.027778in}{0.000000in}}{\pgfqpoint{-0.000000in}{0.000000in}}{%
\pgfpathmoveto{\pgfqpoint{-0.000000in}{0.000000in}}%
\pgfpathlineto{\pgfqpoint{-0.027778in}{0.000000in}}%
\pgfusepath{stroke,fill}%
}%
\begin{pgfscope}%
\pgfsys@transformshift{5.977941in}{2.797617in}%
\pgfsys@useobject{currentmarker}{}%
\end{pgfscope}%
\end{pgfscope}%
\begin{pgfscope}%
\pgfsetbuttcap%
\pgfsetroundjoin%
\definecolor{currentfill}{rgb}{0.000000,0.000000,0.000000}%
\pgfsetfillcolor{currentfill}%
\pgfsetlinewidth{0.602250pt}%
\definecolor{currentstroke}{rgb}{0.000000,0.000000,0.000000}%
\pgfsetstrokecolor{currentstroke}%
\pgfsetdash{}{0pt}%
\pgfsys@defobject{currentmarker}{\pgfqpoint{-0.027778in}{0.000000in}}{\pgfqpoint{-0.000000in}{0.000000in}}{%
\pgfpathmoveto{\pgfqpoint{-0.000000in}{0.000000in}}%
\pgfpathlineto{\pgfqpoint{-0.027778in}{0.000000in}}%
\pgfusepath{stroke,fill}%
}%
\begin{pgfscope}%
\pgfsys@transformshift{5.977941in}{2.951553in}%
\pgfsys@useobject{currentmarker}{}%
\end{pgfscope}%
\end{pgfscope}%
\begin{pgfscope}%
\pgfsetbuttcap%
\pgfsetroundjoin%
\definecolor{currentfill}{rgb}{0.000000,0.000000,0.000000}%
\pgfsetfillcolor{currentfill}%
\pgfsetlinewidth{0.602250pt}%
\definecolor{currentstroke}{rgb}{0.000000,0.000000,0.000000}%
\pgfsetstrokecolor{currentstroke}%
\pgfsetdash{}{0pt}%
\pgfsys@defobject{currentmarker}{\pgfqpoint{-0.027778in}{0.000000in}}{\pgfqpoint{-0.000000in}{0.000000in}}{%
\pgfpathmoveto{\pgfqpoint{-0.000000in}{0.000000in}}%
\pgfpathlineto{\pgfqpoint{-0.027778in}{0.000000in}}%
\pgfusepath{stroke,fill}%
}%
\begin{pgfscope}%
\pgfsys@transformshift{5.977941in}{3.060771in}%
\pgfsys@useobject{currentmarker}{}%
\end{pgfscope}%
\end{pgfscope}%
\begin{pgfscope}%
\pgfsetbuttcap%
\pgfsetroundjoin%
\definecolor{currentfill}{rgb}{0.000000,0.000000,0.000000}%
\pgfsetfillcolor{currentfill}%
\pgfsetlinewidth{0.602250pt}%
\definecolor{currentstroke}{rgb}{0.000000,0.000000,0.000000}%
\pgfsetstrokecolor{currentstroke}%
\pgfsetdash{}{0pt}%
\pgfsys@defobject{currentmarker}{\pgfqpoint{-0.027778in}{0.000000in}}{\pgfqpoint{-0.000000in}{0.000000in}}{%
\pgfpathmoveto{\pgfqpoint{-0.000000in}{0.000000in}}%
\pgfpathlineto{\pgfqpoint{-0.027778in}{0.000000in}}%
\pgfusepath{stroke,fill}%
}%
\begin{pgfscope}%
\pgfsys@transformshift{5.977941in}{3.145488in}%
\pgfsys@useobject{currentmarker}{}%
\end{pgfscope}%
\end{pgfscope}%
\begin{pgfscope}%
\pgfsetbuttcap%
\pgfsetroundjoin%
\definecolor{currentfill}{rgb}{0.000000,0.000000,0.000000}%
\pgfsetfillcolor{currentfill}%
\pgfsetlinewidth{0.602250pt}%
\definecolor{currentstroke}{rgb}{0.000000,0.000000,0.000000}%
\pgfsetstrokecolor{currentstroke}%
\pgfsetdash{}{0pt}%
\pgfsys@defobject{currentmarker}{\pgfqpoint{-0.027778in}{0.000000in}}{\pgfqpoint{-0.000000in}{0.000000in}}{%
\pgfpathmoveto{\pgfqpoint{-0.000000in}{0.000000in}}%
\pgfpathlineto{\pgfqpoint{-0.027778in}{0.000000in}}%
\pgfusepath{stroke,fill}%
}%
\begin{pgfscope}%
\pgfsys@transformshift{5.977941in}{3.214707in}%
\pgfsys@useobject{currentmarker}{}%
\end{pgfscope}%
\end{pgfscope}%
\begin{pgfscope}%
\pgfsetbuttcap%
\pgfsetroundjoin%
\definecolor{currentfill}{rgb}{0.000000,0.000000,0.000000}%
\pgfsetfillcolor{currentfill}%
\pgfsetlinewidth{0.602250pt}%
\definecolor{currentstroke}{rgb}{0.000000,0.000000,0.000000}%
\pgfsetstrokecolor{currentstroke}%
\pgfsetdash{}{0pt}%
\pgfsys@defobject{currentmarker}{\pgfqpoint{-0.027778in}{0.000000in}}{\pgfqpoint{-0.000000in}{0.000000in}}{%
\pgfpathmoveto{\pgfqpoint{-0.000000in}{0.000000in}}%
\pgfpathlineto{\pgfqpoint{-0.027778in}{0.000000in}}%
\pgfusepath{stroke,fill}%
}%
\begin{pgfscope}%
\pgfsys@transformshift{5.977941in}{3.273230in}%
\pgfsys@useobject{currentmarker}{}%
\end{pgfscope}%
\end{pgfscope}%
\begin{pgfscope}%
\pgfsetbuttcap%
\pgfsetroundjoin%
\definecolor{currentfill}{rgb}{0.000000,0.000000,0.000000}%
\pgfsetfillcolor{currentfill}%
\pgfsetlinewidth{0.602250pt}%
\definecolor{currentstroke}{rgb}{0.000000,0.000000,0.000000}%
\pgfsetstrokecolor{currentstroke}%
\pgfsetdash{}{0pt}%
\pgfsys@defobject{currentmarker}{\pgfqpoint{-0.027778in}{0.000000in}}{\pgfqpoint{-0.000000in}{0.000000in}}{%
\pgfpathmoveto{\pgfqpoint{-0.000000in}{0.000000in}}%
\pgfpathlineto{\pgfqpoint{-0.027778in}{0.000000in}}%
\pgfusepath{stroke,fill}%
}%
\begin{pgfscope}%
\pgfsys@transformshift{5.977941in}{3.323926in}%
\pgfsys@useobject{currentmarker}{}%
\end{pgfscope}%
\end{pgfscope}%
\begin{pgfscope}%
\pgfsetbuttcap%
\pgfsetroundjoin%
\definecolor{currentfill}{rgb}{0.000000,0.000000,0.000000}%
\pgfsetfillcolor{currentfill}%
\pgfsetlinewidth{0.602250pt}%
\definecolor{currentstroke}{rgb}{0.000000,0.000000,0.000000}%
\pgfsetstrokecolor{currentstroke}%
\pgfsetdash{}{0pt}%
\pgfsys@defobject{currentmarker}{\pgfqpoint{-0.027778in}{0.000000in}}{\pgfqpoint{-0.000000in}{0.000000in}}{%
\pgfpathmoveto{\pgfqpoint{-0.000000in}{0.000000in}}%
\pgfpathlineto{\pgfqpoint{-0.027778in}{0.000000in}}%
\pgfusepath{stroke,fill}%
}%
\begin{pgfscope}%
\pgfsys@transformshift{5.977941in}{3.368642in}%
\pgfsys@useobject{currentmarker}{}%
\end{pgfscope}%
\end{pgfscope}%
\begin{pgfscope}%
\pgfsetbuttcap%
\pgfsetroundjoin%
\definecolor{currentfill}{rgb}{0.000000,0.000000,0.000000}%
\pgfsetfillcolor{currentfill}%
\pgfsetlinewidth{0.602250pt}%
\definecolor{currentstroke}{rgb}{0.000000,0.000000,0.000000}%
\pgfsetstrokecolor{currentstroke}%
\pgfsetdash{}{0pt}%
\pgfsys@defobject{currentmarker}{\pgfqpoint{-0.027778in}{0.000000in}}{\pgfqpoint{-0.000000in}{0.000000in}}{%
\pgfpathmoveto{\pgfqpoint{-0.000000in}{0.000000in}}%
\pgfpathlineto{\pgfqpoint{-0.027778in}{0.000000in}}%
\pgfusepath{stroke,fill}%
}%
\begin{pgfscope}%
\pgfsys@transformshift{5.977941in}{3.671797in}%
\pgfsys@useobject{currentmarker}{}%
\end{pgfscope}%
\end{pgfscope}%
\begin{pgfscope}%
\pgfsetbuttcap%
\pgfsetroundjoin%
\definecolor{currentfill}{rgb}{0.000000,0.000000,0.000000}%
\pgfsetfillcolor{currentfill}%
\pgfsetlinewidth{0.602250pt}%
\definecolor{currentstroke}{rgb}{0.000000,0.000000,0.000000}%
\pgfsetstrokecolor{currentstroke}%
\pgfsetdash{}{0pt}%
\pgfsys@defobject{currentmarker}{\pgfqpoint{-0.027778in}{0.000000in}}{\pgfqpoint{-0.000000in}{0.000000in}}{%
\pgfpathmoveto{\pgfqpoint{-0.000000in}{0.000000in}}%
\pgfpathlineto{\pgfqpoint{-0.027778in}{0.000000in}}%
\pgfusepath{stroke,fill}%
}%
\begin{pgfscope}%
\pgfsys@transformshift{5.977941in}{3.825732in}%
\pgfsys@useobject{currentmarker}{}%
\end{pgfscope}%
\end{pgfscope}%
\begin{pgfscope}%
\pgfsetbuttcap%
\pgfsetroundjoin%
\definecolor{currentfill}{rgb}{0.000000,0.000000,0.000000}%
\pgfsetfillcolor{currentfill}%
\pgfsetlinewidth{0.602250pt}%
\definecolor{currentstroke}{rgb}{0.000000,0.000000,0.000000}%
\pgfsetstrokecolor{currentstroke}%
\pgfsetdash{}{0pt}%
\pgfsys@defobject{currentmarker}{\pgfqpoint{-0.027778in}{0.000000in}}{\pgfqpoint{-0.000000in}{0.000000in}}{%
\pgfpathmoveto{\pgfqpoint{-0.000000in}{0.000000in}}%
\pgfpathlineto{\pgfqpoint{-0.027778in}{0.000000in}}%
\pgfusepath{stroke,fill}%
}%
\begin{pgfscope}%
\pgfsys@transformshift{5.977941in}{3.934951in}%
\pgfsys@useobject{currentmarker}{}%
\end{pgfscope}%
\end{pgfscope}%
\begin{pgfscope}%
\pgfsetbuttcap%
\pgfsetroundjoin%
\definecolor{currentfill}{rgb}{0.000000,0.000000,0.000000}%
\pgfsetfillcolor{currentfill}%
\pgfsetlinewidth{0.602250pt}%
\definecolor{currentstroke}{rgb}{0.000000,0.000000,0.000000}%
\pgfsetstrokecolor{currentstroke}%
\pgfsetdash{}{0pt}%
\pgfsys@defobject{currentmarker}{\pgfqpoint{-0.027778in}{0.000000in}}{\pgfqpoint{-0.000000in}{0.000000in}}{%
\pgfpathmoveto{\pgfqpoint{-0.000000in}{0.000000in}}%
\pgfpathlineto{\pgfqpoint{-0.027778in}{0.000000in}}%
\pgfusepath{stroke,fill}%
}%
\begin{pgfscope}%
\pgfsys@transformshift{5.977941in}{4.019667in}%
\pgfsys@useobject{currentmarker}{}%
\end{pgfscope}%
\end{pgfscope}%
\begin{pgfscope}%
\pgfsetbuttcap%
\pgfsetroundjoin%
\definecolor{currentfill}{rgb}{0.000000,0.000000,0.000000}%
\pgfsetfillcolor{currentfill}%
\pgfsetlinewidth{0.602250pt}%
\definecolor{currentstroke}{rgb}{0.000000,0.000000,0.000000}%
\pgfsetstrokecolor{currentstroke}%
\pgfsetdash{}{0pt}%
\pgfsys@defobject{currentmarker}{\pgfqpoint{-0.027778in}{0.000000in}}{\pgfqpoint{-0.000000in}{0.000000in}}{%
\pgfpathmoveto{\pgfqpoint{-0.000000in}{0.000000in}}%
\pgfpathlineto{\pgfqpoint{-0.027778in}{0.000000in}}%
\pgfusepath{stroke,fill}%
}%
\begin{pgfscope}%
\pgfsys@transformshift{5.977941in}{4.088886in}%
\pgfsys@useobject{currentmarker}{}%
\end{pgfscope}%
\end{pgfscope}%
\begin{pgfscope}%
\pgfsetbuttcap%
\pgfsetroundjoin%
\definecolor{currentfill}{rgb}{0.000000,0.000000,0.000000}%
\pgfsetfillcolor{currentfill}%
\pgfsetlinewidth{0.602250pt}%
\definecolor{currentstroke}{rgb}{0.000000,0.000000,0.000000}%
\pgfsetstrokecolor{currentstroke}%
\pgfsetdash{}{0pt}%
\pgfsys@defobject{currentmarker}{\pgfqpoint{-0.027778in}{0.000000in}}{\pgfqpoint{-0.000000in}{0.000000in}}{%
\pgfpathmoveto{\pgfqpoint{-0.000000in}{0.000000in}}%
\pgfpathlineto{\pgfqpoint{-0.027778in}{0.000000in}}%
\pgfusepath{stroke,fill}%
}%
\begin{pgfscope}%
\pgfsys@transformshift{5.977941in}{4.147409in}%
\pgfsys@useobject{currentmarker}{}%
\end{pgfscope}%
\end{pgfscope}%
\begin{pgfscope}%
\pgfsetbuttcap%
\pgfsetroundjoin%
\definecolor{currentfill}{rgb}{0.000000,0.000000,0.000000}%
\pgfsetfillcolor{currentfill}%
\pgfsetlinewidth{0.602250pt}%
\definecolor{currentstroke}{rgb}{0.000000,0.000000,0.000000}%
\pgfsetstrokecolor{currentstroke}%
\pgfsetdash{}{0pt}%
\pgfsys@defobject{currentmarker}{\pgfqpoint{-0.027778in}{0.000000in}}{\pgfqpoint{-0.000000in}{0.000000in}}{%
\pgfpathmoveto{\pgfqpoint{-0.000000in}{0.000000in}}%
\pgfpathlineto{\pgfqpoint{-0.027778in}{0.000000in}}%
\pgfusepath{stroke,fill}%
}%
\begin{pgfscope}%
\pgfsys@transformshift{5.977941in}{4.198105in}%
\pgfsys@useobject{currentmarker}{}%
\end{pgfscope}%
\end{pgfscope}%
\begin{pgfscope}%
\pgfsetbuttcap%
\pgfsetroundjoin%
\definecolor{currentfill}{rgb}{0.000000,0.000000,0.000000}%
\pgfsetfillcolor{currentfill}%
\pgfsetlinewidth{0.602250pt}%
\definecolor{currentstroke}{rgb}{0.000000,0.000000,0.000000}%
\pgfsetstrokecolor{currentstroke}%
\pgfsetdash{}{0pt}%
\pgfsys@defobject{currentmarker}{\pgfqpoint{-0.027778in}{0.000000in}}{\pgfqpoint{-0.000000in}{0.000000in}}{%
\pgfpathmoveto{\pgfqpoint{-0.000000in}{0.000000in}}%
\pgfpathlineto{\pgfqpoint{-0.027778in}{0.000000in}}%
\pgfusepath{stroke,fill}%
}%
\begin{pgfscope}%
\pgfsys@transformshift{5.977941in}{4.242821in}%
\pgfsys@useobject{currentmarker}{}%
\end{pgfscope}%
\end{pgfscope}%
\begin{pgfscope}%
\pgfsetbuttcap%
\pgfsetroundjoin%
\definecolor{currentfill}{rgb}{0.000000,0.000000,0.000000}%
\pgfsetfillcolor{currentfill}%
\pgfsetlinewidth{0.602250pt}%
\definecolor{currentstroke}{rgb}{0.000000,0.000000,0.000000}%
\pgfsetstrokecolor{currentstroke}%
\pgfsetdash{}{0pt}%
\pgfsys@defobject{currentmarker}{\pgfqpoint{-0.027778in}{0.000000in}}{\pgfqpoint{-0.000000in}{0.000000in}}{%
\pgfpathmoveto{\pgfqpoint{-0.000000in}{0.000000in}}%
\pgfpathlineto{\pgfqpoint{-0.027778in}{0.000000in}}%
\pgfusepath{stroke,fill}%
}%
\begin{pgfscope}%
\pgfsys@transformshift{5.977941in}{4.545976in}%
\pgfsys@useobject{currentmarker}{}%
\end{pgfscope}%
\end{pgfscope}%
\begin{pgfscope}%
\pgfpathrectangle{\pgfqpoint{5.977941in}{1.100000in}}{\pgfqpoint{3.419118in}{3.500000in}}%
\pgfusepath{clip}%
\pgfsetrectcap%
\pgfsetroundjoin%
\pgfsetlinewidth{1.505625pt}%
\definecolor{currentstroke}{rgb}{0.121569,0.466667,0.705882}%
\pgfsetstrokecolor{currentstroke}%
\pgfsetdash{}{0pt}%
\pgfpathmoveto{\pgfqpoint{6.133356in}{4.440909in}}%
\pgfpathlineto{\pgfqpoint{6.237312in}{4.367257in}}%
\pgfpathlineto{\pgfqpoint{6.382850in}{4.267187in}}%
\pgfpathlineto{\pgfqpoint{6.580367in}{4.134244in}}%
\pgfpathlineto{\pgfqpoint{6.871444in}{3.941191in}}%
\pgfpathlineto{\pgfqpoint{7.401621in}{3.592671in}}%
\pgfpathlineto{\pgfqpoint{9.241644in}{2.387067in}}%
\pgfpathlineto{\pgfqpoint{9.241644in}{2.387067in}}%
\pgfusepath{stroke}%
\end{pgfscope}%
\begin{pgfscope}%
\pgfpathrectangle{\pgfqpoint{5.977941in}{1.100000in}}{\pgfqpoint{3.419118in}{3.500000in}}%
\pgfusepath{clip}%
\pgfsetrectcap%
\pgfsetroundjoin%
\pgfsetlinewidth{1.505625pt}%
\definecolor{currentstroke}{rgb}{1.000000,0.498039,0.054902}%
\pgfsetstrokecolor{currentstroke}%
\pgfsetdash{}{0pt}%
\pgfpathmoveto{\pgfqpoint{6.133356in}{4.380219in}}%
\pgfpathlineto{\pgfqpoint{6.216521in}{4.313606in}}%
\pgfpathlineto{\pgfqpoint{6.310081in}{4.241920in}}%
\pgfpathlineto{\pgfqpoint{6.372455in}{4.195607in}}%
\pgfpathlineto{\pgfqpoint{6.382850in}{3.733953in}}%
\pgfpathlineto{\pgfqpoint{6.393246in}{3.700788in}}%
\pgfpathlineto{\pgfqpoint{6.403642in}{3.677088in}}%
\pgfpathlineto{\pgfqpoint{6.424433in}{3.638380in}}%
\pgfpathlineto{\pgfqpoint{6.445224in}{3.605886in}}%
\pgfpathlineto{\pgfqpoint{6.466015in}{3.577024in}}%
\pgfpathlineto{\pgfqpoint{6.497202in}{3.538137in}}%
\pgfpathlineto{\pgfqpoint{6.528389in}{3.502895in}}%
\pgfpathlineto{\pgfqpoint{6.569971in}{3.459868in}}%
\pgfpathlineto{\pgfqpoint{6.611554in}{3.420126in}}%
\pgfpathlineto{\pgfqpoint{6.632345in}{3.401201in}}%
\pgfpathlineto{\pgfqpoint{6.642741in}{3.625550in}}%
\pgfpathlineto{\pgfqpoint{6.653136in}{3.600667in}}%
\pgfpathlineto{\pgfqpoint{6.663532in}{3.586193in}}%
\pgfpathlineto{\pgfqpoint{6.684323in}{3.563037in}}%
\pgfpathlineto{\pgfqpoint{6.715510in}{3.533224in}}%
\pgfpathlineto{\pgfqpoint{6.757092in}{3.497671in}}%
\pgfpathlineto{\pgfqpoint{6.809071in}{3.456888in}}%
\pgfpathlineto{\pgfqpoint{6.871444in}{3.410997in}}%
\pgfpathlineto{\pgfqpoint{6.892235in}{3.396173in}}%
\pgfpathlineto{\pgfqpoint{6.902631in}{3.619199in}}%
\pgfpathlineto{\pgfqpoint{6.913027in}{3.592686in}}%
\pgfpathlineto{\pgfqpoint{6.923422in}{3.577944in}}%
\pgfpathlineto{\pgfqpoint{6.944214in}{3.554984in}}%
\pgfpathlineto{\pgfqpoint{6.975400in}{3.525913in}}%
\pgfpathlineto{\pgfqpoint{7.016983in}{3.491441in}}%
\pgfpathlineto{\pgfqpoint{7.068961in}{3.451847in}}%
\pgfpathlineto{\pgfqpoint{7.141730in}{3.399815in}}%
\pgfpathlineto{\pgfqpoint{7.152126in}{3.392581in}}%
\pgfpathlineto{\pgfqpoint{7.162521in}{3.618652in}}%
\pgfpathlineto{\pgfqpoint{7.172917in}{3.591999in}}%
\pgfpathlineto{\pgfqpoint{7.183313in}{3.577225in}}%
\pgfpathlineto{\pgfqpoint{7.204104in}{3.554273in}}%
\pgfpathlineto{\pgfqpoint{7.235291in}{3.525269in}}%
\pgfpathlineto{\pgfqpoint{7.276873in}{3.490914in}}%
\pgfpathlineto{\pgfqpoint{7.328851in}{3.451465in}}%
\pgfpathlineto{\pgfqpoint{7.401621in}{3.399605in}}%
\pgfpathlineto{\pgfqpoint{7.412016in}{3.392393in}}%
\pgfpathlineto{\pgfqpoint{7.422412in}{3.618579in}}%
\pgfpathlineto{\pgfqpoint{7.432807in}{3.591921in}}%
\pgfpathlineto{\pgfqpoint{7.443203in}{3.577142in}}%
\pgfpathlineto{\pgfqpoint{7.463994in}{3.554188in}}%
\pgfpathlineto{\pgfqpoint{7.495181in}{3.525188in}}%
\pgfpathlineto{\pgfqpoint{7.536764in}{3.490846in}}%
\pgfpathlineto{\pgfqpoint{7.588742in}{3.451417in}}%
\pgfpathlineto{\pgfqpoint{7.661511in}{3.399586in}}%
\pgfpathlineto{\pgfqpoint{7.671907in}{3.392377in}}%
\pgfpathlineto{\pgfqpoint{7.682302in}{3.618564in}}%
\pgfpathlineto{\pgfqpoint{7.692698in}{3.591909in}}%
\pgfpathlineto{\pgfqpoint{7.703093in}{3.577130in}}%
\pgfpathlineto{\pgfqpoint{7.723885in}{3.554174in}}%
\pgfpathlineto{\pgfqpoint{7.755071in}{3.525175in}}%
\pgfpathlineto{\pgfqpoint{7.796654in}{3.490835in}}%
\pgfpathlineto{\pgfqpoint{7.848632in}{3.451410in}}%
\pgfpathlineto{\pgfqpoint{7.921401in}{3.399585in}}%
\pgfpathlineto{\pgfqpoint{7.931797in}{3.392377in}}%
\pgfpathlineto{\pgfqpoint{7.942193in}{3.618561in}}%
\pgfpathlineto{\pgfqpoint{7.952588in}{3.591907in}}%
\pgfpathlineto{\pgfqpoint{7.962984in}{3.577127in}}%
\pgfpathlineto{\pgfqpoint{7.983775in}{3.554172in}}%
\pgfpathlineto{\pgfqpoint{8.014962in}{3.525173in}}%
\pgfpathlineto{\pgfqpoint{8.056544in}{3.490833in}}%
\pgfpathlineto{\pgfqpoint{8.108522in}{3.451408in}}%
\pgfpathlineto{\pgfqpoint{8.181292in}{3.399585in}}%
\pgfpathlineto{\pgfqpoint{8.191687in}{3.392377in}}%
\pgfpathlineto{\pgfqpoint{8.202083in}{3.618567in}}%
\pgfpathlineto{\pgfqpoint{8.212479in}{3.591913in}}%
\pgfpathlineto{\pgfqpoint{8.222874in}{3.577134in}}%
\pgfpathlineto{\pgfqpoint{8.243665in}{3.554178in}}%
\pgfpathlineto{\pgfqpoint{8.274852in}{3.525179in}}%
\pgfpathlineto{\pgfqpoint{8.316435in}{3.490839in}}%
\pgfpathlineto{\pgfqpoint{8.368413in}{3.451415in}}%
\pgfpathlineto{\pgfqpoint{8.441182in}{3.399592in}}%
\pgfpathlineto{\pgfqpoint{8.451578in}{3.392385in}}%
\pgfpathlineto{\pgfqpoint{8.461973in}{3.618563in}}%
\pgfpathlineto{\pgfqpoint{8.472369in}{3.591910in}}%
\pgfpathlineto{\pgfqpoint{8.482765in}{3.577130in}}%
\pgfpathlineto{\pgfqpoint{8.503556in}{3.554174in}}%
\pgfpathlineto{\pgfqpoint{8.534743in}{3.525175in}}%
\pgfpathlineto{\pgfqpoint{8.576325in}{3.490836in}}%
\pgfpathlineto{\pgfqpoint{8.628303in}{3.451411in}}%
\pgfpathlineto{\pgfqpoint{8.701072in}{3.399588in}}%
\pgfpathlineto{\pgfqpoint{8.711468in}{3.392381in}}%
\pgfpathlineto{\pgfqpoint{8.721864in}{3.618567in}}%
\pgfpathlineto{\pgfqpoint{8.732259in}{3.591913in}}%
\pgfpathlineto{\pgfqpoint{8.742655in}{3.577133in}}%
\pgfpathlineto{\pgfqpoint{8.763446in}{3.554177in}}%
\pgfpathlineto{\pgfqpoint{8.794633in}{3.525178in}}%
\pgfpathlineto{\pgfqpoint{8.836215in}{3.490839in}}%
\pgfpathlineto{\pgfqpoint{8.888193in}{3.451414in}}%
\pgfpathlineto{\pgfqpoint{8.960963in}{3.399591in}}%
\pgfpathlineto{\pgfqpoint{8.971358in}{3.392384in}}%
\pgfpathlineto{\pgfqpoint{8.981754in}{3.618566in}}%
\pgfpathlineto{\pgfqpoint{8.992150in}{3.591913in}}%
\pgfpathlineto{\pgfqpoint{9.002545in}{3.577133in}}%
\pgfpathlineto{\pgfqpoint{9.023336in}{3.554177in}}%
\pgfpathlineto{\pgfqpoint{9.054523in}{3.525178in}}%
\pgfpathlineto{\pgfqpoint{9.096106in}{3.490838in}}%
\pgfpathlineto{\pgfqpoint{9.148084in}{3.451414in}}%
\pgfpathlineto{\pgfqpoint{9.220853in}{3.399591in}}%
\pgfpathlineto{\pgfqpoint{9.231249in}{3.392384in}}%
\pgfpathlineto{\pgfqpoint{9.241644in}{3.618566in}}%
\pgfpathlineto{\pgfqpoint{9.241644in}{3.618566in}}%
\pgfusepath{stroke}%
\end{pgfscope}%
\begin{pgfscope}%
\pgfpathrectangle{\pgfqpoint{5.977941in}{1.100000in}}{\pgfqpoint{3.419118in}{3.500000in}}%
\pgfusepath{clip}%
\pgfsetrectcap%
\pgfsetroundjoin%
\pgfsetlinewidth{1.505625pt}%
\definecolor{currentstroke}{rgb}{0.172549,0.627451,0.172549}%
\pgfsetstrokecolor{currentstroke}%
\pgfsetdash{}{0pt}%
\pgfpathmoveto{\pgfqpoint{6.133356in}{3.887113in}}%
\pgfpathlineto{\pgfqpoint{6.143751in}{3.822134in}}%
\pgfpathlineto{\pgfqpoint{6.154147in}{3.785586in}}%
\pgfpathlineto{\pgfqpoint{6.164542in}{3.756679in}}%
\pgfpathlineto{\pgfqpoint{6.185334in}{3.708736in}}%
\pgfpathlineto{\pgfqpoint{6.206125in}{3.668391in}}%
\pgfpathlineto{\pgfqpoint{6.226916in}{3.632939in}}%
\pgfpathlineto{\pgfqpoint{6.247707in}{3.601064in}}%
\pgfpathlineto{\pgfqpoint{6.278894in}{3.558266in}}%
\pgfpathlineto{\pgfqpoint{6.310081in}{3.519969in}}%
\pgfpathlineto{\pgfqpoint{6.341268in}{3.485002in}}%
\pgfpathlineto{\pgfqpoint{6.382850in}{3.442179in}}%
\pgfpathlineto{\pgfqpoint{6.424433in}{3.402512in}}%
\pgfpathlineto{\pgfqpoint{6.476411in}{3.356062in}}%
\pgfpathlineto{\pgfqpoint{6.538785in}{3.303539in}}%
\pgfpathlineto{\pgfqpoint{6.611554in}{3.245334in}}%
\pgfpathlineto{\pgfqpoint{6.694719in}{3.181656in}}%
\pgfpathlineto{\pgfqpoint{6.798675in}{3.105085in}}%
\pgfpathlineto{\pgfqpoint{6.923422in}{3.016316in}}%
\pgfpathlineto{\pgfqpoint{7.068961in}{2.915651in}}%
\pgfpathlineto{\pgfqpoint{7.256082in}{2.789075in}}%
\pgfpathlineto{\pgfqpoint{7.515972in}{2.616142in}}%
\pgfpathlineto{\pgfqpoint{7.952588in}{2.328643in}}%
\pgfpathlineto{\pgfqpoint{9.002545in}{1.640636in}}%
\pgfpathlineto{\pgfqpoint{9.241644in}{1.484189in}}%
\pgfpathlineto{\pgfqpoint{9.241644in}{1.484189in}}%
\pgfusepath{stroke}%
\end{pgfscope}%
\begin{pgfscope}%
\pgfpathrectangle{\pgfqpoint{5.977941in}{1.100000in}}{\pgfqpoint{3.419118in}{3.500000in}}%
\pgfusepath{clip}%
\pgfsetrectcap%
\pgfsetroundjoin%
\pgfsetlinewidth{1.505625pt}%
\definecolor{currentstroke}{rgb}{0.839216,0.152941,0.156863}%
\pgfsetstrokecolor{currentstroke}%
\pgfsetdash{}{0pt}%
\pgfpathmoveto{\pgfqpoint{6.133356in}{4.294613in}}%
\pgfpathlineto{\pgfqpoint{6.206125in}{4.239466in}}%
\pgfpathlineto{\pgfqpoint{6.299685in}{4.171648in}}%
\pgfpathlineto{\pgfqpoint{6.424433in}{4.084265in}}%
\pgfpathlineto{\pgfqpoint{6.590763in}{3.970808in}}%
\pgfpathlineto{\pgfqpoint{6.819466in}{3.817820in}}%
\pgfpathlineto{\pgfqpoint{7.172917in}{3.584396in}}%
\pgfpathlineto{\pgfqpoint{7.827841in}{3.154628in}}%
\pgfpathlineto{\pgfqpoint{9.241644in}{2.685850in}}%
\pgfpathlineto{\pgfqpoint{9.241644in}{2.685850in}}%
\pgfusepath{stroke}%
\end{pgfscope}%
\begin{pgfscope}%
\pgfsetrectcap%
\pgfsetmiterjoin%
\pgfsetlinewidth{0.803000pt}%
\definecolor{currentstroke}{rgb}{0.000000,0.000000,0.000000}%
\pgfsetstrokecolor{currentstroke}%
\pgfsetdash{}{0pt}%
\pgfpathmoveto{\pgfqpoint{5.977941in}{1.100000in}}%
\pgfpathlineto{\pgfqpoint{5.977941in}{4.600000in}}%
\pgfusepath{stroke}%
\end{pgfscope}%
\begin{pgfscope}%
\pgfsetrectcap%
\pgfsetmiterjoin%
\pgfsetlinewidth{0.803000pt}%
\definecolor{currentstroke}{rgb}{0.000000,0.000000,0.000000}%
\pgfsetstrokecolor{currentstroke}%
\pgfsetdash{}{0pt}%
\pgfpathmoveto{\pgfqpoint{9.397059in}{1.100000in}}%
\pgfpathlineto{\pgfqpoint{9.397059in}{4.600000in}}%
\pgfusepath{stroke}%
\end{pgfscope}%
\begin{pgfscope}%
\pgfsetrectcap%
\pgfsetmiterjoin%
\pgfsetlinewidth{0.803000pt}%
\definecolor{currentstroke}{rgb}{0.000000,0.000000,0.000000}%
\pgfsetstrokecolor{currentstroke}%
\pgfsetdash{}{0pt}%
\pgfpathmoveto{\pgfqpoint{5.977941in}{1.100000in}}%
\pgfpathlineto{\pgfqpoint{9.397059in}{1.100000in}}%
\pgfusepath{stroke}%
\end{pgfscope}%
\begin{pgfscope}%
\pgfsetrectcap%
\pgfsetmiterjoin%
\pgfsetlinewidth{0.803000pt}%
\definecolor{currentstroke}{rgb}{0.000000,0.000000,0.000000}%
\pgfsetstrokecolor{currentstroke}%
\pgfsetdash{}{0pt}%
\pgfpathmoveto{\pgfqpoint{5.977941in}{4.600000in}}%
\pgfpathlineto{\pgfqpoint{9.397059in}{4.600000in}}%
\pgfusepath{stroke}%
\end{pgfscope}%
\begin{pgfscope}%
\definecolor{textcolor}{rgb}{0.000000,0.000000,0.000000}%
\pgfsetstrokecolor{textcolor}%
\pgfsetfillcolor{textcolor}%
\pgftext[x=7.687500in,y=4.683333in,,base]{\color{textcolor}{\sffamily\fontsize{12.000000}{14.400000}\selectfont\catcode`\^=\active\def^{\ifmmode\sp\else\^{}\fi}\catcode`\%=\active\def%{\%}GS}}%
\end{pgfscope}%
\begin{pgfscope}%
\pgfsetbuttcap%
\pgfsetmiterjoin%
\definecolor{currentfill}{rgb}{1.000000,1.000000,1.000000}%
\pgfsetfillcolor{currentfill}%
\pgfsetlinewidth{0.000000pt}%
\definecolor{currentstroke}{rgb}{0.000000,0.000000,0.000000}%
\pgfsetstrokecolor{currentstroke}%
\pgfsetstrokeopacity{0.000000}%
\pgfsetdash{}{0pt}%
\pgfpathmoveto{\pgfqpoint{10.080882in}{1.100000in}}%
\pgfpathlineto{\pgfqpoint{13.500000in}{1.100000in}}%
\pgfpathlineto{\pgfqpoint{13.500000in}{4.600000in}}%
\pgfpathlineto{\pgfqpoint{10.080882in}{4.600000in}}%
\pgfpathlineto{\pgfqpoint{10.080882in}{1.100000in}}%
\pgfpathclose%
\pgfusepath{fill}%
\end{pgfscope}%
\begin{pgfscope}%
\pgfsetbuttcap%
\pgfsetroundjoin%
\definecolor{currentfill}{rgb}{0.000000,0.000000,0.000000}%
\pgfsetfillcolor{currentfill}%
\pgfsetlinewidth{0.803000pt}%
\definecolor{currentstroke}{rgb}{0.000000,0.000000,0.000000}%
\pgfsetstrokecolor{currentstroke}%
\pgfsetdash{}{0pt}%
\pgfsys@defobject{currentmarker}{\pgfqpoint{0.000000in}{-0.048611in}}{\pgfqpoint{0.000000in}{0.000000in}}{%
\pgfpathmoveto{\pgfqpoint{0.000000in}{0.000000in}}%
\pgfpathlineto{\pgfqpoint{0.000000in}{-0.048611in}}%
\pgfusepath{stroke,fill}%
}%
\begin{pgfscope}%
\pgfsys@transformshift{10.236297in}{1.100000in}%
\pgfsys@useobject{currentmarker}{}%
\end{pgfscope}%
\end{pgfscope}%
\begin{pgfscope}%
\definecolor{textcolor}{rgb}{0.000000,0.000000,0.000000}%
\pgfsetstrokecolor{textcolor}%
\pgfsetfillcolor{textcolor}%
\pgftext[x=10.236297in,y=1.002778in,,top]{\color{textcolor}{\sffamily\fontsize{10.000000}{12.000000}\selectfont\catcode`\^=\active\def^{\ifmmode\sp\else\^{}\fi}\catcode`\%=\active\def%{\%}0}}%
\end{pgfscope}%
\begin{pgfscope}%
\pgfsetbuttcap%
\pgfsetroundjoin%
\definecolor{currentfill}{rgb}{0.000000,0.000000,0.000000}%
\pgfsetfillcolor{currentfill}%
\pgfsetlinewidth{0.803000pt}%
\definecolor{currentstroke}{rgb}{0.000000,0.000000,0.000000}%
\pgfsetstrokecolor{currentstroke}%
\pgfsetdash{}{0pt}%
\pgfsys@defobject{currentmarker}{\pgfqpoint{0.000000in}{-0.048611in}}{\pgfqpoint{0.000000in}{0.000000in}}{%
\pgfpathmoveto{\pgfqpoint{0.000000in}{0.000000in}}%
\pgfpathlineto{\pgfqpoint{0.000000in}{-0.048611in}}%
\pgfusepath{stroke,fill}%
}%
\begin{pgfscope}%
\pgfsys@transformshift{10.864234in}{1.100000in}%
\pgfsys@useobject{currentmarker}{}%
\end{pgfscope}%
\end{pgfscope}%
\begin{pgfscope}%
\definecolor{textcolor}{rgb}{0.000000,0.000000,0.000000}%
\pgfsetstrokecolor{textcolor}%
\pgfsetfillcolor{textcolor}%
\pgftext[x=10.864234in,y=1.002778in,,top]{\color{textcolor}{\sffamily\fontsize{10.000000}{12.000000}\selectfont\catcode`\^=\active\def^{\ifmmode\sp\else\^{}\fi}\catcode`\%=\active\def%{\%}20}}%
\end{pgfscope}%
\begin{pgfscope}%
\pgfsetbuttcap%
\pgfsetroundjoin%
\definecolor{currentfill}{rgb}{0.000000,0.000000,0.000000}%
\pgfsetfillcolor{currentfill}%
\pgfsetlinewidth{0.803000pt}%
\definecolor{currentstroke}{rgb}{0.000000,0.000000,0.000000}%
\pgfsetstrokecolor{currentstroke}%
\pgfsetdash{}{0pt}%
\pgfsys@defobject{currentmarker}{\pgfqpoint{0.000000in}{-0.048611in}}{\pgfqpoint{0.000000in}{0.000000in}}{%
\pgfpathmoveto{\pgfqpoint{0.000000in}{0.000000in}}%
\pgfpathlineto{\pgfqpoint{0.000000in}{-0.048611in}}%
\pgfusepath{stroke,fill}%
}%
\begin{pgfscope}%
\pgfsys@transformshift{11.492171in}{1.100000in}%
\pgfsys@useobject{currentmarker}{}%
\end{pgfscope}%
\end{pgfscope}%
\begin{pgfscope}%
\definecolor{textcolor}{rgb}{0.000000,0.000000,0.000000}%
\pgfsetstrokecolor{textcolor}%
\pgfsetfillcolor{textcolor}%
\pgftext[x=11.492171in,y=1.002778in,,top]{\color{textcolor}{\sffamily\fontsize{10.000000}{12.000000}\selectfont\catcode`\^=\active\def^{\ifmmode\sp\else\^{}\fi}\catcode`\%=\active\def%{\%}40}}%
\end{pgfscope}%
\begin{pgfscope}%
\pgfsetbuttcap%
\pgfsetroundjoin%
\definecolor{currentfill}{rgb}{0.000000,0.000000,0.000000}%
\pgfsetfillcolor{currentfill}%
\pgfsetlinewidth{0.803000pt}%
\definecolor{currentstroke}{rgb}{0.000000,0.000000,0.000000}%
\pgfsetstrokecolor{currentstroke}%
\pgfsetdash{}{0pt}%
\pgfsys@defobject{currentmarker}{\pgfqpoint{0.000000in}{-0.048611in}}{\pgfqpoint{0.000000in}{0.000000in}}{%
\pgfpathmoveto{\pgfqpoint{0.000000in}{0.000000in}}%
\pgfpathlineto{\pgfqpoint{0.000000in}{-0.048611in}}%
\pgfusepath{stroke,fill}%
}%
\begin{pgfscope}%
\pgfsys@transformshift{12.120108in}{1.100000in}%
\pgfsys@useobject{currentmarker}{}%
\end{pgfscope}%
\end{pgfscope}%
\begin{pgfscope}%
\definecolor{textcolor}{rgb}{0.000000,0.000000,0.000000}%
\pgfsetstrokecolor{textcolor}%
\pgfsetfillcolor{textcolor}%
\pgftext[x=12.120108in,y=1.002778in,,top]{\color{textcolor}{\sffamily\fontsize{10.000000}{12.000000}\selectfont\catcode`\^=\active\def^{\ifmmode\sp\else\^{}\fi}\catcode`\%=\active\def%{\%}60}}%
\end{pgfscope}%
\begin{pgfscope}%
\pgfsetbuttcap%
\pgfsetroundjoin%
\definecolor{currentfill}{rgb}{0.000000,0.000000,0.000000}%
\pgfsetfillcolor{currentfill}%
\pgfsetlinewidth{0.803000pt}%
\definecolor{currentstroke}{rgb}{0.000000,0.000000,0.000000}%
\pgfsetstrokecolor{currentstroke}%
\pgfsetdash{}{0pt}%
\pgfsys@defobject{currentmarker}{\pgfqpoint{0.000000in}{-0.048611in}}{\pgfqpoint{0.000000in}{0.000000in}}{%
\pgfpathmoveto{\pgfqpoint{0.000000in}{0.000000in}}%
\pgfpathlineto{\pgfqpoint{0.000000in}{-0.048611in}}%
\pgfusepath{stroke,fill}%
}%
\begin{pgfscope}%
\pgfsys@transformshift{12.748045in}{1.100000in}%
\pgfsys@useobject{currentmarker}{}%
\end{pgfscope}%
\end{pgfscope}%
\begin{pgfscope}%
\definecolor{textcolor}{rgb}{0.000000,0.000000,0.000000}%
\pgfsetstrokecolor{textcolor}%
\pgfsetfillcolor{textcolor}%
\pgftext[x=12.748045in,y=1.002778in,,top]{\color{textcolor}{\sffamily\fontsize{10.000000}{12.000000}\selectfont\catcode`\^=\active\def^{\ifmmode\sp\else\^{}\fi}\catcode`\%=\active\def%{\%}80}}%
\end{pgfscope}%
\begin{pgfscope}%
\pgfsetbuttcap%
\pgfsetroundjoin%
\definecolor{currentfill}{rgb}{0.000000,0.000000,0.000000}%
\pgfsetfillcolor{currentfill}%
\pgfsetlinewidth{0.803000pt}%
\definecolor{currentstroke}{rgb}{0.000000,0.000000,0.000000}%
\pgfsetstrokecolor{currentstroke}%
\pgfsetdash{}{0pt}%
\pgfsys@defobject{currentmarker}{\pgfqpoint{0.000000in}{-0.048611in}}{\pgfqpoint{0.000000in}{0.000000in}}{%
\pgfpathmoveto{\pgfqpoint{0.000000in}{0.000000in}}%
\pgfpathlineto{\pgfqpoint{0.000000in}{-0.048611in}}%
\pgfusepath{stroke,fill}%
}%
\begin{pgfscope}%
\pgfsys@transformshift{13.375982in}{1.100000in}%
\pgfsys@useobject{currentmarker}{}%
\end{pgfscope}%
\end{pgfscope}%
\begin{pgfscope}%
\definecolor{textcolor}{rgb}{0.000000,0.000000,0.000000}%
\pgfsetstrokecolor{textcolor}%
\pgfsetfillcolor{textcolor}%
\pgftext[x=13.375982in,y=1.002778in,,top]{\color{textcolor}{\sffamily\fontsize{10.000000}{12.000000}\selectfont\catcode`\^=\active\def^{\ifmmode\sp\else\^{}\fi}\catcode`\%=\active\def%{\%}100}}%
\end{pgfscope}%
\begin{pgfscope}%
\definecolor{textcolor}{rgb}{0.000000,0.000000,0.000000}%
\pgfsetstrokecolor{textcolor}%
\pgfsetfillcolor{textcolor}%
\pgftext[x=11.790441in,y=0.812809in,,top]{\color{textcolor}{\sffamily\fontsize{10.000000}{12.000000}\selectfont\catcode`\^=\active\def^{\ifmmode\sp\else\^{}\fi}\catcode`\%=\active\def%{\%}V cycles}}%
\end{pgfscope}%
\begin{pgfscope}%
\pgfsetbuttcap%
\pgfsetroundjoin%
\definecolor{currentfill}{rgb}{0.000000,0.000000,0.000000}%
\pgfsetfillcolor{currentfill}%
\pgfsetlinewidth{0.803000pt}%
\definecolor{currentstroke}{rgb}{0.000000,0.000000,0.000000}%
\pgfsetstrokecolor{currentstroke}%
\pgfsetdash{}{0pt}%
\pgfsys@defobject{currentmarker}{\pgfqpoint{-0.048611in}{0.000000in}}{\pgfqpoint{-0.000000in}{0.000000in}}{%
\pgfpathmoveto{\pgfqpoint{-0.000000in}{0.000000in}}%
\pgfpathlineto{\pgfqpoint{-0.048611in}{0.000000in}}%
\pgfusepath{stroke,fill}%
}%
\begin{pgfscope}%
\pgfsys@transformshift{10.080882in}{1.660284in}%
\pgfsys@useobject{currentmarker}{}%
\end{pgfscope}%
\end{pgfscope}%
\begin{pgfscope}%
\pgfsetbuttcap%
\pgfsetroundjoin%
\definecolor{currentfill}{rgb}{0.000000,0.000000,0.000000}%
\pgfsetfillcolor{currentfill}%
\pgfsetlinewidth{0.803000pt}%
\definecolor{currentstroke}{rgb}{0.000000,0.000000,0.000000}%
\pgfsetstrokecolor{currentstroke}%
\pgfsetdash{}{0pt}%
\pgfsys@defobject{currentmarker}{\pgfqpoint{-0.048611in}{0.000000in}}{\pgfqpoint{-0.000000in}{0.000000in}}{%
\pgfpathmoveto{\pgfqpoint{-0.000000in}{0.000000in}}%
\pgfpathlineto{\pgfqpoint{-0.048611in}{0.000000in}}%
\pgfusepath{stroke,fill}%
}%
\begin{pgfscope}%
\pgfsys@transformshift{10.080882in}{2.534463in}%
\pgfsys@useobject{currentmarker}{}%
\end{pgfscope}%
\end{pgfscope}%
\begin{pgfscope}%
\pgfsetbuttcap%
\pgfsetroundjoin%
\definecolor{currentfill}{rgb}{0.000000,0.000000,0.000000}%
\pgfsetfillcolor{currentfill}%
\pgfsetlinewidth{0.803000pt}%
\definecolor{currentstroke}{rgb}{0.000000,0.000000,0.000000}%
\pgfsetstrokecolor{currentstroke}%
\pgfsetdash{}{0pt}%
\pgfsys@defobject{currentmarker}{\pgfqpoint{-0.048611in}{0.000000in}}{\pgfqpoint{-0.000000in}{0.000000in}}{%
\pgfpathmoveto{\pgfqpoint{-0.000000in}{0.000000in}}%
\pgfpathlineto{\pgfqpoint{-0.048611in}{0.000000in}}%
\pgfusepath{stroke,fill}%
}%
\begin{pgfscope}%
\pgfsys@transformshift{10.080882in}{3.408642in}%
\pgfsys@useobject{currentmarker}{}%
\end{pgfscope}%
\end{pgfscope}%
\begin{pgfscope}%
\pgfsetbuttcap%
\pgfsetroundjoin%
\definecolor{currentfill}{rgb}{0.000000,0.000000,0.000000}%
\pgfsetfillcolor{currentfill}%
\pgfsetlinewidth{0.803000pt}%
\definecolor{currentstroke}{rgb}{0.000000,0.000000,0.000000}%
\pgfsetstrokecolor{currentstroke}%
\pgfsetdash{}{0pt}%
\pgfsys@defobject{currentmarker}{\pgfqpoint{-0.048611in}{0.000000in}}{\pgfqpoint{-0.000000in}{0.000000in}}{%
\pgfpathmoveto{\pgfqpoint{-0.000000in}{0.000000in}}%
\pgfpathlineto{\pgfqpoint{-0.048611in}{0.000000in}}%
\pgfusepath{stroke,fill}%
}%
\begin{pgfscope}%
\pgfsys@transformshift{10.080882in}{4.282822in}%
\pgfsys@useobject{currentmarker}{}%
\end{pgfscope}%
\end{pgfscope}%
\begin{pgfscope}%
\pgfsetbuttcap%
\pgfsetroundjoin%
\definecolor{currentfill}{rgb}{0.000000,0.000000,0.000000}%
\pgfsetfillcolor{currentfill}%
\pgfsetlinewidth{0.602250pt}%
\definecolor{currentstroke}{rgb}{0.000000,0.000000,0.000000}%
\pgfsetstrokecolor{currentstroke}%
\pgfsetdash{}{0pt}%
\pgfsys@defobject{currentmarker}{\pgfqpoint{-0.027778in}{0.000000in}}{\pgfqpoint{-0.000000in}{0.000000in}}{%
\pgfpathmoveto{\pgfqpoint{-0.000000in}{0.000000in}}%
\pgfpathlineto{\pgfqpoint{-0.027778in}{0.000000in}}%
\pgfusepath{stroke,fill}%
}%
\begin{pgfscope}%
\pgfsys@transformshift{10.080882in}{1.203194in}%
\pgfsys@useobject{currentmarker}{}%
\end{pgfscope}%
\end{pgfscope}%
\begin{pgfscope}%
\pgfsetbuttcap%
\pgfsetroundjoin%
\definecolor{currentfill}{rgb}{0.000000,0.000000,0.000000}%
\pgfsetfillcolor{currentfill}%
\pgfsetlinewidth{0.602250pt}%
\definecolor{currentstroke}{rgb}{0.000000,0.000000,0.000000}%
\pgfsetstrokecolor{currentstroke}%
\pgfsetdash{}{0pt}%
\pgfsys@defobject{currentmarker}{\pgfqpoint{-0.027778in}{0.000000in}}{\pgfqpoint{-0.000000in}{0.000000in}}{%
\pgfpathmoveto{\pgfqpoint{-0.000000in}{0.000000in}}%
\pgfpathlineto{\pgfqpoint{-0.027778in}{0.000000in}}%
\pgfusepath{stroke,fill}%
}%
\begin{pgfscope}%
\pgfsys@transformshift{10.080882in}{1.312413in}%
\pgfsys@useobject{currentmarker}{}%
\end{pgfscope}%
\end{pgfscope}%
\begin{pgfscope}%
\pgfsetbuttcap%
\pgfsetroundjoin%
\definecolor{currentfill}{rgb}{0.000000,0.000000,0.000000}%
\pgfsetfillcolor{currentfill}%
\pgfsetlinewidth{0.602250pt}%
\definecolor{currentstroke}{rgb}{0.000000,0.000000,0.000000}%
\pgfsetstrokecolor{currentstroke}%
\pgfsetdash{}{0pt}%
\pgfsys@defobject{currentmarker}{\pgfqpoint{-0.027778in}{0.000000in}}{\pgfqpoint{-0.000000in}{0.000000in}}{%
\pgfpathmoveto{\pgfqpoint{-0.000000in}{0.000000in}}%
\pgfpathlineto{\pgfqpoint{-0.027778in}{0.000000in}}%
\pgfusepath{stroke,fill}%
}%
\begin{pgfscope}%
\pgfsys@transformshift{10.080882in}{1.397130in}%
\pgfsys@useobject{currentmarker}{}%
\end{pgfscope}%
\end{pgfscope}%
\begin{pgfscope}%
\pgfsetbuttcap%
\pgfsetroundjoin%
\definecolor{currentfill}{rgb}{0.000000,0.000000,0.000000}%
\pgfsetfillcolor{currentfill}%
\pgfsetlinewidth{0.602250pt}%
\definecolor{currentstroke}{rgb}{0.000000,0.000000,0.000000}%
\pgfsetstrokecolor{currentstroke}%
\pgfsetdash{}{0pt}%
\pgfsys@defobject{currentmarker}{\pgfqpoint{-0.027778in}{0.000000in}}{\pgfqpoint{-0.000000in}{0.000000in}}{%
\pgfpathmoveto{\pgfqpoint{-0.000000in}{0.000000in}}%
\pgfpathlineto{\pgfqpoint{-0.027778in}{0.000000in}}%
\pgfusepath{stroke,fill}%
}%
\begin{pgfscope}%
\pgfsys@transformshift{10.080882in}{1.466348in}%
\pgfsys@useobject{currentmarker}{}%
\end{pgfscope}%
\end{pgfscope}%
\begin{pgfscope}%
\pgfsetbuttcap%
\pgfsetroundjoin%
\definecolor{currentfill}{rgb}{0.000000,0.000000,0.000000}%
\pgfsetfillcolor{currentfill}%
\pgfsetlinewidth{0.602250pt}%
\definecolor{currentstroke}{rgb}{0.000000,0.000000,0.000000}%
\pgfsetstrokecolor{currentstroke}%
\pgfsetdash{}{0pt}%
\pgfsys@defobject{currentmarker}{\pgfqpoint{-0.027778in}{0.000000in}}{\pgfqpoint{-0.000000in}{0.000000in}}{%
\pgfpathmoveto{\pgfqpoint{-0.000000in}{0.000000in}}%
\pgfpathlineto{\pgfqpoint{-0.027778in}{0.000000in}}%
\pgfusepath{stroke,fill}%
}%
\begin{pgfscope}%
\pgfsys@transformshift{10.080882in}{1.524872in}%
\pgfsys@useobject{currentmarker}{}%
\end{pgfscope}%
\end{pgfscope}%
\begin{pgfscope}%
\pgfsetbuttcap%
\pgfsetroundjoin%
\definecolor{currentfill}{rgb}{0.000000,0.000000,0.000000}%
\pgfsetfillcolor{currentfill}%
\pgfsetlinewidth{0.602250pt}%
\definecolor{currentstroke}{rgb}{0.000000,0.000000,0.000000}%
\pgfsetstrokecolor{currentstroke}%
\pgfsetdash{}{0pt}%
\pgfsys@defobject{currentmarker}{\pgfqpoint{-0.027778in}{0.000000in}}{\pgfqpoint{-0.000000in}{0.000000in}}{%
\pgfpathmoveto{\pgfqpoint{-0.000000in}{0.000000in}}%
\pgfpathlineto{\pgfqpoint{-0.027778in}{0.000000in}}%
\pgfusepath{stroke,fill}%
}%
\begin{pgfscope}%
\pgfsys@transformshift{10.080882in}{1.575567in}%
\pgfsys@useobject{currentmarker}{}%
\end{pgfscope}%
\end{pgfscope}%
\begin{pgfscope}%
\pgfsetbuttcap%
\pgfsetroundjoin%
\definecolor{currentfill}{rgb}{0.000000,0.000000,0.000000}%
\pgfsetfillcolor{currentfill}%
\pgfsetlinewidth{0.602250pt}%
\definecolor{currentstroke}{rgb}{0.000000,0.000000,0.000000}%
\pgfsetstrokecolor{currentstroke}%
\pgfsetdash{}{0pt}%
\pgfsys@defobject{currentmarker}{\pgfqpoint{-0.027778in}{0.000000in}}{\pgfqpoint{-0.000000in}{0.000000in}}{%
\pgfpathmoveto{\pgfqpoint{-0.000000in}{0.000000in}}%
\pgfpathlineto{\pgfqpoint{-0.027778in}{0.000000in}}%
\pgfusepath{stroke,fill}%
}%
\begin{pgfscope}%
\pgfsys@transformshift{10.080882in}{1.620284in}%
\pgfsys@useobject{currentmarker}{}%
\end{pgfscope}%
\end{pgfscope}%
\begin{pgfscope}%
\pgfsetbuttcap%
\pgfsetroundjoin%
\definecolor{currentfill}{rgb}{0.000000,0.000000,0.000000}%
\pgfsetfillcolor{currentfill}%
\pgfsetlinewidth{0.602250pt}%
\definecolor{currentstroke}{rgb}{0.000000,0.000000,0.000000}%
\pgfsetstrokecolor{currentstroke}%
\pgfsetdash{}{0pt}%
\pgfsys@defobject{currentmarker}{\pgfqpoint{-0.027778in}{0.000000in}}{\pgfqpoint{-0.000000in}{0.000000in}}{%
\pgfpathmoveto{\pgfqpoint{-0.000000in}{0.000000in}}%
\pgfpathlineto{\pgfqpoint{-0.027778in}{0.000000in}}%
\pgfusepath{stroke,fill}%
}%
\begin{pgfscope}%
\pgfsys@transformshift{10.080882in}{1.923438in}%
\pgfsys@useobject{currentmarker}{}%
\end{pgfscope}%
\end{pgfscope}%
\begin{pgfscope}%
\pgfsetbuttcap%
\pgfsetroundjoin%
\definecolor{currentfill}{rgb}{0.000000,0.000000,0.000000}%
\pgfsetfillcolor{currentfill}%
\pgfsetlinewidth{0.602250pt}%
\definecolor{currentstroke}{rgb}{0.000000,0.000000,0.000000}%
\pgfsetstrokecolor{currentstroke}%
\pgfsetdash{}{0pt}%
\pgfsys@defobject{currentmarker}{\pgfqpoint{-0.027778in}{0.000000in}}{\pgfqpoint{-0.000000in}{0.000000in}}{%
\pgfpathmoveto{\pgfqpoint{-0.000000in}{0.000000in}}%
\pgfpathlineto{\pgfqpoint{-0.027778in}{0.000000in}}%
\pgfusepath{stroke,fill}%
}%
\begin{pgfscope}%
\pgfsys@transformshift{10.080882in}{2.077373in}%
\pgfsys@useobject{currentmarker}{}%
\end{pgfscope}%
\end{pgfscope}%
\begin{pgfscope}%
\pgfsetbuttcap%
\pgfsetroundjoin%
\definecolor{currentfill}{rgb}{0.000000,0.000000,0.000000}%
\pgfsetfillcolor{currentfill}%
\pgfsetlinewidth{0.602250pt}%
\definecolor{currentstroke}{rgb}{0.000000,0.000000,0.000000}%
\pgfsetstrokecolor{currentstroke}%
\pgfsetdash{}{0pt}%
\pgfsys@defobject{currentmarker}{\pgfqpoint{-0.027778in}{0.000000in}}{\pgfqpoint{-0.000000in}{0.000000in}}{%
\pgfpathmoveto{\pgfqpoint{-0.000000in}{0.000000in}}%
\pgfpathlineto{\pgfqpoint{-0.027778in}{0.000000in}}%
\pgfusepath{stroke,fill}%
}%
\begin{pgfscope}%
\pgfsys@transformshift{10.080882in}{2.186592in}%
\pgfsys@useobject{currentmarker}{}%
\end{pgfscope}%
\end{pgfscope}%
\begin{pgfscope}%
\pgfsetbuttcap%
\pgfsetroundjoin%
\definecolor{currentfill}{rgb}{0.000000,0.000000,0.000000}%
\pgfsetfillcolor{currentfill}%
\pgfsetlinewidth{0.602250pt}%
\definecolor{currentstroke}{rgb}{0.000000,0.000000,0.000000}%
\pgfsetstrokecolor{currentstroke}%
\pgfsetdash{}{0pt}%
\pgfsys@defobject{currentmarker}{\pgfqpoint{-0.027778in}{0.000000in}}{\pgfqpoint{-0.000000in}{0.000000in}}{%
\pgfpathmoveto{\pgfqpoint{-0.000000in}{0.000000in}}%
\pgfpathlineto{\pgfqpoint{-0.027778in}{0.000000in}}%
\pgfusepath{stroke,fill}%
}%
\begin{pgfscope}%
\pgfsys@transformshift{10.080882in}{2.271309in}%
\pgfsys@useobject{currentmarker}{}%
\end{pgfscope}%
\end{pgfscope}%
\begin{pgfscope}%
\pgfsetbuttcap%
\pgfsetroundjoin%
\definecolor{currentfill}{rgb}{0.000000,0.000000,0.000000}%
\pgfsetfillcolor{currentfill}%
\pgfsetlinewidth{0.602250pt}%
\definecolor{currentstroke}{rgb}{0.000000,0.000000,0.000000}%
\pgfsetstrokecolor{currentstroke}%
\pgfsetdash{}{0pt}%
\pgfsys@defobject{currentmarker}{\pgfqpoint{-0.027778in}{0.000000in}}{\pgfqpoint{-0.000000in}{0.000000in}}{%
\pgfpathmoveto{\pgfqpoint{-0.000000in}{0.000000in}}%
\pgfpathlineto{\pgfqpoint{-0.027778in}{0.000000in}}%
\pgfusepath{stroke,fill}%
}%
\begin{pgfscope}%
\pgfsys@transformshift{10.080882in}{2.340528in}%
\pgfsys@useobject{currentmarker}{}%
\end{pgfscope}%
\end{pgfscope}%
\begin{pgfscope}%
\pgfsetbuttcap%
\pgfsetroundjoin%
\definecolor{currentfill}{rgb}{0.000000,0.000000,0.000000}%
\pgfsetfillcolor{currentfill}%
\pgfsetlinewidth{0.602250pt}%
\definecolor{currentstroke}{rgb}{0.000000,0.000000,0.000000}%
\pgfsetstrokecolor{currentstroke}%
\pgfsetdash{}{0pt}%
\pgfsys@defobject{currentmarker}{\pgfqpoint{-0.027778in}{0.000000in}}{\pgfqpoint{-0.000000in}{0.000000in}}{%
\pgfpathmoveto{\pgfqpoint{-0.000000in}{0.000000in}}%
\pgfpathlineto{\pgfqpoint{-0.027778in}{0.000000in}}%
\pgfusepath{stroke,fill}%
}%
\begin{pgfscope}%
\pgfsys@transformshift{10.080882in}{2.399051in}%
\pgfsys@useobject{currentmarker}{}%
\end{pgfscope}%
\end{pgfscope}%
\begin{pgfscope}%
\pgfsetbuttcap%
\pgfsetroundjoin%
\definecolor{currentfill}{rgb}{0.000000,0.000000,0.000000}%
\pgfsetfillcolor{currentfill}%
\pgfsetlinewidth{0.602250pt}%
\definecolor{currentstroke}{rgb}{0.000000,0.000000,0.000000}%
\pgfsetstrokecolor{currentstroke}%
\pgfsetdash{}{0pt}%
\pgfsys@defobject{currentmarker}{\pgfqpoint{-0.027778in}{0.000000in}}{\pgfqpoint{-0.000000in}{0.000000in}}{%
\pgfpathmoveto{\pgfqpoint{-0.000000in}{0.000000in}}%
\pgfpathlineto{\pgfqpoint{-0.027778in}{0.000000in}}%
\pgfusepath{stroke,fill}%
}%
\begin{pgfscope}%
\pgfsys@transformshift{10.080882in}{2.449746in}%
\pgfsys@useobject{currentmarker}{}%
\end{pgfscope}%
\end{pgfscope}%
\begin{pgfscope}%
\pgfsetbuttcap%
\pgfsetroundjoin%
\definecolor{currentfill}{rgb}{0.000000,0.000000,0.000000}%
\pgfsetfillcolor{currentfill}%
\pgfsetlinewidth{0.602250pt}%
\definecolor{currentstroke}{rgb}{0.000000,0.000000,0.000000}%
\pgfsetstrokecolor{currentstroke}%
\pgfsetdash{}{0pt}%
\pgfsys@defobject{currentmarker}{\pgfqpoint{-0.027778in}{0.000000in}}{\pgfqpoint{-0.000000in}{0.000000in}}{%
\pgfpathmoveto{\pgfqpoint{-0.000000in}{0.000000in}}%
\pgfpathlineto{\pgfqpoint{-0.027778in}{0.000000in}}%
\pgfusepath{stroke,fill}%
}%
\begin{pgfscope}%
\pgfsys@transformshift{10.080882in}{2.494463in}%
\pgfsys@useobject{currentmarker}{}%
\end{pgfscope}%
\end{pgfscope}%
\begin{pgfscope}%
\pgfsetbuttcap%
\pgfsetroundjoin%
\definecolor{currentfill}{rgb}{0.000000,0.000000,0.000000}%
\pgfsetfillcolor{currentfill}%
\pgfsetlinewidth{0.602250pt}%
\definecolor{currentstroke}{rgb}{0.000000,0.000000,0.000000}%
\pgfsetstrokecolor{currentstroke}%
\pgfsetdash{}{0pt}%
\pgfsys@defobject{currentmarker}{\pgfqpoint{-0.027778in}{0.000000in}}{\pgfqpoint{-0.000000in}{0.000000in}}{%
\pgfpathmoveto{\pgfqpoint{-0.000000in}{0.000000in}}%
\pgfpathlineto{\pgfqpoint{-0.027778in}{0.000000in}}%
\pgfusepath{stroke,fill}%
}%
\begin{pgfscope}%
\pgfsys@transformshift{10.080882in}{2.797617in}%
\pgfsys@useobject{currentmarker}{}%
\end{pgfscope}%
\end{pgfscope}%
\begin{pgfscope}%
\pgfsetbuttcap%
\pgfsetroundjoin%
\definecolor{currentfill}{rgb}{0.000000,0.000000,0.000000}%
\pgfsetfillcolor{currentfill}%
\pgfsetlinewidth{0.602250pt}%
\definecolor{currentstroke}{rgb}{0.000000,0.000000,0.000000}%
\pgfsetstrokecolor{currentstroke}%
\pgfsetdash{}{0pt}%
\pgfsys@defobject{currentmarker}{\pgfqpoint{-0.027778in}{0.000000in}}{\pgfqpoint{-0.000000in}{0.000000in}}{%
\pgfpathmoveto{\pgfqpoint{-0.000000in}{0.000000in}}%
\pgfpathlineto{\pgfqpoint{-0.027778in}{0.000000in}}%
\pgfusepath{stroke,fill}%
}%
\begin{pgfscope}%
\pgfsys@transformshift{10.080882in}{2.951553in}%
\pgfsys@useobject{currentmarker}{}%
\end{pgfscope}%
\end{pgfscope}%
\begin{pgfscope}%
\pgfsetbuttcap%
\pgfsetroundjoin%
\definecolor{currentfill}{rgb}{0.000000,0.000000,0.000000}%
\pgfsetfillcolor{currentfill}%
\pgfsetlinewidth{0.602250pt}%
\definecolor{currentstroke}{rgb}{0.000000,0.000000,0.000000}%
\pgfsetstrokecolor{currentstroke}%
\pgfsetdash{}{0pt}%
\pgfsys@defobject{currentmarker}{\pgfqpoint{-0.027778in}{0.000000in}}{\pgfqpoint{-0.000000in}{0.000000in}}{%
\pgfpathmoveto{\pgfqpoint{-0.000000in}{0.000000in}}%
\pgfpathlineto{\pgfqpoint{-0.027778in}{0.000000in}}%
\pgfusepath{stroke,fill}%
}%
\begin{pgfscope}%
\pgfsys@transformshift{10.080882in}{3.060771in}%
\pgfsys@useobject{currentmarker}{}%
\end{pgfscope}%
\end{pgfscope}%
\begin{pgfscope}%
\pgfsetbuttcap%
\pgfsetroundjoin%
\definecolor{currentfill}{rgb}{0.000000,0.000000,0.000000}%
\pgfsetfillcolor{currentfill}%
\pgfsetlinewidth{0.602250pt}%
\definecolor{currentstroke}{rgb}{0.000000,0.000000,0.000000}%
\pgfsetstrokecolor{currentstroke}%
\pgfsetdash{}{0pt}%
\pgfsys@defobject{currentmarker}{\pgfqpoint{-0.027778in}{0.000000in}}{\pgfqpoint{-0.000000in}{0.000000in}}{%
\pgfpathmoveto{\pgfqpoint{-0.000000in}{0.000000in}}%
\pgfpathlineto{\pgfqpoint{-0.027778in}{0.000000in}}%
\pgfusepath{stroke,fill}%
}%
\begin{pgfscope}%
\pgfsys@transformshift{10.080882in}{3.145488in}%
\pgfsys@useobject{currentmarker}{}%
\end{pgfscope}%
\end{pgfscope}%
\begin{pgfscope}%
\pgfsetbuttcap%
\pgfsetroundjoin%
\definecolor{currentfill}{rgb}{0.000000,0.000000,0.000000}%
\pgfsetfillcolor{currentfill}%
\pgfsetlinewidth{0.602250pt}%
\definecolor{currentstroke}{rgb}{0.000000,0.000000,0.000000}%
\pgfsetstrokecolor{currentstroke}%
\pgfsetdash{}{0pt}%
\pgfsys@defobject{currentmarker}{\pgfqpoint{-0.027778in}{0.000000in}}{\pgfqpoint{-0.000000in}{0.000000in}}{%
\pgfpathmoveto{\pgfqpoint{-0.000000in}{0.000000in}}%
\pgfpathlineto{\pgfqpoint{-0.027778in}{0.000000in}}%
\pgfusepath{stroke,fill}%
}%
\begin{pgfscope}%
\pgfsys@transformshift{10.080882in}{3.214707in}%
\pgfsys@useobject{currentmarker}{}%
\end{pgfscope}%
\end{pgfscope}%
\begin{pgfscope}%
\pgfsetbuttcap%
\pgfsetroundjoin%
\definecolor{currentfill}{rgb}{0.000000,0.000000,0.000000}%
\pgfsetfillcolor{currentfill}%
\pgfsetlinewidth{0.602250pt}%
\definecolor{currentstroke}{rgb}{0.000000,0.000000,0.000000}%
\pgfsetstrokecolor{currentstroke}%
\pgfsetdash{}{0pt}%
\pgfsys@defobject{currentmarker}{\pgfqpoint{-0.027778in}{0.000000in}}{\pgfqpoint{-0.000000in}{0.000000in}}{%
\pgfpathmoveto{\pgfqpoint{-0.000000in}{0.000000in}}%
\pgfpathlineto{\pgfqpoint{-0.027778in}{0.000000in}}%
\pgfusepath{stroke,fill}%
}%
\begin{pgfscope}%
\pgfsys@transformshift{10.080882in}{3.273230in}%
\pgfsys@useobject{currentmarker}{}%
\end{pgfscope}%
\end{pgfscope}%
\begin{pgfscope}%
\pgfsetbuttcap%
\pgfsetroundjoin%
\definecolor{currentfill}{rgb}{0.000000,0.000000,0.000000}%
\pgfsetfillcolor{currentfill}%
\pgfsetlinewidth{0.602250pt}%
\definecolor{currentstroke}{rgb}{0.000000,0.000000,0.000000}%
\pgfsetstrokecolor{currentstroke}%
\pgfsetdash{}{0pt}%
\pgfsys@defobject{currentmarker}{\pgfqpoint{-0.027778in}{0.000000in}}{\pgfqpoint{-0.000000in}{0.000000in}}{%
\pgfpathmoveto{\pgfqpoint{-0.000000in}{0.000000in}}%
\pgfpathlineto{\pgfqpoint{-0.027778in}{0.000000in}}%
\pgfusepath{stroke,fill}%
}%
\begin{pgfscope}%
\pgfsys@transformshift{10.080882in}{3.323926in}%
\pgfsys@useobject{currentmarker}{}%
\end{pgfscope}%
\end{pgfscope}%
\begin{pgfscope}%
\pgfsetbuttcap%
\pgfsetroundjoin%
\definecolor{currentfill}{rgb}{0.000000,0.000000,0.000000}%
\pgfsetfillcolor{currentfill}%
\pgfsetlinewidth{0.602250pt}%
\definecolor{currentstroke}{rgb}{0.000000,0.000000,0.000000}%
\pgfsetstrokecolor{currentstroke}%
\pgfsetdash{}{0pt}%
\pgfsys@defobject{currentmarker}{\pgfqpoint{-0.027778in}{0.000000in}}{\pgfqpoint{-0.000000in}{0.000000in}}{%
\pgfpathmoveto{\pgfqpoint{-0.000000in}{0.000000in}}%
\pgfpathlineto{\pgfqpoint{-0.027778in}{0.000000in}}%
\pgfusepath{stroke,fill}%
}%
\begin{pgfscope}%
\pgfsys@transformshift{10.080882in}{3.368642in}%
\pgfsys@useobject{currentmarker}{}%
\end{pgfscope}%
\end{pgfscope}%
\begin{pgfscope}%
\pgfsetbuttcap%
\pgfsetroundjoin%
\definecolor{currentfill}{rgb}{0.000000,0.000000,0.000000}%
\pgfsetfillcolor{currentfill}%
\pgfsetlinewidth{0.602250pt}%
\definecolor{currentstroke}{rgb}{0.000000,0.000000,0.000000}%
\pgfsetstrokecolor{currentstroke}%
\pgfsetdash{}{0pt}%
\pgfsys@defobject{currentmarker}{\pgfqpoint{-0.027778in}{0.000000in}}{\pgfqpoint{-0.000000in}{0.000000in}}{%
\pgfpathmoveto{\pgfqpoint{-0.000000in}{0.000000in}}%
\pgfpathlineto{\pgfqpoint{-0.027778in}{0.000000in}}%
\pgfusepath{stroke,fill}%
}%
\begin{pgfscope}%
\pgfsys@transformshift{10.080882in}{3.671797in}%
\pgfsys@useobject{currentmarker}{}%
\end{pgfscope}%
\end{pgfscope}%
\begin{pgfscope}%
\pgfsetbuttcap%
\pgfsetroundjoin%
\definecolor{currentfill}{rgb}{0.000000,0.000000,0.000000}%
\pgfsetfillcolor{currentfill}%
\pgfsetlinewidth{0.602250pt}%
\definecolor{currentstroke}{rgb}{0.000000,0.000000,0.000000}%
\pgfsetstrokecolor{currentstroke}%
\pgfsetdash{}{0pt}%
\pgfsys@defobject{currentmarker}{\pgfqpoint{-0.027778in}{0.000000in}}{\pgfqpoint{-0.000000in}{0.000000in}}{%
\pgfpathmoveto{\pgfqpoint{-0.000000in}{0.000000in}}%
\pgfpathlineto{\pgfqpoint{-0.027778in}{0.000000in}}%
\pgfusepath{stroke,fill}%
}%
\begin{pgfscope}%
\pgfsys@transformshift{10.080882in}{3.825732in}%
\pgfsys@useobject{currentmarker}{}%
\end{pgfscope}%
\end{pgfscope}%
\begin{pgfscope}%
\pgfsetbuttcap%
\pgfsetroundjoin%
\definecolor{currentfill}{rgb}{0.000000,0.000000,0.000000}%
\pgfsetfillcolor{currentfill}%
\pgfsetlinewidth{0.602250pt}%
\definecolor{currentstroke}{rgb}{0.000000,0.000000,0.000000}%
\pgfsetstrokecolor{currentstroke}%
\pgfsetdash{}{0pt}%
\pgfsys@defobject{currentmarker}{\pgfqpoint{-0.027778in}{0.000000in}}{\pgfqpoint{-0.000000in}{0.000000in}}{%
\pgfpathmoveto{\pgfqpoint{-0.000000in}{0.000000in}}%
\pgfpathlineto{\pgfqpoint{-0.027778in}{0.000000in}}%
\pgfusepath{stroke,fill}%
}%
\begin{pgfscope}%
\pgfsys@transformshift{10.080882in}{3.934951in}%
\pgfsys@useobject{currentmarker}{}%
\end{pgfscope}%
\end{pgfscope}%
\begin{pgfscope}%
\pgfsetbuttcap%
\pgfsetroundjoin%
\definecolor{currentfill}{rgb}{0.000000,0.000000,0.000000}%
\pgfsetfillcolor{currentfill}%
\pgfsetlinewidth{0.602250pt}%
\definecolor{currentstroke}{rgb}{0.000000,0.000000,0.000000}%
\pgfsetstrokecolor{currentstroke}%
\pgfsetdash{}{0pt}%
\pgfsys@defobject{currentmarker}{\pgfqpoint{-0.027778in}{0.000000in}}{\pgfqpoint{-0.000000in}{0.000000in}}{%
\pgfpathmoveto{\pgfqpoint{-0.000000in}{0.000000in}}%
\pgfpathlineto{\pgfqpoint{-0.027778in}{0.000000in}}%
\pgfusepath{stroke,fill}%
}%
\begin{pgfscope}%
\pgfsys@transformshift{10.080882in}{4.019667in}%
\pgfsys@useobject{currentmarker}{}%
\end{pgfscope}%
\end{pgfscope}%
\begin{pgfscope}%
\pgfsetbuttcap%
\pgfsetroundjoin%
\definecolor{currentfill}{rgb}{0.000000,0.000000,0.000000}%
\pgfsetfillcolor{currentfill}%
\pgfsetlinewidth{0.602250pt}%
\definecolor{currentstroke}{rgb}{0.000000,0.000000,0.000000}%
\pgfsetstrokecolor{currentstroke}%
\pgfsetdash{}{0pt}%
\pgfsys@defobject{currentmarker}{\pgfqpoint{-0.027778in}{0.000000in}}{\pgfqpoint{-0.000000in}{0.000000in}}{%
\pgfpathmoveto{\pgfqpoint{-0.000000in}{0.000000in}}%
\pgfpathlineto{\pgfqpoint{-0.027778in}{0.000000in}}%
\pgfusepath{stroke,fill}%
}%
\begin{pgfscope}%
\pgfsys@transformshift{10.080882in}{4.088886in}%
\pgfsys@useobject{currentmarker}{}%
\end{pgfscope}%
\end{pgfscope}%
\begin{pgfscope}%
\pgfsetbuttcap%
\pgfsetroundjoin%
\definecolor{currentfill}{rgb}{0.000000,0.000000,0.000000}%
\pgfsetfillcolor{currentfill}%
\pgfsetlinewidth{0.602250pt}%
\definecolor{currentstroke}{rgb}{0.000000,0.000000,0.000000}%
\pgfsetstrokecolor{currentstroke}%
\pgfsetdash{}{0pt}%
\pgfsys@defobject{currentmarker}{\pgfqpoint{-0.027778in}{0.000000in}}{\pgfqpoint{-0.000000in}{0.000000in}}{%
\pgfpathmoveto{\pgfqpoint{-0.000000in}{0.000000in}}%
\pgfpathlineto{\pgfqpoint{-0.027778in}{0.000000in}}%
\pgfusepath{stroke,fill}%
}%
\begin{pgfscope}%
\pgfsys@transformshift{10.080882in}{4.147409in}%
\pgfsys@useobject{currentmarker}{}%
\end{pgfscope}%
\end{pgfscope}%
\begin{pgfscope}%
\pgfsetbuttcap%
\pgfsetroundjoin%
\definecolor{currentfill}{rgb}{0.000000,0.000000,0.000000}%
\pgfsetfillcolor{currentfill}%
\pgfsetlinewidth{0.602250pt}%
\definecolor{currentstroke}{rgb}{0.000000,0.000000,0.000000}%
\pgfsetstrokecolor{currentstroke}%
\pgfsetdash{}{0pt}%
\pgfsys@defobject{currentmarker}{\pgfqpoint{-0.027778in}{0.000000in}}{\pgfqpoint{-0.000000in}{0.000000in}}{%
\pgfpathmoveto{\pgfqpoint{-0.000000in}{0.000000in}}%
\pgfpathlineto{\pgfqpoint{-0.027778in}{0.000000in}}%
\pgfusepath{stroke,fill}%
}%
\begin{pgfscope}%
\pgfsys@transformshift{10.080882in}{4.198105in}%
\pgfsys@useobject{currentmarker}{}%
\end{pgfscope}%
\end{pgfscope}%
\begin{pgfscope}%
\pgfsetbuttcap%
\pgfsetroundjoin%
\definecolor{currentfill}{rgb}{0.000000,0.000000,0.000000}%
\pgfsetfillcolor{currentfill}%
\pgfsetlinewidth{0.602250pt}%
\definecolor{currentstroke}{rgb}{0.000000,0.000000,0.000000}%
\pgfsetstrokecolor{currentstroke}%
\pgfsetdash{}{0pt}%
\pgfsys@defobject{currentmarker}{\pgfqpoint{-0.027778in}{0.000000in}}{\pgfqpoint{-0.000000in}{0.000000in}}{%
\pgfpathmoveto{\pgfqpoint{-0.000000in}{0.000000in}}%
\pgfpathlineto{\pgfqpoint{-0.027778in}{0.000000in}}%
\pgfusepath{stroke,fill}%
}%
\begin{pgfscope}%
\pgfsys@transformshift{10.080882in}{4.242821in}%
\pgfsys@useobject{currentmarker}{}%
\end{pgfscope}%
\end{pgfscope}%
\begin{pgfscope}%
\pgfsetbuttcap%
\pgfsetroundjoin%
\definecolor{currentfill}{rgb}{0.000000,0.000000,0.000000}%
\pgfsetfillcolor{currentfill}%
\pgfsetlinewidth{0.602250pt}%
\definecolor{currentstroke}{rgb}{0.000000,0.000000,0.000000}%
\pgfsetstrokecolor{currentstroke}%
\pgfsetdash{}{0pt}%
\pgfsys@defobject{currentmarker}{\pgfqpoint{-0.027778in}{0.000000in}}{\pgfqpoint{-0.000000in}{0.000000in}}{%
\pgfpathmoveto{\pgfqpoint{-0.000000in}{0.000000in}}%
\pgfpathlineto{\pgfqpoint{-0.027778in}{0.000000in}}%
\pgfusepath{stroke,fill}%
}%
\begin{pgfscope}%
\pgfsys@transformshift{10.080882in}{4.545976in}%
\pgfsys@useobject{currentmarker}{}%
\end{pgfscope}%
\end{pgfscope}%
\begin{pgfscope}%
\pgfpathrectangle{\pgfqpoint{10.080882in}{1.100000in}}{\pgfqpoint{3.419118in}{3.500000in}}%
\pgfusepath{clip}%
\pgfsetrectcap%
\pgfsetroundjoin%
\pgfsetlinewidth{1.505625pt}%
\definecolor{currentstroke}{rgb}{0.121569,0.466667,0.705882}%
\pgfsetstrokecolor{currentstroke}%
\pgfsetdash{}{0pt}%
\pgfpathmoveto{\pgfqpoint{10.236297in}{4.161702in}}%
\pgfpathlineto{\pgfqpoint{10.267694in}{4.099679in}}%
\pgfpathlineto{\pgfqpoint{10.299091in}{4.040242in}}%
\pgfpathlineto{\pgfqpoint{10.330487in}{3.982696in}}%
\pgfpathlineto{\pgfqpoint{10.361884in}{3.926661in}}%
\pgfpathlineto{\pgfqpoint{10.393281in}{3.871835in}}%
\pgfpathlineto{\pgfqpoint{10.424678in}{3.849543in}}%
\pgfpathlineto{\pgfqpoint{10.456075in}{3.827373in}}%
\pgfpathlineto{\pgfqpoint{10.487472in}{3.805342in}}%
\pgfpathlineto{\pgfqpoint{10.518868in}{3.783439in}}%
\pgfpathlineto{\pgfqpoint{10.550265in}{3.761650in}}%
\pgfpathlineto{\pgfqpoint{10.581662in}{3.739962in}}%
\pgfpathlineto{\pgfqpoint{10.613059in}{3.718366in}}%
\pgfpathlineto{\pgfqpoint{10.644456in}{3.696852in}}%
\pgfpathlineto{\pgfqpoint{10.675853in}{3.675412in}}%
\pgfpathlineto{\pgfqpoint{10.707250in}{3.654039in}}%
\pgfpathlineto{\pgfqpoint{10.738646in}{3.632725in}}%
\pgfpathlineto{\pgfqpoint{10.770043in}{3.611464in}}%
\pgfpathlineto{\pgfqpoint{10.801440in}{3.590252in}}%
\pgfpathlineto{\pgfqpoint{10.832837in}{3.569083in}}%
\pgfpathlineto{\pgfqpoint{10.864234in}{3.547953in}}%
\pgfpathlineto{\pgfqpoint{10.895631in}{3.526858in}}%
\pgfpathlineto{\pgfqpoint{10.927028in}{3.505794in}}%
\pgfpathlineto{\pgfqpoint{10.958424in}{3.484759in}}%
\pgfpathlineto{\pgfqpoint{10.989821in}{3.463748in}}%
\pgfpathlineto{\pgfqpoint{11.021218in}{3.442760in}}%
\pgfpathlineto{\pgfqpoint{11.052615in}{3.421793in}}%
\pgfpathlineto{\pgfqpoint{11.084012in}{3.400843in}}%
\pgfpathlineto{\pgfqpoint{11.115409in}{3.379910in}}%
\pgfpathlineto{\pgfqpoint{11.146806in}{3.358992in}}%
\pgfpathlineto{\pgfqpoint{11.178202in}{3.338087in}}%
\pgfpathlineto{\pgfqpoint{11.209599in}{3.317193in}}%
\pgfpathlineto{\pgfqpoint{11.240996in}{3.296311in}}%
\pgfpathlineto{\pgfqpoint{11.272393in}{3.275438in}}%
\pgfpathlineto{\pgfqpoint{11.303790in}{3.254573in}}%
\pgfpathlineto{\pgfqpoint{11.335187in}{3.233716in}}%
\pgfpathlineto{\pgfqpoint{11.366584in}{3.212867in}}%
\pgfpathlineto{\pgfqpoint{11.397980in}{3.192023in}}%
\pgfpathlineto{\pgfqpoint{11.429377in}{3.171185in}}%
\pgfpathlineto{\pgfqpoint{11.460774in}{3.150352in}}%
\pgfpathlineto{\pgfqpoint{11.492171in}{3.129524in}}%
\pgfpathlineto{\pgfqpoint{11.523568in}{3.108700in}}%
\pgfpathlineto{\pgfqpoint{11.554965in}{3.087880in}}%
\pgfpathlineto{\pgfqpoint{11.586362in}{3.067063in}}%
\pgfpathlineto{\pgfqpoint{11.617758in}{3.046249in}}%
\pgfpathlineto{\pgfqpoint{11.649155in}{3.025438in}}%
\pgfpathlineto{\pgfqpoint{11.680552in}{3.004629in}}%
\pgfpathlineto{\pgfqpoint{11.711949in}{2.983823in}}%
\pgfpathlineto{\pgfqpoint{11.743346in}{2.963019in}}%
\pgfpathlineto{\pgfqpoint{11.774743in}{2.942217in}}%
\pgfpathlineto{\pgfqpoint{11.806140in}{2.921417in}}%
\pgfpathlineto{\pgfqpoint{11.837536in}{2.900618in}}%
\pgfpathlineto{\pgfqpoint{11.868933in}{2.879820in}}%
\pgfpathlineto{\pgfqpoint{11.900330in}{2.859024in}}%
\pgfpathlineto{\pgfqpoint{11.931727in}{2.838229in}}%
\pgfpathlineto{\pgfqpoint{11.963124in}{2.817436in}}%
\pgfpathlineto{\pgfqpoint{11.994521in}{2.796643in}}%
\pgfpathlineto{\pgfqpoint{12.025918in}{2.775851in}}%
\pgfpathlineto{\pgfqpoint{12.057314in}{2.755061in}}%
\pgfpathlineto{\pgfqpoint{12.088711in}{2.734271in}}%
\pgfpathlineto{\pgfqpoint{12.120108in}{2.713481in}}%
\pgfpathlineto{\pgfqpoint{12.151505in}{2.692692in}}%
\pgfpathlineto{\pgfqpoint{12.182902in}{2.671904in}}%
\pgfpathlineto{\pgfqpoint{12.214299in}{2.651117in}}%
\pgfpathlineto{\pgfqpoint{12.245696in}{2.630329in}}%
\pgfpathlineto{\pgfqpoint{12.277092in}{2.609543in}}%
\pgfpathlineto{\pgfqpoint{12.308489in}{2.588757in}}%
\pgfpathlineto{\pgfqpoint{12.339886in}{2.567971in}}%
\pgfpathlineto{\pgfqpoint{12.371283in}{2.547185in}}%
\pgfpathlineto{\pgfqpoint{12.402680in}{2.526400in}}%
\pgfpathlineto{\pgfqpoint{12.434077in}{2.505615in}}%
\pgfpathlineto{\pgfqpoint{12.465474in}{2.484830in}}%
\pgfpathlineto{\pgfqpoint{12.496870in}{2.464046in}}%
\pgfpathlineto{\pgfqpoint{12.528267in}{2.443261in}}%
\pgfpathlineto{\pgfqpoint{12.559664in}{2.422477in}}%
\pgfpathlineto{\pgfqpoint{12.591061in}{2.401694in}}%
\pgfpathlineto{\pgfqpoint{12.622458in}{2.380909in}}%
\pgfpathlineto{\pgfqpoint{12.653855in}{2.360124in}}%
\pgfpathlineto{\pgfqpoint{12.685252in}{2.339339in}}%
\pgfpathlineto{\pgfqpoint{12.716648in}{2.318555in}}%
\pgfpathlineto{\pgfqpoint{12.748045in}{2.297770in}}%
\pgfpathlineto{\pgfqpoint{12.779442in}{2.276986in}}%
\pgfpathlineto{\pgfqpoint{12.810839in}{2.256200in}}%
\pgfpathlineto{\pgfqpoint{12.842236in}{2.235414in}}%
\pgfpathlineto{\pgfqpoint{12.873633in}{2.214628in}}%
\pgfpathlineto{\pgfqpoint{12.905030in}{2.193842in}}%
\pgfpathlineto{\pgfqpoint{12.936426in}{2.173055in}}%
\pgfpathlineto{\pgfqpoint{12.967823in}{2.152266in}}%
\pgfpathlineto{\pgfqpoint{12.999220in}{2.131478in}}%
\pgfpathlineto{\pgfqpoint{13.030617in}{2.110688in}}%
\pgfpathlineto{\pgfqpoint{13.062014in}{2.089897in}}%
\pgfpathlineto{\pgfqpoint{13.093411in}{2.069106in}}%
\pgfpathlineto{\pgfqpoint{13.124808in}{2.048315in}}%
\pgfpathlineto{\pgfqpoint{13.156204in}{2.027522in}}%
\pgfpathlineto{\pgfqpoint{13.187601in}{2.006728in}}%
\pgfpathlineto{\pgfqpoint{13.218998in}{1.985931in}}%
\pgfpathlineto{\pgfqpoint{13.250395in}{1.965134in}}%
\pgfpathlineto{\pgfqpoint{13.281792in}{1.944335in}}%
\pgfpathlineto{\pgfqpoint{13.313189in}{1.923532in}}%
\pgfpathlineto{\pgfqpoint{13.344586in}{1.902730in}}%
\pgfusepath{stroke}%
\end{pgfscope}%
\begin{pgfscope}%
\pgfpathrectangle{\pgfqpoint{10.080882in}{1.100000in}}{\pgfqpoint{3.419118in}{3.500000in}}%
\pgfusepath{clip}%
\pgfsetrectcap%
\pgfsetroundjoin%
\pgfsetlinewidth{1.505625pt}%
\definecolor{currentstroke}{rgb}{1.000000,0.498039,0.054902}%
\pgfsetstrokecolor{currentstroke}%
\pgfsetdash{}{0pt}%
\pgfpathmoveto{\pgfqpoint{10.236297in}{4.270754in}}%
\pgfpathlineto{\pgfqpoint{10.267694in}{4.218026in}}%
\pgfpathlineto{\pgfqpoint{10.299091in}{4.166120in}}%
\pgfpathlineto{\pgfqpoint{10.330487in}{4.114704in}}%
\pgfpathlineto{\pgfqpoint{10.361884in}{4.063640in}}%
\pgfpathlineto{\pgfqpoint{10.393281in}{4.012837in}}%
\pgfpathlineto{\pgfqpoint{10.424678in}{3.962234in}}%
\pgfpathlineto{\pgfqpoint{10.456075in}{3.941528in}}%
\pgfpathlineto{\pgfqpoint{10.487472in}{3.920833in}}%
\pgfpathlineto{\pgfqpoint{10.518868in}{3.900162in}}%
\pgfpathlineto{\pgfqpoint{10.550265in}{3.879513in}}%
\pgfpathlineto{\pgfqpoint{10.581662in}{3.858883in}}%
\pgfpathlineto{\pgfqpoint{10.613059in}{3.838270in}}%
\pgfpathlineto{\pgfqpoint{10.644456in}{3.817671in}}%
\pgfpathlineto{\pgfqpoint{10.675853in}{3.630624in}}%
\pgfpathlineto{\pgfqpoint{10.707250in}{3.531435in}}%
\pgfpathlineto{\pgfqpoint{10.738646in}{3.505741in}}%
\pgfpathlineto{\pgfqpoint{10.770043in}{3.481596in}}%
\pgfpathlineto{\pgfqpoint{10.801440in}{3.458459in}}%
\pgfpathlineto{\pgfqpoint{10.832837in}{3.436018in}}%
\pgfpathlineto{\pgfqpoint{10.864234in}{3.414072in}}%
\pgfpathlineto{\pgfqpoint{10.895631in}{3.392485in}}%
\pgfpathlineto{\pgfqpoint{10.927028in}{3.371161in}}%
\pgfpathlineto{\pgfqpoint{10.958424in}{3.350030in}}%
\pgfpathlineto{\pgfqpoint{10.989821in}{3.329042in}}%
\pgfpathlineto{\pgfqpoint{11.021218in}{3.308160in}}%
\pgfpathlineto{\pgfqpoint{11.052615in}{3.287357in}}%
\pgfpathlineto{\pgfqpoint{11.084012in}{3.266614in}}%
\pgfpathlineto{\pgfqpoint{11.115409in}{3.245916in}}%
\pgfpathlineto{\pgfqpoint{11.146806in}{3.616861in}}%
\pgfpathlineto{\pgfqpoint{11.178202in}{3.518872in}}%
\pgfpathlineto{\pgfqpoint{11.209599in}{3.493975in}}%
\pgfpathlineto{\pgfqpoint{11.240996in}{3.470459in}}%
\pgfpathlineto{\pgfqpoint{11.272393in}{3.447806in}}%
\pgfpathlineto{\pgfqpoint{11.303790in}{3.425733in}}%
\pgfpathlineto{\pgfqpoint{11.335187in}{3.404068in}}%
\pgfpathlineto{\pgfqpoint{11.366584in}{3.382695in}}%
\pgfpathlineto{\pgfqpoint{11.397980in}{3.361534in}}%
\pgfpathlineto{\pgfqpoint{11.429377in}{3.340529in}}%
\pgfpathlineto{\pgfqpoint{11.460774in}{3.319638in}}%
\pgfpathlineto{\pgfqpoint{11.492171in}{3.298833in}}%
\pgfpathlineto{\pgfqpoint{11.523568in}{3.278091in}}%
\pgfpathlineto{\pgfqpoint{11.554965in}{3.257397in}}%
\pgfpathlineto{\pgfqpoint{11.586362in}{3.236739in}}%
\pgfpathlineto{\pgfqpoint{11.617758in}{3.616756in}}%
\pgfpathlineto{\pgfqpoint{11.649155in}{3.518650in}}%
\pgfpathlineto{\pgfqpoint{11.680552in}{3.493778in}}%
\pgfpathlineto{\pgfqpoint{11.711949in}{3.470285in}}%
\pgfpathlineto{\pgfqpoint{11.743346in}{3.447650in}}%
\pgfpathlineto{\pgfqpoint{11.774743in}{3.425592in}}%
\pgfpathlineto{\pgfqpoint{11.806140in}{3.403937in}}%
\pgfpathlineto{\pgfqpoint{11.837536in}{3.382572in}}%
\pgfpathlineto{\pgfqpoint{11.868933in}{3.361417in}}%
\pgfpathlineto{\pgfqpoint{11.900330in}{3.340416in}}%
\pgfpathlineto{\pgfqpoint{11.931727in}{3.319529in}}%
\pgfpathlineto{\pgfqpoint{11.963124in}{3.298726in}}%
\pgfpathlineto{\pgfqpoint{11.994521in}{3.277986in}}%
\pgfpathlineto{\pgfqpoint{12.025918in}{3.257293in}}%
\pgfpathlineto{\pgfqpoint{12.057314in}{3.236636in}}%
\pgfpathlineto{\pgfqpoint{12.088711in}{3.616753in}}%
\pgfpathlineto{\pgfqpoint{12.120108in}{3.518641in}}%
\pgfpathlineto{\pgfqpoint{12.151505in}{3.493770in}}%
\pgfpathlineto{\pgfqpoint{12.182902in}{3.470278in}}%
\pgfpathlineto{\pgfqpoint{12.214299in}{3.447644in}}%
\pgfpathlineto{\pgfqpoint{12.245696in}{3.425586in}}%
\pgfpathlineto{\pgfqpoint{12.277092in}{3.403932in}}%
\pgfpathlineto{\pgfqpoint{12.308489in}{3.382567in}}%
\pgfpathlineto{\pgfqpoint{12.339886in}{3.361412in}}%
\pgfpathlineto{\pgfqpoint{12.371283in}{3.340411in}}%
\pgfpathlineto{\pgfqpoint{12.402680in}{3.319524in}}%
\pgfpathlineto{\pgfqpoint{12.434077in}{3.298721in}}%
\pgfpathlineto{\pgfqpoint{12.465474in}{3.277981in}}%
\pgfpathlineto{\pgfqpoint{12.496870in}{3.257288in}}%
\pgfpathlineto{\pgfqpoint{12.528267in}{3.236631in}}%
\pgfpathlineto{\pgfqpoint{12.559664in}{3.616761in}}%
\pgfpathlineto{\pgfqpoint{12.591061in}{3.518648in}}%
\pgfpathlineto{\pgfqpoint{12.622458in}{3.493777in}}%
\pgfpathlineto{\pgfqpoint{12.653855in}{3.470285in}}%
\pgfpathlineto{\pgfqpoint{12.685252in}{3.447652in}}%
\pgfpathlineto{\pgfqpoint{12.716648in}{3.425594in}}%
\pgfpathlineto{\pgfqpoint{12.748045in}{3.403939in}}%
\pgfpathlineto{\pgfqpoint{12.779442in}{3.382574in}}%
\pgfpathlineto{\pgfqpoint{12.810839in}{3.361420in}}%
\pgfpathlineto{\pgfqpoint{12.842236in}{3.340419in}}%
\pgfpathlineto{\pgfqpoint{12.873633in}{3.319532in}}%
\pgfpathlineto{\pgfqpoint{12.905030in}{3.298729in}}%
\pgfpathlineto{\pgfqpoint{12.936426in}{3.277989in}}%
\pgfpathlineto{\pgfqpoint{12.967823in}{3.257296in}}%
\pgfpathlineto{\pgfqpoint{12.999220in}{3.236638in}}%
\pgfpathlineto{\pgfqpoint{13.030617in}{3.616758in}}%
\pgfpathlineto{\pgfqpoint{13.062014in}{3.518646in}}%
\pgfpathlineto{\pgfqpoint{13.093411in}{3.493774in}}%
\pgfpathlineto{\pgfqpoint{13.124808in}{3.470282in}}%
\pgfpathlineto{\pgfqpoint{13.156204in}{3.447649in}}%
\pgfpathlineto{\pgfqpoint{13.187601in}{3.425591in}}%
\pgfpathlineto{\pgfqpoint{13.218998in}{3.403937in}}%
\pgfpathlineto{\pgfqpoint{13.250395in}{3.382572in}}%
\pgfpathlineto{\pgfqpoint{13.281792in}{3.361417in}}%
\pgfpathlineto{\pgfqpoint{13.313189in}{3.340416in}}%
\pgfpathlineto{\pgfqpoint{13.344586in}{3.319529in}}%
\pgfusepath{stroke}%
\end{pgfscope}%
\begin{pgfscope}%
\pgfpathrectangle{\pgfqpoint{10.080882in}{1.100000in}}{\pgfqpoint{3.419118in}{3.500000in}}%
\pgfusepath{clip}%
\pgfsetrectcap%
\pgfsetroundjoin%
\pgfsetlinewidth{1.505625pt}%
\definecolor{currentstroke}{rgb}{0.172549,0.627451,0.172549}%
\pgfsetstrokecolor{currentstroke}%
\pgfsetdash{}{0pt}%
\pgfpathmoveto{\pgfqpoint{10.236297in}{4.026395in}}%
\pgfpathlineto{\pgfqpoint{10.267694in}{3.740113in}}%
\pgfpathlineto{\pgfqpoint{10.299091in}{3.609609in}}%
\pgfpathlineto{\pgfqpoint{10.330487in}{3.503235in}}%
\pgfpathlineto{\pgfqpoint{10.361884in}{3.462691in}}%
\pgfpathlineto{\pgfqpoint{10.393281in}{3.423693in}}%
\pgfpathlineto{\pgfqpoint{10.424678in}{3.386008in}}%
\pgfpathlineto{\pgfqpoint{10.456075in}{3.349450in}}%
\pgfpathlineto{\pgfqpoint{10.487472in}{3.313870in}}%
\pgfpathlineto{\pgfqpoint{10.518868in}{3.279175in}}%
\pgfpathlineto{\pgfqpoint{10.550265in}{3.245293in}}%
\pgfpathlineto{\pgfqpoint{10.581662in}{3.212177in}}%
\pgfpathlineto{\pgfqpoint{10.613059in}{3.179792in}}%
\pgfpathlineto{\pgfqpoint{10.644456in}{3.148110in}}%
\pgfpathlineto{\pgfqpoint{10.675853in}{3.117108in}}%
\pgfpathlineto{\pgfqpoint{10.707250in}{3.086767in}}%
\pgfpathlineto{\pgfqpoint{10.738646in}{3.057066in}}%
\pgfpathlineto{\pgfqpoint{10.770043in}{3.027986in}}%
\pgfpathlineto{\pgfqpoint{10.801440in}{2.999506in}}%
\pgfpathlineto{\pgfqpoint{10.832837in}{2.971603in}}%
\pgfpathlineto{\pgfqpoint{10.864234in}{2.944254in}}%
\pgfpathlineto{\pgfqpoint{10.895631in}{2.917435in}}%
\pgfpathlineto{\pgfqpoint{10.927028in}{2.891118in}}%
\pgfpathlineto{\pgfqpoint{10.958424in}{2.865278in}}%
\pgfpathlineto{\pgfqpoint{10.989821in}{2.839887in}}%
\pgfpathlineto{\pgfqpoint{11.021218in}{2.814916in}}%
\pgfpathlineto{\pgfqpoint{11.052615in}{2.790339in}}%
\pgfpathlineto{\pgfqpoint{11.084012in}{2.766129in}}%
\pgfpathlineto{\pgfqpoint{11.115409in}{2.742258in}}%
\pgfpathlineto{\pgfqpoint{11.146806in}{2.718702in}}%
\pgfpathlineto{\pgfqpoint{11.178202in}{2.695435in}}%
\pgfpathlineto{\pgfqpoint{11.209599in}{2.672434in}}%
\pgfpathlineto{\pgfqpoint{11.240996in}{2.649678in}}%
\pgfpathlineto{\pgfqpoint{11.272393in}{2.627144in}}%
\pgfpathlineto{\pgfqpoint{11.303790in}{2.604814in}}%
\pgfpathlineto{\pgfqpoint{11.335187in}{2.582668in}}%
\pgfpathlineto{\pgfqpoint{11.366584in}{2.560690in}}%
\pgfpathlineto{\pgfqpoint{11.397980in}{2.538864in}}%
\pgfpathlineto{\pgfqpoint{11.429377in}{2.517176in}}%
\pgfpathlineto{\pgfqpoint{11.460774in}{2.495613in}}%
\pgfpathlineto{\pgfqpoint{11.492171in}{2.474163in}}%
\pgfpathlineto{\pgfqpoint{11.523568in}{2.452814in}}%
\pgfpathlineto{\pgfqpoint{11.554965in}{2.431557in}}%
\pgfpathlineto{\pgfqpoint{11.586362in}{2.410382in}}%
\pgfpathlineto{\pgfqpoint{11.617758in}{2.389281in}}%
\pgfpathlineto{\pgfqpoint{11.649155in}{2.368245in}}%
\pgfpathlineto{\pgfqpoint{11.680552in}{2.347270in}}%
\pgfpathlineto{\pgfqpoint{11.711949in}{2.326349in}}%
\pgfpathlineto{\pgfqpoint{11.743346in}{2.305476in}}%
\pgfpathlineto{\pgfqpoint{11.774743in}{2.284646in}}%
\pgfpathlineto{\pgfqpoint{11.806140in}{2.263854in}}%
\pgfpathlineto{\pgfqpoint{11.837536in}{2.243097in}}%
\pgfpathlineto{\pgfqpoint{11.868933in}{2.222371in}}%
\pgfpathlineto{\pgfqpoint{11.900330in}{2.201672in}}%
\pgfpathlineto{\pgfqpoint{11.931727in}{2.180998in}}%
\pgfpathlineto{\pgfqpoint{11.963124in}{2.160347in}}%
\pgfpathlineto{\pgfqpoint{11.994521in}{2.139717in}}%
\pgfpathlineto{\pgfqpoint{12.025918in}{2.119103in}}%
\pgfpathlineto{\pgfqpoint{12.057314in}{2.098506in}}%
\pgfpathlineto{\pgfqpoint{12.088711in}{2.077923in}}%
\pgfpathlineto{\pgfqpoint{12.120108in}{2.057353in}}%
\pgfpathlineto{\pgfqpoint{12.151505in}{2.036796in}}%
\pgfpathlineto{\pgfqpoint{12.182902in}{2.016248in}}%
\pgfpathlineto{\pgfqpoint{12.214299in}{1.995713in}}%
\pgfpathlineto{\pgfqpoint{12.245696in}{1.975181in}}%
\pgfpathlineto{\pgfqpoint{12.277092in}{1.954658in}}%
\pgfpathlineto{\pgfqpoint{12.308489in}{1.934141in}}%
\pgfpathlineto{\pgfqpoint{12.339886in}{1.913631in}}%
\pgfpathlineto{\pgfqpoint{12.371283in}{1.893131in}}%
\pgfpathlineto{\pgfqpoint{12.402680in}{1.872632in}}%
\pgfpathlineto{\pgfqpoint{12.434077in}{1.852144in}}%
\pgfpathlineto{\pgfqpoint{12.465474in}{1.831656in}}%
\pgfpathlineto{\pgfqpoint{12.496870in}{1.811165in}}%
\pgfpathlineto{\pgfqpoint{12.528267in}{1.790684in}}%
\pgfpathlineto{\pgfqpoint{12.559664in}{1.770204in}}%
\pgfpathlineto{\pgfqpoint{12.591061in}{1.749725in}}%
\pgfpathlineto{\pgfqpoint{12.622458in}{1.729255in}}%
\pgfpathlineto{\pgfqpoint{12.653855in}{1.708788in}}%
\pgfpathlineto{\pgfqpoint{12.685252in}{1.688318in}}%
\pgfpathlineto{\pgfqpoint{12.716648in}{1.667854in}}%
\pgfpathlineto{\pgfqpoint{12.748045in}{1.647395in}}%
\pgfpathlineto{\pgfqpoint{12.779442in}{1.626935in}}%
\pgfpathlineto{\pgfqpoint{12.810839in}{1.606479in}}%
\pgfpathlineto{\pgfqpoint{12.842236in}{1.586021in}}%
\pgfpathlineto{\pgfqpoint{12.873633in}{1.565575in}}%
\pgfpathlineto{\pgfqpoint{12.905030in}{1.545121in}}%
\pgfpathlineto{\pgfqpoint{12.936426in}{1.524670in}}%
\pgfpathlineto{\pgfqpoint{12.967823in}{1.504216in}}%
\pgfpathlineto{\pgfqpoint{12.999220in}{1.483768in}}%
\pgfpathlineto{\pgfqpoint{13.030617in}{1.463328in}}%
\pgfpathlineto{\pgfqpoint{13.062014in}{1.442893in}}%
\pgfpathlineto{\pgfqpoint{13.093411in}{1.422451in}}%
\pgfpathlineto{\pgfqpoint{13.124808in}{1.402027in}}%
\pgfpathlineto{\pgfqpoint{13.156204in}{1.381589in}}%
\pgfpathlineto{\pgfqpoint{13.187601in}{1.361172in}}%
\pgfpathlineto{\pgfqpoint{13.218998in}{1.340738in}}%
\pgfpathlineto{\pgfqpoint{13.250395in}{1.320322in}}%
\pgfpathlineto{\pgfqpoint{13.281792in}{1.299892in}}%
\pgfpathlineto{\pgfqpoint{13.313189in}{1.279481in}}%
\pgfpathlineto{\pgfqpoint{13.344586in}{1.259091in}}%
\pgfusepath{stroke}%
\end{pgfscope}%
\begin{pgfscope}%
\pgfsetrectcap%
\pgfsetmiterjoin%
\pgfsetlinewidth{0.803000pt}%
\definecolor{currentstroke}{rgb}{0.000000,0.000000,0.000000}%
\pgfsetstrokecolor{currentstroke}%
\pgfsetdash{}{0pt}%
\pgfpathmoveto{\pgfqpoint{10.080882in}{1.100000in}}%
\pgfpathlineto{\pgfqpoint{10.080882in}{4.600000in}}%
\pgfusepath{stroke}%
\end{pgfscope}%
\begin{pgfscope}%
\pgfsetrectcap%
\pgfsetmiterjoin%
\pgfsetlinewidth{0.803000pt}%
\definecolor{currentstroke}{rgb}{0.000000,0.000000,0.000000}%
\pgfsetstrokecolor{currentstroke}%
\pgfsetdash{}{0pt}%
\pgfpathmoveto{\pgfqpoint{13.500000in}{1.100000in}}%
\pgfpathlineto{\pgfqpoint{13.500000in}{4.600000in}}%
\pgfusepath{stroke}%
\end{pgfscope}%
\begin{pgfscope}%
\pgfsetrectcap%
\pgfsetmiterjoin%
\pgfsetlinewidth{0.803000pt}%
\definecolor{currentstroke}{rgb}{0.000000,0.000000,0.000000}%
\pgfsetstrokecolor{currentstroke}%
\pgfsetdash{}{0pt}%
\pgfpathmoveto{\pgfqpoint{10.080882in}{1.100000in}}%
\pgfpathlineto{\pgfqpoint{13.500000in}{1.100000in}}%
\pgfusepath{stroke}%
\end{pgfscope}%
\begin{pgfscope}%
\pgfsetrectcap%
\pgfsetmiterjoin%
\pgfsetlinewidth{0.803000pt}%
\definecolor{currentstroke}{rgb}{0.000000,0.000000,0.000000}%
\pgfsetstrokecolor{currentstroke}%
\pgfsetdash{}{0pt}%
\pgfpathmoveto{\pgfqpoint{10.080882in}{4.600000in}}%
\pgfpathlineto{\pgfqpoint{13.500000in}{4.600000in}}%
\pgfusepath{stroke}%
\end{pgfscope}%
\begin{pgfscope}%
\definecolor{textcolor}{rgb}{0.000000,0.000000,0.000000}%
\pgfsetstrokecolor{textcolor}%
\pgfsetfillcolor{textcolor}%
\pgftext[x=11.790441in,y=4.683333in,,base]{\color{textcolor}{\sffamily\fontsize{12.000000}{14.400000}\selectfont\catcode`\^=\active\def^{\ifmmode\sp\else\^{}\fi}\catcode`\%=\active\def%{\%}MG}}%
\end{pgfscope}%
\begin{pgfscope}%
\definecolor{textcolor}{rgb}{0.000000,0.000000,0.000000}%
\pgfsetstrokecolor{textcolor}%
\pgfsetfillcolor{textcolor}%
\pgftext[x=1.118970in, y=2.216406in, left, base,rotate=90.000000]{\color{textcolor}{\sffamily\fontsize{18.000000}{21.600000}\selectfont\catcode`\^=\active\def^{\ifmmode\sp\else\^{}\fi}\catcode`\%=\active\def%{\%}2D Poisson}}%
\end{pgfscope}%
\begin{pgfscope}%
\definecolor{textcolor}{rgb}{0.000000,0.000000,0.000000}%
\pgfsetstrokecolor{textcolor}%
\pgfsetfillcolor{textcolor}%
\pgftext[x=1.118970in, y=6.316406in, left, base,rotate=90.000000]{\color{textcolor}{\sffamily\fontsize{18.000000}{21.600000}\selectfont\catcode`\^=\active\def^{\ifmmode\sp\else\^{}\fi}\catcode`\%=\active\def%{\%}1D Poisson}}%
\end{pgfscope}%
\begin{pgfscope}%
\definecolor{textcolor}{rgb}{0.000000,0.000000,0.000000}%
\pgfsetstrokecolor{textcolor}%
\pgfsetfillcolor{textcolor}%
\pgftext[x=7.500000in,y=9.800000in,,top]{\color{textcolor}{\sffamily\fontsize{18.000000}{21.600000}\selectfont\catcode`\^=\active\def^{\ifmmode\sp\else\^{}\fi}\catcode`\%=\active\def%{\%}Error Comparison of Different Routing Strategies}}%
\end{pgfscope}%
\begin{pgfscope}%
\pgfsetbuttcap%
\pgfsetmiterjoin%
\definecolor{currentfill}{rgb}{1.000000,1.000000,1.000000}%
\pgfsetfillcolor{currentfill}%
\pgfsetfillopacity{0.800000}%
\pgfsetlinewidth{1.003750pt}%
\definecolor{currentstroke}{rgb}{0.800000,0.800000,0.800000}%
\pgfsetstrokecolor{currentstroke}%
\pgfsetstrokeopacity{0.800000}%
\pgfsetdash{}{0pt}%
\pgfpathmoveto{\pgfqpoint{14.097388in}{4.152227in}}%
\pgfpathlineto{\pgfqpoint{17.075000in}{4.152227in}}%
\pgfpathquadraticcurveto{\pgfqpoint{17.125000in}{4.152227in}}{\pgfqpoint{17.125000in}{4.202227in}}%
\pgfpathlineto{\pgfqpoint{17.125000in}{5.645000in}}%
\pgfpathquadraticcurveto{\pgfqpoint{17.125000in}{5.695000in}}{\pgfqpoint{17.075000in}{5.695000in}}%
\pgfpathlineto{\pgfqpoint{14.097388in}{5.695000in}}%
\pgfpathquadraticcurveto{\pgfqpoint{14.047388in}{5.695000in}}{\pgfqpoint{14.047388in}{5.645000in}}%
\pgfpathlineto{\pgfqpoint{14.047388in}{4.202227in}}%
\pgfpathquadraticcurveto{\pgfqpoint{14.047388in}{4.152227in}}{\pgfqpoint{14.097388in}{4.152227in}}%
\pgfpathlineto{\pgfqpoint{14.097388in}{4.152227in}}%
\pgfpathclose%
\pgfusepath{stroke,fill}%
\end{pgfscope}%
\begin{pgfscope}%
\pgfsetrectcap%
\pgfsetroundjoin%
\pgfsetlinewidth{1.505625pt}%
\definecolor{currentstroke}{rgb}{0.121569,0.466667,0.705882}%
\pgfsetstrokecolor{currentstroke}%
\pgfsetdash{}{0pt}%
\pgfpathmoveto{\pgfqpoint{14.147388in}{5.492559in}}%
\pgfpathlineto{\pgfqpoint{14.397388in}{5.492559in}}%
\pgfpathlineto{\pgfqpoint{14.647388in}{5.492559in}}%
\pgfusepath{stroke}%
\end{pgfscope}%
\begin{pgfscope}%
\definecolor{textcolor}{rgb}{0.000000,0.000000,0.000000}%
\pgfsetstrokecolor{textcolor}%
\pgfsetfillcolor{textcolor}%
\pgftext[x=14.847388in,y=5.405059in,left,base]{\color{textcolor}{\sffamily\fontsize{18.000000}{21.600000}\selectfont\catcode`\^=\active\def^{\ifmmode\sp\else\^{}\fi}\catcode`\%=\active\def%{\%}Solver Only}}%
\end{pgfscope}%
\begin{pgfscope}%
\pgfsetrectcap%
\pgfsetroundjoin%
\pgfsetlinewidth{1.505625pt}%
\definecolor{currentstroke}{rgb}{1.000000,0.498039,0.054902}%
\pgfsetstrokecolor{currentstroke}%
\pgfsetdash{}{0pt}%
\pgfpathmoveto{\pgfqpoint{14.147388in}{5.125615in}}%
\pgfpathlineto{\pgfqpoint{14.397388in}{5.125615in}}%
\pgfpathlineto{\pgfqpoint{14.647388in}{5.125615in}}%
\pgfusepath{stroke}%
\end{pgfscope}%
\begin{pgfscope}%
\definecolor{textcolor}{rgb}{0.000000,0.000000,0.000000}%
\pgfsetstrokecolor{textcolor}%
\pgfsetfillcolor{textcolor}%
\pgftext[x=14.847388in,y=5.038115in,left,base]{\color{textcolor}{\sffamily\fontsize{18.000000}{21.600000}\selectfont\catcode`\^=\active\def^{\ifmmode\sp\else\^{}\fi}\catcode`\%=\active\def%{\%}HINTS}}%
\end{pgfscope}%
\begin{pgfscope}%
\pgfsetrectcap%
\pgfsetroundjoin%
\pgfsetlinewidth{1.505625pt}%
\definecolor{currentstroke}{rgb}{0.172549,0.627451,0.172549}%
\pgfsetstrokecolor{currentstroke}%
\pgfsetdash{}{0pt}%
\pgfpathmoveto{\pgfqpoint{14.147388in}{4.758672in}}%
\pgfpathlineto{\pgfqpoint{14.397388in}{4.758672in}}%
\pgfpathlineto{\pgfqpoint{14.647388in}{4.758672in}}%
\pgfusepath{stroke}%
\end{pgfscope}%
\begin{pgfscope}%
\definecolor{textcolor}{rgb}{0.000000,0.000000,0.000000}%
\pgfsetstrokecolor{textcolor}%
\pgfsetfillcolor{textcolor}%
\pgftext[x=14.847388in,y=4.671172in,left,base]{\color{textcolor}{\sffamily\fontsize{18.000000}{21.600000}\selectfont\catcode`\^=\active\def^{\ifmmode\sp\else\^{}\fi}\catcode`\%=\active\def%{\%}Greedy}}%
\end{pgfscope}%
\begin{pgfscope}%
\pgfsetrectcap%
\pgfsetroundjoin%
\pgfsetlinewidth{1.505625pt}%
\definecolor{currentstroke}{rgb}{0.839216,0.152941,0.156863}%
\pgfsetstrokecolor{currentstroke}%
\pgfsetdash{}{0pt}%
\pgfpathmoveto{\pgfqpoint{14.147388in}{4.391729in}}%
\pgfpathlineto{\pgfqpoint{14.397388in}{4.391729in}}%
\pgfpathlineto{\pgfqpoint{14.647388in}{4.391729in}}%
\pgfusepath{stroke}%
\end{pgfscope}%
\begin{pgfscope}%
\definecolor{textcolor}{rgb}{0.000000,0.000000,0.000000}%
\pgfsetstrokecolor{textcolor}%
\pgfsetfillcolor{textcolor}%
\pgftext[x=14.847388in,y=4.304229in,left,base]{\color{textcolor}{\sffamily\fontsize{18.000000}{21.600000}\selectfont\catcode`\^=\active\def^{\ifmmode\sp\else\^{}\fi}\catcode`\%=\active\def%{\%}Greedy-GS,Jacobi}}%
\end{pgfscope}%
\end{pgfpicture}%
\makeatother%
\endgroup%

    \includegraphics[width=0.7\linewidth]{images/errorcomparison.png}
    \caption{Convergence histories for representative test instances. Rows: 1D Poisson (top) and 2D Poisson (bottom). Columns: Jacobi, Gauss–Seidel (GS), and two-grid multigrid (MG). Greedy yields near-monotone decay and the lowest errors, whereas HINTS shows characteristic sawtooth transients due to periodic DeepONet corrections.}
    \label{fig:errorcomparison}
\end{figure}


\section{Discussion}

In this paper, we introduce an adaptive method for selecting solvers that efficiently minimize the final error when solving a PDE iteratively. This opens several directions for future work. Our current DeepONet is trained in isolation, without accounting for its downstream role as a correction term predictor. Jointly training the ML solver and the router could yield larger gains. Another avenue is to employ reinforcement learning to learn a cost-aware routing strategy that optimizes the terminal error under compute budgets. Finally, it would be interesting to extend this routing scheme to broader optimization settings in which, at each iteration, the router selects from a suite of iterative solvers.




% \section{Random Notes}

% \begin{itemize}
%     \item thoughts on title
%     \item Difficulty of finite sample guarantees: If we know that $f(S \ \oplus r(S)) - \min_{\omega \in \Omega} f(S \ \oplus \omega) < \epsilon$ (can obtain this lower bound from the consistency bound), this doesn't guarantee that $f(Z^T)$ where $Z^T$ is constructed using $r$ such that $Z^T = r(\emptyset) \ \oplus r(Z_1) \ \oplus r(Z_2) \ \oplus \ \cdots \oplus r(Z_{T - 1})$ is lower bounded by the performance of the true greedy sequence $S^T$,
%     \item Budget constraints: There is some work on how to 
%     \item $f$ is not needed for the router
%     \item can neural operator be a matrix? I thought maybe something something lipschitz function $I_N - C_j \circ \mathcal{L}_h^a$ is L then all the conditions blah blah the norm of the error propagation matrix needs to be less than 1 will translate to $L < 1$. BUT this only makes sense if $\mathcal{N}(0) = 0$. Is it normal to train a neural operator without any biases?
%     \item check if the error norms are not spectral radius. l2 norm
%     \item TODO: include proof of consistency, generalization bounds (maybe)
%     \item TODO: look up this upper bound weirdness for the surrogate
%     \item TODO: Appendix A
%     \item \Cref{sec:approxgreedy} I am not sure if neural operators mathematically act on function spaces. I have the routers recieving discretized inputs mathmatically so not sure if it's consistent with neural operator literature. f but we do need to connect \Cref{sec:greedy} somehow. feels disruptive to talk aboutdiscrete things and then moving to functional space.
%     \item what is $\mu$. do we need $\mu$. I feel invertibility is not that deep
%     \item should we have a background section on neural operators when the discretization invariance property is not leveraged by the greedy router. 
%     \item title idea
% \end{itemize}


% $$h(S) = \left\|\circ_{t = |S|}^1\left(I_N - C_{S_t} \circ \mathcal{L}_h^a\right)  e^{(0)}_h \right\|^2_2$$


% $g(0) = 0$ and $f(0) = 0$

% $ \|(f\circ g) (x)\| = L_fL_g\|x\|$

% $\|ABx\| = \|ABA^{-1}Ax\| \leq \|A\|\|B\|\|A^{-1}\|\|Ax\|$

% $\|(f\circ g)(x)\| = \|(f\circ g \circ f^{-1} \circ f)(x)\| \leq L_fL_g\frac{1}{L_f}\|f(x)\|$




\bibliography{iclr2026_conference}
\bibliographystyle{iclr2026_conference}

\appendix
\section{Background} \label{sec:appbackground}

\subsection{Multigrid} \label{sec:multigrid}

Let $A_h, A_{2h}$ represent the coefficient matrix on a fine grid with discretization parameter $h$ and $2h$, $u_h$ and $f_h$ represent the PDE solution and constant vector on a grid discretized by $h$, $R^{2h}_h$ denote the restriction matrix which transfers vectors from a fine grid to a coarse one, and $I^{h}_{2h}$ is the interpolation matrix transfers vectors from a coarse grid to a fine one. In a 2-grid method, a few iterations of the smoother (e.x. Jacobi or Gauss-Seidel) are first applied on the fine grid to approximate the solution of $A_hu_h = f_h$. The residual is then computed as $r_h = f_h - A_hu_h$ and restricted to the coarse grid via $r_{2h} = R^{2h}_hr_h$. The error equation, $A_{2h}e_{2h} = r_{2h}$, is solved on the coarse grid. The resulting estimate of the error is interpolated in the fine grid by $e_h = I^{h}_{2h}e_{2h}$, and the fine grid solution is updated by adding this correction, $u_h = u_h + e_h$. Finally, additional smoothing steps are performed on the fine grid to further reduce any high frequency errors. The preconditioning matrix for two-grid solver is $C_\mathrm{2G} = I^{h}_{2h}A_{2h}^{-1}R^{2h}_h$.

More complex strategies for multigrid like V-cycle and W-cycle compute error corrections recursively across multiple grids of varying coarseness %(See \Cref{sec:multigrid} for pseudocode of these methods).

\section{Proofs for \Cref{sec:greedy}} \label{sec:proofgreedy}

\subsection{Proof of Proposition \ref{th:alpha=1}} \label{sec:proofalpha=1}

\begin{proposition} \label{th:alpha=1}
    Any prefix monotonically non-increasing sequence supermodular function $g$ is weakly-$\alpha$-supermodular with $\alpha = 1$
\end{proposition}

\begin{proof}

This proof is adapted from \citet{liberty2017greedy}.

If $g$ is sequence supermodular then,
\begin{align*}
    &g(S) - g(S \ \oplus S' ) \\
    &= \sum_{i = 1}^{|S'|} g(S \ \oplus \left(S'_1, \dots, S'_{i - 1}\right)) - g(S \ \oplus \left(S'_1, \dots, S'_{i}\right)) \\
    &\overset{(a)}{\leq} \sum_{i = 1}^{|S'|} g(S) - g(S \ \oplus  S'_{i}) \\
    &\leq |S'|\max_{i \in [|S'|]} g(S) - g(S \ \oplus  S'_{i})
\end{align*}

(a) by supermodularity


\end{proof}

\subsection{Proof of \Cref{th:suboptgreedy}} \label{sec:proofsuboptgreedy}



\begin{customtheorem}{\ref{th:suboptgreedy}}
    Let $g: \Omega^* \to \mathbb{R}$ be a weakly supermodular function with respect to the optimal solution $O = \argmin_{S \in \Omega^T} h(S)$ with a supermodularity ratio of $\alpha(O)$ and weak-$\mu$-postfix monotonicity. Let the greedy solution of length $T$ be $S^T$. Then,
    \begin{equation*}
        g(S^T) \leq  \mu\left(1 - \left(1 - \frac{1}{\alpha(O) T}\right)^T\right)
        g(O) +  \left(1 - \frac{1}{\alpha(O) T}\right)^Tf(\emptyset)
    \end{equation*}
    Equivalently, 
    \begin{equation*}
        \frac{g(\emptyset) - g(S^T)}{g(\emptyset) - g(O)} \geq \mu\left(1 - \left(1 - \frac{1}{\alpha T}\right)^T\right) \geq \mu\left(1 - \frac{1}{e^{\alpha}}\right)
    \end{equation*}
\end{customtheorem}

\begin{proof}

This proof strategy is inspired by \citet{streeter2008online} and \citet{liberty2017greedy}

\begin{align*}
    g(S^t) -g(O) &\overset{(a)}{\leq} g(S^t) - \frac{1}{\mu}g(S^t\ \oplus O) \\
    &= \left(1 - \frac{1}{\mu}\right)g(S^t) + \frac{1}{\mu}\left(g(S^t) - g(S^t\ \oplus O)\right) \\
    &=\left(1 - \frac{1}{\mu}\right)g(S^t) + \frac{1}{\mu}\sum_{i = 1}^{|O|} g(S^t \oplus \left(o_1, \dots o_{i - 1}\right)) -g(S^t\ \oplus \left(o_1, \dots, o_{i}\right))\\
    &\overset{(b)}{\leq} \left(1 - \frac{1}{\mu}\right)g(S^t) + \frac{\alpha(O)}{\mu}\sum_{i = 1}^{|O|} g(S^t) -g(S^t\ \oplus  o_{i}) \\
    &\leq \left(1 - \frac{1}{\mu}\right)g(S^t) + \frac{\alpha(O)|O|}{\mu}\max_{i \in [|O|]} g(S^t) -g(S^t\ \oplus  o_{i}) \\
    &\leq  \left(1 - \frac{1 - \alpha(O) T}{\mu}\right)g(S^t) -  \frac{\alpha(O) T}{\mu}\min_{\omega \in \Omega} g(S^t\ \oplus  \omega) \\
    &=  \left(1 - \frac{1 - \alpha(O) T}{\mu}\right)g(S^t) -  \frac{\alpha(O) T}{\mu}g(S^{t + 1})
    % &= Tg(S^t) -  Tg(S^{t + 1}) \\
    % &= T \left(g(S^t) - g(O)\right) - T\left(g(S^{t + 1}) -g(O)\right) \\
    % g(S^{t + 1}) -g(O) &\leq \left(1 - \frac{1}{T}\right)\left(g(S^t) - g(O)\right)
\end{align*}
(a) by $\mu$- postfix monotonicity, (b) by supermodularity. 

After rearranging the inequality, we get:

\begin{align*}
    g(S^{t + 1}) &\leq \frac{\mu}{\alpha(O) T}\left(g(O) -  \frac{1 - \alpha(O) T}{\mu}g(S^t)\right) \\
    &= \frac{\mu}{\alpha(O) T}g(O) + \left(1 - \frac{1}{\alpha(O) T}\right)g(S^t)
\end{align*}

When recursively applying this inequality, we get:

\begin{align*}
     g(S^{T}) &\leq \frac{\mu}{\alpha(O) T}g(O)\sum_{i = 0}^{T - 1}\left(1 - \frac{1}{\alpha(O) T}\right)^i + \left(1 - \frac{1}{\alpha(O) T}\right)^Tg(\emptyset) \\
     &= \mu \left(1 - \left(1 - \frac{1}{\alpha(O) T}\right)^T\right)g(O) + \left(1 - \frac{1}{\alpha(O) T}\right)^Tg(\emptyset)
\end{align*}
\end{proof}

\subsection{Proof of Proposition \ref{th:weaklyalphasupermodular}} \label{sec:proofweaklyalphasupermodular}


% \begin{customproposition}{\ref{th:weaklyalphasupermodular}}
%     If $\|I_N - C_j\mathcal{L}_h^a\| \leq 1$ for all $j \in [K]$, $h$ is (i) prefix monotonically non-increasing and (ii) f is weakly-$\alpha$-supermodular with $$\alpha = \max_{S \in [K]^*: |S| = T}\frac{\left\|I_N - \prod_{t = T}^{1}\left(I_N - C_{S_t}\mathcal{L}_h^a\right)\right\|}{T\left(1 - \max_{t \in [T]}\|I_N - C_{S_t}\mathcal{L}_h^a\|\right) }$$Furthermore, if $I_N - C_j\mathcal{L}_h^a$ is invertible for all $j \in [K]$, $h$ is also postfix monotonically non-increasing.
% \end{customproposition}

\begin{customproposition}{\ref{th:weaklyalphasupermodular}}
     Suppose that for all $j \in [K]$, the error propagation function $I - C_j \circ \mathcal{L}_h^a$ is $\rho_j$-Lipschitz continuous with $\rho_j < 1$, and that $(I - C_j \circ \mathcal{L}_h^a) (0_N) = 0_N$. Then, the function $h$ is weakly supermodular with respect to the optimal solution $O$, with $$\alpha(O) = \max\left\{\frac{4}{T - \sum_{i = 1}^{T} \rho_{O_i}^2} , 1\right\}  $$Furthermore, if $I - C_j \circ \mathcal{L}_h^a$ is invertible for all $j \in [K]$, $h$ is also postfix monotonically non-increasing.
\end{customproposition}

\begin{proof}
For brevity, we use the notation $\left(g_1 \circ \dots \circ g_T\right)(x) = \circ_{t = 1}^{T} g_t(x)$, where composition is applied from right to left so that $g_T$ acts first. In this proof, we use a few properties of Lipschitz continuous functions:
\begin{itemize}
    \item \textbf{Property 1:} If $g$ is $\rho$-Lipschitz continuous and $g(0) = 0$, then $\|g(x)\|_2 = \|g(x) - 0\|_2 = \|g(x) - g(0)\|_2 \leq \rho \|x - 0\|_2 = \rho\|x\|_2$
    \item \textbf{Property 2:} If $g_1$ and $g_2$ are Lipschitz continuous functions with Lipschitz constants of $\rho_1$ and $\rho_2$ respectively, the Lipschitz constant of $g_1 + g_2$ and $g_1 - g_2$ is $\rho_1 + \rho_2$.
    \item \textbf{Property 3:} If $g_1$ and $g_2$ are Lipschitz continuous functions with Lipschitz constants of $\rho_1$ and $\rho_2$ respectively, the Lipschitz constant of $g_1 \circ g_2$ is $\rho_1\rho_2$.
\end{itemize}

In order to prove weakly-$\alpha$-supermodularity, we must first prove prefix monotonicity.

\textbf{Prefix monotonicity: } Let $S \preceq S'$ where $S' = S \ \oplus N$.
\begin{align*}
    h(S') &= \left \|\circ_{t = |S'|}^1\left(I - C_{S'_t} \circ \mathcal{L}_h^a\right) \left(e^{(0)}_h\right)\right\|^2_2 \\
    &= \left \|\circ_{t = |S'|}^{|S| + 1}\left(I - C_{S'_t}  \circ \mathcal{L}_h^a\right)\circ \circ_{t = |S|}^1\left(I - C_{S'_t}  \circ \mathcal{L}_h^a\right) \left(e^{(0)}_h\right)\right\|^2_2 \\
    &\overset{(a)}{\leq}\left( \prod_{t = |S| + 1 }^{|S'|} \rho_{S'_t}^2 \right)\left\|\circ_{t = |S|}^1\left(I - C_{S_t} \circ \mathcal{L}_h^a\right) \left(e^{(0)}_h\right)\right\|^2_2 \\
    &\leq h(S)
\end{align*}
(a) by Property 1


\textbf{Weak supermodularity: } 
% \paragraph{ } 
We will upper bound $\alpha(O)$ by providing an upper bound for $ h(S) - h(S\ \oplus O) $ and a lower bound for $\sum_{i = 1}^T h(S) - h(S\ \oplus O_i) $. 
\begin{align*}
    h(S) - h(S\ \oplus O) &=  \left \|\circ_{t = |S|}^1\left(I - C_{S_t} \circ \mathcal{L}_h^a\right) \left(e^{(0)}_h\right)\right\|^2_2 -  \left \|\circ_{t = |O|}^{1}\left(I- C_{O_t} \circ\mathcal{L}_h^a\right) \circ \circ_{t = |S|}^1\left(I - C_{S_t} \circ \mathcal{L}_h^a\right)\left(e^{(0)}_h\right)\right\|^2_2 \\
    &\overset{(a)}{\leq} \left \|\circ_{t = |S|}^1\left(I - C_{S_t} \circ \mathcal{L}_h^a\right) \left(e^{(0)}_h\right)-\circ_{t = |O|}^{1}\left(I - C_{O_t} \circ \mathcal{L}_h^a\right) \circ \circ_{t = |S|}^1\left(I - C_{S_t} \circ \mathcal{L}_h^a\right)\left(e^{(0)}_h\right)\right\|^2_2 \\
    &=  \left \| \left(I -\circ_{t = |O|}^{1}\left(I - C_{O_t}\circ\mathcal{L}_h^a\right) \right)\circ  \circ_{t = |S|}^1\left(I - C_{S_t}\circ\mathcal{L}_h^a\right) \left(e^{(0)}_h\right)\right\|^2_2 \\
    &\overset{(b)}{\leq}  \left(1 + \prod_{t = 1}^{|O|}\rho_{O_t}\right)^2\left\| \circ_{t = |S|}^1\left(I - C_{S_t} \circ \mathcal{L}_h^a\right) \left(e^{(0)}_h\right)\right\|^2_2 \\
    &\overset{(c)}{\leq } 4 \left\| \circ_{t = |S|}^1\left(I - C_{S_t} \circ \mathcal{L}_h^a\right) \left(e^{(0)}_h\right)\right\|^2_2
\end{align*}
% \begin{align*}
%     h(S) - h(S\ \oplus S') &=  \left \|\circ_{t = |S|}^1\left(I - C_{S_t} \circ \mathcal{L}_h^a\right) \left(e^{(0)}_h\right)\right\|^2_2 -  \left \|\circ_{t = |S'|}^{1}\left(I- C_{S'_t} \circ\mathcal{L}_h^a\right) \circ \circ_{t = |S|}^1\left(I - C_{S_t} \circ \mathcal{L}_h^a\right)\left(e^{(0)}_h\right)\right\|^2_2 \\
%     &\overset{(a)}{\leq} \left \|\circ_{t = |S|}^1\left(I - C_{S_t} \circ \mathcal{L}_h^a\right) \left(e^{(0)}_h\right)-\circ_{t = |S'|}^{1}\left(I - C_{S'_t} \circ \mathcal{L}_h^a\right) \circ \circ_{t = |S|}^1\left(I - C_{S_t} \circ \mathcal{L}_h^a\right)\left(e^{(0)}_h\right)\right\|^2_2 \\
%     &=  \left \| \left(I -\circ_{t = |S'|}^{1}\left(I - C_{S'_t}\circ\mathcal{L}_h^a\right) \right)\circ  \circ_{t = |S|}^1\left(I - C_{S_t}\circ\mathcal{L}_h^a\right) \left(e^{(0)}_h\right)\right\|^2_2 \\
%     &\overset{(b)}{\leq}  \left(1 + \prod_{t = 1}^{|S'|}\rho_{S'_t}\right)^2\left\| \circ_{t = |S|}^1\left(I - C_{S_t} \circ \mathcal{L}_h^a\right) \left(e^{(0)}_h\right)\right\|^2_2 \\
%     &\overset{(c)}{\leq } 4 \left\| \circ_{t = |S|}^1\left(I - C_{S_t} \circ \mathcal{L}_h^a\right) \left(e^{(0)}_h\right)\right\|^2_2
% \end{align*}

(a) by reverse triangle property and $h(S) - h(S\ \oplus S') > 0$ by prefix monotonicity, (b) since the Lipschitz constant of $I -\circ_{t = |S'|}^{1}\left(I - C_{S'_t}\circ\mathcal{L}_h^a\right)$ is $1 + \prod_{t = 1}^{|S'|}\rho_{S'_t}$ by Property 2 and 3, (c) since $\rho_j < 1$


To lower bound $\sum_{i = 1}^{|O|} h(S) - h(S\ \oplus O_i) $
\begin{align*}
    \sum_{i = 1}^{|O|} h(S) - h(S\ \oplus O_i) &= \sum_{i = 1}^{|O|} \left \|\circ_{t = |S|}^1\left(I - C_{S_t} \circ \mathcal{L}_h^a\right) \left(e^{(0)}_h\right)\right\|^2_2 - \left \|\left(I - C_{O_i} \circ \mathcal{L}_h^a\right) \circ \circ_{t = |S|}^1\left(I - C_{S_t} \circ\mathcal{L}_h^a\right) \left(e^{(0)}_h\right)\right\|^2_2  \\
    &\geq \sum_{i = 1}^{|O|} \left \|\circ_{t = |S|}^1\left(I - C_{S_t} \circ\mathcal{L}_h^a\right) \left(e^{(0)}_h\right)\right\|^2_2 - \rho_{O_i}^2\left\|\circ_{t = |S|}^1\left(I - C_{S_t}\circ\mathcal{L}_h^a\right) \left(e^{(0)}_h\right)\right\|^2_2 \\
    &= \left \|\circ_{t = |S|}^1\left(I_N - C_{S_t}\circ\mathcal{L}_h^a\right) \left(e^{(0)}_h\right)\right\|^2_2 \left(T - \sum_{i = 1}^{T} \rho_{O_i}^2\right) 
\end{align*}
% \begin{align*}
%     \max_i h(S) - h(S\ \oplus S'_i) &= \left \|\circ_{t = |S|}^1\left(I - C_{S_t} \circ \mathcal{L}_h^a\right) \left(e^{(0)}_h\right)\right\|^2_2 - \min_i \left \|\left(I - C_{S_i'} \circ \mathcal{L}_h^a\right) \circ \circ_{t = |S|}^1\left(I - C_{S_t} \circ\mathcal{L}_h^a\right) \left(e^{(0)}_h\right)\right\|^2_2  \\
%     &\geq \left \|\circ_{t = |S|}^1\left(I - C_{S_t} \circ\mathcal{L}_h^a\right) \left(e^{(0)}_h\right)\right\|^2_2 - \min_i \rho_{S_i'}^2\left\|\circ_{t = |S|}^1\left(I - C_{S_t}\circ\mathcal{L}_h^a\right) \left(e^{(0)}_h\right)\right\|^2_2 \\
%     &= \left(1 - \min_i \rho_{S_i'}^2\right) \left \|\circ_{t = |S|}^1\left(I_N - C_{S_t}\circ\mathcal{L}_h^a\right) \left(e^{(0)}_h\right)\right\|^2_2 
% \end{align*}

Finally,

\begin{align*}
    \frac{h(S) - h(S\ \oplus O)}{T\max_i h(S) - h(S\ \oplus O_i)} &\leq \max\left\{ \frac{4\left\| \circ_{t = |S|}^1\left(I - C_{S_t} \circ \mathcal{L}_h^a\right) \left(e^{(0)}_h\right)\right\|^2_2}{\left \|\circ_{t = |S|}^1\left(I_N - C_{S_t}\circ\mathcal{L}_h^a\right) \left(e^{(0)}_h\right)\right\|^2_2 \left(T - \sum_{i = 1}^{T} \rho_{O_i}^2\right) }, 1\right\}\\
    &= \max\left\{\frac{4}{T - \sum_{i = 1}^{T} \rho_{O_i}^2} , 1\right\} 
\end{align*}

\textbf{Postfix monotonicity: } Let $S' = S \ \oplus N$. 

% Consider a function $\tilde{h}^{-1}(z)$ for $z \in \mathbb{R}^N$ be a function such that $\tilde{h}^{-1}(h(N)) = $

%Let $b$ a function such that $b \ \circ  \circ_{t = |N|}^{1}\left(I - C_{N_t}\circ\mathcal{L}_h\right) e^{(0)} = e^{(0)}$

\begin{align*}
    h(S') &= \left \|\circ_{t = |S'|}^1\left(I - C_{S'_t}\circ\mathcal{L}_h\right) e^{(0)}_h\right\|^2_2 \\
    &= \left \|\circ_{t = |N|}^{1}\left(I - C_{N_t}\circ \mathcal{L}_h\right)\circ \circ_{t = |S|}^1\left(I - C_{S_t}\circ \mathcal{L}_h\right) e^{(0)}_h\right\|^2_2 \\
    &\overset{(a)}{=}  \left \|\circ_{t = |N|}^{1}\left(I - C_{N_t} \circ \mathcal{L}_h\right)\circ \circ_{t = |S|}^1\left(I - C_{S_t}\circ\mathcal{L}_h\right)\circ\left(\circ_{t = |N|}^{1}\left(I - C_{N_t}\circ\mathcal{L}_h\right)\right)^{-1}\circ \circ_{t = |N|}^{1}\left(I - C_{N_t}\circ\mathcal{L}_h\right) e^{(0)}_h\right\|^2_2 \\
    &\leq \prod_{t = 1}^{|N|} \rho^2_{N_t} \prod_{t = 1}^{|S|} \rho^2_{S_t} \prod_{t = 1}^{|N|} \rho^{-2}_{N_t}  \left\|\circ_{t = |N|}^{1}\left(I - C_{N_t}\circ\mathcal{L}_h\right) e^{(0)}_h\right\|^2_2 \\
    % &\leq \left \|\prod_{t = |N|}^{1}\left(I_N - C_{N_t}\mathcal{L}_h\right)\right\|^2_2\left \|\prod_{t = |S|}^1\left(I_N - C_{S_t}\mathcal{L}_h\right)\right\|^2_2\left \|\prod_{t = |N|}^{1}\left(I_N - C_{N_t}\mathcal{L}_h\right)\right\|^{-2}_2\left \| \prod_{t = |N|}^{1}\left(I_N - C_{N_t}\mathcal{L}_h\right)e^{(0)}\right\|^2_2 \\
    % &\leq \circ_{t = |S|}^{1}\left \|\left(I - C_{S_t}\circ\mathcal{L}_h\right) \right\|^2_2 \left\|\circ_{t = |N|}^1\left(I - C_{N_t}\circ\mathcal{L}_h\right) e^{(0)}\right\|^2_2 \\
    &\leq h(N)
\end{align*}

% \begin{align*}
%     h(S') &= \left \|\circ_{t = |S'|}^1\left(I - C_{S'_t}\circ\mathcal{L}_h\right) e^{(0)}_h\right\|^2_2 \\
%     &= \left \|\circ_{t = |N|}^{1}\left(I - C_{N_t}\circ \mathcal{L}_h\right)\circ \circ_{t = |S|}^1\left(I - C_{S_t}\circ \mathcal{L}_h\right) e^{(0)}_h\right\|^2_2 \\
%     &\overset{(a)}{=}  \left \|\circ_{t = |N|}^{1}\left(I - C_{N_t} \circ \mathcal{L}_h\right)\circ \circ_{t = |S|}^1\left(I - C_{S_t}\circ\mathcal{L}_h\right)\circ\left(\circ_{t = |N|}^{1}\left(I - C_{N_t}\circ\mathcal{L}_h\right)\right)^{-1}\circ \circ_{t = |N|}^{1}\left(I - C_{N_t}\circ\mathcal{L}_h\right) e^{(0)}\right\|^2_2 \\
%     &\leq \prod_{t = 1}^{|N|} \rho^2_{N_t} \prod_{t = 1}^{|S|} \rho^2_{S_t} \rho^2(g)  \left\|\circ_{t = |N|}^{1}\left(I - C_{N_t}\circ\mathcal{L}_h\right) e^{(0)}\right\|^2_2 \\
%     % &\leq \left \|\prod_{t = |N|}^{1}\left(I_N - C_{N_t}\mathcal{L}_h\right)\right\|^2_2\left \|\prod_{t = |S|}^1\left(I_N - C_{S_t}\mathcal{L}_h\right)\right\|^2_2\left \|\prod_{t = |N|}^{1}\left(I_N - C_{N_t}\mathcal{L}_h\right)\right\|^{-2}_2\left \| \prod_{t = |N|}^{1}\left(I_N - C_{N_t}\mathcal{L}_h\right)e^{(0)}\right\|^2_2 \\
%     % &\leq \circ_{t = |S|}^{1}\left \|\left(I - C_{S_t}\circ\mathcal{L}_h\right) \right\|^2_2 \left\|\circ_{t = |N|}^1\left(I - C_{N_t}\circ\mathcal{L}_h\right) e^{(0)}\right\|^2_2 \\
%     &\leq h(N)
% \end{align*}

(a) due the invertibility of $I_N - C_j\mathcal{L}_h$

\end{proof}

\subsection{Proof of \Cref{th:simultaneousallocation}} \label{sec:proofsimultaneousallocation}

\begin{lemma} \label{th:simultaneousallocation}
     Let $\left(I - C_j \mathcal{L}_{h}^a\right) = P\Lambda_jP^{-1}$ where $P$ is an orthogonal matrix and $\Lambda_j = \mathrm{diag}(\lambda_{j1}, \dots, \lambda_{jN})$. If $P^{-1}e^{(0)}_h = z$, then the following equality holds:
     \begin{equation} \label{eq:simultaneous}
          h(S) = \sum_{i = 1}^N z_i^2 \prod_{j = 1}^K \lambda_{ji}^{2m_j(S)}
          % \qquad\mathrm{s.t}\qquad \sum_{k=1}^{K} m_k = T 
     \end{equation}
     where $m_j(S) = \sum_{t= 1}^{|S|}\1_{S_t = j}$
\end{lemma}

\begin{proof}

\begin{align*}
    h(S) &= \left\|\prod_{t = |S|}^1\left(I_N - C_{S_t} \mathcal{L}_{h}^a\right)  e^{(0)}_h \right\|^2 \\
    &=  \left\|\prod_{t = |S|}^1\left(P\Lambda_{S_t}P^{-1}\right)  e^{(0)}_h \right\|^2 \\
    &=  \left\|P\prod_{t = |S|}^1\Lambda_{S_t}P^{-1}  e^{(0)}_h \right\|^2 \\
    &= \left\|\prod_{t = |S|}^1\Lambda_{S_t}P^{-1}  e^{(0)}_h \right\|^2\\
    &=  \sum_{i = 1}^N z_i^2 \prod_{t = T}^1 \lambda_{S_ti}^2 \\
    &=  \sum_{i = 1}^N z_i^2 \prod_{j = 1}^K \lambda_{ji}^{2m_j(S)} 
    % &= \min_{\mathbf{m} \in\mathbb{N}^{K}}\quad \sum_{j = 1}^N z_i^2 \exp \left\{\sum_{k = 1}^K m_k \log \lambda_{kj}^{2}\right\}
    % \qquad\mathrm{s.t}\qquad \sum_{k=1}^{K} m_k = T
\end{align*}

\end{proof}

\subsection{Proof of \Cref{th:simultaneoussupermodular}} \label{sec:proofsimultaneoussupermodular}
\begin{customproposition}{\ref{th:simultaneoussupermodular}}
    Let $\|I_N - C_j\mathcal{L}_h^a\| \leq 1$ for all $j \in [K]$ and $\left(I - C_j \mathcal{L}_{h}^a\right) = P\Lambda_jP^{-1}$. Then, $h$ is supermodular. 
\end{customproposition}

\begin{proof}

Let $S \preceq S'$ where $S' = S \ \oplus B$. By \Cref{th:simultaneousallocation},
\begin{align*}
    h(S) &= \sum_{i = 1}^N z_i^2 \prod_{j = 1}^K \lambda_{ji}^{2m_j(S)} 
\end{align*}
where $m_j(S) = \sum_{t = 1}^{|S|} \1_{S_t = j}$ be the number of times a sequence $S$ calls the solver $j$. Recall that $h$ is considered sequence supermodular if $\forall S', S \in \Omega^*$ such that $S \preceq S'$, it holds that
\begin{align*}
    h(S) - h(S \ \oplus \omega ) \geq h(S') - h(S' \ \oplus \omega )
\end{align*}
% Without loss of generality, let's assume that $\omega$ that will be appended to $S, S'$ is $1$.
\begin{align*}
    h(S) - h(S \ \oplus \omega) &= \sum_{i = 1}^N z_i^2 \prod_{j = 1}^K \lambda_{ji}^{2m_j(S)} - \sum_{i = 1}^N z_i^2 \lambda_{\omega i}^{2}\prod_{k = 1}^K \lambda_{ji}^{2m_j(S)} \\
    &=  \sum_{i = 1}^N \left(1 - \lambda_{\omega i}^2\right) z_i^2 \prod_{j = 1}^K \lambda_{ji}^{2m_j(S)}
\end{align*}
Similarly, 
\begin{align*}
    h(S') - h(S \ \oplus \omega) &= \sum_{i = 1}^N \left(1 - \lambda_{\omega i}^2\right) z_i^2 \prod_{j = 1}^K \lambda_{ji}^{2m_j(S')} \\
    &\overset{(a)}{=}  \sum_{i = 1}^N \left(1 - \lambda_{\omega i}^2\right) z_i^2 \prod_{j = 1}^K \lambda_{ji}^{2\left(m_j(S) + m_j(B)\right)} \\
    &\overset{(b)}{\leq } \sum_{i = 1}^N \left(1 - \lambda_{\omega i}^2\right) z_i^2 \prod_{j = 1}^K \lambda_{ji}^{2m_j(S)} \\
    &= h(S) - h(S \ \oplus \omega)
\end{align*}


(a) Since $m_j(S') = \sum_{t = 1}^{|S'|} \1_{S_t' = j} = \sum_{t = 1}^{|S|} \1_{S_t' = j} + \sum_{t = |S|}^{|S'|} \1_{S_t' = j} = \sum_{t = 1}^{|S|} \1_{S_t = k} + \sum_{t = 1}^{|B|} 1_{B_t = j} = m_j(S) + m_j(B)$, (b) since $\rho(I_N - C_j \mathcal{L}_h^a) < 1$ and $m_j(B) \geq 0$ for all $j \in [K]$

\end{proof}


\section{Proofs for \Cref{sec:approxgreedy}}

\subsection{Proof of \Cref{th:routingequivalence}} \label{sec:proofroutingequivalence}

\begin{lemma} \label{th:routingequivalence}

For any set of preconditioning functions $\mathcal{C}$, any discrete operator $\mathcal{L}_h^a$, any router $r$, any $a_h, f_h, u_h^{(t)}, u_h \in \mathcal{A} \times \mathcal{F} \times \mathcal{U} \times \mathcal{U}$, the following equality holds true:
\begin{align*}
     l_{\text{route}}\left(r, a_h, f_h, u^{(t)}_h, u_h\right) = &\sum_{j = 1}^K \sum_{k = 1}^K\left\|\left( I - C_{k} \circ \mathcal{L}_{h}^a\right)\left(u_h - u^{(t)}_h\right)\right\|^2_2\1_{k\neq j}\1_{r(a_h, f_h, u^{(t)}_h) \neq j} \\
     &-\left(K - 2\right)\sum_{j = 1}^K\left\|\left( I - C_{j} \circ \mathcal{L}_{h}^a\right)\left(u_h - u^{(t)}_h\right)\right\|^2_2
\end{align*}
    
\end{lemma}

\begin{proof}
Note that $\sum_{j = 1}^K \1_{r(a_h, f_h, u^{(t)}_h) \neq j}= K - 1$
\begin{align*}
    l_{\text{route}}\left(r, a_h, f_h, u^{(t)}_h, u_h\right) &= \sum_{j = 1}^K \left\|\left( I - C_{j} \circ \mathcal{L}_{h}^a\right)\left(u_h - u^{(t)}_h\right)\right\|^2_2\1_{r(a_h, f_h, u^{(t)}_h) = j} \\
    &= \sum_{j = 1}^K \left\|\left( I - C_{j} \circ \mathcal{L}_{h}^a\right)\left(u_h - u^{(t)}_h\right)\right\|^2_2 - \sum_{j = 1}^K \left\|\left( I - C_{j} \circ \mathcal{L}_{h}^a\right)\left(u_h - u^{(t)}_h\right)\right\|^2_2\1_{r(a_h, f_h, u^{(t)}_h) \neq j} \\
    &= \sum_{j = 1}^K \left\|\left( I - C_{j} \circ \mathcal{L}_{h}^a\right)\left(u_h - u^{(t)}_h\right)\right\|^2_2 - \sum_{j = 1}^K \left\|\left( I - C_{j} \circ \mathcal{L}_{h}^a\right)\left(u_h - u^{(t)}_h\right)\right\|^2_2\1_{r(a_h, f_h, u^{(t)}_h) \neq j} \\
    &+ \left(K - 1\right)\sum_{k = 1}^K\left\|\left( I - C_{k} \circ \mathcal{L}_{h}^a\right)\left(u_h - u^{(t)}_h\right)\right\|^2_2 - \left(K - 1\right)\sum_{k = 1}^K\left\|\left( I - C_{k} \circ \mathcal{L}_{h}^a\right)\left(u_h - u^{(t)}_h\right)\right\|^2_2 \\
    &= \sum_{j = 1}^K \left\|\left( I - C_{j} \circ \mathcal{L}_{h}^a\right)\left(u_h - u^{(t)}_h\right)\right\|^2_2 - \sum_{j = 1}^K \left\|\left( I - C_{j} \circ \mathcal{L}_{h}^a\right)\left(u_h - u^{(t)}_h\right)\right\|^2_2\1_{r(a_h, f_h, u^{(t)}_h) \neq j} \\
    &+ \sum_{j = 1}^K\1_{r(a_h, f_h, u^{(t)}_h) \neq j}\sum_{k = 1}^K\left\|\left( I - C_{k} \circ \mathcal{L}_{h}^a\right)\left(u_h - u^{(t)}_h\right)\right\|^2_2 \\
    &- \left(K - 1\right)\sum_{k = 1}^K\left\|\left( I - C_{k} \circ \mathcal{L}_{h}^a\right)\left(u_h - u^{(t)}_h\right)\right\|^2_2 \\
    &= \sum_{j = 1}^K\left(\sum_{k = 1}^K\left\|\left( I - C_{k} \circ \mathcal{L}_{h}^a\right)\left(u_h - u^{(t)}_h\right)\right\|^2_2  - \left\|\left( I - C_{j} \circ \mathcal{L}_{h}^a\right)\left(u_h - u^{(t)}_h\right)\right\|^2_2\right)\1_{r(a_h, f_h, u^{(t)}_h) \neq j} \\
    &- \left(K - 2\right)\sum_{k = 1}^K\left\|\left( I - C_{k} \circ \mathcal{L}_{h}^a\right)\left(u_h - u^{(t)}_h\right)\right\|^2_2 \\
    &=\sum_{j = 1}^K\sum_{k = 1}^K\left\|\left( I - C_{k} \circ \mathcal{L}_{h}^a\right)\left(u_h - u^{(t)}_h\right)\right\|^2_2\1_{k\neq j}\1_{r(a_h, f_h, u^{(t)}_h) \neq j} \\
    &- \left(K - 2\right)\sum_{k = 1}^K\left\|\left( I - C_{k} \circ \mathcal{L}_{h}^a\right)\left(u_h - u^{(t)}_h\right)\right\|^2_2
 \end{align*}
\end{proof}


\subsection{Proof of \Cref{th:upperbound}} \label{sec:proofupperbound}

\begin{proposition} \label{th:upperbound}
    For any router $r$ defined by $r(a, f, u^{(t)}) = \text{argmax}_{j \in [K]} g_j(a, f, u^{(t)})$, any $a_h \in \mathcal{A}$, $f_h \in \mathcal{F}$, and $u_h^{(t)}, u_h \in \mathcal{U}$, the routing loss $l_{\text{route}}$ satisfies:
    $$\log(2)l_{\text{route}}\left(r, a_h, f_h, u^{(t)}_h, u_h\right)\leq \Psi(\mathbf{g}, a_h, f_h, u_h^{(t)}, u_h) $$
\end{proposition}


\begin{proof}

By \Cref{th:routingequivalence}, we know that

\begin{align*}
    \log(2)l_{\text{route}}\left(r, a_h, f_h, u^{(t)}_h, u_h\right) = &\log(2)\sum_{j = 1}^K \sum_{k = 1}^K\left\|\left( I - C_{k} \circ \mathcal{L}_{h}^a\right)\left(u_h - u^{(t)}_h\right)\right\|^2_2\1_{k\neq j}\1_{r(a_h, f_h, u^{(t)}_h) \neq j} \\
     &-\log(2)\left(K - 2\right)\sum_{j = 1}^K\left\|\left( I - C_{j} \circ \mathcal{L}_{h}^a\right)\left(u_h - u^{(t)}_h\right)\right\|^2_2 \\
     &\leq \log(2)\sum_{j = 1}^K \sum_{k = 1}^K\left\|\left( I - C_{k} \circ \mathcal{L}_{h}^a\right)\left(u_h - u^{(t)}_h\right)\right\|^2_2\1_{k\neq j}\1_{r(a_h, f_h, u^{(t)}_h) \neq j} \\
     &\overset{(a)}{\leq} -\sum_{j = 1}^K \sum_{k = 1}^K\left\|\left( I - C_{k} \circ \mathcal{L}_{h}^a\right)\left(u_h - u^{(t)}_h\right)\right\|^2_2\1_{k\neq j}\log\left(\frac{\exp\left(g_j(a, f, u^{(t)})\right)}{\sum_{k = 1}^K\exp\left(g_k(a, f, u^{(t)})\right)}\right) \\
     &= \Psi(\mathbf{g}, a_h, f_h, u_h^{(t)}, u_h)
\end{align*}
(a) if $r(a_h, f_h, u^{(t)}_h) \neq j$, $\frac{\exp\left(g_j(a, f, u^{(t)})\right)}{\sum_{k = 1}^K\exp\left(g_k(a, f, u^{(t)})\right)} < 0.5$ which implies that $-\log\left(\frac{\exp\left(g_j(a, f, u^{(t)})\right)}{\sum_{k = 1}^K\exp\left(g_k(a, f, u^{(t)})\right)}\right)  \geq \log(2)\1_{r(a_h, f_h, u^{(t)}_h) \neq j}$

\end{proof}


\subsection{Proof of \Cref{th:consistency}} \label{sec:proofconsistency}

\begin{customtheorem}{\ref{th:consistency}}
    Let $\left\|\left( I - C_{j} \circ \mathcal{L}_{h}^a\right)\left(e^{(t)}_h\right)\right\|^2_2 < \Bar{E} < \infty$  for all $j \in [K]$. If there exists $j\in [K]$ such that $\left\|\left( I - C_{j} \circ \mathcal{L}_{h}^a\right)\left(e^{(t)}_h\right)\right\|^2_2 > E_{min} > 0$, then, for any collection of solvers $\{C_j\}_{j= 1}^K$ and linear discrete operator $\mathcal{L}_h^a$, $\Psi$ is Bayes consistent surrogate for $l_{\text{route}}$.
\end{customtheorem}

\begin{proof}

For a given $a_h, f_h$, let $u_h$ be $\mathcal{G}_h\left(a_h, f_h\right)$ where $\mathcal{G}_h$ denotes the solution operator acting on the grid $G_h$. Furthermore, let's consider routers of the form $$r(a, f, u^{(t)}) = \text{argmax}_{j \in [K]} g_j(a, f, u^{(t)})$$

For a given $a_h, f_h, u_h^{(t)} \in \mathcal{A} \times \mathcal{F} \times \mathcal{U}$, let the optimal loss under $l_{\text{route}}$ be $l_{\text{route}}^*\left(a_h, f_h, u^{(t)}_h\right) = \inf_{\Tilde{r}} l_{\text{route}}\left(\tilde{r}, a_h, f_h, u^{(t)}_h, \mathcal{G}_h\left(a_h, f_h\right)\right)$. Similarly, let the optimal loss under $\Psi$ be $\Psi^*\left(a_h, f_h, u^{(t)}_h\right) = \inf_{\Tilde{\mathbf{g}}} \Psi\left(\tilde{\mathbf{g}}, a_h, f_h, u^{(t)}_h, \mathcal{G}_h\left(a_h, f_h\right)\right)$. Let $B_j(a_h, f_h, u^{(t)}_h) = \sum_{k = 1}^K\left\|\left( I - C_{k} \circ \mathcal{L}_{h}^a\right)\left(\mathcal{G}_h\left(a_h, f_h\right) - u^{(t)}_h\right)\right\|^2_2\1_{k\neq j}$.

\begin{align*}
    &l_{\text{route}}\left(r, a_h, f_h, u^{(t)}_h, \mathcal{G}_h\left(a_h, f_h\right)\right) - l_{\text{route}}^*\left(a_h, f_h, u^{(t)}_h\right) \\
    &\overset{(a)}{=} \sum_{j = 1}^K\sum_{k = 1}^K\left\|\left( I - C_{k} \circ \mathcal{L}_{h}^a\right)\left(u_h - u^{(t)}_h\right)\right\|^2_2\1_{k\neq j}\1_{r(a_h, f_h, u^{(t)}_h) \neq j}  -\left(K - 2\right)\sum_{j = 1}^K\left\|\left( I - C_{j} \circ \mathcal{L}_{h}^a\right)\left(u_h - u^{(t)}_h\right)\right\|^2_2 \\
    &\quad- \inf_{\Tilde{r}} \sum_{j = 1}^K\sum_{k = 1}^K\left\|\left( I - C_{k} \circ \mathcal{L}_{h}^a\right)\left(u_h - u^{(t)}_h\right)\right\|^2_2\1_{k\neq j}\1_{\tilde{r}(a_h, f_h, u^{(t)}_h) \neq j}  +\left(K - 2\right)\sum_{j = 1}^K\left\|\left( I - C_{j} \circ \mathcal{L}_{h}^a\right)\left(u_h - u^{(t)}_h\right)\right\|^2_2 \\
    &= \sum_{j = 1}^K B_j(a_h, f_h, u^{(t)}_h)\1_{r(a_h, f_h, u^{(t)}_h) \neq j} - \inf_{\Tilde{r}} \sum_{j = 1}^KB_j(a_h, f_h, u^{(t)}_h)\1_{\tilde{r}(a_h, f_h, u^{(t)}_h) \neq j}  \\
    &= \sum_{k = 1}^K B_k(a_h, f_h, u^{(t)}_h)\left(\sum_{j = 1}^K \frac{B_j(a_h, f_h, u^{(t)}_h)}{\sum_{k = 1}^K B_k(a_h, f_h, u^{(t)}_h)}\1_{r(a_h, f_h, u^{(t)}_h) \neq j} - \inf_{\Tilde{r}} \sum_{j = 1}^K \frac{B_j(a_h, f_h, u^{(t)}_h)}{ \sum_{k = 1}^K B_k(a_h, f_h, u^{(t)}_h)}\1_{\tilde{r}(a_h, f_h, u^{(t)}_h) \neq j}\right)
\end{align*}
(a) by \Cref{th:routingequivalence}

Let $\mathcal{X} = \mathcal{A} \times \mathcal{F} \times \mathcal{U}$ and $\mathcal{Y} = [K]$. Let $\mathcal{P}_{\mathcal{X}}$ denote the degenerate distribution suported at the point $(a_h, f_h, u_h^{(t)})$. We define the conditional distribtion - $P(Y = j \mid X = (a_h, f_h, u_h^{(t)})) = \frac{B_j(a_h, f_h, u^{(t)}_h)}{ \sum_{k = 1}^K B_k(a_h, f_h, u^{(t)}_h)}$ for $j \in [K]$. The risk and optimal risk of $0-1$ loss under this distribution can be written as:
\begin{align*}
    \mathcal{R}_{0-1}(r) &= \sum_{j = 1}^K \frac{B_j(a_h, f_h, u^{(t)}_h)}{\sum_{k = 1}^K B_k(a_h, f_h, u^{(t)}_h)}\1_{r(a_h, f_h, u^{(t)}_h) \neq j} \\
    \mathcal{R}_{0-1}^* &= \inf_r \sum_{j = 1}^K \frac{B_j(a_h, f_h, u^{(t)}_h)}{\sum_{k = 1}^K B_k(a_h, f_h, u^{(t)}_h)}\1_{r(a_h, f_h, u^{(t)}_h) \neq j} 
    % \mathcal{R}_{ce}(r) &= \sum_{j = 1}^K \frac{B_j(a_h, f_h, u^{(t)}_h)}{\sum_{k = 1}^K B_k(a_h, f_h, u^{(t)}_h)}\1_{f(a_h, f_h, u^{(t)}_h) \neq j} \\
    % \mathcal{R}_{ce}^* &= \inf_f \sum_{j = 1}^K \frac{B_j(a_h, f_h, u^{(t)}_h)}{\sum_{k = 1}^K B_k(a_h, f_h, u^{(t)}_h)}\1_{f(a_h, f_h, u^{(t)}_h) \neq j}
\end{align*}

If $r(a_h, f_h, u^{(t)}_h) = \argmax_{j \in [k]} g_j(a_h, f_h, u^{(t)}_h)$ for all $x \in \mathcal{X}$, then the he risk and optimal risk of cross entropy loss ($l_{ce}(\mathbf{g}, x, y)-\log\left(\frac{\exp\left(g_y(x)\right)}{\sum_{k = 1}^K\exp\left(g_k(x)\right)}\right)$) under this distribution can be written as:

\begin{align*}
    \mathcal{R}_{ce}(\mathbf{g}) &= -\sum_{j = 1}^K \frac{B_j(a_h, f_h, u^{(t)}_h)}{\sum_{k = 1}^K B_k(a_h, f_h, u^{(t)}_h)}\log\left(\frac{\exp\left(g_j(a_h, f_h, u^{(t)}_h)\right)}{\sum_{k = 1}^K\exp\left(g_k(a_h, f_h, u^{(t)}_h)\right)}\right) \\
    \mathcal{R}_{ce}^* &= \inf_{\mathbf{g}} -\sum_{j = 1}^K \frac{B_j(a_h, f_h, u^{(t)}_h)}{\sum_{k = 1}^K B_k(a_h, f_h, u^{(t)}_h)}\log\left(\frac{\exp\left(g_j(a_h, f_h, u^{(t)}_h)\right)}{\sum_{k = 1}^K\exp\left(g_k(a_h, f_h, u^{(t)}_h)\right)}\right) 
\end{align*}

From Theorem 3.1 of \citep{mao2023cross}, $\mathcal{R}_{0-1}(r) - \mathcal{R}_{0-1}^* \leq \Gamma^{-1}\left(\mathcal{R}_{ce}(\mathbf{g}) - \mathcal{R}_{ce}^*\right)$ if $r(a_h, f_h, u^{(t)}_h) = \argmax_{j \in [k]} g_j(a_h, f_h, u^{(t)}_h)$ where $\Gamma(z) = \frac{1 + z}{2}\log(1 + z) + \frac{1 - z}{2}\log(1 - z)$. Then,

\begin{align*}
    &l_{\text{route}}\left(r, a_h, f_h, u^{(t)}_h, \mathcal{G}_h\left(a_h, f_h\right)\right) - l_{\text{route}}^*\left(a_h, f_h, u^{(t)}_h\right) \\
    &= \sum_{k = 1}^K B_k(a_h, f_h, u^{(t)}_h)\left(\sum_{j = 1}^K \frac{B_j(a_h, f_h, u^{(t)}_h)}{\sum_{k = 1}^K B_k(a_h, f_h, u^{(t)}_h)}\1_{r(a_h, f_h, u^{(t)}_h) \neq j} - \inf_{\Tilde{r}} \sum_{j = 1}^K \frac{B_j(a_h, f_h, u^{(t)}_h)}{ \sum_{k = 1}^K B_k(a_h, f_h, u^{(t)}_h)}\1_{\tilde{r}(a_h, f_h, u^{(t)}_h) \neq j}\right) \\
    &\leq \sum_{k = 1}^K B_k(a_h, f_h, u^{(t)}_h)\Gamma^{-1}\left(-\sum_{j = 1}^K \frac{B_j(a_h, f_h, u^{(t)}_h)}{\sum_{k = 1}^K B_k(a_h, f_h, u^{(t)}_h)}\log\left(\frac{\exp\left(g_j(a_h, f_h, u^{(t)}_h)\right)}{\sum_{k = 1}^K\exp\left(g_k(a_h, f_h, u^{(t)}_h)\right)}\right) \right.\\&\quad\quad\quad\quad\quad\quad\quad\quad\quad\quad\quad\quad- \left.\inf_{\mathbf{g}} -\sum_{j = 1}^K \frac{B_j(a_h, f_h, u^{(t)}_h)}{\sum_{k = 1}^K B_k(a_h, f_h, u^{(t)}_h)}\log\left(\frac{\exp\left(g_j(a_h, f_h, u^{(t)}_h)\right)}{\sum_{k = 1}^K\exp\left(g_k(a_h, f_h, u^{(t)}_h)\right)}\right) \right) \\
    &\overset{(a)}{\leq} \Bar{E}K\left(K- 1\right)\Gamma^{-1}\left(-\sum_{j = 1}^K \frac{B_j(a_h, f_h, u^{(t)}_h)}{\sum_{k = 1}^K B_k(a_h, f_h, u^{(t)}_h)}\log\left(\frac{\exp\left(g_j(a_h, f_h, u^{(t)}_h)\right)}{\sum_{k = 1}^K\exp\left(g_k(a_h, f_h, u^{(t)}_h)\right)}\right) \right.\\&\quad\quad\quad\quad\quad\quad\quad\quad\quad\quad\quad\quad- \left.\inf_{\mathbf{g}} -\sum_{j = 1}^K \frac{B_j(a_h, f_h, u^{(t)}_h)}{\sum_{k = 1}^K B_k(a_h, f_h, u^{(t)}_h)}\log\left(\frac{\exp\left(g_j(a_h, f_h, u^{(t)}_h)\right)}{\sum_{k = 1}^K\exp\left(g_k(a_h, f_h, u^{(t)}_h)\right)}\right) \right) \\
    &\overset{(b)}{\leq} \Bar{E}K\left(K- 1\right)\Gamma^{-1}\left(-\sum_{j = 1}^K \frac{B_j(a_h, f_h, u^{(t)}_h)}{(K - 1)E_{min}}\log\left(\frac{\exp\left(g_j(a_h, f_h, u^{(t)}_h)\right)}{\sum_{k = 1}^K\exp\left(g_k(a_h, f_h, u^{(t)}_h)\right)}\right) \right.\\&\quad\quad\quad\quad\quad\quad\quad\quad\quad\quad\quad\quad- \left.\inf_{\mathbf{g}} -\sum_{j = 1}^K \frac{B_j(a_h, f_h, u^{(t)}_h)}{(K - 1)E_{min}}\log\left(\frac{\exp\left(g_j(a_h, f_h, u^{(t)}_h)\right)}{\sum_{k = 1}^K\exp\left(g_k(a_h, f_h, u^{(t)}_h)\right)}\right) \right) \\
    &= \Bar{E}K\left(K- 1\right)\Gamma^{-1}\left(\frac{\Psi\left(\mathbf{g}, a_h, f_h, u^{(t)}_h, \mathcal{G}_h\left(a_h, f_h\right)\right)- \Psi^*\left(a_h, f_h, u^{(t)}_h\right)}{(K - 1)E_{min}} \right)
\end{align*}
(a) since $\left\|\left( I - C_{j} \circ \mathcal{L}_{h}^a\right)\left(e^{(t)}_h\right)\right\|^2_2 < \Bar{E}$ for all $j \in [K]$, (b) since $\Gamma^{-1}$ is non-decreasing and $\exists j \in [K]$ such that $\left\|\left( I - C_{j} \circ \mathcal{L}_{h}^a\right)\left(e^{(t)}_h\right)\right\|^2_2 > E_{min}$

Finally,

\begin{align*}
    &\lim_{n \to \infty }\mathcal{R}_{\text{route}}\left(r_n\right) - \mathcal{R}_{\text{route}}^* \\
    &\overset{(a)}{=} \lim_{n \to \infty } \mathbb{E}_{a_h, f_h \sim \mathcal{P}_{\mathcal{A} \times \mathcal{F}}}\left[l_{\text{route}}\left(r_n, a_h, f_h, u^{(t)}_h, \mathcal{G}_h\left(a_h, f_h\right)\right) - l_{\text{route}}^*\left(a_h, f_h, u^{(t)}_h\right)\right] \\
    &\leq \lim_{n \to \infty } \mathbb{E}_{a_h, f_h \sim \mathcal{P}_{\mathcal{A} \times \mathcal{F}}}\left[\Bar{E}K\left(K- 1\right)\Gamma^{-1}\left(\frac{\Psi\left(\tilde{\mathbf{g}}, a_h, f_h, u^{(t)}_h, \mathcal{G}_h\left(a_h, f_h\right)\right)- \Psi^*\left(a_h, f_h, u^{(t)}_h\right)}{(K - 1)E_{min}} \right)\right] \\
    &\overset{(b)}{\leq} \lim_{n \to \infty } \Bar{E}K\left(K- 1\right)\Gamma^{-1}\left(\frac{\mathbb{E}_{a_h, f_h \sim \mathcal{P}_{\mathcal{A} \times \mathcal{F}}}\left[\Psi\left(\mathbf{g}_n, a_h, f_h, u^{(t)}_h, \mathcal{G}_h\left(a_h, f_h\right)\right)- \Psi^*\left(a_h, f_h, u^{(t)}_h\right) \right]}{(K - 1)E_{min}}\right) \\
    &= \lim_{n \to \infty } \Bar{E}K\left(K- 1\right)\Gamma^{-1}\left(\frac{\mathcal{R}_{\Psi}\left(\mathbf{g}_n\right) - \mathcal{R}_{\Psi}^*}{(K - 1)E_{min}}\right) \\
    &\overset{(c)}{=} \Bar{E}K\left(K- 1\right)\Gamma^{-1}\left(\frac{\lim_{n \to \infty } \mathcal{R}_{\Psi}\left(\mathbf{g}_n\right) - \mathcal{R}_{\Psi}^*}{(K - 1)E_{min}}\right) \\
    &= \Bar{E}K\left(K- 1\right)\Gamma^{-1}\left(0\right) \\
    &\overset{(d)}{=} 0
\end{align*}

(a) $\mathcal{R}_{\text{route}}^*  = \mathbb{E}_{a_h, f_h \sim \mathcal{P}_{\mathcal{A} \times \mathcal{F}}}\left[l_{\text{route}}^*\left(a_h, f_h, u^{(t)}_h\right)\right]$ since the infimum is taken over all measurable functions, (b) by Jensen's inequality since $\Gamma^{-1}$ is concave, (c) by continuity of $\Gamma^{-1}$ at $0$, (d) $\Gamma^{-1}(0) = 0$

% \begin{align*}
%     \mathcal{R}_{\Psi}\left(\mathbf{g}_n\right) = \mathbb{E}_{a, f \sim \mathcal{P}_{\mathcal{A} \times \mathcal{F}}}\left[\Psi\left(\mathbf{g}, a, f, u^{(t)}, u\right)\right]
% \end{align*}


\end{proof}



\section{Training Details} \label{sec:training details}


All the models were implemented using PyTorch and all the models were trained on an Nvidia A40 GPU. \Cref{tab:hpsettings} contains all hyperparameter details for the DeepONet.

\begin{table}[h]
    \centering
    \begin{tabular}{|c|c|}
        \hline
        Hyperparameter & Value\\
        \hline \hline
        Learning rate & 1e-3 \\
        Branch Dimension & 64 \\
        Hidden dimension for branch net & 128 \\
        No. of hidden layers in branch net & 2\\
        Hidden dimension for trunk net & 128 \\
        No. of hidden layers in trunk net & 2\\
        Gradient Clipping Norm & 1.0  \\
        Weight Decay & 0.005 \\
        Batch size & 128 \\
        Training samples & 15000 \\
        Epochs & 100 \\
        % Hidden layers & 1 & 3 & 1  \\
        % Hidden units & 8 & 128 & 8   \\
        % Dropout & 0.2 & 0.4 & 0.2   \\
        % Epochs & 200 & 100 & 200  \\
        % Training Samples & 2400 & 800 & 2400  \\
        % Gradient Clipping Norm & 1.0 & 1.0 & 1.0  \\
        % Weight Decay & 0.005 & 0.001 & 0.005  \\
        \hline
    \end{tabular}
    \caption{Hyperparameter settings for DeepONet}
    \label{tab:hpsettings}
\end{table}

The routers are long short-term memory (LSTM) models which were trained using Scheduled sampling. 


\section{Additional Experimental Results} \label{sec:moreexperimentalresults}
\subsection{Residual Comparison} \label{sec:residualcomparison}

\begin{table}[]
    \centering
    \begin{tabular}{|c|c|c|c|}
        \hline
         Methods/ Equation & 1D Poisson & 2D Poisson & 2D Helmholtz\\
         \hline
         \multicolumn{4}{|c|}{Jacobi-related Solvers} \\
        \hline
        Jacobi Only & 5156.608 (2729.501) & 4979.462 (1570.207) & NaN \\
        HINTS-Jacobi & 2511.313 (1113.55) & 6194.836 (480.136) & NaN \\
        Greedy-Jacobi & 2505.155 (1638.618) & 2097.546 (468.004) & NaN \\
        \hline
        \multicolumn{4}{|c|}{GS-related Solvers} \\
        \hline
        GS Only & 3272.362 (1741.517) & 2686.598 (858.228) & NaN \\
        HINTS-GS & 904.028 (394.014) & 4945.064 (339.648) & NaN \\
        Greedy-GS & 170.604 (57.765) & 998.784 (193.159) & NaN \\
        Greedy-GS,Jacobi & 389.426 (188.019) & 2721.092 (874.031) & NaN \\
         \hline
         \multicolumn{4}{|c|}{Multigrid methods} \\
         \hline
        MG Only & 819.828 (292.076) & 460.981 (111.146) & NaN \\
        HINTS-MG & 354.601 (149.921) & 1414.353 (73.994) & NaN \\
        Greedy-MG & 61.195 (17.887) & 284.36 (51.462) & NaN \\
         \hline
    \end{tabular}
    \caption{AUC of Residuals in $\times 10^{-3}$ over 300 iterations}
    \label{tab:residualcomparison}
\end{table}


\begin{figure}
    \centering
    \includegraphics[width=\linewidth]{images/residualcomparison1Dpoissonfinal.png}
    \caption{Residual Comparison}
    \label{fig:residualcomparison}
\end{figure}



\subsection{Fourier mode-wise error comparison} \label{sec:fouriermodeerrorcomparison}


\begin{table}[]
    \centering
    \begin{tabular}{|c|c|c|c|}
        \hline
         Methods/ Mode & Mode 1 Error & Mode 5 Error & Mode 10 Error\\
         \hline
         \multicolumn{4}{|c|}{Jacobi-related Solvers} \\
        \hline
        Jacobi Only & 732.713 (399.296) & 0.076 (0.043) & 0.001 (0.0) \\
        HINTS-Jacobi & 128.076 (64.871) & 0.348 (0.121) & 0.041 (0.021) \\
        Greedy-Jacobi & 4.152 (2.593) & 0.192 (0.094) & 0.025 (0.016) \\
        \hline
        \multicolumn{4}{|c|}{GS-related Solvers} \\
        \hline
        GS Only & 458.834 (250.064) & 0.223 (0.136) & 0.096 (0.072) \\
        HINTS-GS & 100.582 (52.598) & 0.176 (0.049) & 0.049 (0.028) \\
        Greedy-GS & 2.661 (1.595) & 0.12 (0.058) & 0.026 (0.015) \\
        Greedy-GS,Jacobi & 4.162 (2.504) & 0.158 (0.078) & 0.023 (0.015) \\
         \hline
         \multicolumn{4}{|c|}{Multigrid methods} \\
         \hline
        MG Only & 116.885 (42.581) & 0.043 (0.022) & 0.013 (0.007) \\
        HINTS-MG & 43.405 (20.908) & 0.027 (0.006) & 0.009 (0.003) \\
        Greedy-MG & 0.999 (0.622) & 0.029 (0.013) & 0.008 (0.005) \\
         \hline
    \end{tabular}
    \caption{AUC of Poisson 1d mode1, 5, 10 in $\times 10^{-3}$ over 300 iterations}
    \label{tab:residualcomparison}
\end{table}

% \begin{table}[]
%     \centering
%     \begin{tabular}{|c|c|c|c|c|c|c|c|c|c|}
%         \hline
%          Methods/ Equation & \multicolumn{3}{|c|}{1D Poisson} &\multicolumn{3}{|c|}{2D Poisson} & \multicolumn{3}{|c|}{2D Helmholtz}\\
%          \hline
%          &Mode 1 Error & Mode 5 Error  & Mode 10 error&Mode 1 Error & Mode 5 Error  & Mode 10 error&Mode 1 Error & Mode 5 Error  & Mode 10 error\\
%          \hline
         
%          \multicolumn{10}{|c|}{Jacobi-related Solvers} \\
%         \hline
%         Jacobi Only & 732.713 (399.296) & 0.076 (0.043) & 0.001 (0.0) & NaN & NaN & NaN & NaN & NaN & NaN \\
%         HINTS-Jacobi & 128.076 (64.871) & 0.348 (0.121) & 0.041 (0.021) & NaN & NaN & NaN & NaN & NaN & NaN \\
%         Greedy-Jacobi & 4.152 (2.593) & 0.192 (0.094) & 0.025 (0.016) & NaN & NaN & NaN & NaN & NaN & NaN \\
%         \hline
%         \multicolumn{10}{|c|}{GS-related Solvers} \\
%         \hline
%         GS Only & 458.834 (250.064) & 0.223 (0.136) & 0.096 (0.072) & NaN & NaN & NaN & NaN & NaN & NaN \\
%         HINTS-GS & 100.582 (52.598) & 0.176 (0.049) & 0.049 (0.028) & NaN & NaN & NaN & NaN & NaN & NaN \\
%         Greedy-GS & 2.661 (1.595) & 0.12 (0.058) & 0.026 (0.015) & NaN & NaN & NaN & NaN & NaN & NaN \\
%         Greedy-GS,Jacobi & 4.162 (2.504) & 0.158 (0.078) & 0.023 (0.015) & NaN & NaN & NaN & NaN & NaN & NaN \\
%         \hline
%          \multicolumn{10}{|c|}{Multigrid methods} \\
%          \hline
%         MG Only & 116.885 (42.581) & 0.043 (0.022) & 0.013 (0.007) & NaN & NaN & NaN & NaN & NaN & NaN \\
%         HINTS-MG & 43.405 (20.908) & 0.027 (0.006) & 0.009 (0.003) & NaN & NaN & NaN & NaN & NaN & NaN \\
%         Greedy-MG & 0.999 (0.622) & 0.029 (0.013) & 0.008 (0.005) & NaN & NaN & NaN & NaN & NaN & NaN \\
%          \hline
%     \end{tabular}
%     \caption{AUC of Modewise errors in $\times 10^{-3}$ over 300 iterations}
%     \label{tab:modewiseerrorcomparison}
% \end{table}


\begin{figure}
    \centering
    \includegraphics[width=\linewidth]{images/modecomparison1dpoissonfinal.png}
    \caption{Mode-wise Error Comparison}
    \label{fig:modewiseerrorcomparison}
\end{figure}

\end{document}

